% !TEX program = xelatex
\documentclass[12pt,a4paper,fontset=fandol]{ctexbook}

% ==================== 页面 ====================
\usepackage[top=2.5cm,bottom=2.5cm,left=2.5cm,right=2.5cm]{geometry}

% ==================== 字体 ====================
\usepackage{fontspec}
\usepackage{iffont}
% CJK 钉 fandol（两侧都取自 TeX 树里的同一份 FandolSong-Regular.otf）；
% mono 钉 DejaVu Sans Mono：CI 由 fontconfig 命中同名，本机没有该系统字体，
% 就用 MiKTeX dejavu 包里的同一版本 TTF（kpathsea 按文件名命中）。
\IfFontExistsTF{DejaVu Sans Mono}
  {\setmonofont{DejaVu Sans Mono}[Scale=MatchLowercase]}
  {\setmonofont[Extension=.ttf,
      UprightFont=DejaVuSansMono,
      BoldFont=DejaVuSansMono-Bold,
      ItalicFont=DejaVuSansMono-Oblique,
      BoldItalicFont=DejaVuSansMono-BoldOblique]
     {DejaVuSansMono}[Scale=MatchLowercase]}

% ==================== 颜色 ====================
\usepackage[dvipsnames,svgnames,x11names]{xcolor}
\definecolor{codebg}{HTML}{F5F5F5}
\definecolor{codeframe}{HTML}{E0E0E0}
\definecolor{cppkeyword}{HTML}{1A1AFF}
\definecolor{cppcomment}{HTML}{2E8B2E}
\definecolor{cppstring}{HTML}{B22222}
\definecolor{linenumgray}{HTML}{999999}
\definecolor{toolheadbg}{HTML}{2D3E50}
\definecolor{yamlkey}{HTML}{7A3E9D}
\definecolor{diffadd}{HTML}{1A7F37}
\definecolor{diffdel}{HTML}{B62324}

% ==================== listings ====================
\usepackage{listings}

% ---- 所有代码块共用的排版（先定义，各语言 style 组合它）----
\lstdefinestyle{codebase}{
  basicstyle=\small\ttfamily,
  keywordstyle=\color{cppkeyword}\bfseries,
  commentstyle=\color{cppcomment}\itshape,
  stringstyle=\color{cppstring},
  numberstyle=\tiny\color{linenumgray},
  numbers=left,
  frame=single,
  rulecolor=\color{codeframe},
  framesep=6pt,
  breaklines=true,
  tabsize=4,
  columns=flexible,
  keepspaces=true,
  backgroundcolor=\color{codebg},
  showstringspaces=false,
  escapeinside={(*@}{@*)},
}

\lstdefinestyle{cppstyle}{
  style=codebase,
  language=C++,
  morekeywords={constexpr,consteval,constinit,decltype,auto,
    concept,requires,co_await,co_yield,co_return,
    noexcept,override,final,nullptr,static_assert,
    template,typename,namespace,using,mutable,volatile,
    explicit,operator,friend,inline,virtual,
    thread_local,alignas,alignof},
  deletekeywords={int,char,bool,float,double,long,short,unsigned,signed,void,wchar_t},
  emph={size_t,int32_t,uint32_t,int64_t,uint64_t,ssize_t,
    string,vector,map,set,unique_ptr,shared_ptr,optional,variant,any,
    span,string_view,FILE,wchar_t},
  emphstyle=\color{teal!80!black},
}
\lstset{style=cppstyle}

% ---- C：寄存器、RTOS 对象与内核接口的 C ABI 边界 ----
\lstdefinestyle{cstyle}{
  style=codebase,
  language=C,
  morekeywords={bool,FILE,size_t,restrict,inline,volatile,sig_atomic_t,_Atomic},
  emph={size_t,int32_t,uint32_t,int64_t,uint64_t,ssize_t,
    TaskHandle_t,QueueHandle_t,SemaphoreHandle_t,EventGroupHandle_t,TimerHandle_t,
    BaseType_t,UBaseType_t,TickType_t,can_frame_t,sockaddr_can,
    ADC_HandleTypeDef,SPI_HandleTypeDef,UART_HandleTypeDef,I2C_HandleTypeDef,
    TIM_HandleTypeDef,DMA_HandleTypeDef,CAN_HandleTypeDef,GPIO_TypeDef},
  emphstyle=\color{teal!80!black},
}

% ---- YAML：设备树绑定、板级配置与 CI 描述文件 ----
\lstdefinestyle{yamlstyle}{
  style=codebase,
  morecomment=[l]{\#},
  morestring=[b]",
  morestring=[d]',
  alsoletter={-,:},
  morekeywords=[1]{Version,Checks,CheckOptions,HeaderFilterRegex,FormatStyle,
    User,WarningsAsErrors,CompilationDatabase,ExtraArgs,Add,Remove,
    Diagnostics,ClangTidy,UnusedIncludes,MissingIncludes,index,background-index,
    clang-tidy,InlayHints,Style,BasedOnStyle,IndentWidth,TabWidth,UseTab,
    ColumnLimit,BreakBeforeBraces,AlignAfterOpenBracket,AllowShortIfStatementsOnASingleLine,
    IncludeBlocks,SortIncludes,SeparateDefinitionBlocks,SpacesInAngles,
    Language,If,Else,Disable,Key,Name,files,exclude,repo,rev,hooks,id,entry,
    language_types,on,push_pull_request,steps,uses,runs-on,defaults},
  keywordstyle=[1]\color{yamlkey}\bfseries,
}

% ---- JSON：跟踪事件、机读输出与工具配置 ----
\lstdefinestyle{jsonstyle}{
  style=codebase,
  morestring=[b]",
  alsoletter={-},
  morekeywords={directory,file,output,arguments,command,version,
    cmakeMinimumRequired,name,generator,binaryDir,sourceDir,caches,
    toolchainFile,C_CXX_FLAGS,defineConstants,environments,
    editorMode,rulerColumn,tabSize,insertSpaces,files,settings,
    codec_name,codec_type,width,height,pix_fmt,r_frame_rate,avg_frame_rate,
    duration,bit_rate,sample_rate,sample_fmt,channel_layout,nb_frames,pkt_size},
  keywordstyle=\color{yamlkey}\bfseries,
}

% ---- Shell：命令行与工具调用 ----
\lstdefinelanguage{ShellLD}{
  morekeywords={if,then,else,elif,fi,for,in,do,done,while,case,esac,
    function,return,local,export,source,echo,cd,set,exec,trap,read,
    printf,test,exit,shift,eval,alias},
  sensitive=true,
  morecomment=[l]{\#},
  morestring=[b]",
  morestring=[d]',
}
\lstdefinestyle{shellstyle}{
  style=codebase,
  language=ShellLD,
  keywordstyle=\color{cppkeyword}\bfseries,
  emph={arm-none-eabi-gcc,arm-none-eabi-objdump,arm-none-eabi-nm,
    arm-none-eabi-size,arm-none-eabi-gdb,riscv64-unknown-elf-gcc,
    openocd,pyocd,JLinkGDBServer,st-flash,dfu-util,flashrom,dtc,west,idf.py,
    cyclictest,oslat,hwlatdetect,rt-migrate-test,chrt,taskset,
    stress-ng,trace-cmd,ethtool,udevadm,dmesg,modprobe,
    candump,cansend,canecho,i2cdetect,i2cget,i2cset,spidev_test,devmem2,
    lscpu,turbostat,cpupower,qemu-system-arm,renode,
    grep,sed,awk,xargs,jq,python3,pip,pip3,pwsh,bash,git,
    cmake,ninja,make,ccache,lldb,gdb,code,nvim,tar,mkdir,touch,ls,cat,
    pkg-config,ldconfig,strace,ltrace,perf,valgrind,gh,yamllint,wc,diff,find,sort,uniq,tee,rm},
  emphstyle=\color{teal!80!black}\bfseries,
}

% ---- CMake ----
\lstdefinelanguage{CMakeLD}{
  morekeywords={cmake_minimum_required,project,add_executable,
    add_library,target_link_libraries,target_include_directories,
    target_compile_definitions,target_compile_options,set,if,elseif,
    else,endif,find_package,option,message,install,file,
    target_sources,add_subdirectory,enable_language,set_target_properties,
    FetchContent_Declare,FetchContent_MakeAvailable,list,string},
  sensitive=false,
  morecomment=[l]{\#},
  morestring=[b]",
}
\lstdefinestyle{cmakestyle}{
  style=codebase,
  language=CMakeLD,
  keywordstyle=\color{cppkeyword}\bfseries,
}

% ---- Python：抖动统计、跟踪日志解析与离线分析脚本 ----
\lstdefinestyle{pythonstyle}{
  style=codebase,
  language=Python,
  morekeywords={self,lambda,async,await,yield,with,as,pass,del,
    global,nonlocal,is,in,not,and,or,True,False,None},
  keywordstyle=\color{cppkeyword}\bfseries,
  emph={av,subprocess,np,numpy,plt,json,os,sys,pathlib,struct,numpy},
  emphstyle=\color{teal!80!black},
}

% ---- 统一 diff：补丁式改动对比 ----
\lstdefinestyle{diffstyle}{
  style=codebase,
  numbers=none,
  frame=none,
  backgroundcolor=\color{white},
  morecomment=[l][\color{diffdel}]{-},
  morecomment=[l][\color{diffadd}]{+},
  morecomment=[l][\color{blue!70!black}\bfseries]{@@},
}

% ==================== tcolorbox ====================
\usepackage[most]{tcolorbox}

\newtcolorbox{guideline}[1][]{
  enhanced,breakable,
  colback=blue!5!white,colframe=blue!75!black,
  fonttitle=\bfseries,title=要点,#1,
  boxed title style={colback=blue!75!black},
  attach boxed title to top left={yshift=-2mm,xshift=2mm},
  arc=2mm,boxrule=0.8mm,
}
\newtcolorbox{compare}[1][]{
  enhanced,breakable,
  colback=green!3!white,colframe=green!50!black,
  fonttitle=\bfseries,title=对比,#1,
  boxed title style={colback=green!50!black},
  attach boxed title to top left={yshift=-2mm,xshift=2mm},
  arc=2mm,boxrule=0.8mm,
}
\newtcolorbox{warning}[1][]{
  enhanced,breakable,
  colback=red!5!white,colframe=red!75!black,
  fonttitle=\bfseries,title=常见陷阱,#1,
  boxed title style={colback=red!75!black},
  attach boxed title to top left={yshift=-2mm,xshift=2mm},
  arc=2mm,boxrule=0.8mm,
}
\newtcolorbox{note}[1][]{
  enhanced,breakable,
  colback=orange!5!white,colframe=orange!80!black,
  fonttitle=\bfseries,title=注意,#1,
  boxed title style={colback=orange!80!black},
  attach boxed title to top left={yshift=-2mm,xshift=2mm},
  arc=2mm,boxrule=0.8mm,
}
\newtcolorbox{bestpractice}[1][]{
  enhanced,breakable,
  colback=purple!5!white,colframe=purple!60!black,
  fonttitle=\bfseries,title=最佳实践,#1,
  boxed title style={colback=purple!60!black},
  attach boxed title to top left={yshift=-2mm,xshift=2mm},
  arc=2mm,boxrule=0.8mm,
}

% ==================== 章节标题 ====================
\usepackage{titlesec}
\usepackage{fancyhdr}
\usepackage{graphicx}
\usepackage{amssymb}
\usepackage{amsmath}
\usepackage{tabularx}
\usepackage{booktabs}
\usepackage{longtable}
\usepackage{array}
\usepackage{multirow}
\usepackage{enumitem}
\usepackage{seqsplit}
\usepackage{float}

\titleformat{\chapter}[display]
  {\normalfont\huge\bfseries\color{toolheadbg}}
  {\CTEXthechapter}{20pt}{\Huge}
\titleformat{\section}
  {\normalfont\Large\bfseries\color{blue!70!black}}
  {\thesection}{1em}{}
\titleformat{\subsection}
  {\normalfont\large\bfseries\color{blue!60!black}}
  {\thesubsection}{1em}{}

\pagestyle{fancy}
\fancyhf{}
\fancyhead[LE,RO]{\thepage}
\fancyhead[RE]{\leftmark}
\fancyhead[LO]{\rightmark}
\setlength{\headheight}{14.5pt}

\setlist[itemize]{leftmargin=2em,itemsep=2pt}
\setlist[enumerate]{leftmargin=2em,itemsep=2pt}
\emergencystretch=6em
% 本册不可断的框体与长表多，book 双页模式的 \flushbottom 会靠拉伸行距补页
\raggedbottom

% ==================== 超链接 ====================
\usepackage{hyperref}
\hypersetup{
  colorlinks=true,
  linkcolor=blue!70!black,
  urlcolor=blue!60!black,
  bookmarks=true,
  pdfauthor={C++ Reference},
  pdftitle={C/C++ 嵌入式与硬实时参考手册：启动链、RTOS 调度与确定性 I/O},
}

% ==================== 自定义命令 ====================
\newcommand{\cpp}[1]{C++#1}
\newcommand{\Fn}[1]{\texttt{#1()}}
\newcommand{\Tp}[1]{\texttt{#1}}
\newcommand{\Opt}[1]{\texttt{-\/-#1}}
\newcommand{\Bk}{\discretionary{}{}{}}
\newcommand{\Lib}[1]{\textbf{\texttt{#1}}}
\newcommand{\Lang}[1]{\textsf{#1}}
\newcommand{\CL}{\Lang{C}}
\newcommand{\CPP}{\Lang{C++}}
\newcommand{\PY}{\Lang{Python}}
% 域标签：本书讲"两侧对比"时统一用这四个写法，避免同义多形
% \DomBare/{\Bk}\DomRtos 标裸机与 RTOS 两条实现路线，\DomRT/{\Bk}\DomHost 标硬实时与通用主机侧
\newcommand{\DomBare}{\Lang{Bare}}
\newcommand{\DomRtos}{\Lang{RTOS}}
\newcommand{\DomRT}{\Lang{Realtime}}
\newcommand{\DomHost}{\Lang{Host}}

% ====================================================================
\begin{document}

\frontmatter

\begin{titlepage}
\centering
\vspace*{4cm}
{\Huge\bfseries C/C++ 嵌入式与硬实时参考手册\par}
\vspace{1cm}
{\LARGE 启动链、调度模型与确定性输入输出\par}
\vspace{0.4cm}
{\Large 裸机、FreeRTOS、Zephyr、PREEMPT\_RT 与实时以太网上的 C/C++\par}
\vfill
{\large \today\par}
\end{titlepage}

\tableofcontents

\mainmatter

\cleardoublepage

% ====================================================================
% 章节正文由 frag_chNN.tex 合并插入到下面这行之前（merge_frags.py）

% ============================================================
% ---- frag_ch01.tex ----
\part{裸机：启动链、时序与确定性 I/O}

\chapter{交叉工具链、链接脚本与启动时序}

\label{ch:ch01-toolchain-boot}

\DomBare{} \DomRT{} 一本讨论硬实时的书从编译器讲起，是因为实时系统交付的第一个成果物不是某个调度算法，而是一段``上电后必然走到已知状态''的代码。决定这段代码能否兑现时间承诺的东西，绝大部分不在函数体里：它藏在三元组选出的 ABI 里、藏在链接脚本给出的加载地址里、藏在 \Tp{Reset\_Handler} 执行的六个动作的顺序里、也藏在优化等级与 \Tp{volatile} 的相互作用里。本章把这条链条从\emph{构建}一路写到\emph{第一条指令}，并且只在``时序与有界性''这个视角下讨论它；寄存器位域该怎么读写属于硬件访问那一册的范围，本章不重复，只回答一个不同的问题：这段访问被安排在启动链的哪一步、它的时间代价是多少、它破坏了哪个不变量。

阅读顺序上，本章的 \S\ref{sec:ch01-triplet} 到 \S\ref{sec:ch01-size} 是自上而下的构建视角，\S\ref{sec:ch01-debug} 是闭环验证视角，\S\ref{sec:ch01-pitfalls} 是把前面所有机制反向压缩成的排查清单。若手上的板子已经能跑但启动时间不稳定，可以直接从 \S\ref{sec:ch01-startup} 的时钟树次序与 \S\ref{sec:ch01-pitfalls} 读起。

\section{目标三元组与工具链选择}

\label{sec:ch01-triplet}

GCC 的交叉驱动名是一个三元组（或多出 ABI 字段的四段名）：架构、厂商、操作系统或 ABI。\Tp{arm-none-eabi-gcc} 里厂商与 OS 字段都是 \Tp{none}，尾部的 \Tp{eabi} 才是真正有话说的部分：它声明目标遵循 ARM Embedded ABI，即 AAPCS 的调用约定、路径分隔符为 \Tp{/} 的名字形式、以及不含 Linux 系统调用层的裸机运行环境。\Tp{aarch64-none-elf} 是 AArch64 裸机的对应写法（Arm 官方较新的发行版把它改名成 \Tp{aarch64-none-elf-gnu} 一类的形态，用来与 \Tp{aarch64-none-linux-gnu} 区分），\Tp{riscv64-unknown-elf} 则是 RV64 裸机的写法，其中的 \Tp{unknown} 只是历史遗留的占位厂商字段。

真正影响时间行为的是三条开关：\Tp{-mcpu=cortex-m7} 决定微架构调度模型并隐含一个 \Tp{-march}；\Tp{-march=armv7e-m+fp.dp} 决定指令集可用性（未列出的扩展会被编译期拒绝，而不是留到运行期崩溃）；\Tp{-mfpu=fpv5-d16} 与 \Tp{-mfloat-abi=hard} 决定浮点参数走寄存器还是走栈。三者的组合不是自由排列的：\Tp{-mfpu} 必须与 \Tp{-mcpu} 声明的浮点单元一致，\Tp{-mfloat-abi=hard} 要求 \Tp{-march} 带上 \Tp{+fp} 或 \Tp{+fpv4} 一类的浮点扩展标记，否则较新的 GCC 会直接报错，而老版本只是安静地产生一个 ABI 标签不一致的目标文件。

从实时视角看，最需要记住的一条是：\emph{同一 ABI 字段下的目标文件与库文件不能混着链接}。soft-float 对象在属性段 \Tp{.gnu.attributes} 里记录的浮点 ABI 与 hard-float 对象不同，链接器可能只给出 \Tp{incompatible float ABI} 警告并继续产出镜像，在有 FPU 的核上表现为浮点参数被读成垃圾——这类故障的时间特征是``每次都不一样''，而不是``每次都崩在同一处''。诊断手段是固定的：用 \Tp{arm-none-eabi-readelf} \Tp{-A} \Tp{foo.o} 查看 \Tp{Tag\_ABI\_VFP\_args} 是否等于 \Tp{VFP registers}。

\subsection{multilib：一个驱动后面的多份库}

一个交叉工具链通常内置若干组预编译的运行库（不同架构、不同 Thumb/ARM 状态、soft 与 hard float 的笛卡尔积），这组矩阵称为 multilib。选错了不会立即失败：GCC 会根据命令行推导出一个 multilib 子目录，若该目录不存在就回落到默认库；若你手工加了 \Tp{-L} 或 \Opt{specs=}，就可能链接到与编译选项 ABI 不一致的那一份。可诊断的两条命令是 \Tp{-print-multi-directory}（打印本次选择的子目录）与 \Tp{-print-search-dirs}（打印搜索路径）。把这两条输出连同 \Tp{-dumpversion}、\Tp{-dumpmachine} 一起写进构建可复现性的记录，比事后凭记忆复原环境便宜得多。

\subsection{libc 选型与 \Tp{-nostdlib} 的代价}

newlib 是裸机上最常见的一份：它提供 \Fn{printf}、\Fn{malloc}、数学函数，代价是你要为它实现一层系统调用桩（\Tp{\_sbrk}、\Tp{\_read}、\Tp{\_write}、\Tp{\_close}、\Tp{\_exit} 等一整套下划线前缀函数），通常先链接 \Tp{libnosys}（桩全部返回失败）或 \Tp{rdimon}（半主机，见 \S\ref{sec:ch01-debug}）把流程跑通。picolibc 是 newlib 的分支，把桩统一成 \Tp{\_sys\_} 前缀的一套接口、把栈与堆的所有权交还给启动代码、并原生配合链接脚本提供的 \Tp{\_sidata}、\Tp{\_estack} 一类符号，因此在 Cortex-M 与 RISC-V 的新项目中越来越多。musl 面向有 MMU 与 ELF 加载器的 Linux 用户态，不适合裸机；裸机项目里真正可比的选项是 newlib、picolibc、以及完全不带 libc。

完全不带 libc 的写法是 \Tp{-nostdlib}（连启动文件与默认库都不带）或 \Tp{-nostartfiles} 加 \Tp{-nodefaultlibs}（只去掉一部分）。它省下的 flash 很实在，账单也一样清楚：(1) 没有 \Fn{malloc}，需要自己给出静态堆与分配器并证明其最坏情况时间；(2) 没有 \Fn{printf}，格式化调试要么整体关掉，要么自己写定长版本；(3) 没有 \Tp{\_\_libc\_init\_array} 之类的构造调用入口，C++ 全局对象的构造必须自己安排次序；(4) \Tp{assert}、\Tp{errno}、\Fn{setjmp}、浮点与字符串的互相转换全部消失。硬实时项目里常见的折中是：链接 newlib 但只用它的字符串与数学部分，把 \Fn{malloc} 用 \Opt{wrap} 重定向到一个静态池版本，从而在链接期就禁止不确定的分配。

\begin{lstlisting}[style=shellstyle,caption={裸机交叉编译与链接的完整命令行（Cortex-M7、hard float）},label={lst:ch01-gccline}]
arm-none-eabi-gcc -c -mcpu=cortex-m7 -mthumb \
    -mfpu=fpv5-d16 -mfloat-abi=hard \
    -ffunction-sections -fdata-sections -fno-common \
    -fno-strict-aliasing -Wall -O2 -g3 -std=c11 \
    -DUSE_HAL_DRIVER -I../Inc -I../Drivers \
    ../Core/Src/main.c -o build/main.o

arm-none-eabi-gcc -mcpu=cortex-m7 -mthumb -mfpu=fpv5-d16 -mfloat-abi=hard \
    -T../Ld/STM32H743ZITx_FLASH.ld -specs=nano.specs \
    -Wl,-Map=build/fw.map,--cref,--gc-sections,--print-gc-sections \
    -Wl,--wrap=malloc -Wl,--defsym=__flash_size=0x100000 \
    build/*.o -o build/fw.elf

arm-none-eabi-objcopy -O ihex build/fw.elf build/fw.hex
arm-none-eabi-objcopy -O binary --gap-fill 0xFF build/fw.elf build/fw.bin
\end{lstlisting}

清单 \ref{lst:ch01-gccline} 里有三处与时间直接相关。\Tp{-ffunction-sections} 与 \Tp{-fdata-sections} 把每个函数、每个全局对象放进独立段，配合 \Opt{gc-sections} 按段回收，这是裸机压缩 flash 的主力；代价是回收按可达性遍历，只被向量表引用的中断处理函数若链接脚本里没有保护，会被判定为不可达（见 \S\ref{sec:ch01-linker}）。\Tp{-fdata-sections} 还有一个副作用：同一结构体的成员不再保证物理相邻，看 \Tp{objdump -h} 会读不出原来的全局表，需要改看 map 文件。\Opt{gap-fill} 是给 NOR flash 的：擦除态就是 0xFF，填 0x00 会让烧写工具多做一次编程，某些外部 QSPI 器件的烧写时间因此明显变长。

\begin{table}[htbp]
\centering
\caption{三元组、工具链与 libc 的选型对照（实时视角）}
\label{tab:ch01-toolchain}
\begin{tabular}{>{\raggedright\arraybackslash}p{0.22\textwidth}
                >{\raggedright\arraybackslash}p{0.20\textwidth}
                >{\raggedright\arraybackslash}p{0.20\textwidth}
                >{\raggedright\arraybackslash}p{0.21\textwidth}}
\toprule
驱动 / 三元组 & 典型目标 & libc 与系统调用桩 & 实时视角的注意点 \\
\midrule
\Tp{arm-none-eabi-gcc} 加 \Tp{-mfloat-abi=hard} & Cortex-M3/M4F/M7、Cortex-R5 & newlib 加 \Tp{libnosys}，或 picolibc 的 \Tp{\_sys\_} 桩 & hard 与 soft 对象混链接只给警告；multilib 选错时符号能解析但 ABI 不成立 \\
\Tp{aarch64-none-elf} & Cortex-A53/A55 裸机、固件层 & picolibc；跑 Linux 时改 glibc & 计数器读 \Tp{CNTPCT\_EL0}，与内核频率无关，别假设与 SysTick 同源 \\
\Tp{riscv64-{\Bk}unknown-{\Bk}elf} 或 \Tp{riscv32-{\Bk}unknown-{\Bk}elf} & RV32EM/{\Bk}RV64MC 各类核 & newlib 或 picolibc，后者对 RISC-V 的支持已相当完整 & 中断入口由 \Tp{\_{\Bk}\_{\Bk}attribute\_{\Bk}\_{\Bk}({\Bk}({\Bk}interrupt))} 决定栈帧形状，不要手搓保存列表 \\
clang 加 \Opt{target=arm-{\Bk}none-{\Bk}eabi} \Tp{-{\Bk}mcpu=cortex-{\Bk}m4} & 与 GCC 相同的裸机目标 & 复用同一份 newlib 或 picolibc & 链接脚本语法共用；内建函数与内联汇编方言不同，CI 里要双跑对比 \\
musl 或 glibc（Linux 用户态） & Cortex-A 上的 Linux 应用 & 完整系统调用层 & 裸机那套 \Tp{-T} 链接脚本不适用；时间语义见第 2 章的软件时间基 \\
不带 libc：\Tp{-nostdlib} & 极小引导、二级 boot & 全部自己提供，入口用 \Tp{-e} 指定 & 失去 \Fn{malloc}、\Fn{printf}、全局构造；换来可控的段布局与确定的启动路径 \\
\bottomrule
\end{tabular}
\end{table}

\section{语言档位：仅 \CL{} 的裸机、可用 \CPP{} 的裸机与嵌入式系统}

\label{sec:ch01-langtier}

表~\ref{tab:ch01-toolchain} 定的是驱动与 libc，还剩一个在项目早期常被当作口味之争、事后极难回退的决定：这块板子允许写哪一档语言。习惯的分法按运行环境（裸机就是 \CL{}，有系统才谈 \CPP{}），但真正的门槛只有三条，而且都能在构建期查出来：链接的是哪一份 \CPP{} 运行库、有没有一个容量写得出了上界的堆、展开表与全局构造遍历这两笔开销能不能进预算。按这三条排下来，常见目标落在四档上；本节末尾的表~\ref{tab:ch01-langtier} 是速查件，下面逐档说清``为什么是这一档''与``越档写会付出什么''。

\subsection{档位一：只支持 \CL{} 的裸机}

这一档的成因不是芯片不行，而是 \CPP{} 需要的运行期支撑里至少缺一件，缺了是编译链接就断，而不是``写得勉强''。（1）没有 \CPP{} 运行库：\Tp{-nostdlib} 的构建、以及只发 C 库的厂商 SDK 里没有 \Tp{libstdc++} 或 \Tp{libc++} 对应的 multilib 变体，\Tp{new} 表达式找不到 \Tp{operator new}，虚表里引用的 \Tp{\_\_cxa{\Bk}\_pure{\Bk}\_virtual} 无人提供。（2）没有构造遍历入口：全局对象与函数局部静态对象依赖 \Tp{.init\_array} 与 \Tp{.ctors}（\S\ref{sec:ch01-startup}）以及 \Tp{\_\_cxa{\Bk}\_guard{\Bk}\_acquire}；启动代码不去遍历这些段时构造函数根本不会执行，症状是对象在、成员全零、一调用就跑飞。（3）没有可界堆：这一档常常连 \Fn{malloc} 都没有，而标准库容器的默认分配器就是它。（4）论证手段只吃 \CL{}：栈深度与时间界分析工具对间接调用、模板实例与展开路径给不出界，安全鉴定里的``语言子集''条款多数只承认受限的 \CL{} 子集。

典型对象：Cortex-M0/M0+/M3 上的小容量固件、ROM 程序与二级引导、只有 C 头文件的厂商 HAL、8 位机的 avr-libc 工程。这一档的正确写法不是用 \CL{} 模拟 \CPP{}（\Tp{void *} 配函数指针表手搓一张虚表，既拿到运行期开销又拿不到编译期检查），而是把 C 写成可移植的接口形态：定长缓冲、显式宽度类型、句柄式 API、每个结构体在类型里写明自己的上界。第 5 章的乒乓缓冲与第 9 章的单写单读环就是这个形态。

\subsection{档位二：可用 \CPP{} 子集的裸机}

门槛降到一条：链接期能拿到 \CPP{} 运行库。堆可以没有，异常与 RTTI 关掉。清单 \ref{lst:ch01-cppflags} 给出这一档最常见的命令行，四个开关各自删掉的东西不同，别把它们当成一组``裸机模板''照抄。

\begin{lstlisting}[style=shellstyle,caption={裸机 \CPP{} 子集：编译开关、链接形态与产物核对},label={lst:ch01-cppflags}]
arm-none-eabi-g++ -c -mcpu=cortex-m7 -mthumb \
    -mfloat-abi=hard -mfpu=fpv5-d16 -std=c++17 \
    -fno-exceptions -fno-rtti -fno-threadsafe-statics \
    -fno-use-cxa-atexit -ffunction-sections -fdata-sections \
    -Wall -O2 -g3 -I../Inc ../App/main.cpp -o build/main.o

# 常规形态：g++ 驱动自动带上 libstdc++ 与 libsupc++
arm-none-eabi-g++ -mcpu=cortex-m7 -mthumb -mfloat-abi=hard -mfpu=fpv5-d16 \
    -T../Ld/h743.ld -specs=nano.specs -Wl,--gc-sections \
    build/*.o -o build/fw.elf

# -nostdlib 形态：运行库必须手给，顺序不能反
#   -lstdc++ -lsupc++ -lm -lc -lnosys -lgcc

# 开关是否生效要在产物上核对，不是在命令行上核对
arm-none-eabi-readelf --section-headers build/fw.elf | grep -E 'exidx|extab'
arm-none-eabi-nm --undefined-only build/fw.elf | grep -E '__cxa|_Znw|guard'
\end{lstlisting}

\begin{itemize}
\item \Tp{-fno-exceptions}：\Tp{throw} 与 \Tp{try} 变成编译期错误；ARM 上不再为每个函数生成 \Tp{.ARM.exidx} 索引项（每函数 8 字节）与配套的 \Tp{.ARM.extab}。若只是``不主动抛但仍要能回溯''，用 \Tp{-funwind-tables} 单独把表留下。
\item \Tp{-fno-rtti}：\Tp{dynamic\_cast} 与对多态对象的 \Tp{typeid} 在编译期就失败；非多态类型的 \Tp{typeid} 是编译期结果，仍然可用。省下的是每个多态类的 \Tp{typeinfo} 对象与虚表里的类型信息指针。
\item \Tp{-fno-{\Bk}threadsafe-{\Bk}statics}：函数局部静态不再走守卫量，去掉对 \Tp{\_\_cxa{\Bk}\_guard\_*} 的依赖。单线程下语义不变；多核共享对象的初始化得自己保证只完成一次。
\item \Tp{-fno-use-cxa-atexit}：全局析构改登记到 \Fn{atexit}。裸机通常没有清理阶段，多数工程接受``析构不跑''这个后果，换来少一层运行库依赖。
\end{itemize}

这一档能直接用的一直是``不做分配、不擦除类型''的那部分标准库，而它恰好覆盖日常写作的大多数：\Tp{std::array} 与 \Tp{std::bitset}（长度写在类型里，常量表进 \Tp{.rodata}）、\cpp{17} 的 \Tp{string\_view} 与 \cpp{20} 的 \Tp{span}（视图，不持有内存）、\Tp{optional}、\Tp{variant}、\Tp{tuple}、\Tp{chrono} 的 duration 与换算（第 2 章的定点换算可以原样用它表达，不引入运行期调用）、\Tp{<algorithm>} 里在既有区间上就地工作的那些。最后一条有个容易记反的点：\Tp{std::sort} 不分配，但同一个 \Tp{std::transform} 配上 \Tp{back\_inserter} 就会往输出容器上要内存---决定代价的是输出迭代器，不是算法名。模板与 CRTP、无捕获 lambda（可隐式转函数指针，中断向量表靠静态表初始化就是这个形态）同样在这一档内。

不在这一档的是：容器与 \Tp{std::string}（默认分配器最终走到 \Tp{operator new}）、\Tp{std::function}（目标大于内联缓冲就上堆，还多一次间接跳转，理由见第 9 章）、异常（连同``只在最外层抛一次''的写法一起禁掉，别留这一层例外）、\Tp{std::regex}（代码量以几十 KB 起跳，递归深度随输入增长）、\Tp{iostream}（首次使用拉入 \Tp{std::ios{\Bk}\_base} 的静态初始化与几十 KB 的格式化实现，\Tp{sync{\Bk}\_with{\Bk}\_stdio} 默认还把它接到 stdio 缓冲上，裸机上等价于一次无界的外部调用）。

确实要容器时的两种换法：给容器换静态池分配器（把第 9 章那个池套一层即可，或用 \cpp{17} 的 \Tp{pmr} 把 \Tp{monotonic{\Bk}\_buffer{\Bk}\_resource} 钉在一块定长缓冲上），或者直接用定长实现；把全局 \Tp{operator new} 重定向到池，则把``能不能分配''变成``池里还剩几槽''。档位是否名副其实只能在产物上看：段表里没有 \Tp{.ARM.exidx} 与 \Tp{.ARM.extab}，未定义符号里没有 \Tp{\_\_cxa\_*} 与 \Tp{\_Znwj}（那是 \Tp{operator new} 的改名形式），才说明真的没要运行库。未定义符号里冒出这类名字，就说明某个头或某个模板实例悄悄越了档，而这一步在链接之前是看不见的。

典型对象：Cortex-M4F/M7 上的 \cpp{11} 到 \cpp{17} 工程、各家的模板工程、Arduino 一类的板级核心（它的运行时默认就带 \Tp{-fno-exceptions} 与 \Tp{-fno-rtti}）。

\subsection{档位三：可用大部分标准库的裸机与 \DomRtos{} 项目}

条件再降一条：有一个上界写得出来的堆，并且分配全部收敛到启动与配置阶段。到这一档，``能不能用标准库''已经不是编译器支持问题，而是时间语义问题：容器可用，但它们的代价是\emph{摊销有界}而不是\emph{最坏有界}。\Tp{std::vector} 的翻倍扩容会在某一次插入上做整表拷贝加一次释放；\Tp{std::map} 每次插入分配一个节点；\Tp{unordered\_map} 的再哈希一次搬运全部元素。于是实时路径上只有两种合法形态：\Tp{reserve} 到最大容量之后不再增长，或者容量由类型给出（第 9 章的池与环）。``启动期分配完、稳态只读''是使用这一档的唯一前提，它不是风格偏好，是把摊销代价从周期路径挪到配置阶段的那一步。

典型对象：带外部存储的大容量 Cortex-M7 采集与通信固件、Cortex-A/R 裸机里应用那一段（同一颗芯片的启动阶段仍属档位一，两档在一份固件里共存是常态）、FreeRTOS（\Tp{configTOTAL{\Bk}\_HEAP{\Bk}\_SIZE} 就是堆的上界，把全局 \Tp{operator new} 打到 \Fn{pvPortMalloc} 是标准做法）、Zephyr（\Tp{CONFIG{\Bk}\_CPP} 打开语言支持，\Tp{CONFIG{\Bk}\_CPP{\Bk}\_EXCEPTIONS} 与 \Tp{CONFIG{\Bk}\_CPP{\Bk}\_RTTI} 各自要求带堆的完整 libstdc++ 变体，最小化构建里没有展开器，见第 7 章）。

这一档仍然要付三笔账，而且都在构建期可见：\Tp{.init\_array} 的构造顺序与它在启动链里的位置（\S\ref{sec:ch01-startup}，跨翻译单元的依赖不可见）、展开表即使一次都不抛也照样进镜像并改变段布局（第 13 章 \S\ref{sec:ch13-compiler} 已从时间侧算过这笔）、以及看着很轻的对象其实不小（\Tp{mt19937} 的状态就是 2.5 KB；\Tp{std::cout} 的首次使用包含 iostream 与 locale 初始化，量级到毫秒，绝不能落在上电第一个周期里）。判据一句话：这一档用标准库换来的是类型安全与开发速度，失去的是``每个周期最坏做什么''那张表的可写性；两张表都必须在项目里同时存在---容器只在配置阶段出现，稳态路径继续按第 9 章的纪律走。

\subsection{档位四：嵌入式系统（带 \Tp{MMU} 的用户态）}

有 MMU、有完整 libc（glibc 或 musl）与系统调用层的目标---网关、域控制器、工控机上的用户态进程---语言侧不存在子集问题：异常、RTTI、全部容器与算法都能用。约束换了位置：时间上界不再由语言决定，而由内核抢占模型、页面驻留、CPU 隔离与优先级配额决定（第 8 章）。这一档要说清的不是``能不能''，而是四条跨过之后才知道代价的边界。

\begin{enumerate}
\item 交叉 \CPP{} 工具链必须与目标 sysroot 成对使用（发行版或构建系统给出的 SDK，或 \Opt{sysroot=} 指向的那一份）。拿主机的 \CPP{} 头去配目标上的运行库，失败发生在 ABI 层，症状是随机崩溃而不是链接错误。
\item 异常不能穿过 \CL{} 的 ABI 边界：从 C 侧注册的回调（中断线程、中间件的消息回调、驱动的辅助函数）往回抛是未定义行为；跨动态库边界抛还要求两侧的展开器与异常表一致。
\item 稳态循环里的分配在这里不再表现为分配失败，而表现为缺页、碎片整理与配额回收带来的毫秒级停顿；这正是第 8 章内存一节的内容，与语言无关。
\item 反过来有一条纪律常被忽略：项目规范里``禁用异常与 RTTI''是裸机档的规矩，套到全量标准库的目标上只会让评审把合规写法当缺陷。每条禁令都要连同它所属的档位一起写。
\end{enumerate}

\begin{longtable}{>{\raggedright\arraybackslash\footnotesize}p{0.12\textwidth}>{\raggedright\arraybackslash\footnotesize}p{0.22\textwidth}>{\raggedright\arraybackslash\footnotesize}p{0.24\textwidth}>{\raggedright\arraybackslash\footnotesize}p{0.25\textwidth}}
\caption{语言档位速查：三条门槛、典型对象与稳态路径上允许的件}\label{tab:ch01-langtier} \\
\toprule
档位 & 三条门槛上的形态 & 典型对象 & 稳态路径上允许的标准库件 \\
\midrule
\endfirsthead
\toprule
档位 & 三条门槛上的形态 & 典型对象 & 稳态路径上允许的标准库件 \\
\midrule
\endhead
\bottomrule
\endfoot
\endlastfoot
档位一：仅 \CL{} & 无 \CPP{} 运行库；无堆；无构造遍历入口 & M0/M0+/M3 小容量固件、ROM 与二级引导、只有 C 头文件的厂商 HAL & 无。定长结构体、显式宽度类型、句柄式 API、自管的定长环 \\
档位二：\CPP{} 子集 & 有 \CPP{} 运行库；无堆或堆不可界；异常与 RTTI 关 & M4F/M7 工程、厂商模板工程、板级核心 & \Tp{array}、\Tp{bitset}、\Tp{string\_view}、\Tp{span}、\Tp{optional}、\Tp{variant}、\Tp{chrono}、就地的 \Tp{algorithm}、模板与 CRTP \\
档位三：大部分标准库 & 有上界写得出的堆，分配收敛到配置阶段 & 大容量 M7 加外部存储、Cortex-A/R 的应用段、FreeRTOS 与 Zephyr 的 \Tp{CONFIG{\Bk}\_CPP} 配置 & 上一档全部，加 \Tp{vector}、\Tp{string}、\Tp{map}：\Tp{reserve} 到顶后不再增长，或换定长容器 \\
档位四：嵌入式系统 & \Tp{MMU} 加完整 libc，异常与 RTTI 全开 & 网关、域控、工控机上的用户态进程 & 全量。上界来自缺页、隔离与配额而非语言 \\
\end{longtable}

\begin{guideline}[title={定档三问，按顺序问}]
\begin{enumerate}
\item 链接的是哪一份运行库？交叉 \Tp{g++} 用 \Opt{print-multi-directory} 找不到对应的 \Tp{libstdc++} 变体，就先别写 \CPP{}；产物未定义符号里出现 \Tp{\_\_cxa\_*}，就说明已经越档。
\item 有没有堆，堆的上界写在哪里？写在 \Tp{configTOTAL{\Bk}\_HEAP{\Bk}\_SIZE}、静态池的字节数或一段预留映射上都算写了；写不出来的项目里容器只能在配置阶段出现。
\item 展开表、构造遍历与 iostream 这三笔开销，能同时进 flash 预算、启动时间预算与栈深度预算吗？三处任一处无人签字，就降一档，而不是``先写着，回头再瘦''。
\end{enumerate}
\end{guideline}

\begin{warning}[title={两个反方向的误判}]
一个是把``裸机就是 \CL{}''当公理：于是错过编译期多态（CRTP、模板参数加函数指针表）省下的那次间接跳转与它带来的栈深度不确定性，也错过用 \Tp{optional} 与 \Tp{span} 表达``可缺省''与``只读视图''的能力，代码退回到手搓虚表。另一个是反过来：在只有几十 KB flash、个位数 KB RAM 的 M0+ 上直接上 \Tp{std::string}、异常与 \Tp{iostream}，于是第一次分配与第一次输出稳稳地出现在最坏路径上，而镜像里那几十 KB 的展开表与格式化实现挤掉的正是功能段。两条来自同一件事---档位是目标环境的事实，不是团队的偏好；规范里每条语言禁令都要写明它管的是哪一档。
\end{warning}

\section{链接脚本：\Tp{MEMORY}、\Tp{SECTIONS} 与双地址模型}

\label{sec:ch01-linker}

链接脚本在实时项目里的地位接近硬件描述：它决定向量表落在哪个地址（这是硬件要求）、决定 \Tp{.data} 的加载地址与运行地址（这决定复位代码必须复制它）、也决定堆栈的预留上限（这决定最坏栈深度能否在链接期被证明）。语法骨架只有两层：\Tp{MEMORY} 声明可用的物理地址区间，\Tp{SECTIONS} 声明每个段被放到哪个区间。

\begin{lstlisting}[style=cstyle,caption={链接脚本片段：向量表、双地址 \Tp{.data}、带断言的堆栈预留},label={lst:ch01-linker}]
/* 地址与长度必须与器件手册一致；写错不会报错，只会随机跑飞 */
MEMORY
{
  FLASH  (rx)  : ORIGIN = 0x08000000, LENGTH = 1024K  /* 代码与只读数据 */
  RAM    (xrw) : ORIGIN = 0x20000000, LENGTH =   56K   /* SRAM1 + SRAM2 */
  DTCMRAM(xrw) : ORIGIN = 0x2001D000, LENGTH =   32K   /* 单周期访问 */
  ITCMRAM(xrw) : ORIGIN = 0x00000000, LENGTH =   64K
}

_estack = ORIGIN(RAM) + LENGTH(RAM);   /* 初始栈顶：栈向低地址生长 */

SECTIONS
{
  .isr_vector : {
    . = ALIGN(4);
    KEEP(*(.isr_vector))               /* 只被硬件引用，必须 KEEP */
    . = ALIGN(4);
  } >FLASH

  .text : {
    . = ALIGN(4);
    *(.text.after_vector)              /* 紧接向量表的早期处理函数 */
    *(.text) *(.text*) *(.glue_7) *(.glue_9)
    *(.init) *(.fini)
    . = ALIGN(4);
  } >FLASH

  .ramfunc : {
    _sramfunc = .;
    *(.ramfunc) *(.ramfunc.*)
    . = ALIGN(4);
    _eramfunc = .;
  } >RAM AT>FLASH                       /* 运行地址在 RAM，加载地址在 FLASH */

  .rodata : { *(.rodata) *(.rodata.*) } >FLASH
  .ARM.exidx : { *(.ARM.exidx*) } >FLASH  /* 有异常展开需求时才需要 */

  .data : {
    _sdata = .;                         /* 运行地址起点（VMA） */
    *(.data) *(.data.*)
    *(.dtcmram) *(.fast)
    . = ALIGN(4);
    _edata = .;
  } >RAM AT>FLASH                       /* 内容仍烧在 FLASH 里 */
  _sidata = LOADADDR(.data);            /* 加载地址起点（LMA） */

  .bss (NOLOAD) : {
    . = ALIGN(4);
    _sbss = .;
    *(.bss) *(.bss.*) *(COMMON)
    . = ALIGN(4);
    _ebss = .;
  } >RAM

  .noinit (NOLOAD) : { *(.noinit) *(.noinit.*) } >RAM

  ._user_heap_stack (NOLOAD) : {
    . = ALIGN(8);
    PROVIDE(end = .); PROVIDE(_end = .);  /* 运行库堆的起点 */
    . = _Min_Heap_Size;
    . = _Min_Stack_Size;
    . = ALIGN(8);
  } >RAM
}
\end{lstlisting}

\subsection{加载地址与运行地址：\Tp{AT>} 才是关键}

清单 \ref{lst:ch01-linker} 中的 \Tp{.ramfunc} 与 \Tp{.data} 体现同一个模式：\emph{VMA（运行地址）在 RAM，LMA（加载地址）在 FLASH}。写法 \Tp{>RAM AT>FLASH} 里前一个决定符号被链接到哪个地址（代码执行时用它访问），后一个决定字节实际被烧在镜像的哪个偏移上（复位代码复制时用它做源）。这两个地址混用是裸机最古老的事故来源之一：把 LMA 当运行地址读，读到的是那份还没被搬走的副本；反过来把运行地址当加载地址交给烧写工具，内容会被写到错误的偏移上。核对方法是固定的：\Tp{objdump -h} 每行的 LMA 与 VMA 两列应按意图成对出现，对 \Tp{.data}、\Tp{.ramfunc}、\Tp{.itcm} 这类段而言两列相等就说明 \Tp{AT>} 写漏了。

\subsection{对齐、页与 \Tp{KEEP}}

Cortex-M 要求向量表按 128 字节或按向量个数向上取整的 2 的幂对齐，\Tp{. = ALIGN(4)} 只保证字对齐，因此向量表段的起始地址要靠 \Tp{ORIGIN} 本身对齐，或在该段显式写 \Tp{. = ALIGN(0x400)}。flash 侧还有第二条约束：编程按页或扇区粒度进行（内部 flash 常见 128 字节页、1 到 4 KB 扇区，外部 QSPI NOR 常见 256 字节页、4 KB 扇区、64 KB 块），一个跨越页边界的大常量表会让烧写与 OTA 差分都要处理``同一页被两次改写''的情况；把大的只读表用 \Tp{SUBALIGN(256)} 或单独成段对齐，可以让差分包更贴合擦除边界。

\Tp{KEEP} 的作用是阻止回收：向量表 \Tp{.isr\_vector}、\Tp{*(.init)} 与 \Tp{*(.fini)}、\Tp{.init\_array} 与 \Tp{.ctors}（全局构造表）都必须 \Tp{KEEP}，因为它们只被硬件或运行时按地址引用，链接器的可达性分析看不到引用方。C 侧的可移植补充手段是 \Tp{\_\_attribute\_\_((used))}，它防的是编译器层的丢弃，而 \Tp{KEEP} 防的是链接器层的丢弃，两者不可互相替代。

\subsection{链接期符号在 C 侧的声明形态}

链接脚本给出的 \Tp{\_estack}、\Tp{\_sidata}、\Tp{\_sdata}、\Tp{\_edata}、\Tp{\_sbss}、\Tp{\_ebss} 是地址而不是值。C 侧必须按``数组首元素取址''的形态拿它们，绝不能声明成普通变量后直接参与运算：

\begin{lstlisting}[style=cstyle,caption={链接期符号的正确声明形态：取地址，不要取值},label={lst:ch01-syms}]
#include <stdint.h>

/* 这些符号只是链接器给出的地址；声明成数组后用下标 0 的取址 */
extern uint32_t _sidata[];   /* .data 在 FLASH 中的镜像起点（LMA） */
extern uint32_t _sdata[];    /* .data 在 RAM 中的起点（VMA） */
extern uint32_t _edata[];    /* .data 终点（不含） */
extern uint32_t _sbss[];
extern uint32_t _ebss[];
extern uint32_t _estack[];   /* 初始栈顶 */

/* 反例：extern uint32_t _edata; 然后把 _edata 当数用，拿到的是内容而不是地址 */
void reset_copy_and_clear(void)
{
    const uint32_t *src = _sidata;
    uint32_t *dst = _sdata;
    while (dst < _edata) *dst++ = *src++;          /* 复制 .data */
    for (dst = _sbss; dst < _ebss; ) *dst++ = 0U;  /* 清零 .bss  */
}
\end{lstlisting}

\begin{table}[htbp]
\centering
\caption{段属性与复位动作对照表}
\label{tab:ch01-sections}
\begin{tabular}{>{\raggedright\arraybackslash}p{0.17\textwidth}
                >{\raggedright\arraybackslash}p{0.18\textwidth}
                >{\raggedright\arraybackslash}p{0.26\textwidth}
                >{\raggedright\arraybackslash}p{0.22\textwidth}}
\toprule
段 & 运行地址 & 复位代码必须做的动作 & 漏做后的典型症状 \\
\midrule
\Tp{.isr\_vector} & FLASH（或 RAM 副本） & 无；但需 \Tp{KEEP} 与对齐，跳转进入时更新 \Tp{SCB->VTOR} & 取中断号读到旧表；进硬错误，栈上第一个返回地址是垃圾 \\
\Tp{.text} 与 \Tp{.rodata} & FLASH & 无 & —— \\
\Tp{.data} & RAM，LMA 在 FLASH & 按 \Tp{\_sidata} 复制到 \Tp{\_edata}；复制完成前不得读任何非零初值全局量 & 全局变量读到 0 或旧内容；状态标志恒为假 \\
\Tp{.bss} 与 \Tp{COMMON} & RAM & 全部写 0 & 未初始化 RAM 的随机值被当状态机；行为随温度与上电顺序变 \\
\Tp{.ramfunc}、\Tp{.itcm}、\Tp{.dtcm} & RAM 或 TCM & 与 \Tp{.data} 同：复制；有缓存的核还要做指令缓存失效（\Tp{ICIALLU} 一类） & 取指错误、执行到 FLASH 里的旧副本、测得的时序与实机不一致 \\
\Tp{.\_{\Bk}user\_{\Bk}heap\_{\Bk}stack} & RAM & 不初始化；只提供堆起点与链接期断言 & 堆起点落进 \Tp{.bss}；两次分配返回同一块内存 \\
\Tp{.noinit} 与 \Tp{.\_\_no\_\_init} & RAM & 明确不清零：\Tp{NOLOAD} 且不进清零循环 & 复位前保存的时间戳、日志环被清掉，取证信息消失 \\
\Tp{.init\_{\Bk}array} 与 \Tp{.preinit\_{\Bk}array} & FLASH 中的表加 RAM 中的对象 & 在时钟与 MPU 就绪后由运行时遍历调用 & C++ 全局对象未构造；调用到纯虚或半初始化的对象 \\
\bottomrule
\end{tabular}
\end{table}

\section{启动文件与 \Tp{Reset\_Handler} 的动作序列}

\label{sec:ch01-startup}

Cortex-M 上电后硬件做的事很少：从 \Tp{0x00000000}（或映射后的向量表基址）取第一个字作为初始 MSP，取第二个字作为复位入口直接跳过去，剩下的全是软件的债。由此推出一个必须写进设计文档的次序，之后任何一次改动都要按这张表复核。

\begin{enumerate}
  \item \textbf{装入 MSP。} 硬件已经装了默认值，但从 bootloader 进来的第二次启动必须显式重装，因为 bootloader 可能已经把主栈用得很深。
  \item \textbf{时钟先于一切。}\Fn{SystemInit} 之类的早期函数按``先设 flash 等待周期、再切频率''的次序执行：频率高于 flash 能支持的读取速度而没有相应等待周期时，取指会间歇性出错，症状是``代码在某条地址上随机跑飞''，看起来与硬件有关。等待周期之后是预取与指令加速器、分支预测的重新配置。
  \item \textbf{复制 \Tp{.data}、清零 \Tp{.bss}。} 顺序固定，在此之前任何 C 代码都不得读写全局对象；这条约束靠把复制放在任何函数调用之前来保证。
  \item \textbf{FPU 使能。} 对 Cortex-M4F/M7，写协处理器访问控制寄存器（CP10、CP11 设为完全访问，安全扩展下还有 NSACR）并置 \Tp{CONTROL.FPEN}；某些器件还需要在复位控制里保持浮点单元上电。跳过这一步的症状是执行到第一条浮点指令触发 \Tp{UsageFault}（NOCP 位置位）。
  \item \textbf{MPU 与故障使能先于全局构造。} 次序理由不是安全而是时间：构造函数一旦运行就可能读外设，MPU 设在其后就有一段无保护窗口；而把 MPU 设在时钟未稳定时，又会让总线错误报告的地址不可信。可用次序是（1）时钟（2）段初始化（3）MPU 与 fault 使能（4）构造。
  \item \textbf{运行时初始化与 C++ 静态构造。} newlib 与 picolibc 的 \Tp{\_\_libc\_init\_array} 遍历 \Tp{.preinit\_array}、\Tp{.init\_array}、\Tp{.ctors}；每一个构造函数都是一个``你不知道它会不会读寄存器''的调用。
  \item \textbf{进入 \Fn{main}。} 返回后不允许落到运行库的退出路径（它可能触发半主机），而是写成 \Tp{for (;;)} 加有节制的看门狗处理，或触发一次受控复位。
\end{enumerate}

\begin{lstlisting}[style=cstyle,caption={复位例程：向量表尾部入口、按序六个动作、\Fn{main} 返回后的死循环},label={lst:ch01-reset}]
/* --- 向量表：段名必须被链接脚本 KEEP，槽位顺序由器件手册给出 --- */
__attribute__((section(".isr_vector"), used))
const vector_table_t g_vectors = {
    .initial_sp  = (uint32_t)&_estack,
    .reset       = (uint32_t)&Reset_Handler,
    .nmi         = (uint32_t)&Default_Handler,
    .hardfault   = (uint32_t)&HardFault_Handler,
    /* ... 其余槽位同样指向 Default_Handler ... */
};

/* naked 保证编译器不插入 prologue，否则会改写刚装入的 MSP */
__attribute__((section(".text.after_vector"), naked))
void Reset_Handler(void)
{
    __asm volatile ("msr msp, %[sp]\n"
                    "bx %[early]\n"
                    :
                    : [sp] "r" (&_estack), [early] "r" (&early_main_stage)
                    : "memory");
}

void early_main_stage(void)
{
    const uint32_t *s;
    uint32_t *d;

    clock_early_init();          /* (2) flash 等待周期 -> PLL -> 回读校验频率 */
    for (s = _sidata, d = _sdata; d < _edata; ) *d++ = *s++;   /* (3) .data */
    for (           d = _sbss;   d < _ebss;  ) *d++ = 0U;      /* (3) .bss  */
    fpu_enable();                /* (4) CPACR 的 CP10/CP11 全访问 + FPEN */
    mpu_config();                /* (5) 区域表来自只读数据 */
    fault_enable();              /*     BusFault/MemManage/UsageFault 全开 */

    __libc_init_array();         /* (6) 全局构造：此刻才允许 I/O */
    main();                      /* (7) */
    for (;;) { /* 受控停住或复位 */ }
}
\end{lstlisting}

\subsection{tick 与中断优先级在启动期的安排}

SysTick 与 PendSV 的优先级必须在任何 RTOS 或时间敏感中断使能之前落定。一个常见事故是启动期用默认优先级 0 跑 SysTick，随后板级支持包把所有中断优先级改成``0 为最高''的自有约定，于是实时性最好的一段恰好落在 bootloader 的时间窗里。第 2 章会展开 tick 与计数器的一致性，这里只强调启动链上的一个动作：把优先级分组（\Tp{AIRCR.PRIGROUP}）与所有内核异常的优先级设定放在 \Tp{.bss} 清零之后、全局中断使能之前，并在 map 文件里留下这段代码独立成段的记录，避免被回收或被拆散。

\subsection{\Tp{-ffreestanding} 与 \Fn{main} 返回}

\Tp{-ffreestanding} 声明运行环境不保证 hosted 语义：没有 \Fn{main} 返回之后的清理、不能假设存在标准库、除 \Tp{sizeof(char)} 之外的类型宽度按实现定义。它通常与 \Tp{-fno-builtin}、\Tp{-fno-exceptions}、\Tp{-fno-rtti}、\Tp{-fno-common} 一组出现。裸机项目里把 \Fn{main} 写成返回 \Tp{int} 并真的返回，等于把控制权交给一个不该存在的地方（很多启动文件会掉进半主机的 \Tp{SYS\_EXIT}，无探针时表现为永久挂起），因此正确形态是 \Tp{void main(void)}，或者返回 \Tp{int} 但永不返回。

\begin{warning}[title={bootloader 跳转后没重装栈指针与 \Tp{VTOR}}]
从 DFU 或一级引导跳应用时，最容易漏的两个动作是：按应用的 \Tp{\_estack} 重装 MSP，以及把 \Tp{SCB->VTOR} 改成应用的向量表地址。漏前一项时应用沿用引导程序的栈区，深调用后写穿到别的模块；漏后一项时中断号取自旧表，症状是``主流程完全正常，一开中断就硬错误''。跳转前的收尾还应包括：关闭并清挂起的中断、停用 SysTick、复位用到的外设、把时钟切到应用假定的配置。
\end{warning}

\begin{guideline}[title={启动代码的优化等级与 \Tp{volatile}}]
启动段（复制、清零、时钟配置）用 \Tp{-O2} 而不是 \Tp{-O0}：这些循环全靠常量折叠与展开，\Tp{-O0} 时边界判断会生成真实比较指令，反而把栈深度与时间都拉长。但 \Tp{-O0} 会掩盖另一类问题，即 \Tp{volatile} 缺失：寄存器忙等循环在 \Tp{-O2} 下可能被整段删掉，\Tp{-O0} 下``能跑''，一旦切回 \Tp{-O2} 就少了一次 PLL 锁定等待。结论是把忙等里的状态读全部标成 \Tp{volatile}，并用 \Tp{-Wall} \Tp{-Wextra} 让优化等级的差异尽早暴露。
\end{guideline}

\section{体积、段与符号的分析}

\label{sec:ch01-size}

\Tp{arm-none-eabi-size} \Tp{-A} \Tp{-x} \Tp{fw.elf} 给出每个段的长度与汇总：\Tp{.text} 与 \Tp{.rodata} 影响 flash 占用，\Tp{.data} 加 \Tp{.bss} 之和与 RAM 容量相减就是可用的堆栈余量。\Tp{-A} 打印全部段，\Tp{-x} 用十六进制显示，\Tp{-m} 按成员文件细分，是``这个模块到底占了多少''的快速答案。要留意 size 报告的段大小是 VMA 上的字节数：\Tp{.ramfunc} 这类段的代码副本同样占 flash，而 \Tp{-A} 只报它运行时的大小；真实占用要看 \Tp{objdump -h} 的 LMA 列，或用 \Tp{objcopy -O binary} 之后按地址区间核算。

map 文件（用 \Tp{-Wl,-Map=} 生成）是唯一能回答``为什么它在这里''的材料。三个区域最值得读：头部的链接脚本回显（确认 \Tp{MEMORY} 的 \Tp{ORIGIN} 与预期一致，以及实际用到的库搜索路径）、中部的段布局（每个段的 VMA 与 LMA、以及并入该段的每个目标文件）、尾部的交叉引用表 \Tp{\# Cross Reference Table}（符号到定义文件的映射，用来回答``谁引用了这个 \Fn{printf}''）。\Opt{cref} 开启交叉引用表，\Opt{print-gc-sections} 打印被回收的段名，后者是``我的处理函数为什么变成了 \Tp{Default\_Handler}''的直接答案。

\subsection{调试符号分离与烧写产物}

带 \Tp{-g3} 的 ELF 可能有几 MB，而 \Tp{.text} 只有几十 KB；调试与烧写需要两个产物而不是一个。标准做法是：

\begin{lstlisting}[style=shellstyle,caption={符号分离：调试 ELF 与烧写镜像各留一份},label={lst:ch01-strip}]
# 1) 保留完整调试符号的副本；调试器只认这一个
arm-none-eabi-objcopy --only-keep-debug build/fw.elf build/fw-dbg.elf
# 2) 去掉调试段与符号，得到烧写用的紧凑 ELF
arm-none-eabi-objcopy --strip-debug --strip-unneeded build/fw.elf build/fw-stripped.elf
# 3) 把调试文件名与 CRC 写进 .gnu_debuglink，GDB 能自动关联
arm-none-eabi-objcopy --add-gnu-debuglink=build/fw-dbg.elf build/fw-stripped.elf
# 4) 只取代码相关段，供 OTA 差分或槽位校验使用
arm-none-eabi-objcopy --only-section=.isr_vector --only-section=.text \
                      --only-section=.rodata --only-section=.ramfunc \
                      build/fw.elf build/slot.bin
\end{lstlisting}

\Opt{only-keep-debug} 会把所有可加载段保留在文件中但清空其内容（只留符号），因此这个副本绝不能拿去烧写；GDB 加载 \Tp{fw-stripped.elf} 时顺着 \Tp{.gnu\_debuglink} 去找 \Tp{fw-dbg.elf}，找不到就退化成没有符号的会话。OTA 与量产场景里把 \Tp{.isr\_vector}、\Tp{.text}、\Tp{.rodata} 单独抽出，是为了让差分包与 flash 槽位边界对齐；\Opt{only-section} 列表必须显式包含 \Tp{.ramfunc}，否则 RAM 中执行的函数会静默缺失，而这类缺失要到第一次调用才暴露。

\begin{note}[title={flash 页与对齐对 OTA 分区的影响}]
分区表按擦除块（4 KB 或 64 KB，视器件而定）对齐只是前提；链接脚本里的对齐还要考虑编程粒度（常为 128 或 256 字节）。把一个三百字节的常量表放在页尾，会让该页在下一次差分写入时被整页重写，擦写寿命与写入时间都变差。给大的只读表单独成段并写 \Tp{SUBALIGN(256)}，是 OTA 项目里成本最低的改进。
\end{note}

\section{仿真与调试闭环}

\label{sec:ch01-debug}

四条路径按``能不能覆盖时间行为''排序，覆盖度从低到高：纯软件仿真（QEMU、Renode）、探针加 GDB server（OpenOCD、厂商 server）、片上计数打印（SWO/ITM），以及 GPIO 翻转配逻辑分析仪（第 2 章的打点一节）。

QEMU 的 Cortex-M 机型有限，常见可用的是 \Tp{mps2-an385}（Cortex-M3）、\Tp{mps2-an505}（Cortex-M33）、\Tp{microbit}（nRF51 系列）、以及 STM32 系列的 \Tp{netduino2} 与 \Tp{olimex--{\Bk}stm32--{\Bk}h405}；RISC-V 侧用 \Tp{virt}。加载 ELF 用 \Tp{-kernel}，附加镜像片段用 \Tp{-device loader}；调试用 \Tp{-s}（等价于在 1234 端口开 GDB 通道）与 \Tp{-S}（暂停等待接入）。它的价值在逻辑与状态机层面的回归，而不是时序：中断延迟、外设时序、flash 等待周期在 QEMU 里不存在或不真实，任何以微秒为单位的结论都不能从它得出。Renode 走的是另一条路：用平台描述与外设模型组合出多节点拓扑，可以在没有硬件时把通信协议栈连同每毫秒的调度一起跑起来，因此多用于 CI 中的协议级回归与多板联调。

探针侧的闭环由 GDB server 提供。OpenOCD 的典型启动是分别指定接口配置与目标配置两个文件，默认 GDB 端口 3333、telnet 3332、RPC 端口 6666；优点是脚本化（暂停、读 RAM 做时间戳、比较 \Tp{.noinit} 里的日志环都能自动化），缺点是连接与 flash 编程的稳定性受复位策略影响，通常需要显式的复位配置（例如只依赖硬件复位并保证 \Tp{srst} 可用）。厂商侧的对应物是 JLinkGDBServer 一类的独立 server，用设备名与接口（SWD 或 JTAG）、时钟速率作为启动参数；两者的实际差异主要体现在复位后能否立即接入，以及对半主机的支持程度上。

\subsection{半主机为什么不能上生产}

\Opt{specs=rdimon.specs}（newlib 的半主机实现，需要额外的初始化入口与系统调用层）让 \Fn{printf}、\Fn{malloc} 通过 \Tp{BKPT 0xAB} 陷入探针，由宿主机代做。它不能出现在量产固件里的原因有两条，第二条比第一条更常被忽略：(1) 该陷入依赖调试主机在场，没有探针时 CPU 停在陷入点永久等待，看门狗只能救回一次复位；(2) 一次半主机输出的最坏延迟由 USB 与宿主机决定，量级是毫秒到几十毫秒，而且完全无界。任何``把 printf 留着当诊断''的做法都会在最坏情况下破坏时间预算。量产的替代是 \Opt{specs=nosys.specs}（桩全返回失败，配合自己的环形缓冲）或不带 libc 加一个定长格式化函数。

SWO/ITM 提供一条不占用外设 UART、也不需要探针持续参与的打印通道：置 \Tp{DEMCR.TRCENA}（寄存器地址 \Tp{0xE000EDFC}）、使能 \Tp{ITM.TCR.TSE} 与 \Tp{ITM.TER} 的对应单元，然后向 \Tp{ITM\_PORT(0).DR} 写数据；数据由跟踪端口接口单元按内核时钟的固定分频（分频比在 \Tp{ACPR}）送出，主机端以几十兆比特率量级采样。它的定位要说清楚：它比半主机好得多（不陷入、单次写入几纳秒到几十纳秒），但打印仍然是\emph{观测扰动}——FIFO 会满，满了就静默丢；而且比特率与内核时钟绑定，一旦进入睡眠或降频，主机端采样率就对不上，出现的是乱码而不是丢字。做时序取证时，把 ITM 当作高带宽但不保证不丢的日志流；边界情况要交给 GPIO 翻转与逻辑分析仪。

\begin{compare}[title={四种观测手段的适用边界}]
QEMU 与 Renode：可进 CI、可复现、没有时间真值。OpenOCD 与厂商 GDB server：暂停会改变 flash 预取与总线状态，暂停期间 SysTick 与外设是否继续取决于时钟与调试配置，因此适合单步看状态而不适合测延迟。SWO 与 ITM：扰动小、可连续观测，但会丢，且随内核时钟变化。GPIO 加逻辑分析仪：唯一能给出硬件真实边界的通道，代价是需要引脚与走线，采样率要不低于被测带宽的十倍。
\end{compare}

\section{启动链上的常见陷阱与排查顺序}

\label{sec:ch01-pitfalls}

\begin{itemize}
  \item \textbf{\Tp{.bss} 未清零。} 三种成因：清零循环只覆盖了一个输入段而链接脚本合并了多个（\Tp{*(.bss)}、\Tp{*(.bss.*)}、\Tp{*(COMMON)} 必须都在循环的范围内）、把 \Tp{.noinit} 一起清了、以及把符号声明成值而不是地址（清单 \ref{lst:ch01-syms} 的反例）。核对方法是用 \Tp{objdump -h} 比较 \Tp{.bss} 的大小与清零循环跨越的地址区间。
  \item \textbf{双地址指针混用。} 凡 LMA 与 VMA 不等的段，其内部的函数指针与跳转表保存的都是 VMA；若把镜像里读到的常量当地址用，就指向了 FLASH 中的副本。LLVM 的链接脚本语法提供 \Tp{\_\_load\_start\_\_} 一类的显式加载地址符号，GNU 侧习惯是自造的 \Tp{\_sidata}。规范比机制重要：凡名字里带 \Tp{load} 或 \Tp{lma} 字样的符号，只允许出现在复制代码里。
  \item \textbf{静态构造里做 I/O。} 全局构造函数在 \Tp{\_\_libc\_init\_array} 里执行，此刻外设时钟可能还没开（多数器件的外设时钟默认关闭），写它的寄存器在某些总线上触发精确总线错误，在另一些上则静默等待几百微秒超时。判据很硬：构造函数里出现任何外设地址访问的，一律改成显式的 \Fn{bsp\_init}。
  \item \textbf{\Opt{gc-sections} 删掉了处理函数。} 只被向量表引用的处理函数在链接器看来不可达（表里存的是取址常量），解决办法是 \Tp{KEEP(*(.isr\_vector))} 保住表本身，并在 C 侧把处理函数声明成外部可见；更稳的做法是在链接脚本里显式收集这些段，或用 \Tp{-Wl,--undefined=Reset\_Handler} 强制引入。
  \item \textbf{栈预留只在链接脚本里检查。}\Tp{.\_user\_heap\_stack} 里的 \Tp{ASSERT} 只保证``预留了声明的量''，不保证``实际用量不超过''。补上 \Tp{-fstack-usage} 与 \Tp{-Wstack-usage=}，配合链接期的边界断言，再对最深的中断上下文做一次静态估算（第 3、4 章展开）。
  \item \textbf{优化等级差异导致启动偶发失败。} 现象是 Debug 配置（\Tp{-O0} \Tp{-g}）正常、Release（\Tp{-O2}）在冷启动或低电压下失败。三个常见原因：PLL 锁定的忙等被优化掉（缺 \Tp{volatile}）、\Tp{naked} 函数里漏了必要的寄存器保存、以及 flash 等待周期与频率不匹配（\Tp{-O0} 下取指节奏更慢从而被掩盖）。冷启动、低电压、边界温度这三个条件应当进硬件回归测试项。
\end{itemize}

\begin{bestpractice}[title={把启动动作做成上电自检清单}]
把 \S\ref{sec:ch01-startup} 的七个动作写成可执行的顺序表，每一项写``如何验证''而不是``如何配置''：MSP 落在预期区间内（读 \Tp{msp} 与两个链接符号比较）、实测频率与预期之差在容许的 ppm 内（用第 2 章的外设定时器测 SysTick 周期）、\Tp{.data} 抽样校验和与 FLASH 副本一致、\Tp{.bss} 抽样全零、FPU 使能后可执行一次无副作用的浮点、MPU 命中一次受控越界访问、构造表遍历次数等于预期对象数。量产固件里通常保留抽样版本，时间代价是几微秒到几十微秒。
\end{bestpractice}

排查顺序建议固定为五步：(1) 读 map 与 \Tp{objdump -h}，确认段的 VMA、LMA 与 \Tp{MEMORY} 一致；(2) 用 \Opt{print-gc-sections} 确认该在的段没被回收；(3) 在复位例程第一行用调试器看 \Tp{SP} 与 \Tp{PC}；(4) 段初始化完成后立刻校验抽样内容，把``复制与清零失败''和``后续逻辑错误''分开；(5) 最后才怀疑时序与硬件。裸机启动问题的绝大多数落在前三步，而前三步都不需要逻辑分析仪。

% ============================================================
% ---- frag_ch02.tex ----
\chapter{时间、计数器与时间戳}

\label{ch:ch02-time-counters}

\DomRT{} \DomBare{} \DomHost{} 硬实时系统的正确性有两种表达方式：逻辑上对不对，时间上有没有在期限内完成。前者由代码决定，后者由\emph{你用什么东西量时间}决定。本章讨论的就是这个量具：硬件计数器的位宽、频率、睡眠时的行为、在多核上是否同源、以及读它这个动作本身要花多少时间。一个反复出现的教训是：绝大多数``抖动变大''的调查报告最后指向的不是调度器，而是测量手段——时间戳来源切换了、频率标定过期了、或者把打点自身的开销算进了被测段。

本章先讲源头（\S\ref{sec:ch02-source}），再讲换算与位宽（\S\ref{sec:ch02-scale}、\S\ref{sec:ch02-wrap}），然后是取时间的动作本身（\S\ref{sec:ch02-probe}），最后是软件时间基与跨设备一致性（\S\ref{sec:ch02-swbase}、\S\ref{sec:ch02-sync}），\S\ref{sec:ch02-traps} 是压缩后的排查清单。与第 1 章的分工是：启动链决定计数器在什么时刻开始有意义（SysTick 的装载值、时钟树的最终频率、内核异常优先级的落定），本章则假定那些已经落定。

\section{硬件时间源}

\label{sec:ch02-source}

\subsection{Cortex-M SysTick}

SysTick 是一个 24 位向下计数器，装载值 \Tp{RELOAD} 最大 \Tp{0xFFFFFF}（16777215），计到 0 时重装、置 \Tp{COUNTFLAG}、并可选产生异常。由位宽直接推出三个结论：(1) 它每 $16777216 / f$ 秒溢出一次，180 MHz 时约 93.2 ms，因此任何超过 93 ms 的窗口都不能只看 SysTick 的当前值，必须数溢出次数；(2) 它的步进来自内核时钟或其八分频（\Tp{CTRL.CLKSOURCE} 选择），\emph{随频率变化}，一旦降频或进入节能状态，cycle 数与微秒数的换算就失效；(3) 装载值决定 tick 频率上限，1 kHz 的 tick 在 180 MHz 上需要装载 180000，远在 24 位范围内，而 1 MHz 的 tick 需要装载 180，此时每次中断的固定开销占比已不可忽略。

优先级是 SysTick 上最容易被写错的一项。FreeRTOS 一类的内核要求 tick 与上下文切换异常使用\emph{最低}优先级（在 \Tp{configKERNEL\_INTERRUPT\_PRIORITY} 里以移位后的原始优先级值表达），目的是让 tick 永远不打断临界区，从而保证切换的原子性。这个设计的代价是 tick 会丢：当 \Tp{PRIMASK} 被长时间置起（一段关中断的驱动）、或有更高优先级的中断连续占用 CPU 时，\Tp{COUNTFLAG} 仍会置位但异常被推迟，两个边沿挤在一起时只进一次中断，tick 就少计了一个。第 1 章里``优先级分组要早于任何中断使能''这条规矩，原因就是它。取舍不在于要不要丢 tick，而在于丢了怎么发现、怎么补偿，见 \S\ref{sec:ch02-swbase}。

\subsection{ARM architected timer}

Cortex-A 与部分实现里有架构定时器：64 位的 \Tp{CNTPCT\_EL0} 给出计数值，\Tp{CNTFRQ\_EL0} 给出频率（常见 24、25、27 MHz 或 1 GHz 中的某个，与内核时钟\emph{无关}），比较器 \Tp{CNTP\_CVAL\_EL0} 用来产生定时器中断，\Tp{CNTVCT\_EL0} 是加了虚拟偏移的视图。三个要点：(1) 频率必须运行时读 \Tp{CNTFRQ\_EL0}，不能从设备树字符串里猜，也不能假设等于 CPU 标称频率；(2) 一次读 64 位是原子的（一条读系统寄存器的指令），不存在 \S\ref{sec:ch02-wrap} 里 32 位核的拼接问题；(3) 它通常由常时电源域供电，核睡眠与降频都不影响，因此是``绝对时间''的天然候选。在 EL0 是否允许读它，取决于 \Tp{CNTKCTL} 与 \Tp{CNTHCTL} 的开放设置，裸机固件要显式打开；若需要与其他访存定序，读取前后要有屏障。

\subsection{x86 不变 TSC}

\Fn{\_rdtsc} 返回 \Tp{IA32\_TSC} 的 64 位值（\Tp{EDX} 与 \Tp{EAX} 两段），它不序列化：前面的指令可能还没执行完，后面的指令也可能已经乱序提前执行，所以裸 \Fn{\_rdtsc} 的边界是模糊的。\Fn{\_rdtscp} 会等待此前所有指令退休（它对后续指令定序，但不阻塞之前的乱序），并顺带返回 \Tp{IA32\_TSC\_AUX} 里的值；标准配对是（1）起点先 \Fn{\_mm\_lfence} 再 \Fn{\_rdtsc}，（2）终点用 \Fn{\_rdtscp} 再补一次 \Fn{\_mm\_lfence}。TSC 频率是否恒定由 CPUID 扩展叶 \Tp{0x80000007} 的 EDX 第 2 位声明；现代 Intel 与 AMD 上通常置位，此时计数率与核心频率解耦，但\emph{不等于}它等于标称频率：基准时钟常是 25 MHz 参考乘一个比例，实际值要用 \Tp{IA32\_TSC\_DEADLINE} 一类的定时器或 CPUID 的频率叶去标定。

TSC 有两个在多核与虚拟化下绕不开的缺陷。跨核偏差（skew）：各核的 TSC 在理想情况下同步，但被 \Tp{TSC\_ADJUST}（\Tp{MSR 0x6B1}）修改、深度睡眠唤醒、或热插拔之后可能不同步，内核的 \Tp{tsc} 时钟源只在确认同步之后才启用。虚拟化：客户机里的 TSC 可被宿主偏移与缩放（\Tp{kvm-clock} 一类的机制），迁移到另一物理核后可能重新对齐，读到的值在极端情况下出现回退。因此除非需要低于 vDSO 的打点开销（约 20 到 25 ns 量级），应当用 \Fn{clock\_gettime} 取 \Tp{CLOCK\_MONOTONIC}，把缩放与校正交给内核完成。

\subsection{外设独立定时器、RTC 与 PHC}

外设定时器由独立总线时钟分频而来，可配成自由运行的 16 或 32 位计数，配合捕获通道可以做与内核完全解耦的时间基准；它的价值在于\emph{即使 CPU 卡在硬错误里也仍在计数}，因此适合做外部可观测的时间锚。RTC 走 32.768 kHz 的常时电源域，精度是 ppm 级但跨越复位与掉电，适合绝对时间而不是分辨率。网络侧的真实时间来自网卡的 PHC（IEEE 1588 与 802.1AS 在 PHY 或 MAC 处打硬件时间戳），它是唯一一个``与链路上其他设备可比''的时间源，\S\ref{sec:ch02-sync} 展开。

\begin{longtable}{>{\raggedright\arraybackslash}p{0.15\textwidth}
                >{\raggedright\arraybackslash}p{0.16\textwidth}
                >{\raggedright\arraybackslash}p{0.17\textwidth}
                >{\raggedright\arraybackslash}p{0.17\textwidth}
                >{\raggedright\arraybackslash}p{0.16\textwidth}}
\caption{时间源特性对比：位宽、频率、时钟域、低功耗行为与读取代价}\label{tab:ch02-sources}\\
\toprule
时间源 & 位宽与典型频率 & 与内核时钟的关系 & 睡眠与掉电行为 & 单次读取的代价 \\
\midrule
\endfirsthead
\toprule
时间源 & 位宽与典型频率 & 与内核时钟的关系 & 睡眠与掉电行为 & 单次读取的代价 \\
\midrule
\endhead
SysTick & 24 位，装载上限 16777215 & 直接由内核时钟或其八分频驱动，降频即改变换算 & 内核时钟停则不计数；\Tp{COUNTFLAG} 需软件补偿 & 一条读指令，常数时间，无副作用 \\
外设定时器 & 16 或 32 位，总线时钟分频 & 与内核时钟无关 & 取决于该外设时钟是否被关；可配异步域 & 一次总线访问，跨慢速总线时几个到十几个周期 \\
DWT \Tp{CYCCNT} & 32 位，内核时钟 & 严格等于内核时钟，一个计数等于一个周期 & 调试运行位未置位时睡眠即停；调试电源域掉电后归零 & 一次内存映射读，约几个周期 \\
RTC & 多为 32 位秒计数或日历形式 & 由独立低速晶振驱动 & 常时电源域保持，跨复位与掉电 & 需要影子寄存器握手，慢且不能进中断 \\
ARM architected timer & 64 位，频率读 \Tp{CNTFRQ\_EL0} & 无关，运行时标定 & 通常常时域，核睡眠不影响 & 一条读系统寄存器指令，64 位原子 \\
x86 不变 TSC & 64 位，参考时钟乘固定比例 & 名义无关，实测需标定 & 浅睡仍走，深睡唤醒可能引入偏移 & 几十 ns；加栅栏指令后代价明显上升 \\
网卡 PHC & 多为 64 位纳秒 & 完全独立，需协议同步 & 独立于 CPU & 系统调用或收包控制消息，取决于路径 \\
\Tp{CLOCK\_{\Bk}MONOTONIC} & 64 位纳秒 & vDSO 下由 TSC 或时钟事件设备换算 & 不含挂起时长（与 \Tp{CLOCK\_{\Bk}BOOTTIME} 的区别） & vDSO 路径约 20 到 25 ns；走系统调用则高一个量级 \\
\bottomrule
\endfoot
\endlastfoot
\end{longtable}

\section{周期、频率与定点换算}

\label{sec:ch02-scale}

换算有三条硬规矩：不要用浮点、不要在周期路径上放除法、不要在运行期才决定频率。

先说浮点。单精度尾数是 24 位（含隐含位），能精确表示的整数上限是 16777216：400 MHz 的 cycle 计数在此之后连续整数开始重复（约 41.9 ms 就跨过边界），把 cycle 存进 \Tp{float} 就已经丢掉分辨率。同理，把跨越几分钟的时间戳换算成以秒为单位的 \Tp{float}，得到的分辨率是几十微秒，比计数本身粗几个量级。双精度有 53 位，能覆盖 64 位纳秒计数到几百年，但裸机上常常没有双精度硬件（\Tp{-mfpu=fpv4-sp-d16} 只有单精度），\Tp{double} 退化成软件浮点，一次乘法上千周期。

再说除法。Cortex-M0/M3 没有整数除法指令，\Tp{/} 是库调用，最坏情况下几百周期且随操作数值变化（newlib 的实现含循环，除数越大循环越长），这正是实时路径要避免的``数据依赖的执行时间''。若除数是编译期常量，编译器本来就会把除法换成乘倒数加移位；手写方案的真正价值在频率只有运行期才知道的场景（例如从 \Tp{CNTFRQ\_EL0} 或时钟树算出的 \Tp{sysclk}）：这时预先算一次乘数并缓存，路径上只剩乘法与移位。清单 \ref{lst:ch02-conv} 给出两种形态与误差分析。

\begin{lstlisting}[style=cstyle,caption={cycle 到微秒的定点换算：常量频率与运行期频率两种形态},label={lst:ch02-conv}]
/* 形态 A：频率是编译期常量。编译器自己会把除法变成乘倒数 + 移位。 */
#define SYSCLK_HZ   180000000UL
static inline uint32_t cyc_to_us_const(uint32_t cyc)
{
    /* 先乘后除，中间量必须 64 位 */
    return (uint32_t)(((uint64_t)cyc * 1000000UL) / SYSCLK_HZ);
}
/* 若频率恰是 1 MHz 的整数倍，连乘都可以省掉：180 MHz 时 1 us 等于 180 cycle */
static inline uint32_t cyc_to_us_180m(uint32_t cyc) { return cyc / 180UL; }

/* 形态 B：频率运行期才知道。初始化时算一次乘数，路径上只有乘和移位。
 * 约定：magic = (1000000 << K) / freq，结果 = (cyc * magic) >> K。
 * K 越大精度越好，但 cyc 的可用上限越小。 */
#define CONV_SHIFT  20U
typedef struct { uint32_t magic; uint32_t freq_hz; } us_scale_t;

void us_scale_init(us_scale_t *s, uint32_t freq_hz)
{
    s->freq_hz = freq_hz;
    /* 例：180 MHz -> magic = 5825（真值 5825.4，截断误差约万分之零点七） */
    s->magic = (uint32_t)(((uint64_t)1000000U << CONV_SHIFT) / freq_hz);
}
static inline uint32_t cyc_to_us_scaled(const us_scale_t *s, uint32_t cyc)
{
    /* 64 位中间量：magic 为 5825 时，32 位只能撑到约 4.1 ms 的 cyc */
    return (uint32_t)(((uint64_t)cyc * s->magic) >> CONV_SHIFT);
}
/* 四舍五入版：加半个最低位再移位，消除单调偏小的系统性偏差 */
static inline uint32_t cyc_to_us_round(const us_scale_t *s, uint32_t cyc)
{
    return (uint32_t)((((uint64_t)cyc * s->magic) + (1U << (CONV_SHIFT - 1U)))
                      >> CONV_SHIFT);
}
\end{lstlisting}

清单 \ref{lst:ch02-conv} 里两个数值值得记住量级：180 MHz 与移位 20 位时，截断带来的相对误差约 0.0069\%，在一秒窗口上是约 12 µs 的偏差；同样参数下 32 位中间量只能安全表示约 4.1 ms 的窗口，超出就必须换成 64 位乘法或先分段。若一定要留在 32 位，另一个方案是把频率约分成整数比（72 MHz 时 1 µs 恰为 72 个周期），此时零误差且只有一条移位。

tick 频率与帧率的错频是显示与控制回路里常见的时间噪声来源。\Tp{configTICK\_{\Bk}RATE\_{\Bk}HZ} 取 1000 Hz 而目标帧率取 120 Hz 时，一帧的理论长度是 8.33 个 tick，软件只能取整成 8 或 9，于是每帧出现约一个 tick（1 ms）的周期性抖动，而这是渲染管线上最难查的抖动。可行方案按优先级有三种：把 tick 频率与帧率取成整数倍关系（1 kHz 配 100 Hz、500 Hz 配 125 Hz）、帧截止用硬件定时器而不用 tick、或者保留错频但在显示侧用垂直同步的硬件时间戳重排。选哪种取决于你愿意把不确定性放在哪一段。

最后是把\textbf{抖动统计口径}固定下来，这一条属于流程而不是代码。报告里必须同时给出微秒与 cycle 两种单位（cycle 才反映真实工作量，微秒才反映用户感受），必须注明采样频率与频率来源（\Tp{CNTFRQ\_EL0}、标称 \Tp{SYSCLK\_HCK}、还是实测标定），必须给出分布量而不是均值（最大值、99 分位、中位数三个一组）。均值在双峰分布上是没有意义的：一个每周期正常、偶尔被 100 µs 的中断阻塞的回路，其平均延迟看起来可以完全达标。

\begin{guideline}[title={打点数据的单位与时钟域写进字段名}]
把变量命名成 \Tp{cycles}、\Tp{us\_from\_cntpct}、\Tp{ns\_monotonic}、\Tp{us\_local\_tsc}，让时钟域出现在名字里。跨域换算只允许发生在一个函数里（清单 \ref{lst:ch02-conv} 的那几个），其余位置禁止出现``裸数字乘一千''这类隐含假设。这条约束能在评审时挡掉大部分把两个计数器的值直接相减的错误。
\end{guideline}

\section{回绕、无符号算术与 64 位读取}

\label{sec:ch02-wrap}

用 32 位存微秒，量程是 $2^{32}$ µs，即 4294.97 秒，约 71.6 分钟；用 32 位存毫秒则是 49.7 天。这些数字不是注意事项，而是一定会撞上的边界：现场部署的设备在运行一小时后出现的异常，第一嫌疑就是微秒计数器回绕。

处理办法靠无符号算术的性质而不是分支：只要 \Tp{now} 与 \Tp{then} 的真实间隔不超过半个量程，\Tp{(int32\_t)(now - then)} 的符号就正确表达先后关系。这是 Linux 内核 \Tp{time\_after} 一类的实现形态，也是裸机上应当统一采用的比较范式。

\begin{lstlisting}[style=cstyle,caption={回绕安全的比较范式与超时判断},label={lst:ch02-wrapcode}]
/* 前提：真实间隔的绝对值不超过半个量程。不要在这个前提之外使用。 */
static inline bool ts_after_u32(uint32_t a, uint32_t b)  { return (int32_t)(a - b) > 0; }
static inline bool ts_eq_u32(uint32_t a, uint32_t b)     { return a == b; }
static inline int32_t ts_diff_u32(uint32_t a, uint32_t b) { return (int32_t)(a - b); }

/* 常见写法 now - then >= timeout_us 在无符号下其实也是回绕安全的，
 * 但它把"负数"变成巨大正数：一旦 now 与 then 来自不同时钟域、
 * 先后关系反过来，就会立刻判定为超时。用显式的有符号差更安全。 */
static inline bool deadline_passed_u32(uint32_t now_us, uint32_t start_us,
                                       uint32_t timeout_us)
{
    return ts_after_u32(now_us, start_us + timeout_us);  /* 和式允许自然回绕 */
}

/* 64 位时间戳不要用 32 位差值判断：用有符号 64 位差，量程足够。 */
static inline int64_t ts_diff_i64(uint64_t a, uint64_t b) { return (int64_t)(a - b); }
\end{lstlisting}

\subsection{32 位核上读 64 位计数器}

32 位核读 64 位量必然是两次 32 位访问，中间可能被中断打断，于是读到``高 32 位是旧值、低 32 位是新值''的组合。误差量级是半个量程，不是 1 ns 级别的错，而是几百秒级别的错。标准方案是两次读高位的夹逼校验，清单 \ref{lst:ch02-readq} 给出两种形态：一个 64 位量由两个 32 位寄存器暴露，以及一个 32 位计数器配溢出计数器。Cortex-M 上还要补一条：低半进位与打断同时发生时，单次校验发现不了，所以夹逼必须是``读高、读低、再读高''的三读形态；或者用硬件提供的锁存命令，把两个通道的一致性交给硬件保证。

\begin{lstlisting}[style=cstyle,caption={32 位核上的 64 位计数读取：夹逼校验与关中断两种形态},label={lst:ch02-readq}]
volatile uint32_t g_timer_hi;     /* 由溢出中断递增 */
volatile uint32_t g_timer_lo;     /* 自由运行的低 32 位 */

/* 形态 A：高、低两个 32 位寄存器，读法必须是"高-低-高"夹逼 */
uint64_t timer_read64(void)
{
    uint32_t hi1, lo, hi2;
    do {
        hi1 = g_timer_hi;         /* 三次 volatile 读，任何一次都可能被打断 */
        lo  = g_timer_lo;
        hi2 = g_timer_hi;
    } while (hi1 != hi2);         /* 高半变过 -> 本次结果不可信，重来 */
    return ((uint64_t)hi1 << 32) | lo;
}

/* 形态 B：溢出计数器由中断递增时，关中断窗口比夹逼更简单，也更有界 */
uint64_t timer_read64_crit(void)
{
    uint32_t hi, lo, primask;
    primask = port_disable_irq();           /* 临界区 = 三条读 + 一次恢复 */
    hi = g_timer_hi;
    lo = g_timer_lo;
    port_restore_irq(primask);
    return ((uint64_t)hi << 32) | lo;
}

/* 反例：*(volatile uint64_t *)&reg —— 编译器仍可能拆成两条字读，
 * 先读高还是先读低不确定，顺序不保证，更不要指望它原子。 */
\end{lstlisting}

多计数器同源对齐是比回绕更隐蔽的问题：若中断延迟用 DWT 的 cycle 计数、任务执行时间用 SysTick 的装载计数、日志时间戳用 RTC，那么三者的差值没有意义，因为三个时区不一样。可执行的原则是：同一份分析里只允许一个度量时钟（measurement clock），其余时钟只作参考时间（reference clock）出现，两者之间的换算在启动期标定一次（例如用外设定时器测 SysTick 周期得到真实的 \Tp{SYSCLK\_HCK}），并把换算结果与标定时刻一起存档。频率会随温度漂移，几百 ppm 量级的漂移在长时间窗口上会显出来，所以标定要有有效期。

\section{打点测量}

\label{sec:ch02-probe}

\subsection{DWT 周期计数器}

Cortex-M 的 DWT 提供 \Tp{CYCCNT}：以内核时钟为步进的 32 位周期计数器，读一个内存映射地址即可。启用序列是（1）置 \Tp{DEMCR.TRCENA}（地址 \Tp{0xE000EDFC}），（2）置 \Tp{DWT\_CTRL} 里的计数使能位，（3）部分器件要先向解锁寄存器写 \Tp{0xC5ACCE55} 序列，（4）有些实现复位后 \Tp{CYCCNT} 不为 0，需要显式写 0。四个限制要一并记住：位宽 32（400 MHz 下 10.7 秒回绕一次，长窗口需自己数溢出）；内核时钟停止时计数停止（调试运行位决定睡眠期间是否继续，未设置时深睡一段就``时间蒸发''）；调试电源域掉电后归零；多核各有一个且不同源。此外读 \Tp{CYCCNT} 本身要走总线，代价是若干周期，所以``测一段 20 周期的代码''在这个平台上做不到，能做的是测几十微秒以上的段。

\subsection{GPIO 翻转与逻辑分析仪}

把空闲引脚在段入口拉高、出口拉低，用逻辑分析仪测脉宽，是唯一绕开``测量扰动''的手段：它不消耗被测 CPU 的时间，不受被测代码的缓存状态影响，分辨率由采样率决定（几十到几百 MHz 量级，注意通道延迟与串扰）。代价是需要一个物理引脚与一次接线，并且翻转本身会改变总线活动与功耗。这条路径的定位是边界情况取证：当软件时间戳与直觉不符时，用一次 GPIO 测量判定究竟是测量错了还是系统真的慢。

\subsection{隔离被测段与测量扰动}

被测段要防三件事：被内联进打点函数导致边界模糊（用 \Tp{\_{\Bk}\_{\Bk}attribute\_{\Bk}\_{\Bk}({\Bk}({\Bk}noinline))}，必要时再加 \Tp{used}）、结果被优化成死代码（在消费处放一条空的 \Tp{asm volatile} 加 \Tp{memory} 约束作为 sink）、以及被自动向量化或重排改变工作量。

真正的干扰来源则有四类：(1) 冷缓存，首次运行与稳态运行可能差几倍，报告时必须区分``首次''与``预热后''；(2) 频率爬坡，节能平台在负载突变后需要几毫秒到几十毫秒才升到目标频率，此时的 cycle 数换算成微秒是错的；(3) 共享执行单元与共享末级缓存的邻居，抖动表现为双峰；(4) 打点自身开销，一次带栅栏的读取可以超过被测段本身。结论是：\emph{先测你的量具}。空段测一千次得到的分布就是本次实验的噪声底，任何小于它的差异不值得解释。

\begin{note}[title={低功耗状态下的计数器不可信}]
SysTick 随内核时钟停、DWT 的 \Tp{CYCCNT} 随调试电源域归零、x86 的 TSC 在部分平台深度睡眠后被重新对齐。若被测流程可能进入睡眠（tickless idle、\Tp{WFI}、\Tp{MWAIT} 一类的等待），必须先确认这三个来源在睡眠期间的行为，否则得到的``执行时间''会比实际短，而且短得不固定。
\end{note}

\begin{lstlisting}[style=cppstyle,caption={x86 打点封装：起止定序、核号校验与不变性检查},label={lst:ch02-rdtscp}]
#include <cstdint>
#include <cpuid.h>
#include <x86intrin.h>

struct TscEdge { uint64_t tsc; uint32_t aux; };

/* 起点：先 lfence，阻止后续访存被乱序提前到起点之前执行 */
static inline TscEdge tsc_start()
{
    _mm_lfence();
    uint32_t aux = 0, lo = 0, hi = 0;
    __asm__ __volatile__("rdtscp" : "=a"(lo), "=d"(hi), "=c"(aux) : : "memory");
    return TscEdge{(static_cast<uint64_t>(hi) << 32) | lo, aux};
}

/* 终点：rdtscp 等待此前指令全部退休，随后再补一条 lfence 定序后续 */
static inline TscEdge tsc_end()
{
    uint32_t aux = 0, lo = 0, hi = 0;
    __asm__ __volatile__("rdtscp" : "=a"(lo), "=d"(hi), "=c"(aux) : : "memory");
    _mm_lfence();
    return TscEdge{(static_cast<uint64_t>(hi) << 32) | lo, aux};
}

/* 跨核 skew 的自检：同一样本的起止 aux 不同 -> 这条数据必须整条丢弃 */
static inline bool same_core(const TscEdge& a, const TscEdge& b)
{
    return a.aux == b.aux;
}

/* TSC 频率不保证等于标称核心频率；invariant 位只保证计数率恒定。
 * 标定用固定窗口（例如 timerfd 的一秒）测 cycle 数，多次取中值。 */
static bool tsc_is_invariant()
{
    unsigned a = 0, b = 0, c = 0, d = 0;
    if (!__get_cpuid_count(0x80000007u, 0u, &a, &b, &c, &d)) return false;
    return (d & (1u << 2)) != 0u;   /* EDX bit 2 = invariant TSC */
}
\end{lstlisting}

清单 \ref{lst:ch02-rdtscp} 里有一个容易忽略的语义差别：\Tp{TSC\_AUX} 通常由内核写入当前 CPU 编号，因此在两次读之间可能因为迁移而变化，所以 \Fn{same\_core} 不是优化而是正确性检查。若大量样本跨核，说明线程没有被绑核，这时任何``最坏执行时间''的统计都是混合了两个执行环境的样本。

\section{软件时间基}

\label{sec:ch02-swbase}

裸机的软件时间基通常是``tick 计数加硬件计数器''两级：tick 给粗粒度绝对时间（毫秒），硬件计数器给 tick 内的插值。丢 tick 的检测靠交叉核对：在 tick 服务例程里读一次自由运行的外设定时器（或 SysTick 的当前值），若两次 tick 之间该计数器的增量超过一个周期的 1.5 倍，就说明至少丢了一个 tick，把差值除以周期取整作为补偿量累加。FreeRTOS 的 tickless idle 走的是同一条路的另一侧：\Tp{vPortSuppressTicksAndSleep}（或 \Tp{xPortSuppressTicksAndSleep}）在唤醒时按睡眠长度补记系统节拍，但\emph{SysTick 的重装载值也要相应改写}，否则唤醒后的第一个 tick 来得太早。绝对时间因此有两个来源（软件累加与硬件计数），必须交叉校验，因为睡眠期间外设可能被重新配置。

\subsection{Linux 侧的时钟选择}

\begin{table}[htbp]
\centering
\caption{POSIX 时钟与到期机制的语义差别}
\label{tab:ch02-linuxclocks}
\begin{tabular}{>{\raggedright\arraybackslash}p{0.22\textwidth}
                >{\raggedright\arraybackslash}p{0.26\textwidth}
                >{\raggedright\arraybackslash}p{0.35\textwidth}}
\toprule
接口 & 语义 & 适用与不适用的场景 \\
\midrule
\Tp{CLOCK\_MONOTONIC} & 单调、不含挂起时长、受 NTP 频率校正影响但不被步进 & 一切测量与超时的默认选择；跨重启不可比；不适合写``事件发生的墙上时刻'' \\
\Tp{CLOCK\_{\Bk}MONOTONIC\_{\Bk}RAW} & 单调，且完全不受 NTP 的调整（包括频率校正） & 做频率标定与计数器换算的基准；部分平台上它的实现精度不如 \Tp{MONOTONIC} 稳定 \\
\Tp{CLOCK\_REALTIME} & 墙上时间，会被 NTP 步进、会被手工改 & 只用于日志与对外协议时间戳；用于测量会遇到回拨（见下方 warning） \\
\Tp{CLOCK\_BOOTTIME} & 单调，含挂起时长 & 需要覆盖睡眠期累计时间的场景；唤醒后的补偿计算用它 \\
\Tp{CLOCK\_{\Bk}MONOTONIC\_{\Bk}COARSE} & 只读 tick 粒度，无高分辨率成本 & 毫秒级粗判（时间片、超时）；不要用于微秒级打点 \\
\Tp{timerfd} 系列 & 到期变成一个可读的文件描述符事件，带溢出计数 & 要进事件循环时的首选，是实时应用里最可控的定时入口 \\
\Tp{timer\_settime}（POSIX 定时器） & 三种通知方式：不通知、信号、投递到指定线程 & 需要把信号精确投递到某线程时用最后一种；投递路径比 \Tp{timerfd} 长 \\
\bottomrule
\end{tabular}
\end{table}

到期偏差在任何系统里都不是零：内核定时器的分辨率由时钟事件设备决定（启用高分辨率定时器内核可做到微秒以下，普通内核受 tick 粒度限制），并且存在一个默认的定时器松弛量（Linux 上常为 50 µs 量级，可用 \Fn{prctl} 的 \Tp{PR\_SET\_TIMERSLACK} 修改，对高分辨率定时器不生效）。使用绝对到期（\Tp{TFD\_TIMER\_ABSTIME}、\Tp{TIMER\_ABSTIME}）而不是相对到期，是让误差不累积的唯一办法：相对到期每轮都把上一轮的处理时间计入下一轮，形成与负载相关的漂移；绝对到期把下一轮边界钉在时间轴上，负载高时表现为``错过一次''而不是``越跑越慢''。读 \Tp{timerfd} 返回的 8 字节计数就是错过的次数，把它记进日志是排查调度饥饿最快的线索。

\begin{lstlisting}[style=cppstyle,caption={周期定时器回路：绝对到期、溢出计数与自校正},label={lst:ch02-timerfd}]
#include <cstdint>
#include <ctime>
#include <poll.h>
#include <sys/timerfd.h>
#include <unistd.h>

static inline uint64_t mono_ns()
{
    timespec ts{};
    clock_gettime(CLOCK_MONOTONIC, &ts);         /* vDSO 路径，约 20 ns 量级 */
    return uint64_t(ts.tv_sec) * 1000000000ull + uint64_t(ts.tv_nsec);
}

void run_1ms_loop(int64_t duration_ms)
{
    int fd = timerfd_create(CLOCK_MONOTONIC, TFD_CLOEXEC | TFD_NONBLOCK);
    itimerspec its{};
    its.it_interval = timespec{0, 1000000};      /* 周期 1 ms */
    its.it_value    = timespec{0, 1000000};      /* 首次到期 */
    timerfd_settime(fd, 0, &its, nullptr);

    uint64_t missed = 0, iters = 0;
    uint64_t next = mono_ns() + 1000000;         /* 绝对边界，避免误差累积 */
    pollfd pfd{fd, POLLIN, 0};

    while (int64_t(mono_ns() - next) < duration_ms * 1000000ll) {
        if (::poll(&pfd, 1, 2 /* ms */) == 1) {
            uint64_t exp = 0;
            ssize_t n = read(fd, &exp, sizeof exp);   /* exp 即错过的次数 */
            if (n == static_cast<ssize_t>(sizeof exp) && exp > 1) missed += exp - 1;
        }
        ++iters;
        /* ... 工作项 ... */
        uint64_t now = mono_ns();
        if (int64_t(now - next) > 200000) {           /* 已晚于 200 us */
            next = now;                               /* 放弃追赶，重设基准 */
        }
        next += 1000000;                              /* 下一轮钉在时间轴上 */
        /* 要求更严格的定序时可改用
         *   clock_nanosleep(CLOCK_MONOTONIC, TIMER_ABSTIME, &abstime, nullptr);
         * 它比 timerfd 少一次读，但不能进 poll 集合。 */
    }
    close(fd);
}
\end{lstlisting}

清单 \ref{lst:ch02-timerfd} 里的``放弃追赶''分支值得单独说一句：过载系统有两种策略，一种是补齐所有错过的周期（保持时间轴对齐，但延迟无界增长），一种是重设基准（牺牲周期性以保持延迟有界）。硬实时里通常选后者，但必须把错过次数报出来，否则``看起来稳定''只是把丢失藏起来了。

\section{时间同步与跨设备一致性}

\label{sec:ch02-sync}

IEEE 1588（PTP）与 802.1AS（gPTP，车载与音视频桥接）在 PHY 或 MAC 处对同步报文打硬件时间戳，把从设备的 PHC 与主设备的 PHC 对齐。在同一段局域网、双向路径对称、且硬件打点齐全的条件下，几十纳秒量级是可达的；条件缺一时误差会跳到微秒到毫秒级，所以``能对齐到什么程度''要写成实测结论而不是引用指标。Linux 侧的典型组合是 \Tp{ptp4l} 管理 PHC、\Tp{phc2sys} 在 PHC 与系统时钟之间搬运偏移，应用则用 \Fn{clock\_gettime} 打开对应的动态时钟 id 读取。802.1AS 额外用邻居速率比让从设备把本地时钟的\emph{频率}调到与主设备一致，这一点很关键：只对齐一次时刻而频率留有 100 ppm 偏差，一秒后误差就回到 100 µs。

收发两侧的硬件时间戳通过套接字选项取回。设置侧声明发送打点与接收过滤条件；取回侧用 \Tp{SO\_TIMESTAMPING} 选项，再用 \Fn{recvmsg} 读到 \Tp{SCM\_TIMESTAMPING} 控制消息，其结构里的三个时间戳分别给出（1）原始硬件时间戳（若请求了原始模式则是驱动原生计数域，否则已换算到系统时间域）、（2）墙上时间域的时间戳、（3）软件时间戳。三个值同时存在非常有用：硬件与软件之差就是驱动栈的延迟，可以直接用来定位哪一段抖动变大了。

跨设备排序的原则是：\emph{只有同一时间基的两个时间戳才可以说谁先发生}。三个推论：(1) 设备 A 的 \Tp{CLOCK\_MONOTONIC} 与设备 B 的同名时钟不可比，因为两者起点（上电时刻）与频率都不同；要可比，必须双方都用 PHC，或者在报文里同时携带各自的 PHC 值以便离线换算。(2) 同步不可用时（主从切换、链路断开、守护进程未就绪），日志要显式标记时间戳来源与置信度，而不是静默给出一个精度被夸大的值。(3) 精度、分辨率、准确度是三个不同的量：一个 64 位纳秒字段宣称的纳秒精度，如果来自软件时间戳加往返估计，实际准确度可能是几十微秒。

\begin{bestpractice}[title={跨设备事件排序的三个前置条件}]
(1) 两个事件的时间戳来自同一类时钟（同为 PHC，或同为某设备的本地单调时钟且只做同域排序）；(2) 记录时一并记录该时钟的同步状态（是否已锁定、最后一次偏移与方差）；(3) 排序判定阈值不小于该同步状态给出不确定度。三条都满足时，``事件 1 早于事件 2''才有可辩护的含义；否则请输出``在不确定度范围内无法判定先后''。
\end{bestpractice}

\section{常见陷阱与排查顺序}

\label{sec:ch02-traps}

\begin{warning}[title={裸 \Fn{\_rdtsc} 与墙上时间做单调测量}]
两个高频事故：(1) 线程被迁移到另一个核而两个核的 TSC 有偏差，差值出现几千 ns 甚至负数；(2) 用 \Tp{CLOCK\_REALTIME} 或 \Fn{gettimeofday} 做单调测量，一次时间同步步进（回拨）让差值突然变成负几十秒。对策分别是：每条样本记录核号并丢弃跨核样本（清单 \ref{lst:ch02-rdtscp}），以及把测量一律换成 \Tp{CLOCK\_MONOTONIC} 或 \Tp{CLOCK\_MONOTONIC\_RAW}。
\end{warning}

\begin{itemize}
  \item \textbf{SysTick 优先级最低导致 tick 丢失。} 表现是任务数量与堆使用都正常，但周期任务越来越慢。检测：在 tick 服务里比较外设定时器的增量与期望周期；找出 \Tp{PRIMASK} 被置起的最长窗口（第 4 章的阻塞时间分析）。补偿：按差值补记 tick，并把补偿次数作为诊断计数器暴露出来。
  \item \textbf{频率标定的量化误差。} 用一个一秒窗口测 TSC 或 \Tp{SYSCLK\_HCK}，端点误差会直接换算成标定误差：若窗口由 1 kHz 的 tick 给出，端点不确定度是 1 ms，也就是千分之一的量化误差。要得到 ppm 级标定，窗口要拉长到数秒以上并多次取中值，或者用更长周期的参考（RTC 的秒脉冲）。
  \item \textbf{把时间戳精度当成测量分辨率。} 64 位纳秒字段不等于纳秒级测量：打点自身的读取开销、总线延迟、以及被测量对象的变化速率共同决定有效分辨率。判据是量具的噪声底（\S\ref{sec:ch02-probe}），不是字段宽度。
  \item \textbf{多核各读各的计数器。} 每个核有自己的 DWT \Tp{CYCCNT}（不同源），而架构定时器的计数在实现上通常同源。跨核的两个 cycle 值相减没有定义。多核分析里只允许两类时间：核内 cycle，以及全局的架构定时器或 PHC。
  \item \textbf{回绕与半量程前提被违反。}\Tp{(int32\_t)(now - then)} 的前提是间隔不超过半个量程。把 32 位微秒值当作``到上电时刻的差''用在长跑模块里，回绕点附近会出现一次 71.6 分钟量级的假超时。对策：跨分钟级的时间量用 64 位，或用带溢出中断的计数器拼接。
  \item \textbf{换算与打点进了关键路径。} 若换算含 64 位除法，其最坏时间可能超过被测段。关键路径的正确做法是``只记 cycle，出界后再换算''，把换算放到日志线程或离线后处理里，使被测路径的时间不依赖换算实现。
  \item \textbf{tickless idle 之后绝对时间错位。} 睡眠长度与补记的 tick 数不一致（例如 SysTick 重装值未同步修改），唤醒后第一个 tick 提前或推后。核对方式：把软件累加的 tick 增量与外设定时器测到的实际睡眠时长对表，误差应在一两个 tick 之内。
\end{itemize}

排查顺序建议固定为五步：(1) 先确认时钟源的身份与频率是从哪里\emph{读到}的，而不是从哪里``知道''的；(2) 检查位宽与回绕假设，把所有 32 位时间量按 \S\ref{sec:ch02-wrap} 的范式复核一遍；(3) 测一次量具（空段、固定开销），确定噪声底；(4) 检查被测流程是否可能进入低功耗状态使计数器停摆；(5) 前面都通过之后，才开始怀疑调度器与中断延迟。时间类问题里一半以上落在前三步，而前三步都不需要修改业务代码。

% ============================================================
% ---- frag_ch03.tex ----
\chapter{中断控制器与中断延迟}
\label{ch:ch03-interrupt-latency}

\DomBare{} \DomRtos{} \DomRT{} 中断是硬实时系统里唯一能把控制权从``你安排的顺序''抢走的机制，所以它也是延迟预算里最难压缩的一段。本章只按时序与有界性来看中断：控制器怎么仲裁、一次中断从引脚有效到 ISR 首指令要经过哪几段、每段的量级与可控手段，以及怎么把这段延迟真正测出来。寄存器位域的字节级写法（掩码、置位清位序列）在硬件访问那一册里已经讲透，本章不重复，只在说明延迟来源时点一下某条寄存器的语义。

先给结论：中断控制器自身的仲裁开销几乎从来不是瓶颈。它固定、可数、有界。真正让最坏延迟比平均延迟高两个数量级的，是软件自己制造的不可中断窗口——关中断的临界区、长跑的 ISR、被更高优先级占住的 CPU，以及底半部与调度器之间的唤醒链路。本章后半部分的所有手段都围绕这个窗口。

\section{三种控制器形态}

同一段``从外部信号有效到 ISR 首指令''的延迟，在三类硬件上由完全不同的机制构成。写可移植的实时代码之前，要先确定自己面对的是哪一种。

\begin{table}[htbp]
\centering
\caption{三类中断入口在延迟视角下的对照}\label{tab:ch03-controllers}
\begin{tabular}{>{\raggedright\arraybackslash}p{0.13\textwidth} >{\raggedright\arraybackslash}p{0.22\textwidth} >{\raggedright\arraybackslash}p{0.22\textwidth} >{\raggedright\arraybackslash}p{0.21\textwidth}}
\toprule
维度 & NVIC（Cortex-M） & GIC（Cortex-A） & x86（APIC 加 IDT） \\
\midrule
号空间 & 由向量表索引，实现相关 & 分段固定：软件、私有、共享三类 & 256 个向量，低段留给异常与保留 \\
\addlinespace
优先级 & 数值小者优先，拆成抢占层与子层 & 每中断一个优先级字段，再与运行优先级屏蔽比较 & 由向量号的高位段决定类别，同族内按号比较 \\
\addlinespace
进入方式 & 硬件压栈后直接跳到表项地址 & 软件读 acknowledge 取号再分派，随后跳到入口 & 过 IDT 门、压栈、走统一入口再分派 \\
\addlinespace
收尾 & 异常返回自动完成，无需显式 EOC & 显式写结束寄存器（安全组还需去激活） & 显式写 EOI 寄存器，或依赖入口约定 \\
\addlinespace
共享与路由 & 多个源可被厂商映射到同一号 & 共享中断在某一刻只落到一个核 & 引脚线可多设备共享，消息中断一号一源 \\
\bottomrule
\end{tabular}
\end{table}

\subsection{NVIC：抢占优先级、子优先级与挂起位}

Cortex-M 把中断控制器放进核内（\Tp{NVIC}），与内核同时钟域，因此仲裁延迟是若干个流水线周期，而不是总线往返。它有两个容易被误用的设计点。

第一，优先级是一个``数值越小越优先''的定宽字段，并被拆成两段：抢占优先级决定能否打断正在服务的 ISR，子优先级只决定同一抢占层内谁先被响应。两者如何切分由 \Tp{AIRCR} 里的分组字段（\Tp{PRIGROUP}）决定，而可用位数由芯片实现固定，CMSIS 侧通过 \Tp{\_\_NVIC\_PRIO\_BITS} 暴露（常见 3 位或 4 位）。分组一旦定下，全片所有中断共享同一个切分方式。子优先级也相同的两个中断，按 IRQ 号小的先响应——这属于``碰巧的顺序''，不要拿它当设计依据。故障类（\Tp{HardFault}、\Tp{NMI}、\Tp{BusFault} 等）的优先级高于一切可编程中断，且 \Tp{NMI} 不可屏蔽，所以它不能被用作普通业务通道。

第二，NVIC 里每个中断有独立的使能位、设置挂起位与清挂起位（\Fn{NVIC\_{\Bk}Enable{\Bk}IRQ}、\Fn{NVIC\_{\Bk}Set{\Bk}Pending{\Bk}IRQ}、\Fn{NVIC\_{\Bk}Clear{\Bk}Pending{\Bk}IRQ}）。对边沿触发的外部中断，往挂起位写 1 可以软件造出一次中断，这是自测延迟的正规手段：脉冲宽度可控、不依赖外设、也不引入总线往返。但对电平语义的外设中断源，只清 NVIC 的挂起位没有用——外设自己的标志还在举着，返回后立刻重进。清源必须先清外设标志，这个顺序是常见事故源。

软件中断 \Tp{SVCall}、\Tp{PendSV} 与 \Tp{SysTick} 是 NVIC 的三个特殊位置。\DomRtos{} RTOS 把上下文切换挂在 \Tp{PendSV} 上并给它最低的可编程优先级，目的是``任何 ISR 都能把切换推迟到自己退出之后''。这个选择本身就是延迟设计：它把切换从``打断中断服务''变成``中断服务全部结束后统一做''，代价是多一次退出与再入。

\begin{lstlisting}[style=cstyle,caption={C：NVIC 优先级与分组配置（4 位优先级、2 加 2 切分的例子）},label={lst:ch03-nvic}]
/* 4 位优先级位：抢占 2 位 + 子优先级 2 位。
   数值越小越优先；最高的一档留给故障类，应用不要用。 */
#define PRIORITY_GROUP_2_2   (5u << 8)     /* AIRCR.PRIGROUP 的编码 */

void nvic_setup(void)
{
    /* CMSIS 内部按实际优先级位数做左移，调用方只传 0..15 */
    NVIC_SetPriorityGrouping(PRIORITY_GROUP_2_2);

    /* 采入 DMA：允许被控制环打断，因为它跑得长且不决定时序 */
    NVIC_SetPriority(DMA1_Stream2_IRQn, 6);

    /* 250 us 控制环：整条链上优先级最高的可编程中断 */
    NVIC_SetPriority(TIM2_IRQn, 4);

    /* 编码器门控与控制在同一抢占层：互不打断，只排队 */
    NVIC_SetPriority(EXTI9_5_IRQn, 4);

    /* 调试串口：尽力而为，可以被打断任意多次 */
    NVIC_SetPriority(USART1_IRQn, 8);

    /* 共享线 EXTI9_5 覆盖 5..9 号引脚。宁可同层排队，
       也不要为了少一个入口把两个语义不同的源塞进同一个 ISR。 */
    NVIC_EnableIRQ(EXTI9_5_IRQn);

    /* 自检：软件挂起一次控制环中断，量一次端到端响应 */
    /* NVIC_SetPendingIRQ(TIM2_IRQn); */
}
\end{lstlisting}

\subsection{GIC：Distributor、Redistributor 与 SPI 路由}

Cortex-A 多核的中断由相对内核外部的 \Tp{GIC} 提供，逻辑上分三层。\Tp{Distributor}（\Tp{GICD\_} 前缀的寄存器）持有全局中断号空间、使能、优先级与目标核路由；\Tp{Redistributor}（\Tp{GICR\_}，GICv3 起才有）是每核一份的银行，保存该核私有中断的状态并代表核发出 acknowledge；核内系统寄存器接口（\Tp{ICC\_} 前缀，如 \Tp{ICC\_PMR\_EL1}、\Tp{ICC\_IAR\_EL1}、\Tp{ICC\_EOIR\_EL1}）才是真正的仲裁比较点。

中断号空间按归属切三类。\Tp{SGI}（0 到 15）软件生成中断，用途是核间打点与调度器唤醒。\Tp{PPI}（16 到 31）每核私有外设，典型是私有定时器与看门狗。\Tp{SPI}（32 起）共享外设中断，必须路由到某一个核（GICv2 写 \Tp{GICD\_ITARGETSR}，GICv3 写 \Tp{GICR\_IROUTER} 的亲和域字段），或者允许多个目标位再由控制器选一个。\DomHost{} ``一个 SPI 在某一刻只落到一个核''这个事实决定了它的延迟性质：路由到哪个核、那个核当时是否屏蔽了中断、是否被更高优先级占着，三者共同决定响应时间。把同类高速中断摊到不同核上，能显著降低单核排队；反过来，把一个高频中断和一个低频但长服务的中断路由到同一个核，两者的抖动就会互相绑定——这是实时装配表里最常见的错配。

取号与收尾也和 NVIC 不同：核要读 acknowledge 寄存器才知道是哪个中断，处理完写结束寄存器；带安全分组的实现里，结束与去激活是两个动作，只做前者会让中断停留在活动状态，表现和忘了发 EOI 一样。写中断路径的启动代码时，``取号、服务、结束、去激活''这四步要一次写全，缺后两步的症状通常要很久以后才出现。

GICv3 还带来两个与延迟直接相关的东西：LPI（基于表格的中断，配 \Tp{ITS} 做翻译，是消息中断在服务器与高性能嵌入式上的载体）与优先级屏蔽的分层。若系统里有安全与非安全分组，分组过滤发生在优先级比较之前——实时中断被分到错误的组里，表现为``完全不来''。这类问题的现场在固件或设备树里，不在驱动代码里。

\subsection{x86：IDT、IO-APIC、MSI/MSI-X 与向量分配}

x86 的``中断控制器''其实是一组东西：核内 local APIC（用 TPR 与 PRR 做优先级比较，用 IRR 与 ISR 位图记录状态，用 EOI 寄存器收尾）、片内或片外的 IO-APIC（一张重定向表项对应一个输入引脚，把引脚中断映射成向量），以及中断描述符表 \Tp{IDT}（256 项，向量号即索引）。传统 8259 的向量基址习惯是 0x20，异常占 0 到 31，若干向量号被保留为 spurious，剩下的才是设备可用空间。向量的分配由内核负责，同一向量不能同时代表两个语义不同的源；调优时能改的是``哪个设备用哪个向量、落到哪个核''，而不是号段本身。

延迟上的关键差异有两点。其一，中断进入要过 IDT 门、压栈、走内核统一入口再分派，固定开销比 NVIC 大，而且是软件层的（入口代码、中断栈切换），量级在亚微秒到几微秒。其二，共享引脚中断是抖动的主要来源：一条 IO-APIC 线可以挂多个设备，处理链要把该线上注册的所有 handler 逐个跑一遍，每个都问一句``这是你吗''，于是任何一次中断的最坏耗时与链上设备数成正比。

\Tp{MSI} 与 \Tp{MSI-X} 把中断变成一次内存写周期的``消息''：不再有共享线，每个 \Tp{MSI-X} 表项有独立的向量、目标与屏蔽位，从根本上消灭了``替别人付探测时间''。这让``一个队列一个中断一个核''的摊平布局成为可能。实时调优里最先该确认的一条就是：设备用的是 MSI-X 而不是传统共享引脚，并且每向量屏蔽可用；表项落在设备的地址空间里，但字节级写法不归本册。

\section{延迟的五段构成}

把``设备把线拉有效''到``ISR 第一条业务指令执行''之间的时间切成五段，是分析中断延迟唯一可靠的办法。每段的量级、是否可数、由谁决定，都不一样。

\begin{longtable}{>{\raggedright\arraybackslash}p{0.15\textwidth} >{\raggedright\arraybackslash}p{0.23\textwidth} >{\raggedright\arraybackslash}p{0.42\textwidth}}
\caption{中断延迟的五段构成、典型量级与可控手段}\label{tab:ch03-latency-breakdown}\\
\toprule
段 & 典型量级 & 由什么决定，怎么压 \\
\midrule
\endfirsthead
\toprule
段 & 典型量级 & 由什么决定，怎么压 \\
\midrule
\endhead
\bottomrule
\endfoot
\endlastfoot
响应 & 裸机：一条指令到数十个周期；含长指令的核可到上百周期 & 当前正在执行的指令何时可被截断。乘法、块拷贝、批量位写、进入深度等待态都会推迟采样点。手段：避开长指令、用 DMA 代替 CPU 搬运、别让等待态挡住中断采样。 \\
\addlinespace
识别与仲裁 & NVIC：2 到 6 个核周期；GIC：若干次总线往返；x86：含入口代码约几微秒 & 控制器的优先级比较、该中断在目标核上是否使能、分组过滤。这一段基本固定，压无可压，只能核对配置是否正确。 \\
\addlinespace
上下文保存 & Cortex-M：硬件压 8 个寄存器，浮点惰性压栈另计几十到上百周期；Linux：栈切换加软件保存全部通用寄存器 & 是否使能 FPU（惰性压栈可把成本推到第一条浮点指令上）、栈对齐、是否在中断栈与任务栈之间切换。手段：短 ISR 里不做浮点。 \\
\addlinespace
服务（ISR 本体） & 亚微秒到几十微秒；失控时到毫秒 & 完全由代码决定。这是唯一一段可以逐行审查的延迟，也是最容易悄悄膨胀的一段。手段：非关键工作挪到底半部，ISR 里不循环等待、不分配、不打印。 \\
\addlinespace
退出与恢复 & 与压栈对称；背靠背中断时 tail-chaining 可省一次完整出栈与入栈 & 出栈、用 \Tp{EXC\_RETURN} 值做模式选择返回、x86 的 \Tp{IRET} 与 EOI 写回。EOI 发在 ISR 末尾时，同一中断要等整段跑完才允许再次响应。 \\
\end{longtable}

\DomBare{} 上面的量级是``设计期用来做预算的量级''，不是芯片手册承诺值：手册给的是控制器时序，端到端最坏值必须自己测，见第~\ref{sec:ch03-measure}节。

\subsection{为什么说不可中断段才是主因}

拿一组典型数字对比就清楚了。假设仲裁加上下文保存在 100 兆赫兹核上是 0.3 微秒，一次有界 ISR 是 2 微秒；而某个驱动路径里有一段持有自旋锁并禁止本地中断 200 微秒的操作。中断落在屏蔽段中间时，响应延迟等于该段的剩余时间，期望值约为它的一半。也就是说，一个大多数人不去量的软件段落，比前面四段加起来大两个数量级。

这段延迟有三个讨厌的性质，决定了它必须被单独治理。第一，它与负载相关：屏蔽段只有在并发出现时才变长，空载测量看不到。第二，它与代码位置无关而与锁有关：你无法通过读 ISR 找出它，必须读所有会禁中断的路径。第三，它在最坏值上不可压缩到零：只要系统里存在必须原子且必须较长的操作（大块硬件初始化、固件升级写入），就一定有下界。所以工程上的正确表述是``把可屏蔽段列出来，逐条给出长度并纳入预算''，而不是``把延迟做到最小''。

\begin{compare}[title={同样五段，主因不同}]
裸机上最坏值通常落在第 4 段（ISR 本体）或它前面的屏蔽区；通用系统上主因变成排队——中断线程与目标任务的唤醒等待。于是同一句``最坏 200 微秒''在两边的含义不同：前者多半是服务时间，后者多半是调度等待，两者的修法完全不同。
\end{compare}

\section{把延迟压下去的具体手段}

\textbf{向量表重定位与在 RAM 里执行 ISR。} Cortex-M 的 \Tp{SCB->VTOR} 可以把向量表搬到 SRAM 或 ITCM。搬动有两个收益：取向量不再受 flash 等待周期与预取未命中的影响，抖动消失；ISR 代码常驻紧耦合内存，取指不再带等待周期。代价是启动时必须做一次拷贝，而且深睡眠唤醒之后不要忘了重建。

\textbf{拆长临界区。} 若一段代码必须原子，先问它是``必须屏蔽中断''还是``必须不被打断地看到一致数据''。后者可以用双缓冲、序列号（读侧用序号校验一致性的思路）、或者让 ISR 只写单字而读者校验两次读值来做到，不需要关中断。把一次长屏蔽拆成两段短屏蔽，收益通常大于任何控制器调优。

\textbf{最小化 \Fn{\_\_disable\_irq} 的包络。} 规则：屏蔽区内只做``改状态''，不做``等状态''。任何等待循环都必须放在屏蔽区外，或者改成带上限的重试并在每轮里允许中断。同理，不要在屏蔽区里做除法、平方根、软件浮点或大结构体拷贝——它们都可能让当前指令跨过几十个周期。

\textbf{tail-chaining 与 EOI 时机。} Cortex-M 在两个中断背靠背时会跳过中间那次完整的出栈再入栈，把第二次响应的开销降到几个周期；这是硬件给的免费收益，作者要做的是不要把该做的事塞到``返回后再排一次''的路径上。EOI 早发与晚发是真实的取舍：晚发让同一中断在本次服务完成前不会重入，行为干净但排队延迟变大；早发配合优先级设计可以让更高优先级中断立刻插入，但代价是 ISR 必须可重入或自带保护。共享线上还有第三种坑：源标志未清就早发 EOI，立刻再进。

\textbf{别名向量与免压栈入口。} 部分实现允许同一 ISR 挂到多个向量号上，或在特定向量上跳过部分压栈，从而给极短的高频中断（例如每 20 微秒一次的采样触发）留一条更便宜的路。做法属于寄存器层，本册不重复；判据是：当某中断的服务时间与压栈时间同量级时，压栈成本才值得专门优化。

\textbf{每个 IRQ 一个独立向量。} 在 x86 上意味着用 \Tp{MSI-X}，把每个队列绑到不同向量与不同核；在 MCU 上意味着不把五个语义不同的源挤到同一条共享线上，宁可多占一个向量号。共享线省下的只是向量表空间，代价是最坏服务时间随链上设备数线性增长——在实时系统里这是坏交易。

\begin{warning}[title={在 ISR 里调用 \Tp{printf} 或日志函数}]
一次格式化的耗时与栈深度都不可控，而且多数实现带着互斥或阻塞语义：ISR 里等锁会把整条中断链拖住，最坏情况是同一中断反复被推迟，整机看起来``卡在某条线上''。需要留痕时只写无锁环形缓冲里的固定长度记录（几个 32 位字），事后由底半部或低优先级任务去格式化。
\end{warning}

\begin{note}[title={清挂起位的顺序：先清源，再谈返回}]
电平语义的外设中断源在标志未清前会持续举手。只清控制器侧的挂起位（或只发 EOI），结果是本次服务刚退出就再进一次，表现为``中断频率异常高但没有任何业务进展''。顺序固定为：读外设状态、清外设标志、必要时再清控制器挂起位。
\end{note}

\begin{guideline}[title={优先级先于代码长度}]
在写第一行 ISR 之前，先把抢占层的分布画出来：哪一层只有一个源，哪一层有多个源以及它们的总服务时间。同层内等待时间等于前面所有 ISR 之和，这一条定下来之后，剩下的优化才是``把本层变短''。反过来做（先优化代码再定优先级）通常会推翻已有结构。
\end{guideline}

\section{中断线程化与底半部}
\label{sec:ch03-threading}

\DomHost{} 通用内核把中断拆成上半部与下半部，动机只有一个：让不可中断的时间尽量短。理解这套拆法对实时的后果，需要精确到``谁在什么上下文里跑、被谁打断、按什么优先级排队''。

\subsection{handler 与 thread 的分工}

\Fn{request\_threaded\_irq} 接受两个入口。primary handler 仍跑在硬中断上下文，只允许做``读走关键寄存器、抓快照、清标志''，然后返回 \Tp{IRQ\_WAKE\_THREAD}；threaded handler 跑在每中断一个的内核线程里，可以被调度、可以用互斥、可以做真正的处理。若 primary 传 \Tp{NULL}，内核用自带的默认 primary（直接唤醒线程），此时必须带 \Tp{IRQF\_ONESHOT}：屏蔽该中断直到线程跑完，保证电平型中断不会因为线还举着而空转。

实时视角下有三个后果必须记住。

其一，工作从``不可打断''变成``可调度''，于是它的延迟不再由中断决定，而由线程优先级与就绪队列决定。启用实时抢占的内核把硬中断默认全部线程化，中断线程是 \Tp{SCHED\_FIFO} 且有默认优先级（历史上是 50），可以在 \Tp{/proc/irq} 下每个中断号的目录里调优先级与亲和。如果应用任务的优先级低于中断线程，控制环就会被压在中断路径之后，表现为``中断来了但数据迟迟不上来''；反过来把应用抬得过高，则会饿死中断路径。两个方向都要显式安排，并把结论写进装配表。

其二，\Tp{IRQF\_ONESHOT} 会把该中断线屏蔽到线程处理完毕，屏蔽时长等于线程的处理时间。对高速或共享设备，这就把``几微秒的 ISR''换成``可能几百微秒的线屏蔽''，抖动反而变大。此时要么把线程处理时间压住，要么拆成多个独立向量，要么干脆回到自有的高优先级线程并在应用侧排序。

其三，\Tp{IRQF\_NO\_THREAD} 强行把 handler 留在硬中断上下文。合法用法很窄：计时精度要求极高、或必须拿 raw 自旋锁（\Tp{raw\_spinlock\_t}，在实时抢占下仍是真自旋、真屏蔽）的路径。除此之外，把带锁的代码标成不线程化，就是在给自己制造不可屏蔽段。

\begin{lstlisting}[style=cstyle,caption={C：request\_threaded\_irq 的最小可用形态（独立向量、显式 primary）},label={lst:ch03-threaded-irq}]
/* 内核侧驱动片段：硬中断只抓快照，真正的处理在 threaded handler */
#define REG_DATA        0x02u
#define STAT_VALID      (1u << 11)
#define RING_SIZE       64u
#define RING_MASK       (RING_SIZE - 1u)

static irqreturn_t drdy_hard(int irq, void *dev_id)
{
    struct state *st = dev_id;
    unsigned int  val;

    /* 硬中断上下文：一次寄存器读加一次快照，不等待、不睡眠 */
    if (regmap_read(st->regmap, REG_DATA, &val))
        return IRQ_NONE;          /* 读不通：可能不是我的设备 */

    if ((val & STAT_VALID) == 0)
        return IRQ_NONE;          /* 共享线上别人的中断：原样退回 */

    st->sample = (u16)val;        /* 单字快照，靠对齐保证原子 */
    return IRQ_WAKE_THREAD;       /* 剩下的交给线程 */
}

static irqreturn_t drdy_thread(int irq, void *dev_id)
{
    struct state *st = dev_id;
    u16 s = st->sample;

    mutex_lock(&st->buf_lock);    /* 线程里可以拿会睡眠的锁 */
    st->ring[st->head++ & RING_MASK] = s;
    mutex_unlock(&st->buf_lock);

    return IRQ_HANDLED;
}

static int dev_probe(struct device *dev, struct state *st)
{
    int irq = dev_get_irq(dev, 0);

    /* 第三个参数是中断名，会出现在中断线程名与中断计数表里。
       显式给出 primary，线在硬中断退出时重新开放；若把 primary
       传 NULL，则必须加 IRQF_ONESHOT，否则注册直接失败。 */
    return request_threaded_irq(irq, drdy_hard, drdy_thread,
                                IRQF_TRIGGER_FALLING, dev_name(dev), st);
}
\end{lstlisting}

\subsection{底半部机制的取舍}

\begin{table}[htbp]
\centering
\caption{四种延后处理机制的时序性质对比}\label{tab:ch03-bh-mechanisms}
\begin{tabular}{>{\raggedright\arraybackslash}p{0.14\textwidth} >{\raggedright\arraybackslash}p{0.19\textwidth} >{\raggedright\arraybackslash}p{0.16\textwidth} >{\raggedright\arraybackslash}p{0.30\textwidth}}
\toprule
机制 & 运行上下文 & 典型被推迟时间 & 实时上的性质 \\
\midrule
softirq & 中断返回时的软中断上下文，可被硬中断打断 & 微秒级，通常在同一次退出路径里就做完 & 延迟最低，但同一类型可在多个 CPU 并发执行，需自己保证无冲突；不可睡眠 \\
\addlinespace
tasklet & 软中断上下文，依附于统一的 tasklet 类型 & 微秒到几十微秒 & 同一实例不会并发，写起来简单；仍不可睡眠，且没有优先级概念，只能等前面排队的同类实例做完 \\
\addlinespace
workqueue & 进程上下文（工作线程） & 毫秒级，完全看调度 & 可睡眠、可长做，但延迟无法承诺；实时路径只用于``一定要推到进程侧''的批量工作，可用延迟工作项做粗粒度定时 \\
\addlinespace
中断线程 & 进程上下文（每中断一个线程） & 几十微秒到毫秒，取决于线程优先级 & 唯一能被纳入优先级体系的底半部；实时抢占场景下的默认选择，代价是屏蔽时长受线程处理时间影响 \\
\bottomrule
\end{tabular}
\end{table}

选型的判据可以压缩成三问。第一问：这个工作能不能睡眠？不能睡就只能留在软中断侧（前两种），能睡才有后两种资格。第二问：它的完成时间要不要进承诺？要进就选中断线程，因为只有它有明确优先级；工作队列那种``大概会在几毫秒内''在控制环里不算答案。第三问：同一实例并发会不会坏？软中断类需要自己保证，tasklet 与中断线程替你保证``同一实例串行''，代价是排队时延随负载增长——一个高频数据流用单个 tasklet 串起来，往往就是抖动的第一现场。

\DomRtos{} 在 RTOS 或裸机上，对应物通常是``ISR 置标志加发信号量，任务做处理''。它比工作队列更可预测，因为那个任务有明确的静态优先级；但要点与线程化相同：一旦推到任务侧，承诺就从``中断延迟''变成``端到端响应延迟''，唤醒到执行那一段必须一起计入预算。

\section{中断与调度的交叉}
\label{sec:ch03-scheduling-cross}

中断延迟很少单独出现，它总是和调度耦合：ISR 唤醒了一个更高优先级的任务，于是``响应时间''里多了一次上下文切换。

\subsection{唤醒、切换与排队}

\textbf{唤醒路径的四步。} ISR 写就绪队列、把某任务的等待状态改为就绪、设置需要切换的标记、退出。在 Cortex-M 的 RTOS 里，第四步通常是让 \Tp{PendSV} 挂起，切换发生在全部 ISR 退出之后。好处是切换期间不必屏蔽中断，坏处是多出一次``退出再入''的开销，量级几十周期。判据：如果控制环周期 250 微秒而抖动预算 10 微秒，这几十周期可以忽略；如果周期 50 微秒且路径上有三个中断串起来，它就不行。

\textbf{切换成本的构成。} 一次切换的真实成本很少等于``保存与恢复寄存器''。主体是栈交换之后的缓存与分支预测失效，以及对下一个任务首段代码的指令缓存重新预热。裸机量级在一微秒上下；通用系统里，一个工作集小且热的线程切换在几微秒，带地址空间切换与缓存冷启动可到十几微秒。测量时要区分``切换本身''与``切换后第一条有用指令''，实时预算通常该算后者。把关键任务与它的中断放在同一个核上，往往比分到不同核更省——省的就是这一次预热。

\subsection{风暴、限流与共享线抖动}

\textbf{中断风暴与限流。} 一个持续举线或疯狂自触发的源可以让 CPU 百分之百花在``退出再重进''上，看起来像整机死掉。防线有三层：控制器级（优先级与屏蔽位）、内核级（对同一中断的重复计数与限流，Linux 会统计``没人认领的中断''并在阈值后禁用整条线——这是保命行为，也是隐蔽故障源）、应用级（发现频率异常就降级为抽样处理）。诊断要点：某个中断计数每秒增长接近一个固定高值、而设备实际事件率低得多，先怀疑共享线或没清源的 ISR。

\textbf{共享线上的探测抖动。} 共享意味着每次中断都要付``问一遍所有人''的成本，而这个成本由别人的代码质量决定。两条缓解：把探测顺序按实测频率排序（常触发的源放前面，把最慢的探测放最后）；以及在探测本身可能很慢的驱动上，先用一个便宜的硬件状态位做门控。若发现``链上第五个设备把第一个的抖动放大''，正确的结论是拆开向量，而不是优化那五个 handler。

\textbf{中断路径上的优先级反转。} 常见形式不是``低优先级持锁挡住高优先级任务''，而是``中优先级 ISR 长时间占 CPU，把高优先级中断的响应挤到后面''。多个源落在同一抢占层时，队头的等待时间等于前面所有 ISR 之和。缓解手段是把长 ISR 拆成``短硬中断加线程或任务处理''，或者给关键源更高的抢占优先级。另一种才是真正的反转：中断路径与一条被更高优先级任务压住的路径共用同一把锁，于是持锁时长由无关任务的调度决定——此时解法是把锁换成无锁结构，或者把持锁区完全移出中断路径。

\begin{bestpractice}[title={让关键中断不被排队}]
先定优先级分布，再写 ISR 长度：同一抢占层内的中断数量与它们的总服务时间，共同决定这一层里最慢者的等待。把控制环放最高、把低频批量源放最低、把耗时处理挪出硬中断上下文，这三件事做完，仲裁延迟就已经进入微秒级；继续抠控制器通常只是徒劳。
\end{bestpractice}

\section{测量方法与判据}
\label{sec:ch03-measure}

\subsection{裸机：三段式 ISR 与最坏时间记账}

\DomBare{} 一个可承诺的 ISR 应该写成三段。第一段做``必须此刻做''的事：读外设状态、清源、抓一个足够宽的快照，长度固定且与外部条件无关。第二段做纯计算与单字发布，没有循环等待、没有分配、没有浮点以外的意外操作，最坏路径可以逐条数出来。第三段做尾注：更新计数、写打点脚、必要时挂起一个更低层的活动。三段的价值不在于好看，而在于它把``必须不可打断的部分''压到第一段，第二段与第三段随时可以被更高优先级截走。

最可靠的原始测量手段也是打点：在能控制的触发侧（或用测试逻辑、软件挂起）产生一个脉冲，在 ISR 第一条指令处翻转一个空闲引脚，在 ISR 退出处再翻转一次。脉冲到第一次翻转是延迟，两次翻转之间是服务时间，用逻辑分析仪或示波器直接读这两段。判据要分两端读——看最坏值而不是平均值，并且要跑足量样本让长尾有机会出现。同时把系统里最长的临界区也打点（入口翻转、出口翻转），否则你会把长屏蔽误读成控制器慢。

\begin{lstlisting}[style=cstyle,caption={C：三段式最小 ISR——清源、有界处理、打点尾注（统计在最低优先级任务里做）},label={lst:ch03-isr}]
#define TRACE_PIN        (1u << 3)          /* PA3 接示波器通道 */
#define TRACE_HIGH()     (GPIOA->BSRR = TRACE_PIN)
#define TRACE_LOW()      (GPIOA->BSRR = TRACE_PIN << 16)

static volatile uint32_t isr_count;
static volatile uint32_t isr_cycles_last;
static volatile uint32_t isr_cycles_max;

void TIM2_IRQHandler(void)
{
    uint32_t t0 = DWT->CYCCNT;       /* 周期计数器需在调试单元里先使能 */

    TRACE_HIGH();                    /* 沿一：ISR 首指令 */
    TIM2->SR = 0u;                   /* 第一段：清源，放在任何可能变长的事前 */

    control_step();                  /* 第二段：有界——无等待、无浮点、无分配 */

    /* 第三段：尾注与记账 */
    uint32_t dt = DWT->CYCCNT - t0;
    isr_cycles_last = dt;
    if (dt > isr_cycles_max)
        isr_cycles_max = dt;         /* 单写者，不必屏蔽 */
    isr_count++;
    TRACE_LOW();                     /* 沿二：服务结束 */
}

/* 最低优先级任务负责搬运与判定。两个读数之间允许轻微不一致，
   因为这里要的是分布和趋势，不是快照一致性。 */
void stats_task(void)
{
    for (;;) {
        uint32_t n  = isr_count;
        uint32_t mx = isr_cycles_max;

        budget_publish(n, mx, core_cycles_per_us());
        if (mx > CTRL_BUDGET_CYCLES)
            alarm_raise(ALARM_ISR_BUDGET);
        os_wait_ms(200);             /* 唯一允许写持久日志的上下文 */
    }
}
\end{lstlisting}

跑上一天，你拿到两张表：每个中断的最坏服务时间，以及单位时间的中断次数，后者超预算就是风暴线索。注意统计代码本身必须短，不要用长临界区去保护统计量——那等于为了测量而制造一个新的不可中断段。

\subsection{通用系统：工具各自测哪一段}

\DomHost{} 现成工具很多，但必须知道各自测的是哪一段，否则数字会被拿去回答它没回答过的问题。

\begin{itemize}
\item \Tp{cyclictest}：测``定时器到期到用户态线程真正跑起来''的调度延迟，包含中断响应、唤醒、上下文切换与内核不可抢占段。它不测具体外设中断的响应，因此得到的数字是对整条链路的下界式估计。常用选项：\Opt{priority}（不高优先级就只是在测普通调度）、\Opt{mlockall}（挡住缺页掺进来）、\Opt{histogram}（看分布而不是只看最大值）、\Opt{interval} 与 \Opt{loops}（周期与样本数，决定长尾能否被看到）。
\item \Opt{break-tracepoints} 一类阈值触发的选项可以在超过门限时停住 tracing，用来抓``那一次''的现场，比只读统计数字有用得多。
\item \Tp{oslat}：专测纯轮询下 CPU 的暂停延迟，用来抓系统管理中断、深度睡眠退出、微码更新这类与外设中断无关的来源。它跑不出设备路径的抖动，但能解释为什么同样的代码某天最大延迟突然变大。常用 \Opt{duration}、\Opt{cpu-list}、\Opt{threshold}、\Opt{bucket-size}。
\item \Tp{ftrace} 的 \Tp{irqsoff} 与 \Tp{preemptoff} 追踪器：直接回答``最长的一次关中断、最长的一次不可抢占分别在哪里、多长''，这是定位长临界区的首选。判据：如果 \Tp{irqsoff} 报出的最大屏蔽段远大于你已知的 ISR 时间，问题在锁或驱动路径里，而不在控制器。工作图类追踪器（\Tp{wakeup}、\Tp{wakeup\_rt}）回答另一半问题：从 ISR 唤醒到任务真正跑起来花了多久。
\item 中断计数与亲和：每个中断号在各核上的计数序列可以看出有没有全压在一个核上、以及计数增长速度是否与业务事件相符。这类指标适合长期记录，因为它揭示的是配置漂移而不是单次抖动。
\end{itemize}

通用作法：先关掉省电与频率调节（否则会把唤醒时间算进中断延迟）、把测量线程钉核、固定亲和与内存锁定，并同时记录系统负载——只在空载下测出来的最坏值没有承诺力。测量结论要写进延迟预算表，与表~\ref{tab:ch03-latency-breakdown}的分段一一对上，才知道下一轮该动哪一段。

\section{常见陷阱}

\begin{itemize}
\item \textbf{忘了发 EOI，或者没清外设标志。} 前者表现为``这个中断只进一次''，后者表现为``ISR 疯狂重入''。两者都容易被误诊成硬件损坏。返回前确认清的是源设备的标志位，而不只是控制器的挂起位。
\item \textbf{在 ISR 里用 \Fn{malloc}、\Tp{printf}、日志或任何会睡眠的接口。} 见前面的警告框；顺带一条：也不要调用会拿普通互斥锁的内核函数。
\item \textbf{把浮点放进 ISR 却没使能 FP 单元。} 未使能时第一条浮点指令直接产生故障；使能了但惰性压栈没配好，则每次中断多背一组寄存器。短 ISR 里不用浮点，把计算放到底半部。
\item \textbf{临界区里提前返回，把开中断漏在后面。} \Fn{\_\_enable\_irq} 必须覆盖所有出口；更好的做法是不手写配对，用作用域对象或统一出口。这类 bug 的表现是``跑一阵才开始丢中断''，最难复现。
\item \textbf{共享线上不属于自己却返回 \Tp{IRQ\_HANDLED}，或该处理时返回 \Tp{IRQ\_NONE}。} 前者让真正的受害者被当成已处理而丢事件；后者被内核累计计数后会导致整条线被禁用，故障要很久以后才以``设备彻底不响应''的形式暴露。
\item \textbf{把周期定时器的优先级设得比杂事低。} 于是控制环节的采样被推迟，表现为控制性能莫名其妙地随负载变化。采样与执行路径上的定时器必须是可编程中断里最高的一档。
\item \textbf{把高频中断与长服务中断路由到同一个核。} 单核排队让两者的抖动互相绑定；换核的成本远低于重写两个驱动。
\item \textbf{用``中断计数在增加''当作系统健康的证据。} 计数增加只说明有人在举手，不说明有人服务得过来。健康证据应该是端到端时间戳差值的分布，见第~\ref{sec:ch03-measure}节。
\item \textbf{把中断绑到正在调试的核上又下断点。} 断点停住的是整个核的中断路径，而总线超时与监督逻辑不会等你。调试期需要配套的旁路策略，这个话题属于下一章。
\end{itemize}

本章的落点是一句可执行的检查顺序：先把不可中断窗口测出来，再定优先级层，再决定哪些工作该挪到线程侧，最后才谈控制器本身的配置与向量分配。顺序颠倒的项目，通常会把几周花在优化一个只占几个周期的环节上。

% ============================================================
% ---- frag_ch04.tex ----
\chapter{看门狗、复位源与失效模式}
\label{ch:ch04-watchdog-failsafe}

\DomBare{} \DomRtos{} \DomRT{} 看门狗是唯一一个专门用来回答``这台设备还活着吗''的机制，也是实时系统里最容易做错的一环：做错的形式通常不是``没装''，而是``装了但装在了一个证明不了任何事的位置''。本章按三条线展开——狗本身的时序语义、复位发生时怎么判定与取证、以及失效之后设备该停在哪一档状态上。与前一章一样，寄存器位域的写法不在这里重复；这里只关心时间、有界性与失效行为。

一句先行的判断标准：喂狗动作只有在``它是某个真实成就的后果''时才有意义。如果一条代码路径无论业务是否推进都会执行到喂狗，那条路径就是狗的盲区，你的监督等于零。

\section{看门狗的类别与各自的时序语义}

先给一张横向对照表，后面四小节分别展开每类的时序细节。

\begin{table}[htbp]
\centering
\caption{五类看门狗监督的时钟独立性、可停性与能抓到什么}\label{tab:ch04-types}
\begin{tabular}{>{\raggedright\arraybackslash}p{0.16\textwidth} >{\raggedright\arraybackslash}p{0.19\textwidth} >{\raggedright\arraybackslash}p{0.15\textwidth} >{\raggedright\arraybackslash}p{0.30\textwidth}}
\toprule
类型 & 时钟与独立性 & 运行中能否停 & 真正能抓到的失效 \\
\midrule
独立看门狗 & 片内低频 RC，与主时钟树无关 & 不能，只能等复位 & 主循环停摆、任务再也没走到喂狗点；抓不到``卡在中断里照常喂狗'' \\
\addlinespace
窗口看门狗 & 总线侧时钟，主时钟停则它也停 & 不能 & 卡死与``跑飞进高频循环''两类；分辨率粗，但边界明确 \\
\addlinespace
任务看门狗（多通道） & 由内核或驱动设备提供时基，通道各自订阅 & 一般可由配置停用 & 与业务成就绑定，能定位``哪一路没完成''；依赖被监督者真的上报 \\
\addlinespace
外部监督芯片与 PMIC & 完全独立的芯片与振荡器 & 通常只能靠断电或复位 & 连芯片自身失效、时钟失效、电源异常都能抓；喂入脚默认态要格外小心 \\
\addlinespace
内核软锁与硬锁检测 & 系统时基加不可屏蔽事件，与用户态无关 & 可由参数关闭 & 软件层的``没人调度''与``连中断都进不来''；不覆盖硬件与电源 \\
\bottomrule
\end{tabular}
\end{table}

\subsection{独立看门狗：只准跑、不准停}

独立狗挂在自己的时钟上（通常是片内低频 RC，量级几十千赫兹），与主时钟树无关，因此主晶振不起振、PLL 失锁、软件把系统时钟切错，都不能让它停下来。这既是它的价值，也是它的局限：它的计时精度很差，批次与温度漂移的量级常在几十个百分点，所以超时值必须按``最坏可能更短''来核算。

典型语义是三条：计数器单调递减到零即产生复位或中断；刷新（重载）由一次写关键码完成，重载后计数器回到重装载值；一旦启动就无法停止，只有复位能让它回到未启动状态（部分家族提供只在调试或低功耗下冻结计数的选项，那是给开发用的，不构成运行时监督）。超时时间等于分频与重装载的乘积除以时钟，形式上写作 $t = 4 \times 2^{p} \times R / f$, 其中 $p$ 是分频指数、$R$ 是重装载值、$f$ 是那个不精确的低频时钟。工程读法：把 $f$ 的上下限都代进去，取最短的那个结果当承诺。

两个必须知道的盲区。第一，某些低功耗状态会关掉那个低频时钟（例如待机态），于是``狗在深度睡眠期间不再监督''——如果你的失效模式包含``卡在待机里醒不来''，独立狗抓不到，要靠唤醒定时器或外部芯片。第二，较新的家族把窗口寄存器也加进了独立狗，于是``独立''指的是时钟域而不是行为；配了窗口之后，别按老习惯认为它只能``晚了才复位''。

\subsection{窗口看门狗：早喂与晚喂都复位}

窗口狗除递减到下限复位之外，还定义了一个不允许重载的区间：如果在计数器仍高于窗口值的时候就写刷新码，等于告诉硬件``我在以超出预期的速度绕圈''，直接复位。它因此能抓两类相反的失效：卡死（晚喂）与跑飞到一个高频循环里（早喂）。

时序上有两个量。到复位的剩余时间等于``当前计数值相对下限的余量''乘以``每计一数所需的时间''，后者通常带一个很大的固定分频（几千倍量级）与一个可选预分频，所以窗口狗的分辨率粗、但边界非常明确。第二个量是预警边界：多数实现允许计数器递减到某个特定值时先产生一次中断（早醒或预警），再往下才复位。这一口气是留给遗言的：写备份域、把最后一帧数据落到持久区、或者主动切断输出。

窗口余量的算法只有一句：允许喂狗的区间必须完整包住喂狗时刻的可能散布，并且两端要按时钟漂移与调度抖动拉开。设狗超时为 $T$, 喂动作名义发生在 $0.6\,T$, 那么窗口边界必须显著晚于上一帧最早可能完成的时间，同时 $0.6\,T$ 减去最大抖动之后仍落在窗口内。数字是次要的，顺序是主要的：先确认最坏抖动，再定窗口，最后才定超时值。

\begin{lstlisting}[style=cstyle,caption={C：窗口看门狗的喂狗与预警中断。这些 \Tp{wwdg\_} 前缀函数是本项目的薄封装，寄存器名以家族手册为准},label={lst:ch04-wwdg}]
/* 超时预算约 20 ms。预警中断设在边界值触发，复位在其后。 */
#define WWDG_PRESCALER_DIV   8u
#define WWDG_WINDOW_VALUE    0x50u  /* 计数高于此值时刷新即算喂得太早 */
#define WWDG_REFRESH_VALUE   0x60u  /* 落在窗口内的刷新码 */
#define FRAME_DEADLINE_US    180u   /* 本帧最晚应完成的时刻（打点值） */
#define EARLY_GUARD_US       120u   /* 比最早完成时刻再晚一点 */

void wwdg_init(void)
{
    /* 顺序不能反：先设预分频与窗口，最后才启动计数器。
       窗口值一旦写入立即生效，先启动再设窗口会当场造成早喂复位。 */
    wwdg_set_prescaler(WWDG_PRESCALER_DIV);
    wwdg_set_window(WWDG_WINDOW_VALUE);
    wwdg_enable_early_irq();          /* 早醒/预警中断，处理见下面 */
    wwdg_start(WWDG_REFRESH_VALUE);
}

void wwdg_feed_on_frame(void)
{
    /* 唯一喂狗点：只有本帧的控制输出真的写进硬件之后才允许调用。
       还要检查时间戳，来得太早就宁可让它复位。 */
    uint32_t now = timer_us();
    if (now < FRAME_DEADLINE_US - EARLY_GUARD_US)
        return;                       /* 到得太早：故意不喂 */
    wwdg_write_counter(WWDG_REFRESH_VALUE);
}

/* 预警中断：复位前最后一段可执行时间。预算要按这段最多能做多少事
   来规划，别指望在里面写 flash——块擦除本身是毫秒量级。 */
void WWDG_IRQHandler(void)
{
    gpio_latch_safe_outputs();        /* 第一优先：切断危险 */
    crashlog_arm_reset();             /* 第二：只写备份域，几个字 */
    wwdg_clear_early_flag();          /* 清标志只消中断，救不了复位 */
}
\end{lstlisting}

\subsection{任务看门狗与用户态监督}
\label{sec:ch04-taskwd}

\DomRtos{} 任务狗把``狗''从单一递减计数器变成一张订阅表：每个被监督的任务或通道在自己的完成点上报，监督逻辑检查本周期的订阅者是否全部报齐，未齐则超时。它的表达力远高于单一狗，因为监督粒度绑定在业务成就上（采入完成、解算完成、输出写入完成），不同通道可以有不同超时；它还顺带回答``哪一路没完成''。

Linux 侧的用户态接口是 man 手册第 4 节的设备接口（\Tp{watchdog(4)}）：字符设备 \Tp{/dev/watchdog0}，可能多个实例，配一组 ioctl。\Tp{WDIOC\_{\Bk}GETSUPPORT} 取驱动自述能力，\Tp{WDIOC\_{\Bk}GETTIMEOUT} 与 \Tp{WDIOC\_{\Bk}SETTIMEOUT} 读写超时，\Tp{WDIOC\_{\Bk}GETSTATUS} 与 \Tp{WDIOC\_{\Bk}GETBOOTSTATUS} 取状态位，\Tp{WDIOC\_{\Bk}KEEPALIVE} 负责刷新，带预警能力的驱动还有 \Tp{WDIOC\_{\Bk}SETPRETIMEOUT} 一类接口。能力位与状态位共用一组 \Tp{WDIOF\_} 常量：\Tp{WDIOF\_{\Bk}SETTIMEOUT} 表示可设超时，\Tp{WDIOF\_{\Bk}PRETIMEOUT} 表示支持预警，\Tp{WDIOF\_{\Bk}CARDRESET} 出现在启动状态里就说明上次是狗咬的，\Tp{WDIOF\_{\Bk}MAGICCLOSE} 表示允许按约定写法再关闭设备来停狗。常见封装（例如 \Fn{wd\_{\Bk}keep\_{\Bk}alive}）本质上是把 \Tp{WDIOC\_{\Bk}KEEPALIVE} 打一遍；有的驱动不实现这个 ioctl，只接受对设备写一个字节作为刷新，所以``先把能力位读回来、再决定怎么喂''是必须的写法。

\begin{lstlisting}[style=cppstyle,caption={C++：\Tp{/dev/watchdog} 用户态监督，按帧刷新并保留启动原因},label={lst:ch04-linux-wdt}]
#include <fcntl.h>
#include <unistd.h>
#include <sys/ioctl.h>
#include <linux/watchdog.h>

/* 约定：只有主循环（真正的帧完成点）可以调用 keepAlive()。
   任何工作线程、信号处理函数里都不许出现对它的调用。 */
class FrameSupervisor {
    int  fd_ = -1;
    int  timeout_ms_ = 500;
    int  feed_period_ms_ = 0;
    bool can_set_timeout_ = false;
 public:
    bool attach(const char *devpath)
    {
        fd_ = ::open(devpath, O_WRONLY);
        if (fd_ < 0)
            return false;

        int boot = 0;
        if (::ioctl(fd_, WDIOC_GETBOOTSTATUS, &boot) == 0 &&
            (boot & WDIOF_CARDRESET) != 0) {
            /* 上次是看门狗复位：走降级启动，别直接进自动模式 */
            enter_degraded_startup();
        }

        struct watchdog_info caps;
        if (::ioctl(fd_, WDIOC_GETSUPPORT, &caps) == 0)
            can_set_timeout_ = (caps.options & WDIOF_SETTIMEOUT) != 0;

        if (can_set_timeout_) {
            if (::ioctl(fd_, WDIOC_SETTIMEOUT, &timeout_ms_) != 0)
                ::ioctl(fd_, WDIOC_GETTIMEOUT, &timeout_ms_);
        } else {
            ::ioctl(fd_, WDIOC_GETTIMEOUT, &timeout_ms_);   /* 记录实际值 */
        }

        feed_period_ms_ = timeout_ms_ / 3;   /* 刷新周期明显短于超时 */
        return true;
    }

    void keepAlive() { ::ioctl(fd_, WDIOC_KEEPALIVE, 0); }

    ~FrameSupervisor()
    {
        if (fd_ < 0)
            return;
        /* 只有能力位含 MAGICCLOSE 时这条路才有效。若内核被配成
           不允许中途停狗，关闭也不会停，进程退出即复位。 */
        const char stop = 'V';
        ::write(fd_, &stop, 1);
        ::close(fd_);
    }
};
\end{lstlisting}

\subsection{外部监督芯片与 PMIC：把监督权交出去}

外部监督芯片（常见形态是一只独立 IC，带喂入脚、窗口配置脚、超时设置脚，输出复位与一路可配置的低有效告警）解决的是 MCU 自己监督自己的根本局限：如果失效的是电源、时钟或芯片本身，片内狗可能同时失效。它的超时由外部元件设定，典型值在几十毫秒到几秒之间；有的型号支持窗口加手动复位两档，也有的把第一次超时只作告警、第二次才复位，用来实现``一次自愈、失败才降级''。

用外部芯片时最容易漏的三件事。第一，喂入脚在 MCU 复位期间的默认状态：若它被复用成普通 GPIO 且默认电平恰好符合芯片的``已喂''判据，监督窗口会被无声拉长。第二，芯片超时与被监督软件的帧周期之间要留余量，并且要跨温度实测喂入脉冲的宽度与间隔，而不是用名义值——窗口型芯片对脉宽有要求，太窄不算喂。第三，PMIC 内部的看门狗或使能超时常常一旦触发不可逆（要么复位要么断电），它的时序必须和软件复位策略一起设计，否则会出现``PMIC 先把几路电源切了，CPU 还没来得及把状态写下来''。

\begin{warning}[title={在 ISR 里喂狗}]
这是最常见的失效设计：主循环死掉之后，周期定时器中断照样按点运行、照样把狗喂好，于是系统以完全正常的时序输出错误数据。喂狗只允许出现在业务成就之后——最坏也要放在任务完成点或调度器的收尾钩子里。若必须用中断做检查，那就让它做检查（读心跳表、判断序号是否推进），而不是让它做刷新。
\end{warning}

\section{喂狗点、心跳与启动交接}
\label{sec:ch04-feed-policy}

\textbf{喂狗点的唯一性。} 一个系统应该只有一个真正的刷新点，其余监督逻辑都以``是否达成成就''的形式向它汇报。多个刷新点等于多个盲区，因为它们之间是逻辑或关系，任一多余点生效就会掩盖全部故障。审查方向是反的：从狗的重载写入处往回追调用者，看得到几条路径。

\textbf{按帧完成的合取条件。} 实用的谓词是三条同时成立：本帧采入的时间戳推进了、控制解算没有触发合理性告警、输出写硬件成功了。三条都在才允许喂。缺一条就让它咬，对应宁可复位也不输出可疑值的 fail-safe 姿态；如果业务要求 fail-operational（见第~\ref{sec:ch04-power}节），那就不能靠狗处理，必须交给冗余通道。

\textbf{多核交叉监督。} 每个核每完成一帧就把自己的帧序号写进共享区，另一核在周期性检查里判断``序号是否比上次看到的大''。发现对方停滞时，正确动作不是替它喂，而是停止自己这边的喂狗或主动触发一次整体复位。要点是复位域：如果两个核共享外设、DMA 通道或电源轨，单侧复位会留下不一致的硬件状态，此时``任一侧发现异常就全体复位''通常是更安全的设计。跨核心跳通道本身也要有超时——对端正在重启时心跳会中断一段时间，别把这件事误判成死机后进入错误的降级分支。

\textbf{周期与余量的算法。} 令最坏帧耗时为 $w$, 喂狗时刻的散布为 $j$, 狗超时的最坏值为 $t$（按漂移取较短的那个），则要求 $w + j$ 显著小于 $t$。狗带窗口时，还要求最早可能喂的时刻晚于窗口边界。写下来只是一行，但每一项都必须来自测量而不是估算：$w$ 来自最长执行时间记账，$j$ 来自调度抖动统计，$t$ 来自时钟漂移规格。

\textbf{启动阶段与 bootloader 的责任交接。} 上电到应用真正跑起来之间有一段``有狗但不干活''的时间：镜像校验、解压、搬到 RAM、外设与传感器稳定。三种可选姿态要提前定死：（1）引导阶段不启狗，由应用初始化完成后再启；（2）引导阶段启狗并使用较长超时，应用接入前不应到期；（3）引导阶段启狗且不允许停，那么应用的第一件事必须是接管刷新，并且在时钟切换之前先喂一次。第三条最严格，也最容易踩两个坑：跳转后代码仍在 flash 上执行但等待周期已变，或者时钟切换期间主频暂时很低使同一段代码的执行时间成倍变长，狗就咬在``看起来一切正常''的位置。若引导程序与应用的狗策略不一致，还要在交接处留一条可验证的记录——第~\ref{sec:ch04-forensics}节的复位原因与帧序号就是为此准备的。

\begin{guideline}[title={改变时间标度的事件要单独过一遍狗}]
凡是改变``一个循环要花多久''的事件——换频、切换 flash 等待周期、进入或退出低功耗态、开启或关闭缓存——都要重新核算最坏帧耗时与狗超时。多数偶发的启动复位，根因就在这类边界上，而不是在业务代码里。
\end{guideline}

\section{复位源判定与崩溃取证}
\label{sec:ch04-forensics}

复位之后第一个问题必须是``上一次是怎么结束的''。答案通常在复位状态寄存器里，用一组粘滞标志区分：上电或掉电复位、外部引脚复位、软件复位、独立狗复位、窗口狗复位、低功耗进入触发的复位、欠压复位。各家族的位名不完全一致（\Tp{PINRSTF}、\Tp{BORRSTF}、\Tp{LPWRRSTF} 这一类），但语义结构通用。

关键动作是把读和清分开并固定顺序：先把标志快照存下来，再清标志（常由一个专门的粘滞位清除控制位完成），然后才做分支决策。不清就复位，会让下次读到陈旧标志而误判；先清后读，会丢掉本次证据。

\begin{lstlisting}[style=cstyle,caption={C：复位源读取、清粘滞标志与启动分支。\Tp{rcc\_} 前缀为本项目封装},label={lst:ch04-reset-flags}]
struct reset_report {
    uint8_t  source;      /* 见下面枚举，既进日志也进备份域 */
    uint8_t  was_watchdog;
    uint16_t boot_count;  /* 复位计数，用于识别重启循环 */
};

void classify_reset(struct reset_report *rr)
{
    uint32_t csr = rcc_read_reset_flags();   /* 只读快照，先不动硬件 */

    rr->source = RESET_SRC_UNKNOWN;
    if (csr & RESET_FLAG_POR)       rr->source = RESET_SRC_POWER;
    else if (csr & RESET_FLAG_PIN)  rr->source = RESET_SRC_EXTERNAL;
    else if (csr & RESET_FLAG_WWDG) rr->source = RESET_SRC_WWDG;
    else if (csr & RESET_FLAG_IWDG) rr->source = RESET_SRC_IWDG;
    else if (csr & RESET_FLAG_SOFT) rr->source = RESET_SRC_SOFT;
    else if (csr & RESET_FLAG_LPWR) rr->source = RESET_SRC_LOWPOWER_EXIT;

    rr->was_watchdog = (rr->source == RESET_SRC_WWDG ||
                        rr->source == RESET_SRC_IWDG);

    /* 必须显式清粘滞位，否则下次读到的是这一轮的组合 */
    rcc_clear_reset_flags();

    /* 计数要带时间窗：只在窗口内累加，别把正常停机也算进重启循环 */
    rr->boot_count = bump_boot_count_if_within_window(BOOT_WINDOW_S);
}

void apply_boot_policy(const struct reset_report *rr)
{
    if (rr->boot_count > 4u && rr->was_watchdog) {
        /* 反复被咬：不要再走同一条初始化路径 */
        load_fallback_image();
        gpio_latch_safe_outputs();
        return;
    }
    if (rr->was_watchdog)
        start_in_reduced_mode();     /* 一次自愈：降级但不中断服务 */
    else
        start_full_mode();
}
\end{lstlisting}

\textbf{把证据放在不会被复位抹掉的地方。} 三类位置按易失性与耐久度排：备份域寄存器（几十到上百个 32 位字，主电丢失后靠电池或备用轨维持）、低功耗保留 RAM（容量大一些，几千字节量级，复位保持但断电丢）、以及普通 SRAM 里的非初始化段（暖复位通常保留，冷上电不保证）。放什么：复位原因、最近帧序号、各环节的执行计数、故障状态寄存器快照、若干返回地址（供事后回溯）、以及一个覆盖记录的校验值。校验值的作用是让下次上能判断这块证据是完整的还是半写的——否则你会拿着一份被复位截断的记录得出错误结论。

\begin{lstlisting}[style=cstyle,caption={C：备份 RAM 里的一槽崩溃证据（固定大小、循环覆盖、不做分配）},label={lst:ch04-bkram}]
/* 放在复位后仍保留的备份域：只写固定大小的槽，绝不动态分配。
   EV_SECTION 由构建系统给出（链接脚本里的保留段），不在此定义。 */
#include <stddef.h>

#define EV_MAGIC        0xA5A5C3C3u
#define EV_SLOTS        8u
#define EV_CRC_FIELDS   offsetof(struct crashlog, crc)  /* 校验覆盖此前字段 */

struct crashlog {
    uint32_t magic;
    uint32_t crc;              /* 覆盖 magic 之前的全部字段 */
    uint32_t seq;              /* 槽序号：最后写，代表“这条已完整” */
    uint32_t reset_flags;      /* 复位标志快照 */
    uint32_t frame_last;       /* 最后完成的帧号 */
    uint32_t faults;           /* 故障状态寄存器组合 */
    uint32_t lr[EV_SLOTS];     /* 返回地址样本，事后做回溯 */
    uint32_t cycles_max;       /* 本帧最长环节耗时（以周期计） */
};

static volatile struct crashlog g_log __attribute__((section(EV_SECTION)));

void crashlog_arm_reset(void)
{
    volatile struct crashlog *p = &g_log;

    p->reset_flags = rcc_read_reset_flags();
    p->frame_last  = app_frame_seq();
    p->faults      = fault_status_snapshot();
    p->cycles_max  = app_worst_step_cycles();
    backtrace_sample(p->lr, EV_SLOTS);      /* 扫描栈取若干返回地址 */

    /* 写序固定：字段 -> 魔数 -> 校验 -> 序号。序号最后落地，
       半路被复位截断的记录会因为序号未推进而被读侧丢弃。 */
    p->magic = EV_MAGIC;
    p->crc   = crc32_over(p, EV_CRC_FIELDS);
    p->seq  += 1u;
}

/* 下次启动时取证：先验魔数与校验，再决定是否采信 */
bool crashlog_take(struct crashlog *out)
{
    *out = (struct crashlog)g_log;

    if (out->magic != EV_MAGIC || out->seq == 0u)
        return false;
    return out->crc == crc32_over(out, EV_CRC_FIELDS);
}
\end{lstlisting}

\textbf{崩溃前转储与写放大。} 预警中断到复位之间只有几毫秒，够写备份域，通常不够做 flash 页擦除（擦除本身是毫秒量级，块更大）。所以设计要分层：紧急路径只写备份域或保留 RAM，事后由下次启动的低优先级任务把它搬到持久介质；或者用铁电存储一类无需擦除的介质承担紧急日志。flash 侧还有两条硬约束：反复复位会写坏寿命，因此``每次复位都追加一条记录''必须配循环覆盖与节流（相同原因不重复写、写之前比对上一条内容）；日志不能吃掉预警窗口，否则你把最后一段可执行时间用在记录死亡上，而安全动作没做。

\begin{note}[title={复位计数本身是一种磨损}]
如果每次复位都往 flash 或 EEPROM 记一笔第 N 次启动，重启循环会把介质写穿。把计数放在备份域、加时间窗（只在窗口内累加）、并在持久化时做同因合并，才能既留下重启趋势又保住寿命。
\end{note}

\section{掉电、电源异常与三档失效姿态}
\label{sec:ch04-power}

\textbf{掉电检测与复位门限的相对位置。} 常见的两级门槛是可编程电压检测（阈值可选，产生中断）与欠压复位（产生复位）。设计要求是检测阈值高于复位门限，并且两者间隔足够留出``从检测到电压真正跌穿''的时间——电压是连续下滑的，可用时间取决于负载电容与泄放电流，量级从几十微秒到几毫秒。这段窗口里能做的事要按最坏下滑率来定：把输出钉到安全值、把关键状态写进备份域（几个字，不要大块搬运）、不要启动任何需要擦除的写操作。如果这段窗口比你需要的短，正解是加储能电容或超级电容，而不是把日志写得更省。

\textbf{电压跌落时狗的行为不能想当然。} 独立狗用自己的 RC，主轨跌落时它可能还在计数而 CPU 已经取不到正确指令——这是``狗咬对了但咬在错误的时刻''。窗口狗用总线时钟，低压下时钟失锁时狗本身可能停摆。还有第三种：复位门限与狗超时在低压下都会漂，两者互相竞争，结果取决于器件个体。结论是：电源异常不能交给看门狗处理，它要有自己的检测与动作路径，而且这条路径的完成时间要计入故障容忍时间。

\textbf{上电时序与就绪等待。} 传感器、ADC 参考、隔离电源、通信收发器往往各有稳定时间，上电即开始采样会把一段无效数据送进控制环。做法是显式的就绪状态机：每个部件的可用由一个可验证条件给出（配置寄存器读回一致、内部状态位、一次自检通过、连续若干帧合理性合格），未全部可用的帧不进入控制解算，也不喂任务狗。不要用固定延时替代条件——固定延时在跨温度与跨批次时是最先失效的假设。

\textbf{三档姿态与``谁定义安全状态''。} 三档是三种不同的承诺。fail-operational：出错后仍继续提供功能，通常靠冗余通道，验收标准是单点故障不影响输出。fail-safe：退到不伤害人与设备的状态，例如切断动力、抱闸、输出置零或置到中间值。fail-silent：停止对外输出但仍然自证安静，例如只维持心跳而不参与控制。安全状态不是软件层能自定的：它来自危险分析，写成``在什么时间内必须到达什么物理状态''，也就是故障容忍时间必须大于检测时间加上到达路径的执行时间。到达路径本身也要监督——``安全继电器已吸合''的反馈信号如果没人检查，那只是假设而不是动作。判据式写法：对每一路输出，都能说出默认电平（掉电时物理上是什么）与主动置安全值的指令及其完成证据。

\begin{bestpractice}[title={预警中断只做两件事：钉输出、留小证据}]
把预警到复位这段窗口的用途严格限制为切断危险与留下最小可解释记录。任何格式化、大块写、通信重试都会吃掉这段窗口。日志的详细部分留给下一次启动的后台搬运。
\end{bestpractice}

\section{软失效：跑飞的程序照样喂狗}

\textbf{位翻转与它的检测手段。} 存储单元与寄存器状态可能被电离辐射或毛刺翻掉。可用的检测按粒度排：带 ECC 的存储（能纠一检二，容量与代价最友好）、关键寄存器的影子副本加周期比对、配置参数的校验（CRC 或奇偶）并在每次加载时校验、以及总线侧奇偶。要点是纠正与报告要分开：悄悄纠掉的错误累计起来就是即将失效的预兆，必须计数并上报，否则降级永远没有依据。

\textbf{控制流完整性与签名循环。} 看门狗对逻辑跑飞的失效边界很硬：如果程序跳到某段代码后仍然按节奏执行到喂狗点，狗什么也抓不到。补法是把执行序列本身当数据校验：若干关键区按预定顺序打点，形成短的签名序列，监督逻辑检查签名是否合法（顺序对、无缺失、无重复）。这类做法在功能安全语境里叫控制流监控，代价是每个关键段增加少量代码与几微秒，收益是把``跑飞但时序正常''变成可检测。更硬的是硬件方案：锁步核逐拍比对两条流水线，或用 MPU 把未使用地址区间设为触发故障。

\textbf{三模冗余与投票。} 对无法用签名覆盖的数据路径，用三路独立计算加中值投票；代价是资源与时间变成三倍，且三路必须有真正的独立性——同一份输入、同一张系数表、同一个时钟源不算独立。投票不一致本身要产生事件：两票对一票是纠错，三票各不相同是必须降级而不是随便挑一个。

\textbf{狗看不见的部分要靠合理性。} 传感器断线常表现为读数贴到轨、跳为零或长期不变；总线被拖住表现为事务超时；算法失效表现为输出超出物理量程或变化率。因此每个通道要有三件最低限度的合理性检查：量程、变化率、以及与上一个可信值的偏差，外加长时间完全不动的僵死检测。检查失败就降级（沿用上值、切备份通道、或退出自动模式）并计数。看门狗是活性证据，不是正确性证据；这两件事在验收上要分开写。

\textbf{故障注入怎么验。} 没有注入测试的看门狗设计等于没做过测试。最小一套四类：（1）拔传感器与断线，含开路到轨、短路到地、短路到电源三种；（2）拖住总线，例如时钟拉伸、数据线被拉低后要求走恢复流程、或者从设备完全不响应；（3）注入延迟，人为在控制环里插忙等，验证``晚到''被当作超时处理，而不是被喂狗掩盖；（4）软件故障注入，屏蔽某个任务的完成上报、打断签名序列、把非法返回地址写进栈。每次注入记录三个数：失效发生到被检测的时间、检测到到达安全状态的时间、留下的证据是否能解释现场。前两者之和就是监督开销，通常是验收里被忽略的那一项。

\begin{longtable}{>{\raggedright\arraybackslash}p{0.21\textwidth} >{\raggedright\arraybackslash}p{0.26\textwidth} >{\raggedright\arraybackslash}p{0.31\textwidth}}
\caption{失效模式到检测手段与安全动作的映射}\label{tab:ch04-failure-map}\\
\toprule
失效模式 & 检测手段 & 安全动作与备注 \\
\midrule
\endfirsthead
\toprule
失效模式 & 检测手段 & 安全动作与备注 \\
\midrule
\endhead
\bottomrule
\endfoot
\endlastfoot
主循环不再推进 & 独立狗或窗口狗超时；任务狗的``本帧订阅者未齐'' & 复位到后备固件或降级启动。前提是喂狗点确实在主循环成就之后 \\
\addlinespace
卡在等待外设返回 & 等待循环自带迭代上限；狗超时应覆盖这种情形 & 放弃本次操作并计数，连续数次则退出自动模式。不要用长关中断实现``原子等待'' \\
\addlinespace
跑飞但仍按节奏喂狗 & 程序流签名、锁步核、MPU 命中非法区间 & 拒绝进入自动模式并告警。狗对这一类基本无效，必须靠流监督 \\
\addlinespace
传感器断线或读数贴轨 & 量程、变化率、长时间不变三类合理性检查 & 切备份通道或保持上值并降级；``数据可疑''要作为独立状态上报 \\
\addlinespace
通信总线被拖住 & 每笔事务的超时与重试上限、总线恢复流程 & 丢帧计数并进入 fail-silent，不再参与控制但维持心跳 \\
\addlinespace
电压跌落 & 可编程电压检测中断，阈值高于复位门限 & 立即钉住危险输出、写备份域小记录；禁止启动擦除型写入 \\
\addlinespace
主时钟失效或失锁 & 时钟安全机制（监控参考与时钟输出，可切内部 RC） & 切内部时钟并降频运行，或复位。注意窗口狗此时可能一起停摆 \\
\addlinespace
重启循环 & 带时间窗的复位计数加复位原因组合 & 超过阈值停在需人工复位的状态或加载后备镜像；不要无限重试同一路径 \\
\addlinespace
状态数据位翻转 & ECC、影子副本比对、配置区校验 & 能纠则纠并计数；纠不了则丢弃该数据并要求重新加载 \\
\end{longtable}

\section{与调度器和 RTOS 的接口}

\textbf{先分清错过期限与卡死。} RT 任务超时的定义是``这一帧没在承诺时间内完成''，它不等于``任务不再前进''。前者是性能问题，处理方式应当是降级、丢帧、告警；后者才是复位问题。把两者混成一件事（``只要有一帧超时就复位''）会把可恢复的瞬时过载变成整机重启，反而降低可用度。判据要写成两条独立的量：完成时刻的分布（用于降级与统计）与帧序号是否推进（用于狗）。

\DomRtos{} 裸机与 RTOS 里常见三条做法。其一，空闲钩子喂狗：狗由 idle 任务刷新，任何任务霸占 CPU 不放都会让狗咬。它简单可靠，但前提是低优先级的实时任务都允许 idle 有机会跑；一旦所有任务都是自顶向下的忙等，idle 永远不出现，狗会在正常负载下也咬，此时要改走第二条。其二，任务心跳位图：每个任务在自己的完成点置一位，由周期检查逻辑（定时器中断或专门的监督任务）判断位图是否在本周期内全部置位，全置位才允许刷新狗；这是任务狗的最小实现，顺带定位``哪一路没完成''。其三，切换钩子做记账：在上下文切换钩子里记录每个任务实际占用的时间片与最长连续占用，用来发现``优先级更高的任务把一切都压住了''这类狗抓不到的失序。

带设备模型的系统里，狗是一个驱动设备：通过设备树别名取实例，用统一的``设定超时并选择模式（中断、复位、或先中断后复位）''接口装配，各通道可绑定回调与各自超时；配置里通常还有一个``启动时是否关掉上一级引导留下的狗''的开关，正好对应第~\ref{sec:ch04-feed-policy}节的责任交接。用得上的一条经验是：把通道当成成就的载体，一个通道对应一个业务里程碑，而不是所有任务共用一条线。

\DomHost{} 通用内核侧的两类检测器要分开理解。软锁检测靠每核一个监督线程：一个高精度定时器周期性去踢本核的监督线程，若发现该线程长时间没被调度而中断是开着的，说明核上没有别的任务在跑而它却饿着，报软锁；阈值由一个可调参数决定（默认量级十秒，实际告警在其两倍附近）。硬锁检测靠不可屏蔽事件（性能计数器的 NMI）或者核间伙伴互相观察，判断的是``连中断都进不来''。前者的盲区正是长时间关中断（此时监督线程根本不会被唤醒），后者恰好补这块。两类都可以配置成只打印全核栈回溯或直接 panic，panic 之后由硬件狗或下一段引导逻辑把系统拉回来。用户态监督则按第~\ref{sec:ch04-taskwd}节的设备接口叠加，或者交给 init 系统按服务配置的重启间隔（服务周期性报告存活），形成应用、内核、硬件三层。三层要各管一段：应用层管业务成就，内核层管活性与不可屏蔽段，硬件层管前两层本身失效。

\textbf{调度器钩子里的看门狗位。} 如果内核或 RTOS 提供切换、空闲钩子，最省事的正确做法是把狗刷新放在``确认调度器本身在推进''的位置，并要求至少一个业务成就标记同时成立。只凭``还在切换''就喂，等于把监督降级为``还有一段代码在跑''。

\begin{warning}[title={调试器 halt 触发复位导致无法单步}]
开发期最常见的``狗误咬''是断点或核暂停期间狗继续计数，恢复后立刻复位，于是根本没法单步。正规处理是在调试模式下冻结狗计数（多数实现有对应选项），或让外部芯片在检测到调试探针在线时旁路。要求是：这条旁路只能由调试探针的物理状态触发，不能由软件标志触发，否则它就变成量产固件里的监督后门。
\end{warning}

\section{常见陷阱}

\begin{itemize}
\item \textbf{在中断里喂狗。} 见前面的警告框。判据：从中断上下文往回追狗重载的写入点，只要找得到一条就是错的。
\item \textbf{有多个喂狗点。} 多点是逻辑或关系，任一多余点会让其余监督失效。审查方式是数``谁写了重载''，而不是数``谁声称自己会喂''。
\item \textbf{没使能狗所依赖的时钟却以为有监督。} 独立狗的 RC 可能在某个低功耗态被关，窗口狗的总线时钟可能在切换瞬间停了一下，那段时间是零监督。上电与状态转移序列里要显式确认狗当时确实在计数。
\item \textbf{复位后没重新配置时钟就进入低电压窗口。} 复位把时钟树恢复到默认（内部 RC、主 PLL 未起），同一段代码的执行时间可能成倍变化。顺序是先喂再改频，改完之后按新预算重设超时。
\item \textbf{把日志写入当成最后要做的事。} 预警窗口里做 flash 擦除或大块串口输出，会把到达安全状态的时间用光。顺序永远是：先钉输出，再留证据。
\item \textbf{忘了清复位标志，或者清了但没保存。} 前者导致下次误判复位原因，后者导致这一轮的取证永远缺失。固定成``快照、清除、分支''三步。
\item \textbf{用固定延时代替就绪条件。} 传感器与电源的稳定时间随温度与批次变化；跨条件后``等得不够''会以随机数据进入控制环的形式出现，而这不会被狗抓到。
\item \textbf{把看门狗当错误处理用。} 复位会丢掉全部未持久化状态，并造成一段无人监督的重启窗口。它应该是最后一档，而不是``检测到异常就重启''的默认答案；能降级的一律先降级，并把``为什么没降级成功''写进证据。
\item \textbf{没测过``多久能到安全状态''。} 从失效发生到输出真正停在安全值之间，包含检测、复位、重启、重新初始化四段，往往是几百毫秒到秒级。这个数字要和安全需求给出的故障容忍时间对得上；对不上时正确的结论是换检测手段或加冗余，而不是把狗超时调短。
\item \textbf{只验狗会咬，不验咬完之后能不能回来。} 复位策略、后备镜像、降级启动模式这些``咬之后''的路径同样需要故障注入覆盖，而它们才真正决定可用度。
\end{itemize}

本章与第 3 章的接口只有一处：前一章管``事件来得多快''，本章管``如果它不来了怎么办''，而两者唯一的连接点就是喂狗点必须是某个可测量成就的后果。这一处写对，整套监督才有语义；写错，装再多的狗也只是给自己一个让人放心的数字。

% ============================================================
% ---- frag_ch05.tex ----
\chapter{DMA、缓冲与确定性数据通路}

\label{ch:ch05-dma-datapath}

本章讨论数据如何在外设与内存之间移动，以及这种移动在时间上是否可预测。\DomBare{} 的实现者常把 DMA 当作一种提速手段，但在实时系统里它首先是一种\textbf{时间结构}：搬运交给硬件之后，CPU 侧的最坏情况执行时间确实下降了，代价是数据到达的时刻不再由软件决定，而由总线矩阵仲裁、通道优先级与缓存状态共同决定。因此本章不讲寄存器位域怎么写（字节级写法属于本仓库的另一分册，讲硬件寄存器与内核侧驱动），只讲三件事：搬运形态的选型、一致性的责任划分、以及把搬运时间真正量出来的方法。

\section{取数据的四种形态}

从一个外设把 $N$ 个字节搬进内存，实现上只有四条路，但它们的实时特性差别极大。

\begin{compare}[title={轮询、中断逐字节、中断批量与 DMA 的实时特性差异}]
轮询的延迟下界最好（一次读寄存器的时间），但上界由轮询周期决定，且 CPU 占用与数据率成正比；中断逐字节把 CPU 占用降到了每字节一次进出中断，却把中断开销乘以了字节数；中断批量在帧边界上打断一次，开销摊薄到整帧；DMA 把搬运彻底移出 CPU，中断只出现在半程与终点，抖动来源从指令流转移到了总线仲裁。选型时不要只比平均占用，要比较\textbf{最坏延迟}与\textbf{延迟抖动}这两列。
\end{compare}

\begin{table}[htbp]
\centering
\caption{四种取数据形态的实时特性对照（数值为量级估计，随主频与总线宽度变化）}
\label{tab:ch05-acquire-modes}
\begin{tabular}{>{\raggedright\arraybackslash}p{0.20\textwidth} >{\raggedright\arraybackslash}p{0.22\textwidth} >{\raggedright\arraybackslash}p{0.20\textwidth} >{\raggedright\arraybackslash}p{0.21\textwidth}}
\toprule
形态 & CPU 占用（量级） & 最坏延迟来源 & 抖动主要来源 \\
\midrule
轮询（含忙等状态机） & 与数据率线性正比，每字节数十周期 & 轮询周期本身；被高优先级任务抢占后整轮丢失 & 抢占时长、循环体编译结果变化 \\
中断逐字节 & 每字节一次进出中断，约百周期每字节 & 中断响应延迟加上同优先级中断排队 & 中断嵌套深度、向量表取指缓存状态 \\
中断批量（帧中断） & 每帧一次中断，摊薄后很低 & 一整帧的到达时间加帧尾处理 & FIFO 水位抖动、总线等待状态 \\
DMA（加完成中断） & 接近零，仅配置与收尾 & 传输排队等待总线授权的时间 & 多主仲裁、缓存行被预取或驱逐 \\
\bottomrule
\end{tabular}
\end{table}

\subsection{什么时候 DMA 反而更差}

三种情形下 DMA 是负收益。第一，\textbf{短包}：一次 DMA 事务的启动需要写若干配置寄存器、做内存屏障、等通道握手，这段固定开销在几个字节到十几个字节的搬运上往往超过直接读寄存器再写内存的成本；许多控制器的通道启动延迟达到数十个总线周期。第二，\textbf{总线仲裁密集}：当同一总线矩阵上有刷新、另一路 DMA、以及调试探针在竞争时，被抢占通道的完成时间上界不可控，而 CPU 搬运的指令流反而更稳定。第三，\textbf{强耦合的即时处理}：如果每收到一个字节就要决定是否改发下一个字节（例如协议协商、变长帧头解析），DMA 只是把搬运做完，判断仍要中断参与，于是付出了 DMA 启动开销又付了中断开销。

一个经验判据：当单块搬运字节数除以总线带宽所得的时间，与 DMA 启动加收尾的固定开销之比小于三到五倍时，DMA 的优势不足以覆盖它的复杂度与抖动。

\section{DMA 通道、请求路由与缓冲模式}

\subsection{请求选择器与通道分配}

早期的 DMA 把外设请求固定映射到某条流或某个通道（换引脚或换外设就要换通道），后来的控制器引入了请求选择器（常见名称为 DMAMUX 或请求路由器）：任意通道可以绑到任意请求线，绑定关系在运行时写入。这带来了灵活性，也带来两类实时风险。

其一是\textbf{通道冲突}：两个模块各自把同一个请求线绑到同一通道，或者两条通道被配成同一优先级却在同一时刻请求总线，后到者必须等前者释放。其二是\textbf{优先级反转的静态化}：通道优先级通常是配置好的固定值，若把关键音频流和低速日志口配成同一等级，日志的长传输会推迟音频流的服务窗口。稳妥做法是在设计阶段列一张通道占用表，让每个请求线只有一个主人，并让关键流的通道优先级高于批量流；如果控制器支持突发长度与轮询权重，把批量流的突发切短，可以压缩关键流的最坏等待。

\subsection{循环模式、双缓冲与乒乓}

循环模式让通道在计数到零后自动重装，省去重新配置的时间与窗口。在循环模式上进一步提供两类硬件形态：\textbf{双缓冲}（硬件持有两个目标地址指针，传输完成时自动切换目标并可选产生缓冲中断）与\textbf{乒乓}（软件用半传输中断与传输完成中断把一个缓冲区一分为二）。前者把切换交给硬件，因而没有软件切换窗口；后者只依赖两个中断点，实现简单，但要求 CPU 侧处理时间小于一块的时长。

清单 \ref{lst:ch05-pingpong} 给出乒乓形态的最小骨架。要点有三个：块尺寸直接等于半传输周期；回调里\textbf{绝不}做耗时动作，只认领索引；消费者对块的认领用一次写完成，避免与硬件争同一个变量。

\begin{lstlisting}[style=cstyle,caption={C：I2S 连续接收的乒乓双缓冲与半传输回调},label={lst:ch05-pingpong}]
#include <stdbool.h>
#include <string.h>

#define I2S_BLOCK_WORDS  256U                      /* 每块字数 */
#define I2S_FRAME_BYTES  (I2S_BLOCK_WORDS * 4U)    /* 32bit 槽 = 2 通道 x 16bit */

static uint8_t  sBuf[2U][I2S_FRAME_BYTES] __attribute__((aligned(32)));
static volatile uint8_t sReady = 0xFFU;            /* 0xFF 表示无新块 */

void AudioPath_Init(void)
{
    /* 整段（两块）一次性交给 DMA，循环模式：
       半传输 -> 块0 就绪，完成 -> 块1 就绪 */
    HAL_I2S_Receive_DMA(&hi2s1, &sBuf[0][0], I2S_BLOCK_WORDS * 2U);
}

void HAL_I2S_RxHalfCpltCallback(I2S_HandleTypeDef *h)
{
    (void)h;
    sReady = 0U;                    /* 硬件此刻正在写块1，块0 可读 */
}

void HAL_I2S_RxCpltCallback(I2S_HandleTypeDef *h)
{
    (void)h;
    sReady = 1U;                    /* 硬件此刻回卷重写块0，块1 可读 */
}

/* 消费者：只读已认领的块，拷贝时长必须小于一块的到达周期 */
bool AudioPath_Take(int16_t *out)
{
    uint8_t idx = sReady;
    if (idx == 0xFFU) { return false; }
    sReady = 0xFFU;                                   /* 认领，先于使用 */
    if (idx > 1U)   { return false; }                 /* 防御：值域检查 */
    memcpy(out, sBuf[idx], I2S_FRAME_BYTES);          /* 唯一一次受控拷贝 */
    return true;
}
\end{lstlisting}

\subsection{散列聚集与描述符链}

复杂搬运（多个源拼成一个帧、多个目的地分发、每块长度可变）用描述符链表达。每个描述符含源地址、目的地址、长度、控制字与指向下一个描述符的指针；控制字里通常还有一个所有权位（习惯称为 own 位）表示这一条目当前归硬件还是软件。硬件推进链时会把已完成的描述符写回（清 own 位、更新状态字段），而软件同一时刻可能在重新初始化同一条目，于是出现状态机竞态：软件刚写的 own 位被硬件的旧写回覆盖，通道要么停在原地，要么跳过一块。

规避方式是把链的更新做成\textbf{单写者循环}：硬件只写它正在使用的条目，软件只重写已被自己确认完成的条目，并把条目数取为足够大，使得软件的重写永远追不上硬件。清单 \ref{lst:ch05-descriptor-chain} 展示四条目环形链的初始化，重点在末尾的屏障与头指针写入顺序。

\begin{lstlisting}[style=cstyle,caption={C：四描述符环形链的构造与提交},label={lst:ch05-descriptor-chain}]
/* 描述符字段顺序按本控制器手册填写；此处只表达链结构与提交顺序 */
typedef struct {
    uint32_t src;          /* 源地址：外设数据寄存器或内存 */
    uint32_t dst;          /* 目的地址 */
    uint32_t next;         /* 下一描述符地址；末元素回指首元素，形成环 */
    uint32_t ctrl;         /* 长度、宽度、增量使能、中断使能、所有权位 */
} DmaDesc;

#define DESC_N     4U
#define CHUNK      256U                 /* 每条目搬运字节数 */

static DmaDesc desc[DESC_N] __attribute__((aligned(32)));
static uint8_t sink[DESC_N][CHUNK] __attribute__((aligned(32)));

void Chain_Init(void)
{
    for (uint32_t i = 0U; i < DESC_N; i++) {
        desc[i].src  = ADC_DR_ADDR;
        desc[i].dst  = (uint32_t)&sink[i][0];
        desc[i].next = (uint32_t)&desc[(i + 1U) % DESC_N];
        desc[i].ctrl = CHUNK | DESC_OWNER_HW | DESC_IRQ_ON_DONE;
    }
    __DSB();          /* 链内容必须先于头指针可见，否则硬件可能读到半初始化条目 */
    WRITE_REG(DMA_HEAD, (uint32_t)&desc[0]);
}

/* 软件只重写硬件已经交还的条目：追不上就丢弃而不是改写 */
bool Chain_Recycle(uint32_t idx)
{
    if (idx >= DESC_N) { return false; }
    if ((desc[idx].ctrl & DESC_OWNER_HW) != 0U) { return false; }   /* 仍属硬件 */
    desc[idx].ctrl = CHUNK | DESC_OWNER_HW | DESC_IRQ_ON_DONE;
    __DSB();
    return true;
}
\end{lstlisting}

另一个易被忽略的问题是\textbf{描述符被优化器重排}。描述符数组必须具有易变语义（声明为易失，或至少用屏障把初始化与提交隔开），且必须落在硬件可访问的地址区间；把描述符放在只被 CPU 访问的紧耦合区看似更快，实际上控制器可能根本无权访问，或访问路径带宽极低。

\section{缓存一致性与 DMA 缓冲}

带数据缓存的核心（如 Cortex-M7 与多数应用处理器）默认把内存区域配成写回、写分配。此时 DMA 与 CPU 看到的是两个不同的数据视图，一致性由软件负责。

\begin{guideline}[title={一致性动作顺序：设备读前 clean，设备写前 invalidate，CPU 读后再次 invalidate}]
CPU 准备数据、DMA 即将读取（CPU 到外设）：对缓冲覆盖的整段地址做 clean（把脏行写回内存），随后一个数据内存屏障，再启动通道。DMA 即将写入、CPU 稍后读取（外设到 CPU）：启动前先 invalidate（丢弃可能已被预取的旧行）；传输完成后\textbf{再}做一次 invalidate，因为预取器可能在传输期间把旧数据取进缓存；最后屏障加指令屏障，再解析数据。只做一次 invalidate 是最常见的错误。
\end{guideline}

\subsection{按缓存行对齐与补齐}

clean 与 invalidate 的最小单位是缓存行。如果缓冲起止不对齐，一次操作会牵连相邻数据：invalidate 会把邻居尚未写回的新值丢掉（邻居随后写回的是旧内存内容），clean 会把邻居尚未定稿的半成品推给内存。因此每个 DMA 缓冲都要同时满足\textbf{地址对齐}到缓存行大小与\textbf{尺寸补齐}到缓存行整数倍。

\subsection{不可缓存区与 MPU 属性}

另一条路是把缓冲放在不可缓存的内存域，或用 MPU 把该地址段标记为不可缓存（更严格的还有设备属性），于是硬件视图与内存视图始终一致，无需任何软件动作。代价是 CPU 访问这些地址每次都要走总线、没有行缓冲，读取大块数据时性能显著下降，而且设备属性区域不允许非对齐访问与突发访问。典型折中：把描述符链与少量控制块放不可缓存区（CPU 访问少、状态必须实时可见），把数据大块留在可缓存区并显式做一致性动作。

\subsection{为什么一个 buffer 不能被两个方向共享}

同一块内存如果既被 DMA 写又被 DMA 读（或同一缓冲被两个通道分别读写），一致性动作就无法安全：为一侧做的 invalidate 会抹掉另一侧尚未写回的脏行；为一侧做的 clean 会把另一侧还没写完的半成品推入内存，破坏数据。硬实时路径上的结论是\textbf{每个方向独立缓冲、独立缓存行范围}；如果确实需要就地处理，就把数据显式搬经 CPU 可控的暂存区，绝不让两个主设备共享同一缓存行区间。

清单 \ref{lst:ch05-cache-aligned} 把这套责任封装成一个类型，避免每处手写序列。

\begin{lstlisting}[style=cppstyle,caption={C++：按缓存行对齐的 DMA 缓冲与一致性动作封装},label={lst:ch05-cache-aligned}]
#include <cstddef>
#include <cstdint>

constexpr std::size_t kCacheLine = 32U;      /* Cortex-M7 数据缓存行为 32 字节 */

template <std::size_t N>
class alignas(kCacheLine) DmaBuf {
    static_assert((N % kCacheLine) == 0U, "尺寸必须是缓存行整数倍");
public:
    std::uint8_t *data()        { return mRaw; }
    std::size_t   size() const  { return N; }

    void beforeDeviceRead()  { CleanByAddr(mRaw, N);         __DSB(); }  /* CPU 到外设 */
    void beforeDeviceWrite() { InvalidateByAddr(mRaw, N);    __DSB(); }  /* 启动前 */
    void afterDeviceWrite()  { InvalidateByAddr(mRaw, N);    __DSB(); __ISB(); }

private:
    std::uint8_t mRaw[N];
};

/* 每个方向各一个实例，绝不共用；尺寸补齐到行 */
DmaBuf<512> gAdcSink;        /* 外设到 CPU */
DmaBuf<512> gDacSource;      /* CPU 到外设 */

/* 状态变量单独占一行，避免与缓冲形成伪共享 */
alignas(kCacheLine) volatile std::uint32_t gAdcSeq;
alignas(kCacheLine) volatile std::uint32_t gDacSeq;
\end{lstlisting}

表 \ref{tab:ch05-cache-actions} 把方向、时机与动作对应起来。

\begin{longtable}{>{\raggedright\arraybackslash}p{0.24\textwidth} >{\raggedright\arraybackslash}p{0.19\textwidth} >{\raggedright\arraybackslash}p{0.19\textwidth} >{\raggedright\arraybackslash}p{0.20\textwidth}}
\caption{DMA 搬运方向与 CPU 侧一致性动作对照}\label{tab:ch05-cache-actions}\\
\toprule
阶段 & 搬运方向 & 必需动作 & 漏掉的后果 \\
\midrule
\endfirsthead
\toprule
阶段 & 搬运方向 & 必需动作 & 漏掉的后果 \\
\midrule
\endhead
\bottomrule
\endfoot
\endlastfoot
启动前，数据待设备读取 & CPU 到外设 & clean 覆盖区间，然后屏障 & 设备读到内存里的旧值，长度看似正确但内容是上一帧 \\
启动前，设备即将写入 & 外设到 CPU & invalidate 覆盖区间 & 预取的旧行被当成新数据，出现周期性陈旧帧 \\
完成后，CPU 解析数据 & 外设到 CPU & 再 invalidate 一次，再解析 & 传输期间被预取的行未被丢弃，偶发解析错误 \\
长循环模式运行中 & 两方向各自独立 & 每个块边界重复上述动作 & 只在启动时做一次，第二圈起数据不可信 \\
描述符与状态字 & 双向交替 & 放不可缓存区，或全程不经缓存访问 & own 位与状态字段被缓存掩盖，链看似没推进 \\
\end{longtable}

\subsection{内核态抽象的边界}

在带操作系统的硬件访问层里，常见 \Fn{memcpy\_toio}、\Fn{memcpy\_fromio} 这类地址空间抽象，以及一致性分配接口与显式同步接口（形如 \Fn{dma\_sync\_single\_for\_device}、\Fn{dma\_sync\_single\_for\_cpu}）。需要看清的是：前一类只处理 I/O 地址空间与访问顺序，\textbf{不}等价于缓存维护；真正把方向与时机表达清楚的是后一类的同步接口。而在缓存型微控制器上裸机运行时，没有任何抽象会替你决定这些时机——把 ``免一致性'' 当作默认，在带预取器的控制器上一定会失效。抽象层可以简化写法，但方向与时机这两件事仍由你决定。

\section{等时流与硬实时约束}

\subsection{服务周期的计算}

对连续流（I2S、SAI、PDM、ADC 硬件触发连续转换），FIFO 或乒乓块的深度决定了必须多久服务一次。服务周期等于深度除以速率：$T = D / R$，其中 $D$ 是缓冲中的样本数（若一个 32 位字打包两个 16 位样本，则按样本数折算），$R$ 是每秒样本数。

举例说明数量级：采样率 48 kHz，立体声，每帧两样本；FIFO 深度 8 个 32 位字，即 16 个样本，服务周期为 16 除以 48000，约 333 微秒。若使用每块 256 帧的乒乓块（512 样本），半传输周期为 256 除以 48000，约 5.33 毫秒。这两次计算给出的就是\textbf{必须响应窗口}：块越大，中断越少，但一旦处理超时，损失是整块。

\begin{note}[title={服务周期计算里最容易错的三件事}]
第一，忘记按通道数折算：立体声的每帧样本数是 2，深度 8 字往往意味着 16 或 32 个样本，取决于是否双相打包。第二，把搬运时间当作缓冲深度：搬运本身占用的总线时间也要计入，深缓冲并不能无限推迟服务点。第三，忽略预取与突发：控制器常以突发为单位占用总线，实际服务边界比理论值提前若干个字。
\end{note}

\subsection{溢出与欠载的检测和恢复}

外设侧的溢出（接收方向 FIFO 被填满）与欠载（发送方向 FIFO 被抽干）通常是\textbf{粘滞标志}：置位后保持，直到软件清除。设计上有三点要求。其一，标志要在错误回调里读取并累加到计数器上，只清标志不计数会让丢样在统计里隐身。其二，检测点要覆盖半传输与完成之外的第三类中断（错误中断），否则错误只在下一块才被发现，无法判断丢样发生在哪一块。其三，恢复策略要明确：音频流的常见做法是丢弃当前块、重新对齐到块边界并保持输出静音，而不是尝试插值补样；控制环的 ADC 流则应把该样本标为无效，让上层沿用上一次有效值一段时间后立即进入安全状态。

\subsection{时钟归属与漂移}

等时链路上的时钟必须由一个主人提供。当声卡或编码器是主时钟提供者，你的位时钟与帧时钟就是它的节奏，本地采样计数只用于对齐；当你自己是主，外部设备的时钟偏差会表现为你收到的数据率略高或略低。真实漂移量级：普通晶振几十 ppm，1000 ppm 级别的差异在 48 kHz 上意味着每秒约 48 个样本的偏差，几分钟就会积累出可闻的咔哒声。对策是明确的：要么统一到同一时钟域（同一音频主时钟分发给所有从设备），要么在接收侧做重采样，或者用异步缓冲加水位控制（水位越界时按固定比例增删样本）。永远不要靠加大缓冲来吸收持续漂移，那只是把失败从此刻推迟到一次突发溢出。

\subsection{连续时钟上的间隙}

SPI 类连续流设备靠时钟边界推进状态机。片选抖动、时钟被 DMA 让路、或每次传输之间的软件间隙都会造成设备侧丢位或帧错位。规避办法：用硬件连续传输模式（DMA 循环且外设保持时钟运行），或把帧长配成硬件能自动维持的长度；若必须分段，让分段边界落在帧边界并保持片选全程有效。间隙是否存在、有多长，本身就是选型判据之一：一个不能维持连续时钟的控制器不适合承担等时流。

\section{总线与内存层次}

\subsection{总线矩阵与多主竞争}

现代控制器用交叉开关式的总线矩阵把多个主设备（CPU、多路 DMA、图形加速单元、外部存储控制器）连到多个从设备（SRAM、闪存、外设桥）。竞争发生在从设备侧：同一 SRAM 端口被两个主设备同时请求时，仲裁决定谁先，另一个只能等待。这就解释了关键缓冲的位置为何如此重要：把关键缓冲放在只有它使用的 RAM 域（例如某个外设独占访问的域，或 CPU 很少触碰的那半容量），把批量数据放到另一域，是压缩最坏等待时间最有效的一招。

\subsection{闪存预取与等待状态}

代码从闪存执行时，等待状态与预取器是否使能决定了指令流的实际延迟。提升主频往往需要增加等待状态，而一次缓存未命中会带来几十周期的停顿；临界路径若常驻闪存又没有指令缓存，其最坏时间会因一次取指失败而显著膨胀。硬实时做法：把中断入口与关键路径函数放紧耦合区或指令缓存常驻区，并且实测最坏执行时间时，在预取开启与关闭的两种配置下各验证一遍。

\subsection{带宽预算}

一条搬运链路的可行性用简单加法判断：把每个主设备的峰值字节率相加，与所经从设备端口的字节率比较，并留出不低于三分之一的余量给仲裁损耗与突发对齐。示例数量级：48 kHz 立体声 32 位槽的音频流约 384 KB 每秒；1 MSPS 的 16 位 ADC 流约 2 MB 每秒；图形刷新可能是几十到几百 MB 每秒。后者会明显挤占前者，而前者虽然字节少却是等时的，不容失败。

\begin{bestpractice}[title={把关键 buffer 放到 TCM 或独占 RAM 域的判据}]
按顺序问四句话：这块内存有两个主设备访问吗？它所在的关键路径有缓存未命中的可能吗？它的地址区间与别的批量流重叠吗？它能被 DMA 直接寻址吗？前三句里只要有一句答``是''，就把它迁到紧耦合区或独占域，并补齐缓存行；最后一句若答``不能''，就不要硬塞，改用受控拷贝把数据从可访问域搬过去，用一次代价已知的拷贝换来稳定的最坏时间。
\end{bestpractice}

\subsection{双口 RAM 与共享内存}

多核或协处理器场景常用双口 RAM：两侧各有独立端口，同步由握手位或硬件信号量完成。实时视角的要点是\textbf{冲突窗口}：同一地址被两侧同时读写时，行为依器件而定（有的给出交叉冲突标志，有的静默返回旧值），所以共享区必须划成互不重叠的分区，或者加校验位。另外共享内存一侧往往不带缓存，两侧都要按对方可见性来写顺序，不能假设自己写完对方立刻能读。

\section{缓冲设计与搬运时间的测量}

\subsection{拷贝次数最小化}

外设到应用的数据通路，拷贝次数是可量化指标。理想形态是\textbf{入栈一次、出栈一次}：DMA 直接写入应用可读的块，应用原地读取，不做中间复制。实践中的约束来自两点：应用希望按帧或按记录取数据（需要解析），而 DMA 只接受连续线性地址。折中做法是让 DMA 落到一个大的循环缓冲，解析层用零拷贝视图（指向缓冲的指针加偏移）向上交付，只有在需要跨越缓冲末尾、或需要长时间持有该数据时才真正拷贝一次。

需要划清职责的是：本章只谈外设与内存之间的搬运与缓冲；软件内部的生产者消费者队列、环形缓冲的索引与内存序规则由第 9 章展开，两处的分工不要混。

\subsection{搬运时间与带宽的计算}

单次 DMA 的完成时间可写成 $t = N / B$，$N$ 为字节数，$B$ 为有效字节率。有效字节率不等于总线标称值：突发拆分、地址非对齐、目标外设的流控、以及仲裁都会压低它，工程上要用实测的 $B$ 而不是手册峰值，实测值通常是标称峰值的六成到八五折。反向用法更有价值：给定服务周期 $T$ 与数据率，可容忍的最大块为 $N = B$ 乘以 $T$，用它来定乒乓块尺寸。

\subsection{用打点方式实测搬运延迟}

\begin{warning}[title={只测平均搬运时间会掩盖总线仲裁造成的尖峰}]
搬运延迟的分布是长尾的：绝大多数事务贴着带宽上限完成，少数事务因为遇到刷新、图形突发或另一路 DMA 而拖长数倍。若只看平均值，你会得出``缓冲尺寸足够''的结论，而运行几小时后恰好在那最长尾的一次上溢出。测量必须报最大值与分布尾部（至少记录一百次里的最坏一次），而且要在满负载的整机上测；孤立测试得到的数字普遍偏乐观。
\end{warning}

两种打点手段互补。示波器法：在启动通道前后翻转一个空闲 GPIO，直接看到事务窗口与仲裁造成的拖长。内核计数法：用数据看门狗定时器的周期计数器取周期数，再除以主频换算成时间，精度可到单个总线周期，但没有外部可见波形，且调试使能本身会占用少量资源。清单 \ref{lst:ch05-dwt-timing} 给出周期计数法的骨架。

\begin{lstlisting}[style=cstyle,caption={C：用 DWT 周期计数器与 GPIO 翻转给搬运打点},label={lst:ch05-dwt-timing}]
#include <stdint.h>

static uint32_t gCpuHz;
static uint32_t gMaxCyc, gSumCyc, gCnt;
static uint32_t gDroppedCount;

#define TRACE_HIGH()  HAL_GPIO_WritePin(T_PORT, T_PIN, GPIO_PIN_SET)
#define TRACE_LOW()   HAL_GPIO_WritePin(T_PORT, T_PIN, GPIO_PIN_RESET)

void Timing_Init(uint32_t cpuHz)
{
    gCpuHz = cpuHz; gMaxCyc = 0U; gSumCyc = 0U; gCnt = 0U;
    CoreDebug->DEMCR |= CoreDebug_DEMCR_TRCENA_Msk;   /* 使能跟踪时钟 */
    /* 部分内核要先写解锁寄存器才能启动计数，按手册顺序执行 */
    DWT->CYCCNT = 0U;
    DWT->CTRL  |= DWT_CTRL_CYCCNTENA_Msk;
}

static void Record(uint32_t startCyc)
{
    uint32_t d = DWT->CYCCNT - startCyc;               /* 无符号回绕安全 */
    if (d > gMaxCyc) { gMaxCyc = d; }
    gSumCyc += d; gCnt++;
}

/* 测一次 DMA 事务：CPU 侧开销与设备完成时刻 */
void Bench_Transfer(void *dst, const void *src, uint32_t n)
{
    uint32_t t0;
    TRACE_HIGH();
    t0 = DWT->CYCCNT;
    StartChannel(dst, src, n);      /* 配置、屏障、启动 */
    while (!ChannelDone()) { }      /* 忙等只用于测量，产品代码应中断驱动 */
    Record(t0);
    TRACE_LOW();
}

uint32_t WorstNanos(void)
{
    /* 周期换算成纳秒：先乘后除保住精度 */
    return (uint32_t)(((uint64_t)gMaxCyc * 1000000000ULL) / gCpuHz);
}
\end{lstlisting}

\section{常见陷阱}

\begin{enumerate}
\item \textbf{缓冲与状态标志共用同一缓存行}。CPU 写标志导致整行写回，覆盖 DMA 刚写入的相邻数据；或 DMA 写数据使 CPU 的标志行变脏。解决办法是让标志独占一行（按缓存行对齐并补齐），清单 \ref{lst:ch05-cache-aligned} 末尾即为此设计。
\item \textbf{半传输回调里读写正在被硬件写的那一半}。半传输中断发生时硬件已经在写另一半，若回调越界去动正在被写的部分，得到的是撕裂数据。回调只做索引与计数，块的实际解析放在消费者任务里。
\item \textbf{忘记完成后 invalidate}，CPU 看到旧数据。现象是首块正确、后续偶发正确，因为缓存状态取决于此前有别的工作是否已把这些行取进来。
\item \textbf{循环模式重装载窗口的丢失}。某些控制器在计数回卷到重装载之间存在极短的内部窗口，此时到达的数据可能丢一个样本。若控制器提供双缓冲或链式回卷就优先用它们；否则让块边界不与外部帧边界重合，并把溢出检查计入错误计数。
\item \textbf{外设关闭后 DMA 仍在搬}。停止顺序写错（先关外设再停通道时，通道可能继续读已清空的外设数据寄存器而报总线错误，或反向情况下继续搬运垃圾）。正确顺序是：先停通道并确认通道空闲，再撤掉外设请求，最后清挂起的中断标志。
\item \textbf{描述符链被优化器重排}。链字段没有易变语义、初始化与头指针写入之间缺少屏障，或每条目长度不是编译期常量而被提出循环，都会造成链在观察上正确却不再推进。
\item \textbf{把 DMA 完成等同于数据可用}。完成中断只说明控制器发出了最后一笔写，这笔写可能在桥接层里仍在排队。跨时钟域或经过缓冲桥的路径上，读数据前加一个数据屏障，或用一个读回操作确认落地。
\item \textbf{用加大缓冲代替解决漂移与过载}。加大缓冲只把失败推迟，同时让服务周期与最坏延迟同步变差；先算清速率与余量，再定尺寸。
\end{enumerate}

本章的判据可以归结为一句话：确定性来自\textbf{单一主人、独占地址域、固定服务周期与实测尾延迟}。DMA 提供了低 CPU 占用的搬运，但它把不确定性从指令流搬到了总线与缓存；只有在选型时比较最坏延迟、在实现时把一致性动作固化成封装、在验证时量出分布尾部，数据通路才真正可预测。

% ============================================================
% ---- frag_ch06.tex ----
\part{RTOS 与调度模型}

\chapter{FreeRTOS：任务、队列与同步}

\label{ch:ch06-freertos}

从本章起进入本册的第二条主线：当搬运形态（第 5 章）确定之后，谁来消费这些数据、以什么周期消费、被谁抢占，就交由内核决定。\DomRtos{} 的选型很多，本章以 FreeRTOS 为具体对象，因为它在微控制器上的实时语义最直白：固定优先级抢占、时间片可选、所有原语的开销都能算清。也正因如此，它对使用者的要求更严格——内核不会替你隐藏不确定性，配置宏写错一个数值，最坏延迟就会翻倍。

\DomRT{} 关心的不是能不能跑，而是在最大负载与最坏中断时序下还守不守得住周期。本章每一节都围绕这个问句展开。

\section{任务对象与栈预算}

\subsection{动态与静态创建}

\Fn{xTaskCreate} 从内核堆里同时取出任务控制块与栈，返回句柄；\Fn{xTaskCreate{\Bk}Static} 由调用方提供这两块存储，内核只登记。二者语义等价，区别全在时间性与可验证性上。

\begin{compare}[title={动态创建与静态创建在硬实时项目里的差别}]
动态创建把存储风险推迟到运行期：堆不够时返回空，而堆是否够用取决于此前所有分配的历史，静态分析无法证明；\Fn{pvPortMalloc} 在多数堆实现里要遍历空闲链表，最坏时间不可界。静态创建在链接期就把地址与容量固定，启动阶段一次登记完成，运行期不再有任何分配点，代价是要手工给出 \Tp{StaticTask\_t} 与栈数组，并保证对齐与尺寸正确。硬实时路径上的任务、中断里会用到的队列与信号量，都应走静态形态；只有生命周期确属可选的功能（诊断命令、临时会话）才允许动态创建。
\end{compare}

清单 \ref{lst:ch06-static-periodic} 给出静态创建加绝对周期等待的完整骨架，这也是本章推荐的默认写法。

\begin{lstlisting}[style=cstyle,caption={C：静态创建的周期任务，用绝对唤醒时间避免漂移},label={lst:ch06-static-periodic}]
#include "FreeRTOS.h"
#include "task.h"

#define CTL_STACK_WORDS   256U          /* 单位是 StackType_t 个数，不是字节 */
#define CTL_PRIO          3U
#define CTL_PERIOD_TICKS  5U            /* 时基 1 毫秒时即 5 毫秒周期 */

static StaticTask_t xCtlTcb;
/* 静态对象须与 portBYTE_ALIGNMENT（常见 8 字节）对齐 */
static StackType_t  xCtlStack[CTL_STACK_WORDS] __attribute__((aligned(8)));

static void Ctl_Task(void *pvParams)
{
    (void)pvParams;
    TickType_t xNextWake = xTaskGetTickCount();

    for (;;) {
        Ctl_RunOnce();              /* 最坏执行时间必须实测，并留进周期 */

        /* 基准是上一次唤醒的目标时刻，不是当前时刻，
           因此本轮耗时不会把下一轮推移掉 */
        vTaskDelayUntil(&xNextWake, CTL_PERIOD_TICKS);

        /* 若本轮严重超时，可用 xTaskCatchUpTicks 把落后的若干个 tick
           一次追平，使下一轮回到正确相位，而不是连续补跑数轮 */
    }
}

void Ctl_Start(void)
{
    TaskHandle_t xHandle = xTaskCreateStatic(Ctl_Task, "ctl",
                                             CTL_STACK_WORDS,
                                             NULL,
                                             CTL_PRIO,
                                             &xCtlStack[0],
                                             &xCtlTcb);
    configASSERT(xHandle != NULL);
}
\end{lstlisting}

\subsection{栈深度的单位与实测}

栈深度参数一律以\textbf{字}为单位（\Tp{StackType\_t} 的个数），在 32 位端口上要乘以 4 才是字节。这是最常见的误配来源：把 200 当成 ``200 字节''，实际只有 800 字节空间；或反过来按字节填数字，浪费四倍内存。

真正的预算要靠实测。系统稳定运行一段时间后调用 \Fn{uxTaskGet{\Bk}Stack{\Bk}High{\Bk}Water{\Bk}Mark}（需要 \Tp{configUSE\_{\Bk}TRACE\_{\Bk}FACILITY} 置 1），它返回该任务历史最深使用时\textbf{剩余}的字数；把每个任务的最坏余量记录下来，再按不低于两倍的余量确定最终深度，因为一次新增的调试打印或一次编译器内联决策就能吃掉几十个字。两点提醒：水印只在被执行过的路径上收敛，从未走到的错误分支不会暴露它的真实需求，因此余量要按覆盖测试的结果定；中断上下文使用主栈或独立的中断栈，不进入任务水印统计，中断里的大数组要单独核算。

\subsection{两种栈溢出检测与保护页}

\Tp{configCHECK\_FOR\_STACK\_OVERFLOW} 取 1 时，内核在每次上下文切换检查栈指针是否越出记录的栈区间，额外开销几乎为零，但只有发生切换时才可能发现；若溢出早已破坏了别的对象而栈指针又回来了，就永远看不到。取 2 时，创建阶段在栈底填入固定填充模式，切换时检查模式是否被改写，能捕获写坏了但栈指针没越界的情形，代价是创建时的填充写入与切换时更多比较。若使用带 MPU 的 Cortex-M 端口，还有第三条路：在每个任务栈下方留一个不映射的保护页，越界写入立刻触发存储管理异常，精确停在出错的那条指令上——这是唯一能在破坏发生瞬间抓到的方案，代价是每次切换要换页配置、每任务额外一页内存。

检测开着不等于安全。硬实时项目的默认答案是：静态分配加保护页，开发构建里再把检测设成 2，并实现 \Fn{vApplicationStackOverflowHook} 记录任务名后进入可控失效，绝不带着已知溢出继续运行。

\subsection{任务删除与资源回收}

\Fn{vTaskDelete} 传入空指针表示删除自身，此时该任务的栈与控制块由空闲任务释放；删除别的任务时，被删者的存储同样要等空闲任务在下一轮运行里回收——如果空闲任务被长期饿死，内存不会真正回来。更麻烦的是副作用：内核不会替被删任务释放它持有的互斥量（于是其他任务永久阻塞）、不会关闭它启动的外设流（DMA 仍在搬、中断仍在来）、不会归还它从应用内存池借出的块。可行的规矩是：任务只用\textbf{自删除}，删除前显式完成停外设、还池、释放自己持有的同步对象这三步；静态创建的任务不回收存储，删除只是让它不再被调度。

\subsection{任务通知：最轻量的唤醒}

\Fn{xTaskNotify}、\Fn{xTaskNotifyGive}、\Fn{ulTaskNotifyTake}、\Fn{xTaskNotifyWait} 把唤醒信息挂在任务控制块上，省掉一个独立内核对象与一次队列拷贝。相对二值信号量，通知更快、更省内存，还能携带一个 32 位值或一组位，或者累计成计数。适用边界要看清：通知只有一个接收者，不能被多任务等待，也不参与队列集合；发送方不等待接收方消费，因此不适合表达有界缓冲满则阻塞的语义。把通知当作事件位加唤醒的组合，用来替代简单信号量与小型事件组，是本册推荐的默认做法。

\section{时基、周期与优先级方向}

\subsection{时基频率与帧周期错频}

\Tp{configTICK\_RATE\_HZ} 决定 tick 周期（\Tp{portTICK\_PERIOD\_MS} 是它的倒数换算），典型取 1000 即 1 毫秒，也有取 100 以降低开销、或取更高以提高分辨率。tick 越高，tick 中断与空闲处理的相对开销越大、功耗越差；tick 越低，最短可表达延迟越粗糙。

错频（beat frequency）指任务周期不是 tick 周期的整数倍。以 1 毫秒 tick 为例，5.33 毫秒的音频帧周期只能表达成 5 或 6 个 tick，误差分别是每帧负 0.33 毫秒与正 0.67 毫秒；即便用 \Fn{vTaskDelayUntil} 保证不累积，实际相位仍会缓慢扫过整帧，表现为周期性的连发两帧或隔一轮三帧。三种处理方式：把周期源改成帧率的整数倍关系（例如用外部定时器或 DMA 半传输中断直接驱动，而不是靠 tick 对齐）；用高精度定时器作周期源，任务只被通知唤醒；或者接受相位扫描，但把缓冲水位设计得足以吸收一整帧的相位误差。绝对不要靠让周期任务稍微跑长一点去对齐外部事件，那会把可控的相位扫描变成不可控的漂移。

\subsection{优先级数值方向与时间片}

FreeRTOS 的数值越大优先级越高，\Tp{tskIDLE\_{\Bk}PRIORITY} 为 0，可用上界由 \Tp{configMAX\_{\Bk}PRIORITIES} 决定（留出一档给内核或诊断，别让所有任务挤在同一数值上）。这一方向与多数通用操作系统相反，代码评审时要专门核一遍优先级常量表：一个把优先级写成越小越高的团队，会造出高优先级任务永远抢不到 CPU 的静默故障。

\Tp{configUSE\_TIME\_SLICING} 置 1 时，同优先级的就绪任务在每次 tick 轮流执行；置 0 时，先到者一直运行直到主动阻塞或让出。硬实时路径上倾向关闭时间片：同优先级的两个关键任务本不该存在，若确有，让它们分档比依赖轮转更可预测。反过来，若干条低重要度的批量任务共享一档优先级时，开启时间片比给它们逐个排优先级更安全，因为不会有谁被永久饿死。

\subsection{相对等待与绝对周期}

\Fn{vTaskDelay} 从现在起算，等待时长是执行耗时的函数：本轮多花 3 个 tick，下一轮就整体后移 3 个 tick，误差逐轮传递；若任务耗时随负载变化，输出就带着抖动。相对等待只适合重试退避、等一会儿再看这类不关心相位的场合。所有周期性工作都应使用 \Fn{vTaskDelayUntil}：基准是上一轮的目标唤醒时刻，本轮耗时不会推移下一轮的目标，周期等于设定的间隔而与执行时间无关。需要追平时用 \Fn{xTaskCatchUpTicks} 传入落后的 tick 数，让相位回到当前时刻之后，避免连续补跑。

阻塞参数一律以 tick 计：换算写成 tick 数乘以 tick 周期得到毫秒，别把毫秒直接当 tick 传。\Tp{portMAX\_DELAY} 表示无限等待，用在纯事件驱动的任务上最省 CPU，但一旦上游丢失唤醒就永久挂起，因此必须配套看门狗（第 10 章）。

\subsection{空闲任务与 tickless 下的时间基恢复}

空闲任务负责回收被删除任务的存储，并在启用 tickless 时决定是否可以睡眠。睡眠前内核关掉 tick 中断，醒来后必须\textbf{补足}睡掉的 tick 数：端口的睡眠实现负责调用推进函数（\Fn{vTaskStepTick}），而实际睡眠时长必须由一个仍在运行的低功耗计数器给出。这里有两个坑。其一，用一个会随主频一起停的计数器测睡眠时长，醒来后 tick 与墙钟已经脱钩，所有周期任务相位错乱。其二，睡眠期间到达的中断唤醒了内核，但补记的 tick 与实际睡着的时间不一致，偏差按次累积；因此唤醒源要触发一次真实的时基校正，并用 \Tp{configEXPECTED\_IDLE\_TIME\_BEFORE\_SLEEP} 设门槛，短于门槛的不睡眠。

\section{堆实现、内存池与对齐}

\subsection{五种堆的行为差异}

\begin{longtable}{>{\raggedright\arraybackslash}p{0.13\textwidth} >{\raggedright\arraybackslash}p{0.22\textwidth} >{\raggedright\arraybackslash}p{0.21\textwidth} >{\raggedright\arraybackslash}p{0.25\textwidth}}
\caption{FreeRTOS 堆实现对比与选型}\label{tab:ch06-heap-variants}\\
\toprule
实现 & 分配行为 & 释放行为 & 最坏时间与适用 \\
\midrule
\endfirsthead
\toprule
实现 & 分配行为 & 释放行为 & 最坏时间与适用 \\
\midrule
\endhead
\bottomrule
\endfoot
\endlastfoot
\Tp{heap\_1} & 单调向前切块，分配即指针相加，池满即返回空 & 不支持释放，调用即触发断言 & 分配时间恒定，可写进最坏情况分析；适合所有对象生命周期等于上电周期的产品 \\
\Tp{heap\_2} & 最佳适配，遍历空闲块 & 支持释放，但相邻空闲块不合并 & 遍历长度随历史增长，碎片只增不减；只用于允许释放但不要求长时间连续运行的场景 \\
\Tp{heap\_3} & 直接包装编译器 \Fn{malloc} & 包装 \Fn{free} & 时间由 C 库决定且通常带锁，中断上下文不可用；仅用于原型验证 \\
\Tp{heap\_4} & 首次适配遍历空闲链表，大块可拆分 & 释放时与相邻空闲块合并 & 最坏时间取决于空闲块数量，仍不可界；适合确有动态对象、又不想引入第三方分配器的产品 \\
\Tp{heap\_5} & 与 \Tp{heap\_4} 同算法，但地址空间由多段不连续区域拼成 & 同 \Tp{heap\_4}，跨段边界不合并 & 时间特性同 \Tp{heap\_4}；用于分散的 RAM 域，例如把紧耦合区与外扩 SRAM 合成一个池 \\
\end{longtable}

结论很直接：\Fn{pvPortMalloc} 没有可用的最坏执行时间界，它要遍历空闲结构、可能拆分与合并，循环次数由内存历史决定。因此硬实时路径（周期任务体、中断服务例程、任何有截止期保证的代码）里不得出现分配；分配集中在启动阶段或维护窗口，并用 \Tp{configUSE\_MALLOC\_FAILED\_HOOK} 挂上失败钩子，让分配失败成为一次可观测的失效，而不是空指针解引用。

\subsection{静态对象全家桶与对齐要求}

内核为每类对象提供了静态版本与对应的结构体类型：\Tp{StaticTask\_t}（任务）、\Tp{StaticQueue\_t}（队列、信号量与互斥量）、\Tp{StaticSemaphore\_t}（独立信号量存储）、\Tp{StaticEventGroup\_t}（事件组）、\Tp{StaticTimer\_t}（软件定时器）、\Tp{StaticStreamBuffer\_t} 与 \Tp{StaticMessageBuffer\_t}（流缓冲与消息缓冲）。使用它们要遵守三条：结构体与缓冲区按 \Tp{portBYTE\_ALIGNMENT}（常见为 8 字节）对齐；栈用 \Tp{StackType\_t} 数组而不是字节数组，长度按字计；静态对象初始化后不要靠 memset 重置，要重启某个功能就走内核支持的删除与重建路径。

\subsection{内存池与碎片}

生产代码更稳的做法是自建固定尺寸池：一个数组加一条空闲链表，取与还都是常数时间、无碎片，容量能在编译期证明。池的尺寸按在途最大对象数再加一两个余量来定；取不到块时的行为要显式定义（丢弃并计数、覆盖最旧、或进入安全状态），绝不能退化成一次带等待的分配。若确实需要变长，用少量离散档位（例如 32、64、128 字节的三档池），不要自己写通用分配器。

\section{IPC 与同步对象选型}

\subsection{队列：按值拷贝与传指针}

创建队列时确定两个尺寸：长度（可容纳的元素个数）与元素字节数；动态形态的内存开销近似为固定控制块加上长度乘以元素大小。元素是\textbf{按值拷贝}进队列的，因此小结构直接传值最安全，发送方在 \Fn{xQueueSendToBack} 返回之后可以立刻复用那块内存。\Fn{xQueueSend} 与 \Fn{xQueueSendToBack} 等价，另有 \Fn{xQueueSendToFront} 可把紧急项插到队首，用它要确认消费者不会因为队首插入而错过顺序约定。

传指针（零拷贝）适合大块：队列元素只是 4 字节或 8 字节的地址，搬运成本与块大小无关。规矩是所有权随消息转移——生产者从池取出、发送后不再触碰，消费者收到、处理完必须归还；两边任何一处越权，得到的是偶发数据撕裂且极难复现。中断侧只用带 \Tp{FromISR} 后缀的版本，例如 \Fn{xQueueSendToBackFromISR} 与 \Fn{xQueueReceiveFromISR}；阻塞版本在中断里调用会被内核断言拦下。

清单 \ref{lst:ch06-isr-queue} 展示中断入队、任务消费、归还缓冲这条最小正确路径。

\begin{lstlisting}[style=cstyle,caption={C：中断经队列唤醒消费者，并用空闲块队列传递所有权},label={lst:ch06-isr-queue}]
#include "FreeRTOS.h"
#include "task.h"
#include "queue.h"

#define Q_LEN      32U              /* 覆盖最坏突发：到达率乘以最长不可运行时长 */
#define ITEM_SZ    sizeof(uint32_t *)

static QueueHandle_t xSampleQ;
static QueueHandle_t xFreeQ;                    /* 空闲块队列，实现零拷贝 */
static uint32_t      xPoolStorage[Q_LEN];
static uint32_t      gDroppedCount, gPoolEmptyCount;

void Sample_Isr(void)
{
    BaseType_t xWoken = pdFALSE;
    uint32_t  *pvItem = NULL;
    uint32_t   raw;

    raw = (uint32_t)READ_SAMPLE();      /* 读外设数据寄存器 */
    (void)CLEAR_SAMPLE_PENDING();       /* 先清挂起，避免清晚了多进一次中断 */

    /* 中断里不能阻塞：拿不到空闲块就丢弃并计数 */
    if (xQueueReceiveFromISR(xFreeQ, &pvItem, &xWoken) == pdTRUE) {
        *pvItem = raw;
        if (xQueueSendToBackFromISR(xSampleQ, &pvItem, &xWoken) != pdTRUE) {
            (void)xQueueSendToBackFromISR(xFreeQ, &pvItem, &xWoken);
            gDroppedCount++;
        }
    } else {
        gPoolEmptyCount++;
    }

    portYIELD_FROM_ISR(xWoken);         /* 唤醒了更高优先级任务就在此处切换 */
}

static void Consumer_Task(void *pvParams)
{
    uint32_t *pvItem;
    (void)pvParams;
    for (;;) {
        if (xQueueReceive(xSampleQ, &pvItem, portMAX_DELAY) == pdTRUE) {
            Process(*pvItem);
            (void)xQueueSendToBack(xFreeQ, &pvItem, 0);  /* 归还，所有权结束 */
        }
    }
}

void Ipc_Init(void)
{
    xSampleQ = xQueueCreate(Q_LEN, ITEM_SZ);
    xFreeQ   = xQueueCreate(Q_LEN, ITEM_SZ);
    configASSERT((xSampleQ != NULL) && (xFreeQ != NULL));
    for (uint32_t i = 0U; i < Q_LEN; i++) {
        uint32_t *p = &xPoolStorage[i];
        (void)xQueueSendToBack(xFreeQ, &p, 0);
    }
    (void)xTaskCreate(Consumer_Task, "cons", 192U, NULL, 4U, NULL);
}
\end{lstlisting}

\subsection{信号量、互斥量，以及继承与天花板的取舍}

二值信号量用 \Fn{xSemaphoreCreateBinary} 创建后处于\textbf{空}状态，必须显式给一次才能让第一次等待通过；这一初始语义常造成第一帧丢失。计数信号量（需要 \Tp{configUSE\_{\Bk}COUNTING\_{\Bk}SEMAPHORES} 置 1）表达有 N 个可用资源，或者累计发生了多少个事件，中断里可用 \Fn{xSemaphoreGiveFromISR}。

互斥量与二值信号量的关键差别是\textbf{所有权}：只有持有者能释放，因此内核知道是谁在持锁，并能在高优先级任务来等锁时把低优先级持有者的优先级临时抬到等待者的水平，这就是内建的优先级继承（\Fn{xSemaphoreCreateMutex} 得到的对象即带此行为，递归版本用 \Fn{xSemaphoreCreateRecursiveMutex}，查持有者用 \Fn{xSemaphoreGetMutexHolder}）。继承解决了瞬时反转，但有两个残余问题：多级嵌套锁时抬升会沿链条传播，可能引发连锁抖动；而且它保证不了完全不推迟——高优先级任务仍可能被一个已经进入临界区的更低优先级任务推迟，推迟量等于那段临界区的长度。

优先级天花板协议把共享资源对应的优先级抬到\textbf{所有可能访问者中的最高值}，从而保证访问一旦开始就不会被打断，最坏推迟更容易算；代价是需要系统级维护资源到优先级的映射表，并且每次访问都无条件抬升。FreeRTOS 只提供继承形式的实现，没有开箱即用的天花板原语；要在这套内核上得到天花板的效果，工程上的通行做法是把共享资源改成\textbf{单一拥有者任务}：其他任务只能通过队列委托它做事，从根本上取消竞争。这也是本册的建议——能改设计就不要加协议。

\subsection{事件组、流缓冲与消息缓冲}

事件组（\Fn{xEventGroupCreate}、\Fn{xEventGroupSetBits}、\Fn{xEventGroupWaitBits}，中断侧用带 \Tp{FromISR} 后缀的置位与等待版本）表达若干条件位的组合等待，支持全部置位触发或任一处位置位触发，并能返回触发时的位值。要区分\textbf{置位语义}与\textbf{清零语义}：默认在等待返回时自动清除指定的位，如果选择不自动清零，下一次等待会立刻被历史位满足。要跨任务传某事发生过用置位，要请对方处理一件事用通知更合适。可用位宽受 \Tp{EventBits\_t} 与内核内部标志位的限制（通常低 24 位可自由使用，其余留给内核），别把它当作 32 位自由位图。

流缓冲（\Fn{xStreamBufferCreate} 一类接口）面向字节流，适合串口收发与日志通道；消息缓冲在其上保证一次写入即一次读出的边界，适合小的定长或变长报文。两者都提供中断安全的收发版本，且\textbf{自身不带互斥语义}：每个方向只应有一个写入者或用一个锁保护写入端，多生产者并发写会交错。需要多生产者入队、单消费者出队时，队列仍是首选。

\subsection{软件定时器}

\Tp{configUSE\_{\Bk}TIMERS} 置 1 后内核创建一条定时器任务，其优先级与栈分别由 \Tp{config{\Bk}TIMER\_{\Bk}TASK\_{\Bk}PRIORITY}、\Tp{config{\Bk}TIMER\_{\Bk}TASK\_{\Bk}STACK\_{\Bk}DEPTH} 决定，队列长度由 \Tp{config{\Bk}TIMER\_{\Bk}QUEUE\_{\Bk}LENGTH} 决定，而所有回调都在这条定时器任务里\textbf{串行}执行。推论有三条：回调不得阻塞或长耗时，否则会推迟后面所有定时器；到期时刻相近的两个回调必然互相推迟；定时器精度受时基与定时器任务被抢占的共同影响。有硬截止期的周期不要交给软件定时器，只在低重要度的超时、重试、喂狗上使用。\Fn{xTimerStartFromISR} 一类接口可以在中断里装载超时，但到期处理仍在同一条中优先级任务里排队。

表 \ref{tab:ch06-ipc-choice} 汇总选型。

\begin{table}[htbp]
\centering
\caption{同步与 IPC 对象选型对照}
\label{tab:ch06-ipc-choice}
\begin{tabular}{>{\raggedright\arraybackslash}p{0.20\textwidth} >{\raggedright\arraybackslash}p{0.20\textwidth} >{\raggedright\arraybackslash}p{0.19\textwidth} >{\raggedright\arraybackslash}p{0.23\textwidth}}
\toprule
需求 & 首选对象 & 携带信息 & 实时注意事项 \\
\midrule
唤醒一个已知任务 & 任务通知 & 计数、位集或 32 位值 & 最快最省，只有一个接收者，不能被多任务等待 \\
中断到任务搬数据 & 队列（元素为指针） & 定长副本或块地址 & 长度按最坏突发算，入队只用 \Tp{FromISR} 版本 \\
某事发生了一次的握手 & 二值信号量 & 无 & 初始为空需先给一次，无所有权概念 \\
N 个等价资源 & 计数信号量 & 无 & 需配置宏开启；超过上限时给出非阻塞失败 \\
保护共享短临界区 & 互斥量 & 无 & 带所有权与优先级继承，不可在中断里使用 \\
多条件同时满足 & 事件组 & 位组合 & 注意等待返回时是否自动清零 \\
字节流 & 流缓冲 & 任意长度分片 & 单写入者，不保证消息边界 \\
小报文 & 消息缓冲或队列 & 一条完整报文 & 消息缓冲无互斥语义，需自行约束写入端 \\
跨模块串行化访问 & 委托任务加队列 & 请求与应答 & 比天花板协议更可预测，代价是多一次往返 \\
\bottomrule
\end{tabular}
\end{table}

\subsection{配置宏段落}

清单 \ref{lst:ch06-config-macros} 列出与实时性直接相关的配置段落，数值只是示例量级，需按产品核对。

\begin{lstlisting}[style=cstyle,caption={C：FreeRTOSConfig.h 中与实时性直接相关的宏段落（数值为示例量级）},label={lst:ch06-config-macros}]
/* 时基与调度 */
#define configCPU_CLOCK_HZ                    480000000UL
#define configTICK_RATE_HZ                    1000      /* 1 毫秒 tick */
#define configMAX_PRIORITIES                  12
#define configUSE_PREEMPTION                  1
#define configUSE_TIME_SLICING                0         /* 关键档不靠轮转 */
#define configMINIMAL_STACK_SIZE              128       /* 单位：字 */
#define configMAX_SYSCALL_INTERRUPT_PRIORITY  5         /* 见第 5 节 */
#define configKERNEL_INTERRUPT_PRIORITY       15
#define configUSE_16_BIT_TICKS                0

/* 内存与检测 */
#define configTOTAL_HEAP_SIZE                 (32 * 1024) /* 只服务动态对象 */
#define configUSE_MALLOC_FAILED_HOOK          1
#define configCHECK_FOR_STACK_OVERFLOW        2

/* 原语开关：不用的一律关掉，省闪存也省状态位 */
#define configUSE_MUTEXES                     1
#define configUSE_RECURSIVE_MUTEXES           0
#define configUSE_COUNTING_SEMAPHORES         1
#define configUSE_QUEUE_SETS                  0
#define configUSE_EVENT_GROUPS                1
#define configUSE_STREAM_BUFFERS              1
#define configSUPPORT_DYNAMIC_ALLOCATION      1
#define configSUPPORT_STATIC_ALLOCATION       1

/* 定时器任务 */
#define configUSE_TIMERS                      1
#define configTIMER_TASK_PRIORITY             2
#define configTIMER_QUEUE_LENGTH              16
#define configTIMER_TASK_STACK_DEPTH          256

/* 低功耗与统计 */
#define configUSE_TICKLESS_IDLE               1
#define configEXPECTED_IDLE_TIME_BEFORE_SLEEP 4
#define configUSE_IDLE_HOOK                   1
#define configUSE_TICK_HOOK                   0
#define configUSE_TRACE_FACILITY              1
#define configGENERATE_RUN_TIME_STATS         1

/* 断言：不要直接停机，先留下现场 */
#define configASSERT(x)   do { if ((x) == 0) { AssertCapture(__FILE__, __LINE__); } } while (0)
\end{lstlisting}

\section{中断与内核协作}

\subsection{哪些函数能在中断里调用}

规则简单粗暴：只有带 \Tp{FromISR} 后缀的内核接口可用于中断，其余一律不可。原因是中断上下文既不能阻塞，也不能安全地触发切换；而阻塞版本会假设自己能被调度出去，并操作不应重入的等待列表。除内核函数之外，C 运行时里也有一大片在中断里不安全：\Fn{printf}（尤其在带锁的新lib 上）、\Fn{malloc}、\Fn{strtok}，以及任何使用线程局部存储或内部游标的函数。诊断打印应从 ISR 拷进缓冲，再由低优先级任务输出。

\subsection{临界区的屏蔽代价}

\Fn{task{\Bk}ENTER\_{\Bk}CRITICAL} 在 Cortex-M 端口上把当前屏蔽级别提到 \Tp{config{\Bk}MAX\_{\Bk}SYSCALL\_{\Bk}INTERRUPT\_{\Bk}PRIORITY} 对应的门槛，也就是关掉所有会调用内核 API 的中断，直到 \Fn{task{\Bk}EXIT\_{\Bk}CRITICAL} 才恢复。它只保证不被任务抢占、也不被这些中断打断，\textbf{不}屏蔽优先级更高（数值更小）的中断。代价必须量化：临界区长度就是这些中断的额外延迟。工程上把它压到几十个指令周期量级，并在注释里写下上限；一旦临界区里出现循环、跨内存的批量访问或函数调用，就应改用互斥量、或改成无锁结构。若确实需要完全关掉可屏蔽中断，端口另有底层接口，但那只用于极少数必须原子的硬件序列，不能当作常规锁。

\subsection{在最高优先级任务里自阻塞}

一个高优先级任务如果等待的事件只能由更低优先级的任务产生，那么它一阻塞就把 CPU 交给别人，而那个别人又因为某种共享资源被它占用而无法推进，形成等待环；同样地，如果它等待的事件本应由中断给出而中断被内核 API 的优先级门槛限制住，它可能永远等不到。表现是系统看似活着、周期完全停摆。避免办法有两条：所有会被高优先级任务等待的对象，其释放方必须能在中断上下文完成（用带 \Tp{FromISR} 后缀的接口）；或者把无限等待改成带超时的等待，并在超时分支进入定义好的安全状态。内核会在若干处用断言拦住非法组合，但拦住不等于提示你把设计做对。

\subsection{唤醒更高优先级任务与切换点}

\Tp{FromISR} 接口用最后一个出参告诉调用者是否有更高优先级任务被唤醒；收集到该标志后必须调用 \Fn{portYIELD\_FROM\_ISR}（或端口等价的挂起切换写法）在\textbf{本次}中断出口完成切换，而不是等下一次 tick。忘记这一步会让被唤醒任务多等一个中断返回路径、甚至多等一个 tick，把几微秒的设计变成毫秒级抖动。另一个细节是中断嵌套：如果一个 ISR 唤醒了任务又立刻被另一个 ISR 打断，切换会推迟到最内层退出，所以 ISR 里先清挂起、尽早返回的顺序直接影响延迟；清单 \ref{lst:ch06-isr-queue} 中先清挂起再做后续处理正是为此。

\section{低功耗与实时性的折中}

\subsection{tickless idle}

\Tp{configUSE{\Bk}\_TICKLESS{\Bk}\_IDLE} 让空闲任务在预计睡眠时长超过 \Tp{configEXPECTED{\Bk}\_IDLE{\Bk}\_TIME{\Bk}\_BEFORE{\Bk}\_SLEEP} 时调用端口的睡眠实现，关掉 tick 中断、按需降低主频或进入停止模式，醒来后补记 tick。它与硬实时的冲突在于：睡眠期间的唤醒延迟等于电源与锁相环的恢复时间，而不是中断响应时间；主频降低之后，正在运行的临界代码最坏执行时间同步放大。因此有周期截止期的产品要么关闭 tickless，要么把睡眠门槛设在没有任何在途周期任务的窗口里，并让低功耗钩子（\Tp{configPRE\_SLEEP\_PROCESSING} 与 \Tp{configPOST\_SLEEP\_PROCESSING}）明确记录是否真的停了外设时钟与调节器。

\begin{note}[title={唤醒延迟预算的写法}]
把唤醒路径拆成可测项逐项给数：深度睡眠退出与电源稳定、内部调节器爬升（典型数十到数百微秒）、锁相环重锁（典型几十微秒量级）、闪存预取与缓存重新使能、外设重新配置、以及第一个 tick 到达前的调度。任何一项没有实测数就不要写进总预算，因为漏掉一项往往就是全部余量。最后把总预算和被唤醒任务的截止期一起核对：若唤醒后立刻要处理一块数据，缓冲深度必须覆盖唤醒延迟加最坏处理时间之和。
\end{note}

\subsection{停止调度器做长事务}

\Fn{vTaskSuspendAll} 与 \Fn{xTaskResumeAll} 只禁止上下文切换，\textbf{不}关中断。它适合必须原子但允许中断继续的一段任务代码：状态表批量更新、快照打印、若干队列的组合操作。好处是中断延迟不受影响；坏处是这段时间内所有任务的就绪都不被响应，若中间出现阻塞调用，内核断言会失败；而且比它低优先级的任务被完全饿死。规矩有三条：事务长度上限写进注释并实测；事务内不得调用任何可能阻塞的接口；不得用它替代互斥量保护共享数据，因为别的任务在临界区里仍在运行。

\subsection{外设时钟与恢复时间}

低功耗下最先被关掉的是外设时钟与 PLL，而恢复代价最高的通常是音频编解码、锁相环式时钟源，以及需要充电泵或基准稳定的模拟模块（毫秒级）。硬实时结论：等时流所依赖的时钟域不参与省电；或者把这些模块的唤醒提前量做成固定流程（唤醒、定时等待、再启动传输），绝不允许在周期任务里现场等待锁定。

\section{工程化：可移植层、调试与陷阱}

\subsection{可移植层与自定义断言}

端口相关的一切集中在 \Tp{portmacro.h}：tick 与毫秒的换算、临界区宏、上下文切换汇编、\Tp{FromISR} 的底层实现、睡眠与运行时间计数接口。改这个文件之前要确认哪些宏同时被配置层与内核源码引用，尤其是中断优先级相关的两处：内核 API 可用的优先级门槛，以及 tick 中断自身的优先级。把它们配反会出现两类故障：中断里调用内核函数被断言拦下，或者硬实时中断被内核临界区挡住。

\Tp{configASSERT} 是内核里唯一主动报告非法用法的地点，务必实现成保存现场（任务名、调用点、当前 tick、栈指针）后进入可观测失效，而不是静默忽略或直接复位；开发阶段让它停住并用调试器看栈，产品阶段让它记录之后进入安全状态。

\subsection{运行时统计}

\Fn{vTaskList} 输出每个任务的状态、优先级、剩余栈与编号；\Fn{vTaskGetRunTimeStats} 输出每个任务占用 CPU 的比例，需要 \Tp{configGENERATE\_{\Bk}RUN\_{\Bk}TIME\_{\Bk}STATS} 置 1，并由端口提供两个宏：一个配置一只比 tick 快得多的自由运行计数器（\Tp{portCONFIGURE\_{\Bk}TIMER\_{\Bk}FOR\_{\Bk}RUN\_{\Bk}TIME\_{\Bk}STATS}），一个读取计数值（\Tp{portGET\_{\Bk}RUN\_{\Bk}TIME\_{\Bk}COUNTER\_{\Bk}VALUE}）。经验上这只计数器比 tick 快两到三个数量级（几十千赫兹到 1 兆赫兹），否则分辨不出微秒级的差别。字符串形态的统计还需要 \Tp{configUSE\_{\Bk}TRACE\_{\Bk}FACILITY} 与 \Tp{configUSE\_{\Bk}STATS\_{\Bk}FORMATTING\_{\Bk}FUNCTIONS}，并且调用方提供足够大的字符缓冲区（按每任务几十字节估算）。若要结构化数据而不是字符串，用 \Fn{uxTaskGetSystemState} 取状态数组，便于外部记录。提醒一句：统计代码自己也要吃 CPU 与栈，采证窗口内的实时性会变差，不能把它当成长期开着的量测手段；尾延迟必须在关掉统计的构建里重测。

\begin{warning}[title={硬实时路径上的四类禁用动作}]
一是分配与释放：\Fn{pvPortMalloc}、任务删除引发的回收、以及库内隐式分配（字符串流、动态容器），它们都可能遍历不可界的结构。二是阻塞调用出现在中断上下文，或被约束的任务里。三是在周期任务体里写无界循环：解析外部输入的扫描循环、重试直到成功的等待、复杂度随长度平方增长的整理。四是长临界区与全中断屏蔽。这四类的共同点是编译期不报错、短期运行不显形，只在最坏组合出现的那一次让你丢帧或错过截止期。
\end{warning}

常见陷阱还可以再列具体几条，评审时逐条对照：优先级数值方向写反；把栈深度按字节填；用 \Fn{vTaskDelay} 做周期导致累积漂移；二值信号量忘记初始给一次；队列长度按平均到达率而不是最坏突发设定；传指针之后生产者又复用了那块内存；在 \Tp{FromISR} 调用里忘了传非空的唤醒出参；软件定时器回调里做耗时处理；用互斥量当中断到任务的同步（互斥量不能在中断里释放）；开启 tickless 却没给唤醒延迟留预算；长期开着统计函数；把附加组件当成硬实时依赖。

\subsection{附加组件在硬实时里的位置}

协议栈（TCP/IP）、文件系统、命令行这类附加组件提供便利，但内部普遍存在动态分配、长循环、重试等待与内部串行任务。它们的定位是与硬实时共存的低优先级服务：单独一档优先级，低于控制面与数据面；跑在自己的缓冲上，不与关键路径共享堆；有周期截止期的功能（例如控制指令接收）通过一条独立的窄通路（队列或通知）与实时任务交互，让协议栈只充当数据来源。文件系统写日志尤其危险：一次块擦除可能是毫秒级且高度不等的等待。实时系统应把它放到最外层，上层用环形缓冲在满时丢弃，而不是把背压传回硬实时路径。

\begin{bestpractice}[title={本章的默认工程配方}]
任务与所有被硬实时使用的 IPC 对象一律静态创建；栈深度按水印实测后加倍；周期任务用 \Fn{vTaskDelayUntil}；数据面用指针队列加自建固定尺寸池，块所有权随消息转移；共享短临界区用互斥量，跨模块的资源访问改成委托任务；临界区长度写上限值并实测；tickless 只在没有在途周期任务的窗口启用；硬实时路径不接受任何分配、无界循环与长事务。这套配方并没有把内核变强，只是让每一处不确定性都能被指出、量出并写进预算。
\end{bestpractice}

% ============================================================
% ---- frag_ch07.tex ----
\chapter{Zephyr 与 RTOS 生态}
\label{ch:ch07-zephyr-ecosystem}

本章把 Zephyr 当成一个硬实时内核来读，而不是当成一个功能齐全的开源 RTOS。
顺序是：构建三件套如何决定最终镜像、线程与调度的数值边界、中断与延迟工作的分工、
同步与 IPC 原语的适用面、设备模型在实时路径上的代价、用户模式与 MPU 的保护粒度，
最后是与 FreeRTOS、ThreadX、RTEMS、NuttX、mbedOS、RT-Thread 的选型判断和可移植层写法。
结论只有一句：在 Zephyr 里可维护性与可预测性通常不冲突，
但两者都建立在静态分配、明确的优先级分区与极短的 ISR 这三个前提之上。

\section{工程结构与 west 构建流程}
\label{sec:ch07-build}

\DomRtos{}工程入口是 \DomHost{} \Tp{west}（一个用 \PY{} 写的命令行工具），
它把工作区清单、模块克隆、构建、烧写与调试串成四条命令：
\Tp{west init} 用应用自带的 \Tp{west.yml} 清单建立工作区骨架，
\Tp{west update} 按清单把内核、hal、\Lib{mbedtls} 等模块克隆到位，
\Tp{west build} 生成 CMake 缓存并调用底层构建系统，\Tp{west flash} 或 \Tp{west debug}
把镜像推到板上。清单固定每个模块的修订号，因此同一工作区可逐字节复现构建结果，
这是认证流程里最容易被审计的一条。清单 \ref{lst:ch07-west-flow} 里的
\Opt{build-dir} 与 \Opt{verbose} 在排错时最常用。

\begin{lstlisting}[style=shellstyle,caption={Shell：west 工作区、构建与调试入口},label={lst:ch07-west-flow}]
west init -l ./my-application      # 以应用自带的 west.yml 作为工作区清单
cd my-workspace && west update     # 按清单克隆 zephyr / hal_stm32 / lib 等模块
west boards | grep -i stm32        # 列出当前工作区可用的板级定义

west build -b nucleo_h743zi2 ./app
west build -b nucleo_h743zi2 ./app --build-dir build-dbg   # 独立目录，便于对比
west build -t menuconfig                                  # 交互式改 Kconfig
west build -t guiconfig                                   # 带依赖树搜索的图形界面
west build -t menuconfig -- --merge-configs               # 只看合并后的最终配置
west build -b qemu_cortex_m3 ./app -p auto                # 主机侧先验证逻辑

west flash --runner openocd        # 也支持 jlink、pyocd、qemu 等 runner
west debugserver --runner pyocd    # 先起 GDB server，再接 IDE
west debug                         # 烧写加起 GDB 的一步到位写法
\end{lstlisting}

\subsection{CMake、Kconfig 与设备树三件套的关系}

三件套各回答一个正交问题，任何一个改完都必须让另两个看得到。
\Lang{CMake} 决定编译哪些源文件、链接哪些库，核心是名为 \Tp{app} 的链接目标，
应用用 \Tp{target\_sources} 把自己的文件挂上去。
Kconfig 决定内核与子系统里存在哪些功能，每个 \Tp{CONFIG\_} 符号要么参与源文件取舍，
要么变成宏写进自动生成的 \Tp{autoconf.h}；它是编译期裁剪，不是运行期开关。
设备树决定这颗芯片上有什么、怎么接，描述寄存器地址、中断号、时钟、引脚复用与状态，
编译期展开成 \Tp{devicetree\_generated.h} 里的一组宏。

耦合点在于：设备树节点的 \Tp{status} 属性与 \Tp{DT\_INST\_OKAY} 一类判定宏决定
驱动是否实例化，Kconfig 决定驱动实现是否被编译。
配置开了、节点没使能，结果是编进了代码但没有设备对象；
节点使能了、配置关了，结果是链接期报未定义符号或干脆静默地少一个 \Tp{struct device}。
硬实时项目应在持续集成里同时归档合并后的配置快照与展开后的设备树，逐版对比。

\subsection{boards 与 soc 的分层，以及配置覆盖顺序}

板级定义在 \Tp{boards/} 下按厂商分目录，一个板目录通常有四件东西：
描述厂商、SoC、硬件能力与依赖的 YAML 元数据文件、同名 \Tp{.dts} 设备树、
同名 \Tp{Kconfig.board}（声明板符号并 select SoC）与同名 \Tp{\_defconfig}。
SoC 层在 \Tp{soc/} 下，描述核数与型号、中断控制器、内存布局、外设基址与家族默认配置，
\Tp{dts/} 里另有对应的 \Tp{\_pins.dtsi}。同一颗 SoC 的不同板子只改引脚引出到哪台连接器；
换芯片才动 SoC 层。自定义板的正确做法是复制最接近的参考板再覆盖差异。

配置合并顺序是越靠后被读入的越优先。Kconfig 侧大致是架构默认，
SoC 的 \Tp{\_defconfig}，板的 \Tp{\_defconfig}，应用的 \Tp{prj.conf}
（即 \Tp{APPLICATION\_CONFIG\_FILE}），再依次是 \Tp{EXTRA\_CONF\_FILE} 列出的片段，
最后是命令行上的 \Tp{-DCONFIG\_XXX=y} 与环境变量形式。设备树侧同理：
SoC skeleton 与板级 \Tp{.dts} 先展开，再叠加应用目录里 \Tp{boards/} 下与板名同名的
\Tp{.overlay} 片段、\Tp{EXTRA\_DTC\_OVERLAY\_FILE} 与环境变量指定的覆盖文件。
后写入的属性覆盖先写入的，但节点存在性只由 \Tp{status} 决定。
这条规则是 \ref{sec:ch07-pitfalls} 里第一类陷阱的根源。

\subsection{构建产物与调试服务器}

关键产物：\Tp{zephyr.elf}（带符号表，给 GDB 与静态分析）、\Tp{zephyr.bin}（给产线烧写）、
\Tp{zephyr.hex} 与 SREC（可选目标）、\Tp{zephyr.map}（段布局，排查栈溢出与区域放错的
唯一入口）、展开后的 \Tp{zephyr.dts} 与生成的 \Tp{autoconf.h}。
\Tp{west build} 结束会按内存区域打印使用率，把它归档进版本记录；
任何一次 RAM 突然多出两千字节的变更都值得追问。

调试两条主流路线：OpenOCD（配置形如 \Tp{interface/{\Bk}stlink.cfg} 加
\Tp{target/{\Bk}stm32h7x.cfg}，并在其 RTOS 配置里选 Zephyr，
这样 \Tp{info threads} 列出的是 Zephyr 线程而不是裸任务），
以及 pyOCD（面向 CMSIS-DAP，自带 target pack，起 \Tp{gdbserver} 后用
\Tp{west debugserver} 桥接）。真正决定调试体验的是镜像有没有关优化、有没有留符号，
以及 \Tp{CONFIG\_{\Bk}SIZE\_{\Bk}HINTS}、\Tp{CONFIG\_{\Bk}THREAD\_{\Bk}STACK\_{\Bk}INFO} 这类信息型开关。

\begin{guideline}[title={可复现的实时构建基线}]
每次提交三件事：固定的工作区清单修订号、合并后的最终配置快照、
以及 \Tp{zephyr.map} 与区域使用率文本。换板或换编译器版本时先比对这三份，
再讨论延迟变化；否则无法区分代码变慢与配置漂移。
\end{guideline}

\section{线程模型：静态栈、优先级与时间片}
\label{sec:ch07-threads}

Zephyr 的调度器是固定优先级、位图就绪队列。优先级是一个整数，数值越小优先级越高，
并分成两段：负的 coop（协作）段与非负的 preemptive（可抢占）段，
段宽由 \Tp{CONFIG\_NUM\_COOP\_PRIORITIES} 与 \Tp{CONFIG\_NUM\_PREEMPT\_PRIORITIES} 决定。
两段各自连续，\Tp{K\_PRIO\_COOP} 与 \Tp{K\_PRIO\_PREEMPT} 负责把段内序号换算成
内核接受的数值；不要手算偏移。行为差异见表 \ref{tab:ch07-coop-preempt}。

\begin{table}[htbp]
\centering
\caption{coop 与 preemptive 线程的行为差异}\label{tab:ch07-coop-preempt}
\begin{tabular}{>{\raggedright\arraybackslash}p{0.20\textwidth}
                >{\raggedright\arraybackslash}p{0.29\textwidth}
                >{\raggedright\arraybackslash}p{0.30\textwidth}}
\toprule
维度 & coop（协作）优先级 & preemptive（可抢占）优先级 \\
\midrule
数值段 & 负值，用 \Tp{K\_PRIO\_COOP} 换算 & 非负，用 \Tp{K\_PRIO\_PREEMPT} 换算 \\
时钟中断 & 不因 tick 被抢占，时间片对它无效 & 同优先级可被 tick 轮转 \\
抢占来源 & 更高优先级 coop 线程、中断返回时的重调度 & 任意更高或同优先级线程、阻塞原语唤醒 \\
典型用途 & 极硬约束、运行时间可静态证明的小循环 & 业务线程、驱动线程、系统工作队列 \\
风险 & 一个死循环锁死整层，低优先级全部饿死 & 同优先级轮转让延迟带抖动 \\
\bottomrule
\end{tabular}
\end{table}

\begin{warning}[title={静态栈估小的表现是直接跑飞}]
栈对象必须用 \Tp{K\_THREAD\_STACK\_DEFINE} 单独声明，它的类型与承载控制块的
\Tp{struct k\_thread} 不是一回事，两者不能互换。栈大小在编译期分配，
没开栈保护时溢出的表现是相邻对象被静默改写、偶发复位或行为异常；
开了 MPU 栈保护则是精确的栈溢出故障。压测方法：在低负载与最坏输入
（最大递归深度、最深格式化调用、全通道中断风暴）两组条件下取栈高水位，
再留不低于一半的余量，并把高水位检查做成量产前的例行项。
\end{warning}

面向内核与驱动线程有一组语义宏：\Tp{K\_{\Bk}THREAD\_{\Bk}PRIORITY\_{\Bk}HIGH}、
\Tp{K\_{\Bk}THREAD\_{\Bk}PRIORITY\_{\Bk}ABOVE\_{\Bk}NORMAL}、\Tp{K\_{\Bk}THREAD\_{\Bk}PRIORITY\_{\Bk}NORMAL}、
\Tp{K\_{\Bk}THREAD\_{\Bk}PRIORITY\_{\Bk}BELOW\_{\Bk}NORMAL}、\Tp{K\_{\Bk}THREAD\_{\Bk}PRIORITY\_{\Bk}LOW}、
\Tp{K\_{\Bk}THREAD\_{\Bk}PRIORITY\_{\Bk}IDLE}，以及两个按类型给的 \Tp{K\_{\Bk}PRIO\_{\Bk}EFFICIENT}
（吞吐型后台线程）与 \Tp{K\_{\Bk}PRIO\_{\Bk}TIME\_{\Bk}CRITICAL}（不能错过的硬实时线程）。
用它们而不是裸数字有两个好处：段宽变化时不必重排代码，
阅读者一眼看出线程属于哪一层。

\Tp{CONFIG\_TIME\_SLICING} 打开时同优先级就绪线程按 tick 轮转；
硬实时路径通常建议关掉，让每个关键线程独占优先级，
否则同优先级线程数会直接进入最坏延迟公式。时间基准来自
\Tp{CONFIG\_SYS\_CLOCK\_TICKS\_PER\_SEC}：1 毫秒 tick 意味着 \Fn{k\_sleep}
的实际唤醒时刻向上取整到毫秒边界。开启 \Tp{CONFIG\_TICKLESS\_KERNEL} 后
内核只在必要时编程硬件定时器，事件驱动系统收益明显，代价是时间边界变成 tick 边界对齐，
短周期任务会读到至少一个 tick 的量化误差；追求亚毫秒稳定周期时宁可保留高频 tick。

等待与放弃 CPU 的三种写法取舍：\Fn{k\_thread\_join} 等待目标线程结束，
是启动阶段与错误收敛路径的标准做法，但目标必须确实退出，否则调用者永远挂起；
\Fn{k\_sleep} 让出 CPU，精度受 tick 取整影响，适合毫秒级以上周期；
\Fn{k\_busy\_wait} 在指定微秒数里真正空转，耗 CPU、推高功耗、
把同优先级与更低优先级线程全挡在门外，还推迟 tick 处理与延迟工作。
\Fn{k\_thread\_yield} 只在同优先级之间轮转，对更高优先级无效。
需要亚毫秒定时时，正确工具是硬件定时器加中断，或只轮寄存器位的计数循环。

线程创建时的选项参数携带行为位（用户模式、essential 等），
浮点行为由内核决定：\Tp{CONFIG\_FPU} 决定是否保存浮点单元，
\Tp{CONFIG\_FPU\_SHARING} 决定上下文切换时是否搬运浮点寄存器；
关掉共享而线程执行浮点指令时内核会拒绝执行。
每多一个可能用浮点的线程，切换就多一组寄存器搬运，
在 100 微秒级周期预算里这不是可忽略项。清单 \ref{lst:ch07-thread}
给出线程加信号量加工作队列的最小骨架。

\begin{lstlisting}[style=cstyle,caption={C：线程、信号量与专用工作队列的最小实时骨架},label={lst:ch07-thread}]
#include <zephyr/kernel.h>
#include <zephyr/logging/log.h>

LOG_MODULE_REGISTER(app, CONFIG_LOG_DEFAULT_LEVEL);

#define CTRL_PRIO   K_PRIO_PREEMPT(2)      /* 抢占段内序号 2，越小越高 */

K_THREAD_STACK_DEFINE(ctrl_stack_obj, 1024);
K_THREAD_STACK_DEFINE(ctrl_wq_stack, 1024);

struct k_sem frame_sem;                    /* 静态对象：运行期不再 malloc */
struct k_work sample_work;
struct k_work_q ctrl_wq;

static void sample_work_handler(struct k_work *w)
{
	ARG_UNUSED(w);
	/* 队列线程上下文：可以阻塞，但仍然禁止动态分配 */
}

static void ctrl_entry(void *p1, void *p2, void *p3)
{
	ARG_UNUSED(p1); ARG_UNUSED(p2); ARG_UNUSED(p3);

	while (1) {
		/* 只等，不做算术；这里的超时值就是最坏等待预算 */
		if (k_sem_take(&frame_sem, K_MSEC(5)) != 0) {
			LOG_WRN("frame timeout, enter safe state");
			continue;
		}
		k_work_submit_to_queue(&ctrl_wq, &sample_work);  /* 重活推给队列 */
	}
}

static void trigger_isr(const void *dev)
{
	ARG_UNUSED(dev);
	k_sem_give(&frame_sem);              /* ISR 里只 give，绝不 take */
}

void main(void)
{
	static struct k_thread ctrl_thread;

	k_sem_init(&frame_sem, 0, 1);
	k_work_init(&sample_work, sample_work_handler);

	/* 关键路径单开工作队列：默认系统队列容易被后台任务堆积拖累 */
	k_work_queue_init(&ctrl_wq);
	k_work_queue_start(&ctrl_wq, ctrl_wq_stack,
			   K_THREAD_STACK_SIZEOF(ctrl_wq_stack),
			   K_PRIO_PREEMPT(5), NULL);

	k_thread_create(&ctrl_thread, ctrl_stack_obj,
			K_THREAD_STACK_SIZEOF(ctrl_stack_obj),
			ctrl_entry, NULL, NULL, NULL, CTRL_PRIO, 0, K_NO_WAIT);
	k_thread_name_set(&ctrl_thread, "ctrl");

	/* 中断号与优先级取自设备树，参数顺序见 include/zephyr/irq.h */
	IRQ_CONNECT(DT_NODELABEL(adc0), 6, trigger_isr,
		    DEVICE_DT_GET(DT_NODELABEL(adc0)), 0);
}
\end{lstlisting}

与实时相关的配置片段见清单 \ref{lst:ch07-prjconf}，它表达的是策略而不是数值最优。

\begin{lstlisting}[style=yamlstyle,caption={YAML：prj.conf 中与实时相关的关键项},label={lst:ch07-prjconf}]
# 调度与时间基准
CONFIG_NUM_COOP_PRIORITIES=8
CONFIG_NUM_PREEMPT_PRIORITIES=16
CONFIG_TIME_SLICING=n              # 关键线程各自独占优先级，避免同优先级轮转
CONFIG_SYS_CLOCK_TICKS_PER_SEC=1000
CONFIG_TICKLESS_KERNEL=n           # 亚毫秒周期不接受 tick 量化误差
CONFIG_ISR_STACK_SIZE=2048
CONFIG_THREAD_STACK_INFO=y         # 栈高水位统计，量产前复核

# 工作队列：默认系统队列太容易堆积，关键路径单开队列
CONFIG_SYSTEM_WORKQUEUE_STACK_SIZE=2048
CONFIG_SYSTEM_WORKQUEUE_PRIORITY=-1

# 日志与调试：实时路径上宁可丢日志也不要延迟
CONFIG_LOG=y
CONFIG_LOG_MODE_DEFERRED=y         # 不用 immediate，避免打印阻塞业务
CONFIG_BOOT_BANNER=n

# 内存：全静态，通用堆只留给启动阶段
CONFIG_HEAP_MEM_POOL_SIZE=2048

# 检查器
CONFIG_WATCHDOG=y
CONFIG_THREAD_WATCHDOG=y
CONFIG_RESET_ON_FATAL_ERROR=n
\end{lstlisting}

\begin{warning}[title={默认系统工作队列是延迟链上最容易被忽略的一环}]
\Tp{CONFIG\_SYSTEM\_WORKQUEUE\_PRIORITY} 的默认值很低，用途正是不要抢实时工作的 CPU。
反过来说：业务线程比它高时，队列里的项可能长时间得不到执行，
表现为中断一切正常但数据就是不上来，而且没有任何错误码。
对策按优先级排列：关键项提交到专用队列（清单 \ref{lst:ch07-thread} 的做法）；
把系统队列抬到业务线程之下但仍在 idle 之上；
给每个队列加运行计数与滞留时长统计，作为持续集成的可观测项。
\end{warning}

\subsection{构建脚本的写法}

\Lang{CMake} 侧保持克制：只在必要时加源目录与条件编译，
并把哪些文件参与实时路径在脚本里写清楚，便于按目录做最坏执行时间分析与代码审查。

\begin{lstlisting}[style=cmakestyle,caption={CMake：应用构建脚本片段},label={lst:ch07-cmake}]
cmake_minimum_required(VERSION 3.20.0)

find_package(Zephyr REQUIRED HINTS $ENV{ZEPHYR_BASE})
project(control_loop)

# 实时路径单独成目录，便于按目录审查与计时
target_sources(app PRIVATE
  src/main.c
  src/rt/control_loop.c
  src/rt/pid_fixed.c          # 定点实现，避免浮点上下文成本
)

target_include_directories(app PRIVATE include)

# 追加配置片段：实时选项与调试选项分开，避免误改
list(APPEND EXTRA_CONF_FILE "prj-rt.conf" "overlay-noise.conf")

# 自定义板或驱动由 BOARD_ROOT 引入，不改上游树
list(APPEND BOARD_ROOT ${CMAKE_CURRENT_SOURCE_DIR}/boards)

if(CONFIG_POSIX_API)
  target_sources(app PRIVATE src/posix_compat/pthread_shim.c)
endif()
\end{lstlisting}

\section{中断、工作队列与延迟工作}
\label{sec:ch07-irq-work}

\DomBare{}ISR 的唯一职责是把事实记录下来并唤醒别人。允许的操作只有：
读写寄存器、置标志位、\Fn{k\_sem\_give}、\Fn{k\_msgq\_put}（配 \Tp{K\_NO\_WAIT}）、
\Fn{k\_fifo\_put}、提交工作项，以及明确标注 ISR 安全的原子接口。
禁止的操作：任何带超时或 \Tp{K\_FOREVER} 的等待
（\Fn{k\_sem\_take}、\Fn{k\_mutex\_lock}、\Fn{k\_msgq\_get}）、任何内存分配、
\Fn{k\_sleep} 与 \Fn{k\_busy\_wait}、以及 immediate 模式下的打印。
在 ISR 里调用阻塞接口不一定立刻死机，它可能表现为
断言被编译掉之后系统卡在中断上下文，比死机更糟。

工作项有三层结构，选错代价不同。
\Tp{struct k\_{\Bk}work} 是一次性工作项，按提交顺序执行，重复提交期间只执行一次，
适合事件已记录、动作可推后的场合。
\Tp{struct k\_{\Bk}work\_{\Bk}q} 是执行容器，绑定一个线程、一个栈、一个优先级；
专用队列的意义在隔离：一条 I2C 传感器读数的慢动作不应该拖住控制闭环的工作项。
\Tp{struct k\_{\Bk}work\_{\Bk}delayable}（早期写法是 \Tp{K\_{\Bk}DELAYED\_{\Bk}WORK\_{\Bk}DEFINE}）
用 \Fn{k\_{\Bk}work\_{\Bk}schedule} 按相对时间投递、用 \Fn{k\_{\Bk}work\_{\Bk}reschedule} 改期，
到期后仍在队列线程上下文执行，因此可以做重活。
\Fn{k\_{\Bk}work\_{\Bk}flush} 与 \Fn{k\_{\Bk}work\_{\Bk}cancel\_{\Bk}sync} 会阻塞调用者，只在线程上下文使用。

\Tp{k\_timer} 与可延迟工作项都能表达周期，区别在时钟源与执行上下文：
定时器回调在内核的定时器中断路径里运行，不能阻塞、不能做重活，
所以惯用组合是定时器 ISR 里只提交一个可延迟工作项。
tickless 内核下定时器到期同样按 tick 边界对齐。

\Fn{irq\_offload} 是测试专用工具：把一个函数以真正的中断上下文执行一次，
用来验证 ISR 代码路径不踩内核 API、不破坏栈。限制要说清楚：
不是所有架构都提供；开启用户模式或内存保护的构建里可能被排除；
它占用一次真实中断向量，因此只适合单元测试与上电自检，不是运行时机制。
把中断上下文演练纳入驱动开发流程，能在集成前抓住绝大多数中断里调阻塞接口的问题。

\section{同步原语、IPC 与内存分配}
\label{sec:ch07-sync-ipc}

同步原语按是否需要优先级继承与数据结构形状两个轴来选，见表 \ref{tab:ch07-sync-ipc}。

\begin{table}[htbp]
\centering
\caption{内核同步与 IPC 对象的形状与约束}\label{tab:ch07-sync-ipc}
\begin{tabular}{>{\raggedright\arraybackslash}p{0.20\textwidth}
                >{\raggedright\arraybackslash}p{0.25\textwidth}
                >{\raggedright\arraybackslash}p{0.14\textwidth}
                >{\raggedright\arraybackslash}p{0.21\textwidth}}
\toprule
对象 & 语义与数据形态 & ISR 侧安全 & 典型用途 \\
\midrule
\Tp{k\_sem}（\Fn{k\_sem\_take}） & 计数信号量，不携带数据，无优先级继承 & give 安全 & 事件到达通知 \\
\Tp{k\_mutex}（\Fn{k\_mutex\_lock}） & 互斥，带优先级继承 & 不适用 & 共享资源互斥 \\
\Tp{k\_condvar} & 与互斥配合的条件等待 & 不适用 & 生产者消费者握手 \\
\Tp{k\_msgq} & 定长记录，入队拷贝，容量固定 & 需 \Tp{K\_NO\_WAIT} & 帧描述符传递 \\
\Tp{k\_fifo} / \Tp{k\_lifo} & 传递指针，不拷贝 & put 安全 & 缓冲块归还 \\
\Tp{k\_pipe} & 字节流，环形缓冲 & 需 \Tp{K\_NO\_WAIT} & 流式数据搬运 \\
\Tp{k\_events} & 位集合，可与或语义等待 & 需 \Tp{K\_NO\_WAIT} & 多源条件汇聚 \\
\Tp{k\_poll} & 同时等多个可轮询对象 & 部分支持 & 网关类多线程 \\
\bottomrule
\end{tabular}
\end{table}

两点强调。第一，保护共享数据要用 \Tp{k\_mutex}，不要把 \Tp{k\_sem} 当互斥：
二值信号量没有优先级继承，持有者在临界区被中优先级线程抢占时，
高优先级等待者会连带被挡住，这就是标准优先级反转。
\Fn{k\_mutex\_lock} 在遇到更高优先级等待者时把持有者临时抬到等待者的优先级，
反转窗口被压缩到临界区长度本身；因此临界区必须真的短，
优先级继承不是让长临界区变合法的许可证。第二，需要一次等多个来源时用 \Fn{k\_poll}
（要开 \Tp{CONFIG\_POLL}），但它比只等一件事多一层簿记开销，
关键周期路径上尽量只等一个对象。

内存方面，实时路径的正确姿势是把分配挪到启动阶段并让它有上界。
\Tp{K\_{\Bk}MEM\_{\Bk}SLAB\_{\Bk}DEFINE} 声明固定块大小、固定块数的对象回收池，
\Fn{k\_{\Bk}mem\_{\Bk}slab\_{\Bk}alloc} 在池空时可阻塞等待或立刻返回失败；
所有块编译期就存在，运行期只是链表操作，耗时与历史无关，
这是它优于通用堆的根本原因。早期的 \Tp{K\_{\Bk}MEM\_{\Bk}POOL\_{\Bk}DEFINE}
（内存池）在新版本中已被弃用并移除，迁移方向是 slab 加通用堆，
新代码不应再依赖它。通用堆容量由 \Tp{CONFIG\_{\Bk}HEAP\_{\Bk}MEM\_{\Bk}POOL\_{\Bk}SIZE} 决定，
\Fn{k\_malloc} 与 \Fn{k\_free} 的耗时与碎片历史有关、无法静态定价，
只允许出现在初始化阶段。数据面需要大块连续缓冲时，
用 slab 的块或静态数组加 \Tp{k\_fifo} 做缓冲块池与归还。

\section{设备模型与驱动的实时代价}
\label{sec:ch07-devices}

驱动注册通过 \Tp{DEVICE\_{\Bk}DT\_{\Bk}DEFINE} 与设备树节点绑定（配套
\Tp{DEVICE\_{\Bk}DT\_{\Bk}GET} 与 \Tp{DEVICE\_{\Bk}DT\_{\Bk}GET\_{\Bk}OR\_{\Bk}NULL}），
参数含初始化函数、电源管理动作、数据与配置指针、初始化级别、优先级和 API 表。
初始化级别按顺序是 \Tp{PRE\_{\Bk}KERNEL\_{\Bk}1}、\Tp{PRE\_{\Bk}KERNEL\_{\Bk}2}、\Tp{POST\_{\Bk}KERNEL\_{\Bk}1}、
\Tp{POST\_{\Bk}KERNEL\_{\Bk}2}、\Tp{APPLICATION}，级别内再用整数优先级排序：
SoC 与内核早期服务落在 PRE\_KERNEL，普通外设驱动通常落在 POST\_KERNEL
配 \Tp{CONFIG\_{\Bk}KERNEL\_{\Bk}INIT\_{\Bk}PRIORITY\_{\Bk}DEFAULT}，依赖一堆服务已就绪的应用层设备落在
APPLICATION。初始化阶段本身不是实时阶段，但如果中途打开中断而 ISR
仍依赖尚未就绪的对象，就会得到随机故障；把使能中断放在初始化函数最后一行是通用做法。

console 与日志是实时路径上的第一号隐形成本。串口 console 每字符一次寄存器写，
115200 波特率按每字符 10 位算是约 87 微秒一字节，一条 120 字节的调试打印就是约 10 毫秒
（与第 8 章内核日志的量级估算同源，瓶颈在同一条物理链路上）。
缓解是结构性的而不是少打一点：\Fn{printk} 与 \Tp{LOG} 走 deferred 模式，
打印落到日志线程；实时线程只把计数与状态写进共享内存；
量产固件把 \Tp{CONFIG\_LOG\_LEVEL} 降到 error 或彻底关；
调试期必须看即时输出时就把被调试的闭环降级，不要让即时打印与硬周期共存。
\Fn{k\_busy\_wait} 与串口写组合能在几毫秒内毁掉一个 500 微秒的周期。

子系统抽象按便利性换时间排列，实时预算要逐项记账。
传感器子系统把 \Fn{sensor\_{\Bk}sample\_{\Bk}fetch} 与 \Fn{sensor\_{\Bk}channel\_{\Bk}get} 分离：
一次 fetch 可能触发完整总线事务，get 只读缓存值，
所以周期任务应按数据手册的转换时间决定 fetch 节律，而不是每周期 fetch。
音频与 DMIC 子系统带环形缓冲与工作队列，端到端延迟由缓冲块数与采样率决定，
属于吞吐优化而非延迟优化。PWM 子系统的 \Fn{pwm\_{\Bk}set\_{\Bk}cycles} 是一次薄封装、开销可预期，
但跨 SoC 的分辨率与精度未必满足闭环要求。
CAN 有底层 \Fn{can\_send} 与面向信号的 \Fn{canbus\_{\Bk}send} 两层，后者带来编解码与额外拷贝。
当某层抽象成为延迟主因时，允许用 \Tp{DEVICE\_{\Bk}DT\_{\Bk}GET} 拿到设备后直接调它的 API 表函数
绕过子系统，但要写进架构决策记录并说明代价。

数据搬运与内存位置的组合能收回不少预算。DMA 的分工是：
驱动申请通道与描述符，完成回调在中断上下文执行，只做置标志、give 信号量、提交工作项；
缓冲按外设要求对齐并在启动时一次性分配；带数据缓存的平台要保证一致性，
传输前用 \Fn{sys\_cache\_data\_flush\_range} 清洗、传输后 invalidate，
更稳妥的做法是用 \Tp{CONFIG\_NOCACHE\_MEMORY} 与非缓存内存域，
把一致性责任一次性放到地址空间上。代码位置同样有代价：
开启 \Tp{CONFIG\_XIP} 时从外部闪存取指要经过带预取的控制器，
分支密集的小循环在等待里抖动；支持 TCM 的平台上把关键函数与数据放进紧耦合内存
（Zephyr 在相应架构的头文件里提供段属性，命名随架构而变）能显著降低延迟方差；
没有 TCM 时用链接段属性把热点函数放进内部 SRAM 的专门区域，
而 \Tp{zephyr.map} 是确认它真在里面而不是被优化掉的唯一凭据。

\section{用户模式、MPU 与内存域}
\label{sec:ch07-userspace}

\Tp{CONFIG\_USERSPACE} 打开后，内核对象登记在受保护的元数据里，线程带一个权限属性：
\Tp{K\_THREAD\_DEFINE} 的选项参数或 \Fn{k\_thread\_create} 的选项位给
\Tp{K\_USER\_MODE} 即得用户态线程，默认全部是特权线程。
用户线程访问内核对象前必须被显式授权：按对象授予，或授予到线程启动时刻
（相关接口以 \Tp{k\_thread\_access} 前缀命名，具体名字随版本演进，
在 \Tp{include/zephyr/kernel.h} 里搜比在文档里猜可靠）。
内存域（\Tp{K\_MEM\_DOMAIN\_DEFINE} 与 \Tp{k\_mem\_domain\_} 系列）把一组线程
划进同一地址视图，域内可见、域外不可见，用来隔离第三方代码或做故障不扩散。

保护的代价与风险要摆明：每次从用户态进内核多一次陷入与参数校验，
\Fn{k\_{\Bk}sem\_{\Bk}take} 这类高频调用在微秒预算下不再免费；
MPU 的 region 数量有限，分区越多越容易撞硬件上限，
栈区还常要求对齐到 2 的幂从而推高占用（约束见
\Tp{CONFIG\_{\Bk}MPU\_{\Bk}REQUIRES\_{\Bk}POWER\_{\Bk}OF\_{\Bk}TWO\_{\Bk}ALIGNMENT}）；
外设寄存器默认对内核可见，是否把它们从用户可访问区域剥离，
是安全设计与调试便利性之间的显式取舍
（以 \Tp{CONFIG\_{\Bk}MPU\_{\Bk}PERIPHERALS{\Bk}\_NONSECURE} 一类开关为代表，
不同版本命名有出入，以 \Tp{menuconfig} 搜索结果为准）。
实践上的分界线很清晰：驱动与实时内核路径留在特权侧，
应用逻辑、协议解析、算法库放用户侧，两者之间用消息队列与明确的所有权约定接口，
不要为了安全把每次中断唤醒都变成一次陷入。

\section{生态选型、可移植层与常见陷阱}
\label{sec:ch07-ecosystem}

\DomRtos{}RTOS 选型很少是纯技术判断，它同时是许可证、认证路径、
平台厂商支持承诺与人才供给的判断。表 \ref{tab:ch07-ecosystem}
按硬实时关心的维度并列几个生态。认证维度不在表里给结论：
Zephyr、FreeRTOS、ThreadX、NuttX 都有公开的认证发行版或第三方评估记录，
但认证几乎总是绑定特定版本加特定配置加特定工具链，
用默认配置或升级内核就会脱离证书范围，
任何项目的认证结论都必须以厂商当前发布的认证包与范围声明为准。

\begin{longtable}{>{\raggedright\arraybackslash}p{0.14\textwidth}
                  >{\raggedright\arraybackslash}p{0.14\textwidth}
                  >{\raggedright\arraybackslash}p{0.13\textwidth}
                  >{\raggedright\arraybackslash}p{0.13\textwidth}
                  >{\raggedright\arraybackslash}p{0.15\textwidth}
                  >{\raggedright\arraybackslash}p{0.14\textwidth}}
\caption{RTOS 生态选型矩阵}\label{tab:ch07-ecosystem}\\
\toprule
系统 & 内核形态与调度 & 许可 & 静态分配 & 网络与无线栈 & 典型场景 \\
\midrule
\endfirsthead
\toprule
系统 & 内核形态与调度 & 许可 & 静态分配 & 网络与无线栈 & 典型场景 \\
\midrule
\endhead
\bottomrule
\endfoot
\endlastfoot
Zephyr & 固定优先级加 coop 段，SMP 支持 & Apache-2.0 & 一等公民（编译期栈、slab） & 自带 TCP/IP 与 BLE，设备树统一 & 多芯片家族、需要现代工具链与认证文档的新项目 \\
FreeRTOS & 固定优先级抢占，单优先级表 & MIT & 编译开关启用静态创建 & 依赖第三方协议栈与厂商库 & 极小 RAM、成熟稳定、既有产品延续 \\
ThreadX（Eclipse ThreadX） & 固定优先级抢占，计数信号量原生继承 & MIT & 支持，块由应用提供 & 配套协议栈与文件系统成体系 & 大批量消费电子与有认证要求的产品 \\
RTEMS & 固定优先级加可选轮转，SMP 成熟 & 类 BSD 许可 & 两者皆可 & 完整 BSD 网络栈，POSIX 面最全 & 航空航天、科研、需要 POSIX 的高可靠系统 \\
NuttX & 抢占式加 POSIX 线程语义 & 自定义宽松许可 & 以内核堆为主 & 自带 TCP/IP 与应用框架 & 需要类 POSIX 应用层又不想上 Linux \\
mbedOS & 基于 RTX 的抢占内核 & Apache-2.0 & 部分 & 蓝牙与网络由框架代管 & Arm 生态既有产品；新项目需评估维护活跃度 \\
RT-Thread & 抢占式加优先级继承 & Apache-2.0 & 支持静态对象与多种堆 & 组件化（\Lib{lwIP}、BLE 包） & 国内模块供应链与快速原型 \\
\end{longtable}

把这张表读成三句话：既有资产（驱动、认证、人才）通常压过技术偏好；
能全静态与有 coop 优先级段是硬实时里最难后期改造的差异；
网络与蓝牙栈的成熟度决定你要不要自己维护一份协议栈。
落到场景：多厂商 SoC 且要用设备树统一硬件描述选 Zephyr；
RAM 以千字节计、单一厂商、已有认证证据链就继续用 FreeRTOS 或既有内核；
量产规模大、要成体系的文件系统与 USB 协议栈、希望认证包与厂商支持绑定选 ThreadX 系；
需要 POSIX 接口面与成熟 SMP 但不需要 Linux 生态选 RTEMS；
应用层想要 shell、动态加载与 POSIX 线程而资源不到 Linux 级别选 NuttX；
看重国内模块供应链与中文社区响应选 RT-Thread；
Arm 官方框架的存量产品继续用 mbedOS，新项目要先评估其维护活跃度。

\subsection{把 RTOS 抽象出来的可移植层}

可移植层的目的不是一套代码跑所有内核，而是让换内核成为一次可控的工程活动。

\begin{enumerate}
  \item 只抽象实际用到的语义，不做全量镜像。通常五类够：任务、互斥与信号量、
        定长消息、定时器或延迟工作、临界区。
  \item 头文件里不出现内核类型。句柄用不透明指针或固定大小结构体，内部再转换；
        对象内存由应用侧静态数组提供，这样 Zephyr 的栈对象与 FreeRTOS 的静态缓冲
        能用同一套声明方式。
  \item 优先级用分层枚举（critical、high、normal、background），每个内核一份映射表，
        映射表只放在一个适配文件里，换内核时改一个文件而不是改全部创建点。
  \item 时间单位统一到毫秒与微秒，转换函数负责 tick 向上取整的语义，
        并显式声明不保证亚毫秒精度，逼迫关键路径改用硬件定时器。
  \item 临界区只包裹确定的一小段代码，适配层里加断言检查嵌套深度与最大持有时间。
  \item 能力差异用编译期检查暴露：是否支持静态分配、优先级继承、tickless，
        逐项在适配层用条件编译与静态断言声明，禁止业务代码直接检测宏名。
  \item 测试双目标：同一套业务代码同时编译到主机（加桩与时间注入）与目标；
        主机跑逻辑与错误路径，目标只跑计时与压力，结论冲突时以目标为准。
\end{enumerate}

\begin{bestpractice}[title={实时项目里 Zephyr 的四条默认设置}]
全静态对象（线程栈、slab、消息队列容量都在编译期确定）、
每个硬约束线程独占一个优先级、关键路径使用专用工作队列、
延迟与栈高水位的测量数据随构建产物一起归档。
这四条做到之后，性能问题才是性能问题，不会与生命周期问题混在一起。
\end{bestpractice}

\subsection{常见陷阱清单}
\label{sec:ch07-pitfalls}

\begin{itemize}
  \item 只改 Kconfig 没改设备树（或反之）。症状是功能开了但探测不到设备，
        或节点在但符号未定义。检查展开后的 \Tp{zephyr.dts} 与合并后的配置快照
        （清单 \ref{lst:ch07-west-flow}），两边都进 CI 归档。
  \item 在 ISR 里调用会阻塞的 \Tp{k\_} 接口。轻则断言，
        重则关掉断言的发布版里静默死锁在中断上下文；上电自检加一次中断上下文演练。
  \item 默认系统工作队列优先级低于业务线程导致堆积，见清单 \ref{lst:ch07-thread}
        与 \Tp{CONFIG\_{\Bk}SYSTEM\_{\Bk}WORKQUEUE\_{\Bk}PRIORITY} 的说明。
  \item 静态栈按演示用例估值，最坏输入下直接跑飞。
  \item \Fn{k\_busy\_wait} 放进 coop 线程饿死所有人：coop 段不被 tick 抢占，
        忙等的每一微秒都是低优先级线程与延迟工作的损失。
  \item 用 \Tp{k\_sem} 做互斥引发优先级反转；改用 \Tp{k\_mutex} 并同时缩短临界区。
  \item 关键线程共用一个优先级还留着时间片轮转，最坏延迟变成同优先级线程数乘 tick，
        看起来像随机抖动。
  \item 用 immediate 日志模式或 \Fn{printk} 做周期诊断，把串口时间写进闭环。
\end{itemize}

\begin{note}[title={本册内的分界与外部参考}]
本章处理内核即全部的裸机加 RTOS 形态，寄存器级位域读写由本册的硬件访问分册承担。
Linux 侧的抢占模型、中断线程化与用户态实时循环由下一章处理；
两者唯一共享的部分是预算方法：先列出不可抢占段的来源，
再决定用什么机制把它请出关键路径。
\end{note}

% ============================================================
% ---- frag_ch08.tex ----
\chapter{Linux PREEMPT\_RT 与用户态实时}
\label{ch:ch08-linux-rt}

可抢占与延迟有界是两个不同的性质，本章大部分篇幅在处理它们之间的落差。
先给内核抢占模型谱系与各自的残留不可抢占段，再逐项列出 PREEMPT\_RT 改造了什么，
然后是隔离、调度策略、内存、频率与硬件干扰、I/O 约束，
最后落到可复制的用户态实时循环骨架与验收量级判据。
本册第 4 章的看门狗与故障收敛在这里被实时任务使用；
时间感知整形（ETF、TSN）的完整讨论在以太网与时间同步章，本章只交代套接字选项的使用位置。

\section{抢占模型谱系：从不可预占到延迟有界}
\label{sec:ch08-preempt-models}

Linux 的编译期抢占模型有四个层次，加上运行期动态切换，构成表
\ref{tab:ch08-preempt-models} 的谱系。对应符号是
\Tp{CONFIG\_{\Bk}PREEMPT\_{\Bk}NONE}、\Tp{CONFIG\_{\Bk}PREEMPT\_{\Bk}VOLUNTARY}、\Tp{CONFIG\_{\Bk}PREEMPT}、
\Tp{CONFIG\_{\Bk}PREEMPT\_{\Bk}RT}，以及 \Tp{CONFIG\_{\Bk}PREEMPT\_{\Bk}DYNAMIC}。后者让 none、voluntary、
full 三种模式在启动时用 \Tp{preempt=} 参数选择（注意它不接受 rt：实时抢占是编译期决定，
不能运行时打开），运行期可从 \Tp{/sys/kernel/debug/sched/} 下的接口读取当前模式。

\begin{longtable}{>{\raggedright\arraybackslash}p{0.15\textwidth}
                  >{\raggedright\arraybackslash}p{0.15\textwidth}
                  >{\raggedright\arraybackslash}p{0.16\textwidth}
                  >{\raggedright\arraybackslash}p{0.19\textwidth}
                  >{\raggedright\arraybackslash}p{0.17\textwidth}}
\caption{抢占模型谱系与残留不可抢占段}\label{tab:ch08-preempt-models}\\
\toprule
模型 & 内核配置 & 抢占点 & 主要残留不可抢占段 & 适用判断 \\
\midrule
\endfirsthead
\toprule
模型 & 内核配置 & 抢占点 & 主要残留不可抢占段 & 适用判断 \\
\midrule
\endhead
\bottomrule
\endfoot
\endlastfoot
不抢占 & \Tp{CONFIG\_{\Bk}PREEMPT\_{\Bk}NONE} & 只在返回用户态与显式内核入口检查 & 系统调用内的任意长循环、持锁段、软中断批量执行 & 吞吐优先，延迟要求毫秒级以上 \\
自愿抢占 & \Tp{CONFIG\_{\Bk}PREEMPT\_{\Bk}VOLUNTARY} & 在原有 \Fn{cond\_{\Bk}resched} 类点缀处增加让出点 & 同上，只是把明显可等待的长段缩短 & 桌面与服务器折中，实时性提升有限 \\
完全抢占 & \Tp{CONFIG\_{\Bk}PREEMPT} & 内核态除显式禁抢占与自旋锁段外均可被抢占 & 仍持自旋锁的段、关中断段、\Fn{preempt\_{\Bk}disable} 区间、console 打印 & 交互延迟改善明显，最坏值仍可达毫秒级 \\
实时抢占 & \Tp{CONFIG\_{\Bk}PREEMPT\_{\Bk}RT} & 绝大多数可睡眠锁变可抢占，硬中断线程化 & 真正的 \Fn{raw\_{\Bk}spin\_{\Bk}lock} 段、极短关中断段、每 CPU 线程、初始化与内存回收 & 目标微秒级最坏延迟时的唯一选项 \\
动态切换 & \Tp{CONFIG\_{\Bk}PREEMPT\_{\Bk}DYNAMIC} & 启动参数或调试接口选择三种非 RT 模式 & 由所选模式决定 & 同一份内核服务多种负载，验证收益差异 \\
\end{longtable}

完全抢占把平均响应压得很好看，但分布尾部由持有自旋锁与关中断的代码段决定，
而这些段的长度在源码里没有静态标注。PREEMPT\_RT 做的正是系统性地把这些段改成可睡眠、
可抢占的形式，从而让尾部进入可控范围。反过来，用了 PREEMPT\_RT 仍被长临界区拖累时，
问题不在模型而在具体实现：console 打印、内存回收、未隔离的中断、深度 C-state 唤醒。
这正是后面几节逐项排查的对象。

\section{PREEMPT\_RT 改了什么}
\label{sec:ch08-rt-changes}

\DomRT{}改造清单里对应用可见的有六项。

\begin{enumerate}
  \item 硬中断线程化。中断上下文只剩最小处理（读寄存器、清中断位、唤醒对应线程），
        实际处理逻辑搬进内核线程，线程名形如 \Tp{irq/\%d-name}（例如 \Tp{irq/44-eth0}），
        默认以 \Tp{SCHED\_FIFO} 优先级 50 运行，每 CPU 线程（定时器、函数调用中断等）
        优先级略高。于是中断处理变成可被更高优先级任务抢占的普通线程：
        延迟与特性能从调度里读出来，也能用 \Tp{chrt} 单独调价。
        全量线程化的启动开关是 \Tp{threadirqs}（非 RT 内核也可用它做实验）；
        反向关闭的参数命名随版本变化（\Tp{no\_threadirqs} 一类写法），
        以 \Tp{rt\_checklist} 的判定为准。
  \item 自旋锁语义分裂。普通 \Fn{spin\_lock} 在 RT 下变成基于 rtmutex 的可睡眠锁，
        因此不能再在硬中断上下文使用它；保留真正自旋语义的是 \Fn{raw\_spin\_lock}，
        只有确实需要关抢占加忙等的内核核心段用它。
        对开发者的直接含义：驱动里随手写的持锁段长度，现在直接进延迟预算。
  \item 优先级继承。rtmutex 携带优先级继承，锁持有者被更高优先级等待者临时抬升，
        反转窗口压缩到临界区长度。这与第 7 章 \Tp{k\_mutex} 的机制同构，
        但用户态要显式声明：\Fn{pthread\_mutexattr\_setprotocol} 配
        \Tp{PTHREAD\_MUTEX\_PRIO\_INHERIT}，默认的 normal 互斥不提供继承。
  \item 大核段拆分。串口与 \Fn{printk} 路径被挪进专门的打印内核线程，
        调度时钟、RCU 回调与软中断的执行点被拆细；软中断主要交给每 CPU 的
        \Tp{ksoftirqd/N} 线程处理，从而纳入调度器优先级管辖。这是收益最大
        也最少被理解的一项：它把不确定的长段换成可被抢占的短段加一个可定价的线程。
  \item \Fn{local\_lock} 与迁移控制。原本用关中断或禁抢占保护的每 CPU 数据改用
        \Fn{local\_lock} 一类抽象，语义随内核是否支持抢占而变；
        \Fn{migrate\_disable} 提供比禁抢占更弱的保证（线程不被搬到别的 CPU，但可被抢占）。
        用户态对应物是亲和性设置加显式绑核策略。
  \item 高精度定时器与无 tick。整条实时路径依赖 \Tp{CONFIG\_HIGH\_RES\_TIMERS}
        提供微秒级绝对到期，依赖 \Tp{CONFIG\_NO\_HZ\_FULL} 让隔离 CPU 不跑周期 tick。
        两者缺失时，无论是否 PREEMPT\_RT，最坏延迟都会出现一个 tick 量级的台阶。
\end{enumerate}

\begin{note}[title={全局实时配额会把高优先级任务掐掉}]
\Tp{/proc/sys/kernel/sched\_rt\_runtime\_us} 默认 950000，配合
\Tp{sched\_rt\_period\_us} 默认 1000000，意思是所有 RT 类任务合计每 1 秒最多用
0.95 秒 CPU，留 5\% 给非实时任务，防止一个失控的 \Tp{SCHED\_FIFO} 循环把系统饿死。
单个周期任务不超过这个比例时没有问题；一旦多个 RT 线程加上中断线程合计超过配额，
就会看到周期性 throttle，表现为固定节律的毛刺而不是随机毛刺。
调试阶段可以把它设成 \Tp{-1} 临时关闭，但量产系统应保留配额，
改为给实时组划分 cgroup 预算或减少 RT 线程数。
\end{note}

\section{CPU 隔离、中断归属与验证}
\label{sec:ch08-isolation}

隔离的目标是把实时核变成只有你自己的线程与指定中断的执行环境。手段分四层，
缺任何一层都会在验收数据里留下可辨认的毛刺。

\begin{itemize}
  \item 调度隔离：\Tp{isolcpus=4--{\Bk}7} 把这些 CPU 从通用调度域里摘出去，
        可选带 \Tp{no\_{\Bk}load\_{\Bk}balance}、\Tp{domain} 等子选项；
        管理层 IRQ（多队列网卡的 MSI-X 向量）由 \Tp{managed\_{\Bk}irqs} 参数控制归属。
        调度域层级可在 \Tp{/sys/kernel/debug/sched/domains/} 下查看。
        新写法倾向只用 cpuset cgroup 划集合，但两者语义不等价：
        \Tp{isolcpus} 影响调度域构建，cgroup 只限制可选 CPU。
  \item 时间基准隔离：\Tp{nohz\_{\Bk}full=4--{\Bk}7} 让隔离核在不必要时不再收 tick，
        条件是该核上只有一个可运行任务且系统调用路径干净；
        \Tp{rcu\_nocbs=4-7} 把 RCU 回调卸载到别的 CPU（较新内核里已由
        \Tp{nohz\_full} 隐含），否则隔离核会周期性执行回调。
        每 CPU 的内核工作线程还需逐个迁走或降低优先级。
  \item 中断隔离：\Tp{irqaffinity=0--{\Bk}3} 把所有中断的默认目标赶到家核；
        已分配的中断号逐个写 \Tp{/proc/irq/N/smp\_{\Bk}affinity\_{\Bk}list} 精调；
        验证看 \Tp{/proc/interrupts} 与 \Tp{/proc/irq/default\_{\Bk}smp\_{\Bk}affinity}。
        这一步之后必须关掉 \Tp{irqbalance}，否则它会在几分钟内把你手写的亲和性改回去。
  \item 任务归属：cgroup v2 的 \Tp{cpuset.cpus} 与 \Tp{cpuset.cpus.exclusive}
        声明独占集合；systemd 单元里用 \Tp{CPUAffinity=} 与 \Tp{AllowedCPUs=}
        （后者只在统一 cgroup 层级生效，通常写在 slice 上）划分服务可运行的 CPU。
\end{itemize}

\begin{warning}[title={只绑核没关 irqbalance：最常见的无效加固}]
把实时线程绑到 CPU 6、却留着 \Tp{irqbalance} 运行，结果是网卡中断被随机搬到 CPU 6，
实时线程的唤醒延迟突然多出几十到几百微秒，而 \Tp{cyclictest} 的曲线看起来只是偶发毛刺。
先确认服务已停再做任何延迟测量：\Tp{systemctl stop} 加 \Tp{disable} 之后，
用 \Tp{/proc/interrupts} 的计数分布复验中断确实落在预期核上。
同类自动优化服务还包括 \Tp{tuned}（除非使用它的实时配置档）与容器运行时的 CPU 管理器。
\end{warning}

清单 \ref{lst:ch08-cmdline} 给出完整的启动参数与校验序列。

\begin{lstlisting}[style=shellstyle,caption={Shell：实时启动参数与校验序列},label={lst:ch08-cmdline}]
# 一、内核命令行。CPU 0-3 为家核，4-7 为实时核。
# 写入 /etc/default/grub 的 GRUB_CMDLINE_LINUX，再重新生成 grub.cfg
isolcpus=4-7 nohz_full=4-7 rcu_nocbs=4-7 irqaffinity=0-3 \
threadirqs skew_tick=1 nosoftlockup nowatchdog tsc=reliable \
intel_idle.max_cstate=0 processor.max_cstate=0 iommu=pt \
pcie_aspm=off transparent_hugepage=never nosmt
# 更激进的写法是再加 idle=poll：完全不进 idle，功耗与热代价需与频率一节一起评估

# 二、校验：内核与抢占模型
cat /proc/cmdline
uname -a | grep -i PREEMPT_RT              # 版本串里应出现 PREEMPT_RT
grep -E 'CONFIG_PREEMPT_RT=y|CONFIG_HIGH_RES_TIMERS=y|CONFIG_NO_HZ_FULL=y' \
     /boot/config-$(uname -r)

# 三、校验：调度与隔离现状
cat /sys/kernel/debug/sched/preempt 2>/dev/null || true
cat /sys/devices/system/cpu/isolated
cat /sys/devices/system/cpu/nohz_full

# 四、校验：中断归属与线程优先级
grep -E '^ *[0-9]+' /proc/interrupts | head -20    # 实时核那一列不该有非预期计数
cat /proc/irq/default_smp_affinity
cat /proc/irq/44/smp_affinity_list
ps -eLo psr,rtprio,ni,pri,cls,comm | awk '$1 >= 4' # 实时核上不该有 SCHED_OTHER 线程
for t in $(pgrep 'irq/'); do chrt -p "$t" | tail -1; done

# 五、关掉自动调整与噪声服务
systemctl stop --now irqbalance
systemctl disable irqbalance
tuned-adm profile cpu-partitioning 2>/dev/null || true

# 六、把实时进程塞进独占 cpuset（cgroup v2）
mkdir -p /sys/fs/cgroup/rt.slice
echo "4-7" > /sys/fs/cgroup/rt.slice/cpuset.cpus
echo "4-7" > /sys/fs/cgroup/rt.slice/cpuset.cpus.exclusive
echo $TARGET_PID > /sys/fs/cgroup/rt.slice/cgroup.procs

# 七、冒烟：先测无业务基线，再测带业务的基线
cyclictest -a4-7 -t1 -p99 -i1000 -l100000 -q -m
\end{lstlisting}

\section{调度策略、优先级与全局配额}
\label{sec:ch08-sched}

\DomRT{}Linux 的调度类是一个严格优先的层次：\Tp{SCHED\_DEADLINE} 高于
\Tp{SCHED\_FIFO} 与 \Tp{SCHED\_RR}，两者高于默认的 \Tp{SCHED\_OTHER}（CFS），
CFS 又高于 \Tp{SCHED\_BATCH} 与 \Tp{SCHED\_IDLE}。这就是 \Tp{SCHED\_OTHER}
上实时性总不够的原因：它按权重分配时间片、允许被任何 RT 类线程抢占、
并在负载变化时用虚拟运行时间重新排队，可等待时长没有上界承诺。

\Tp{SCHED\_FIFO} 是同优先级内先到先服务、不轮转，因此每个实时优先级上只放一个线程
才是可靠的定价方式；\Tp{SCHED\_RR} 只多了同优先级轮转，适合等价的多 worker，不适合硬周期。
实时优先级空间是 1 到 99：把最高几档留给中断线程与故障收敛路径，
业务实时线程放中档，这样一次调价不会牵动整个系统。
\Tp{SCHED\_DEADLINE} 用三个参数表达任务：每 \Tp{period} 内需要 \Tp{runtime} 的执行时间，
且必须在 \Tp{deadline} 前完成；内核按声明的总带宽做准入控制，用 EDF 加 CBS 运行。
它比 \Tp{SCHED\_FIFO} 更适合多个实时任务共存又不想手工排优先级的场合，
代价是参数必须真实，超预算会被 throttle。清单 \ref{lst:ch08-deadline} 是最小设置代码。

\begin{lstlisting}[style=cstyle,caption={C：用 sched\_setattr 进入 SCHED\_DEADLINE 并回读校验},label={lst:ch08-deadline}]
#define _GNU_SOURCE
#include <sched.h>
#include <sys/syscall.h>   /* struct sched_attr 定义在 uapi/linux/sched/types.h */
#include <unistd.h>
#include <errno.h>
#include <stdio.h>

/* 10 毫秒周期，单次执行预算 1.5 毫秒，deadline 等于 period */
#define DL_RUNTIME_NS	1500000ULL
#define DL_PERIOD_NS	10000000ULL

static int enter_deadline(void)
{
	struct sched_attr attr = {
		.size = sizeof(attr),
		.sched_policy = SCHED_DEADLINE,
		.sched_flags = SCHED_FLAG_RESET_ON_FORK,  /* 子进程不继承实时参数 */
		.sched_runtime = DL_RUNTIME_NS,
		.sched_deadline = DL_PERIOD_NS,
		.sched_period = DL_PERIOD_NS,
	};

	if (syscall(__NR_sched_setattr, 0, &attr, 0) != 0) {
		/* EPERM：缺 CAP_SYS_NICE 或 RLIMIT_RTPRIO；EINVAL：不满足准入 */
		perror("sched_setattr");
		return -1;
	}
	/* 回读校验：没读回来就不能认为设置生效 */
	struct sched_attr got = { .size = sizeof(got) };
	if (syscall(__NR_sched_getattr, 0, &got, sizeof(got), 0) != 0) {
		perror("sched_getattr");
		return -1;
	}
	printf("policy=%u runtime=%llu period=%llu\n",
	       got.sched_policy, got.sched_runtime, got.sched_period);
	return 0;
}
\end{lstlisting}

权限与配额是两道独立的门。第一道是能力与资源限制：\Tp{chrt}（\Tp{-p} 只作用于已有进程）、
\Fn{sched\_setattr} 与 \Fn{pthread\_setschedparam} 都需要 \Tp{CAP\_SYS\_NICE}，
非 root 运行时要靠 \Tp{/etc/security/limits.d/} 里的 \Tp{rtprio} 与 \Tp{nice} 限额授予；
不满足时返回失败但常被脚本忽略，看起来像优先级设置没生效。
第二道是全局配额，见 \ref{sec:ch08-rt-changes} 的说明框。
表 \ref{tab:ch08-policies} 并列要点。

\begin{table}[htbp]
\centering
\caption{调度策略的实时相关差异}\label{tab:ch08-policies}
\begin{tabular}{>{\raggedright\arraybackslash}p{0.18\textwidth}
                >{\raggedright\arraybackslash}p{0.22\textwidth}
                >{\raggedright\arraybackslash}p{0.16\textwidth}
                >{\raggedright\arraybackslash}p{0.23\textwidth}}
\toprule
策略 & 抢占与轮转语义 & 权限要求 & 实时项目里的用法 \\
\midrule
\Tp{SCHED\_OTHER} & 权重分时，可被 RT 类抢占 & 无 & 家核上的普通服务、日志、shell \\
\Tp{SCHED\_FIFO} & 严格优先，同优先级不轮转 & \Tp{CAP\_SYS\_NICE} 或 rtprio 限额 & 实时主循环、中断线程（默认 50） \\
\Tp{SCHED\_RR} & 严格优先加同优先级时间片 & 同上 & 等价的多 worker，不用于硬周期 \\
\Tp{SCHED\_DEADLINE} & 参数化带宽，EDF 加 CBS，超预算 throttle & 同上 & 多实时任务共存、想显式给预算 \\
\Tp{SCHED\_BATCH}、\Tp{SCHED\_IDLE} & 低于默认类 & 无 & 后台编译、采集与压缩 \\
\bottomrule
\end{tabular}
\end{table}

优先级反转在 Linux 上有三种来源与对应处置：用户态 normal 互斥不提供继承
（改用 \Tp{PRIO\_INHERIT} 属性）；实时线程等待低优先级线程持有的锁
（PREEMPT\_RT 已把它变成可睡眠加继承）；以及中断线程与业务线程之间的手写优先级倒置。
最后一种最常见症状是给网卡中断线程设了 50、业务线程设了 49，
于是每个包都把业务线程压在后面。\Tp{rt\_checklist}
（rt-tests 提供，有时安装名写作 \Tp{rt-checklist}）会检查中断线程与实时任务的相对优先级，
适合纳入开机自检脚本。

\section{内存：锁定、预触碰与抖动源}
\label{sec:ch08-memory}

实时循环里第一次触碰一页内存会触发缺页，一次缺页在页表填充与页面分配路径上的耗时
很容易达到数十到数百微秒，因此所有可能被碰的页必须在启动阶段就碰过。三个动作配合完成：
\Fn{mlockall} 用 \Tp{MCL\_CURRENT} 与 \Tp{MCL\_FUTURE} 锁定现有与将来映射；
\Tp{MAP\_POPULATE} 在映射时预填充（或映射后立即显式写一遍）；
栈预触碰，即启动时把最坏栈深度整段写一遍，让栈的自动扩展在实时化之前完成。
\Tp{RLIMIT\_MEMLOCK} 必须足够大，否则 \Fn{mlockall} 只成功一部分就返回错误。
清单 \ref{lst:ch08-mlock} 是骨架。

\begin{lstlisting}[style=cppstyle,caption={C++：mlockall 与栈、堆的预触碰},label={lst:ch08-mlock}]
#include <sys/mman.h>
#include <sys/resource.h>
#include <cstdio>
#include <cstring>
#include <cstdint>

namespace rt {

constexpr std::size_t kStackWarmBytes = 512u * 1024u;   // 与最坏递归深度一致
constexpr std::size_t kPoolBytes      = 4u * 1024u * 1024u;

// 递归把栈一路写到底：每帧都真实触碰，避免运行期首次深递归时缺页
void warm_stack(std::size_t remaining) noexcept
{
	alignas(64) char frame[1024];
	for (char &c : frame)
		c = static_cast<char>(remaining & 0x7f);
	if (remaining > sizeof(frame))
		warm_stack(remaining - sizeof(frame));
}

// 只在启动阶段调用一次：此后稳态路径不得分配或释放
bool pin_and_warm_up() noexcept
{
	// (1) 放开 RLIMIT_MEMLOCK；容器内还需要具备相应能力
	rlimit lim{};
	if (getrlimit(RLIMIT_MEMLOCK, &lim) == 0) {
		lim.rlim_cur = lim.rlim_max;
		(void)setrlimit(RLIMIT_MEMLOCK, &lim);   // 失败则依赖已有的部署配置
	}

	// (2) 锁定全部当前与将来的映射；失败必须让上层进入安全状态
	if (mlockall(MCL_CURRENT | MCL_FUTURE) != 0) {
		std::perror("mlockall");
		return false;
	}
	warm_stack(kStackWarmBytes);              // (3) 栈预触碰

	char *p = static_cast<char *>(::malloc(kPoolBytes));   // (4) 堆预触碰
	if (p == nullptr) return false;
	std::memset(p, 0, kPoolBytes);                         // 之后只复用

	// (5) 大缓冲：MAP_POPULATE 映射后立即写，再预填充数据结构与虚表
	void *buf = mmap(nullptr, kPoolBytes, PROT_READ | PROT_WRITE,
			 MAP_PRIVATE | MAP_ANONYMOUS | MAP_POPULATE, -1, 0);
	if (buf == MAP_FAILED) return false;
	std::memset(buf, 0, kPoolBytes);
	warm_buffers(buf, kPoolBytes);

	// (6) 关闭回收抖动：swap 必须真关掉，仅设 swappiness 不等于不换页
	//     swapoff -a
	//     echo 0 > /proc/sys/vm/zone_reclaim_mode
	//     echo never > /sys/kernel/mm/transparent_hugepage/enabled
	return true;
}

}  // namespace rt
\end{lstlisting}

\begin{warning}[title={锁定之后仍然缺页的两个隐蔽入口}]
一是共享库延迟绑定：第一次调用某个符号会触发动态链接器取页与写重定位表。
启动时用 \Tp{LD\_BIND\_NOW} 预绑定，或链接时加 \Tp{-Wl,-z,now}；
二进制自身的页也要在实时化之前完整执行过一遍。
二是内存由别的线程建立：\Fn{mmap} 之后首次触碰发生在哪个线程，
缺页就记在哪个线程上，因此预触碰必须由真正使用该内存的线程在实时化之前完成。
此外，与你的进程无关却作用在你进程上的内核活动同样制造延迟：
内存回收（\Tp{pgscan\_kswapd} 计数上升）、透明大页的合并与拆分、
\Tp{zone\_reclaim\_mode} 非零。用 \Tp{vmstat} 与 \Tp{/proc/vmstat} 在测量窗口内采样。
\end{warning}

\section{频率、idle 与硬件层干扰}
\label{sec:ch08-frequency}

调频与 idle 是最容易被忽略的两类延迟源，量级常常压过调度改造的全部收益。

\begin{itemize}
  \item C-state：从深度睡眠唤醒的成本是几十到上百微秒，某些平台超过 1 毫秒。
        手段按侵入性排序：逐项写
        \Tp{/sys/devices/system/cpu/cpu6/cpuidle/state2/disable}
        只关深状态、保留轻量 halt；打开 \Tp{/dev/cpu\_dma\_latency}
        并持有描述符（全局影响，重启即失效，适合短时测试）；
        粗粒度做法是内核参数 \Tp{intel\_idle.max\_cstate=0}
        与 \Tp{processor.max\_cstate=0}；最激进是 \Tp{idle=poll}，彻底不睡，
        但功耗、热代价与 SMT 邻居要一起评估。
  \item 调频：把实时核固定在工作频率。\Tp{cpupower frequency--{\Bk}set --{\Bk}g performance}
        或 \Tp{x86\_{\Bk}energy\_{\Bk}perf\_{\Bk}policy}；Intel 平台用 \Tp{intel\_{\Bk}pstate=no\_{\Bk}hwp}
        关掉硬件自主调频，再写 \Tp{/sys/devices/system/cpu/intel\_{\Bk}pstate/no\_{\Bk}turbo}
        关睿频。理由不是睿频更快，而是睿频时快时慢让执行时间不可静态定价。
  \item SMT：两个超线程共享流水线与一级缓存，邻居负载直接改变你的每周期耗时。
        可用 \Tp{nosmt} 关掉，或把实时任务两两绑定在同一物理核的两个逻辑核上
        并保证负载互补。
  \item 热降频：确认机箱风道与环境温度，读
        \Tp{/{\Bk}sys{\Bk}/devices{\Bk}/system{\Bk}/cpu{\Bk}/cpu6{\Bk}/thermal\_{\Bk}throttle/} 下的降频计数，
        验收数据必须在热稳定之后取。
  \item PCIe 省电：\Tp{pcie\_{\Bk}aspm=off} 与
        \Tp{/sys/{\Bk}module/{\Bk}pcie\_aspm/{\Bk}parameters/{\Bk}policy} 一起确认；
        ASPM 打开时设备唤醒时间会加进中断响应延迟。
  \item 时钟源：内核判定 TSC 不可靠（常见于某些虚拟化或外频不稳的平台）后会切到
        更慢的时钟源，绝对睡眠的分辨率随之退化。确认
        \Tp{/sys/{\Bk}devices/{\Bk}system/{\Bk}clocksource/{\Bk}clocksource0/current\_{\Bk}clocksource}
        与同目录下的 \Tp{available\_{\Bk}clocksource}，必要时用启动参数强制并加
        \Tp{tsc=reliable}（前提是已排除其它抖动来源）。
\end{itemize}

\section{I/O 约束、实时循环范式与验收}
\label{sec:ch08-io-loop}

I/O 在实时路径上的三条原则：不做打开关闭、不做同步系统调用、不给对齐问题留运行期补救。
\Tp{O\_DIRECT} 绕开页缓存与回写线程，代价是缓冲区地址、文件偏移与长度
都要按块设备逻辑扇区对齐（多数平台要求 512 或 4096 字节），
用 \Fn{posix\_memalign} 或 \Fn{aligned\_alloc} 在启动阶段一次性备好缓冲池。
块层调度器（\Tp{/sys/block/sda/queue/scheduler}）在与实时核绑定的设备上
选延迟优先的策略，或直接 \Tp{none} 配合单队列提交。
异步 I/O 的位置要看清：AIO 与 \Lib{io\_uring} 能把等待移出关键路径，
但注册缓冲区、固定文件描述符与提交时的页钉住都是启动阶段动作，稳态只提交已注册的请求。
网络侧 \Tp{SO\_PRIORITY} 决定出队优先级，\Tp{SO\_TXTIME} 让报文携带发送时间戳
交给 ETF 队列规则整形；一旦用到后者，时钟源与 \Tp{CLOCK\_TAI} 的偏差校准
就进入了实时预算，完整的 TSN 讨论见以太网章。
噪声服务必须处理：syslog 与 \Tp{klogd} 的刷盘、\Tp{cron} 与 \Tp{logrotate} 的定时任务、
\Tp{rasdaemon} 与 \Tp{mcelog} 的错误采集、\Tp{systemd-journald} 的同步写；
把它们的执行位置安排到家核或降低触发频率，并在验收窗口内禁止任何定时任务。

\begin{bestpractice}[title={稳态循环的三条铁律}]
第一，一个周期只做一次绝对时间等待：\Fn{clock\_nanosleep} 配
\Tp{CLOCK\_MONOTONIC} 与 \Tp{TIMER\_ABSTIME}，绝对时刻按周期递增；
不用相对睡眠，否则唤醒延迟会被累加进下一周期。
第二，稳态路径不出现 \Fn{malloc}、\Tp{new}、\Tp{std::string}、异常展开、
动态链接解析、日志格式化与任何首次触碰的页；错误处理路径同样要预触碰并锁定，
否则第一次出错就是最慢的一次执行。
第三，实时线程只做纯计算与无锁交接，与外部世界（磁盘、网络、日志、IPC 服务端）
的交互全部交给家核线程，中间用预分配的无锁环形队列传指针。
\end{bestpractice}

清单 \ref{lst:ch08-loop} 是可直接改用的骨架：绑核与进入实时类在等待之前完成，
退出时按相反顺序解除。

\begin{lstlisting}[style=cppstyle,caption={C++：绝对周期实时循环骨架},label={lst:ch08-loop}]
#include <cerrno>
#include <csignal>
#include <cstdint>
#include <ctime>
#include <atomic>
#include <pthread.h>
#include <sched.h>
#include <sys/mman.h>
#include <unistd.h>

#include "rt_buffers.h"     // 启动阶段分配并预触碰的静态缓冲
#include "safe_state.h"     // 第 4 章的看门狗喂狗与故障收敛

namespace {

constexpr std::uint64_t kPeriodNs = 1000000ULL;   // 1 毫秒
constexpr int kFifoPrio = 60;   // 与中断线程的相对位置按设计定，勿随手取高值
std::atomic<bool> g_stop{false};

void on_signal(int sig) { (void)sig; g_stop.store(true, std::memory_order_relaxed); }

bool pin_to_cpu(int cpu_id) noexcept
{
	cpu_set_t set{};
	CPU_ZERO(&set);
	CPU_SET(cpu_id, &set);
	return pthread_setaffinity_np(pthread_self(), sizeof(set), &set) == 0;
}

bool enter_fifo(int prio) noexcept
{
	sched_param sp{};
	sp.sched_priority = prio;          // 需要 CAP_SYS_NICE 或 rtprio 限额
	return pthread_setschedparam(pthread_self(), SCHED_FIFO, &sp) == 0;
}

void add_ns(timespec &ts, std::uint64_t ns) noexcept
{
	ts.tv_nsec += static_cast<long>(ns);
	while (ts.tv_nsec >= 1000000000L) { ts.tv_nsec -= 1000000000L; ++ts.tv_sec; }
}

}  // namespace

int rt_loop_main(int cpu_id) noexcept
{
	// 只放行退出信号；其余信号不进实时线程，避免任意点被异步打断
	sigset_t mask{};
	sigemptyset(&mask);
	sigaddset(&mask, SIGTERM);
	sigaddset(&mask, SIGINT);
	pthread_sigmask(SIG_UNBLOCK, &mask, nullptr);
	std::signal(SIGTERM, on_signal);

	if (mlockall(MCL_CURRENT | MCL_FUTURE) != 0) return 1;  // 启动阶段已做则幂等
	if (!pin_to_cpu(cpu_id) || !enter_fifo(kFifoPrio)) return 2;
	if (!rt::warm_all_buffers_and_error_paths()) return 3;  // 含错误处理分支

	timespec next{};
	clock_gettime(CLOCK_MONOTONIC, &next);    // 也可换成 timerfd 的周期读取
	long overrun = 0;

	while (!g_stop.load(std::memory_order_relaxed)) {
		add_ns(next, kPeriodNs);

		// 唯一允许被信号打断的地方：重新等待同一个绝对时刻
		while (clock_nanosleep(CLOCK_MONOTONIC, TIMER_ABSTIME, &next, nullptr) != 0) {
			if (errno == EINTR) continue;
			return 4;                 // 参数或时钟问题，交给安全状态
		}

		if (rt::sample_and_control() != 0) {   // 全部为预触碰、无分配的操作
			enter_safe_state();
			feed_watchdog();               // 呼应第 4 章的硬件看门狗位
			continue;
		}

		const timespec now_ts{rt::monotonic_now()};
		if (now_ts > next) {                   // 迟到：只统计，不补做
			++overrun;
			if (rt::overrun_budget_exceeded(overrun)) enter_safe_state();
			// next 不回退到 now，保持标称周期恒定；连续迟到由预算判定
			// 触发安全状态，而不是把错过的执行堆到现在
		}

		rt::publish_pointers_only(next);       // 与家核线程无锁交接
	}

	munlockall();     // 优雅退出：停调度、通知看门狗与家核之后才解除锁定
	return 0;
}
\end{lstlisting}

验收方法与 \Tp{cyclictest} 的完整参数在第 11 章，这里给预算分解的量级判断。
以隔离核、\Tp{SCHED\_FIFO} 高优先级、1 毫秒周期、10 万次采样为例：
过关线通常写作最大延迟不超过 50 到 100 微秒量级。分解下来：中断响应与调度唤醒
占几微秒到十几微秒，PREEMPT\_RT 残留的短临界区占几微秒，深 C-state 唤醒
贡献几十到上百微秒（这就是它被列为首要开关的原因），未锁定的页或缺页处理
贡献单次几百微秒的尖峰，一次 console 打印或 \Tp{cron} 触发能留下毫秒级单点。
换句话说，判据不是越小越好，而是每一项毛刺都能被指认到一个来源；
指认不出来就是还没关干净。表 \ref{tab:ch08-jitter-map} 给出映射。

\begin{table}[htbp]
\centering
\caption{抖动来源到开关与验证命令的映射}\label{tab:ch08-jitter-map}
\begin{tabular}{>{\raggedright\arraybackslash}p{0.24\textwidth}
                >{\raggedright\arraybackslash}p{0.26\textwidth}
                >{\raggedright\arraybackslash}p{0.30\textwidth}}
\toprule
抖动来源 & 开关或动作 & 验证命令 \\
\midrule
自动迁移与负载均衡 & \Tp{isolcpus} 加独占 cpuset & \Tp{cat /{\Bk}sys/{\Bk}devices/{\Bk}system/{\Bk}cpu/{\Bk}isolated}；\Tp{ps -{\Bk}eLo psr,cls,comm} \\
周期 tick 与 RCU 回调 & \Tp{nohz\_{\Bk}full} 与 \Tp{rcu\_{\Bk}nocbs} & \Tp{cat /{\Bk}sys/{\Bk}devices/{\Bk}system/{\Bk}cpu/{\Bk}nohz\_{\Bk}full}；\Tp{perf} 记录 tick 事件 \\
中断落到实时核 & \Tp{irqaffinity} 加逐个写亲和性 & \Tp{cat /proc/interrupts}；停用 \Tp{irqbalance} \\
中断线程优先级不当 & \Tp{chrt -p} 调价，或改默认策略 & \Tp{ps -eo comm,rtprio} 查 \Tp{irq/} 线程 \\
实时配额 throttle & 提高 \Tp{sched\_{\Bk}rt\_{\Bk}runtime\_{\Bk}us} 或减少 RT 线程 & \Tp{cat /{\Bk}proc/{\Bk}sys/{\Bk}kernel/{\Bk}sched\_{\Bk}rt\_{\Bk}runtime\_{\Bk}us} \\
缺页与内存回收 & \Fn{mlockall}、预触碰、\Tp{swapoff}、\Tp{zone\_{\Bk}reclaim\_{\Bk}mode=0} & \Tp{grep pgscan /{\Bk}proc/{\Bk}vmstat}；\Tp{ltrace} \\
调频与睿频 & performance 加 \Tp{no\_turbo}、\Tp{no\_hwp} & \Tp{cpupower frequency-info} \\
深度 idle 唤醒 & 逐项 disable state 或 \Tp{idle=poll} & \Tp{cat .../{\Bk}cpuidle/{\Bk}state*/{\Bk}latency} \\
日志与定时任务 & 刷盘移到家核，验收窗口禁 \Tp{cron} & \Tp{journalctl -{\Bk}-{\Bk}since} 比对毛刺时间戳 \\
\bottomrule
\end{tabular}
\end{table}

\begin{warning}[title={chrt 静默失败与两道权限门}]
非 root 用户下 \Tp{chrt} 失败常常只打印一行错误，在脚本里被重定向后就表现为
系统照旧跑、只是没有实时性。两道门都要开：给 \Tp{CAP\_SYS\_NICE}
（或临时以特权设置后再放权），并在 \Tp{/etc/security/limits.d/rt.conf} 写
\Tp{@rtgrp - rtprio 99} 与 \Tp{@rtgrp - memlock unlimited}；PAM 变更需重新登录生效。
设置完必须复验：\Tp{chrt -p} 读回的优先级才作数，
\Tp{ps -eo pid,cls,rtprio,comm} 的 \Tp{cls} 列显示 \Tp{FF} 或 \Tp{RR}
才算真的进入实时类。
\end{warning}

加固与验收本身要按顺序做，每步留一份测量记录，否则无法归因：
(1) 确认内核是 PREEMPT\_RT 且开启高精度定时器；
(2) 关中断自动平衡、把中断赶离实时核、绑定 cpuset；
(3) 关深度 idle、固定频率、关睿频，必要时关 SMT；
(4) 锁定内存并完成栈、堆与错误路径的预触碰；
(5) 把日志与定时任务移到家核；
(6) 先跑 \Tp{cyclictest} 取无业务基线，再取带业务基线，
两者之差就是业务线程与共享资源的干扰量；
(7) 最后跑 \Tp{rt\_checklist} 做一致性检查，把输出与启动参数一起归档。

% ============================================================
% ---- frag_ch09.tex ----
\chapter{无锁、静态分配与内存纪律}
\label{ch:ch09-lockless-memory}

实时路径上的内存与并发不是性能话题，而是可论证性问题。一个控制环能不能在最坏情况下交出结果，取决于三件事：这段代码的运行时间有没有上界，它占用的内存有没有上界，它等待别人放行资源的时间有没有上界。动态分配破坏第二件，无边界自旋破坏第三件，而编译器与处理器的重排可能让第三件变成一件永远不会发生的事件——数据看起来写好了，读者却永远读不到。本章围绕这三条线展开：先说明为什么实时路径要彻底禁止 \Fn{malloc}，再给出替代它的静态结构（对象池、编译期数组、placement 构造），然后把无锁队列的两个主流形态（SPSC 与 MPMC）拆到 cache line 与内存序的粒度，最后落到栈纪律、对齐与陷阱。

外设与内存之间的搬运由 DMA 完成，那部分在第 5 章讨论；寄存器级位域读写属于另一分册的内容，本章不涉及。这里只谈并发数据结构本身：谁写、谁读、需要什么保证、代价是什么。

\section{为什么实时路径禁止动态分配}

\Fn{malloc} 的问题不在“慢”，而在“最坏情况不可界”。多数分配器的快路径是一个大小分档的空闲链表或位图，命中时几十个周期就能返回；一旦需要向操作系统扩段、需要做块分裂与合并、需要遍历空闲链找回避碎片，同一次调用的耗时可能跨越两到三个数量级。更麻烦的是，这条慢路径不是随机出现的，它跟堆的历史状态相关：系统连续运行三小时之后才第一次踩到它，于是测试永远看不到。

碎片是第二个不可界来源。即便平均分配时间可以接受，长时间运行后可用内存被切碎成不满足任何一次请求的小片，分配开始失败。而失败分支恰恰是最不该存在的东西：\Fn{malloc} 返回空指针意味着你的实时任务此刻必须有一条“放弃这一帧”的处理路径，这条路径几乎从不被执行，因此也几乎从来没被验证过。把一段从未执行过的代码放进安全关键程序，等价于把一个未测试的功能放进产品。

正因如此，主流标准对它的态度是禁用而不是限制。\Lang{C} 侧最直接的是 MISRA C:2012 的 Rule 21.3：动态内存分配与释放函数的使用应被禁用。在合规项目里这通常是强制条款，例外必须写偏离说明并附论证。\Lang{C++} 没有标准层面的禁令，但同一条工程理由完全适用，所以本册对 \Fn{new}、\Fn{delete} 与容器按同一标准处理。功能安全侧，EN 50128 与 IEC 61508 要求在高层级软件中提供静态论证：内存使用要有界，运行期不得出现未定义行为，不得有不可回收的资源泄漏。动态分配同时触碰“有界”与“静态论证不可行”两条，评审时一般直接被要求替换。

替代方案不是“少用”，而是把分配这件事移出运行期。四类可行做法：

\begin{itemize}
\item 编译期数组与静态对象：文件作用域的数组，长度是常量表达式，容量在链接期就定了，上界可以直接从符号表读出来。这是最可取的形态。
\item 对象池：初始化时一次性从静态缓冲区切出固定数量、固定大小的槽，运行期只在槽之间搬“在用/空闲”标记。\DomRtos{} 里几乎所有稳态路径都能用池表达。
\item \DomRtos{} 堆实现选型：以 FreeRTOS 为例，\Tp{heap\_1} 只允许分配不允许释放，所有请求从一个大数组里线性推进，耗尽即失败，行为完全可预测；\Tp{heap\_4} 带空闲块合并，会碎片化；\Tp{heap\_3} 只是标准库 \Fn{malloc} 的包装，等于把上述所有不可界性引进来。安全关键的稳态路径应选 \Tp{heap\_1} 形态，或者干脆不经过堆。
\item placement 构造：在已有的静态存储上调用 \Fn{new} 的 placement 形式，只构造不分配。它把 \CPP{} 的构造函数语义保留下来，同时不触碰分配器。
\end{itemize}

\begin{guideline}[title={实时路径的内存准入三条}]
把每一条稳态路径的内存动作归为三类之一：不访问内存（只碰寄存器与局部量）、访问固定地址（静态对象或池槽）、访问栈上定长空间。任何不属于这三类的动作（\Fn{malloc}、\Fn{free}、\Fn{new}、\Fn{delete}、容器扩容、字符串拼接、异常展开所需的堆）都必须出现在设计文档的“非实时路径”清单里，并有明确的触发上下文。评审时只看这张清单，不看代码。
\end{guideline}

\section{SPSC 环形缓冲：最小可用的无锁结构}

单生产者单消费者队列是实时系统里最常用的形态：ISR 生产、任务消费；或者采集任务生产、处理任务消费。它之所以能真正做到无锁，是因为“只有一个写者”这个前提让索引更新退化成普通存储指令——不需要比较交换，不需要自旋，只需要保证可见性与顺序。

基本形态是一个容量为 $2^{k}$ 的定长数组加两个单调递增的计数器 \Tp{head}（写位置）与 \Tp{tail}（读位置）。用掩码取模代替除法，所以容量必须是 2 的幂。计数器不折回、只增长，靠无符号溢出自然回绕，这样满与空可以用纯减法判别，不必维护额外状态位。

\begin{lstlisting}[style=cppstyle,caption={C++：SPSC 无锁环形缓冲（定长槽）},label={lst:ch09-spsc}]
#include <atomic>
#include <cstddef>
#include <cstdint>

template <typename T, std::size_t PowerOfTwo>
class SpscRing {
  static_assert((PowerOfTwo & (PowerOfTwo - 1)) == 0,
                "capacity must be a power of two");
  static_assert(PowerOfTwo >= 2, "capacity must be at least two");

  static constexpr std::size_t kMask = PowerOfTwo - 1;

  // 写者游标与读者游标各自独占一条 cache line，
  // 否则两端会互相打伪共享。
  struct alignas(64) Cursor {
    std::atomic<std::size_t> v{0};
    char pad[64 - sizeof(std::atomic<std::size_t>)];
  };

  Cursor write_;   // 只有生产者写入
  Cursor read_;    // 只有消费者写入
  T slots_[PowerOfTwo];

public:
  // 返回 false 表示缓冲已满，本帧被丢弃。
  // 丢弃策略必须在设计文档里显式声明，不能靠运气。
  bool push(const T& value) noexcept {
    const std::size_t w = write_.v.load(std::memory_order_relaxed);
    const std::size_t r = read_.v.load(std::memory_order_acquire);
    if (w - r >= PowerOfTwo - 1) {
      return false;            // 留一格约定：追上即视为满
    }
    slots_[w & kMask] = value;
    write_.v.store(w + 1, std::memory_order_release);
    return true;
  }

  bool pop(T& out) noexcept {
    const std::size_t r = read_.v.load(std::memory_order_relaxed);
    const std::size_t w = write_.v.load(std::memory_order_acquire);
    if (r == w) {
      return false;            // 空
    }
    out = slots_[r & kMask];
    read_.v.store(r + 1, std::memory_order_release);
    return true;
  }

  // 容量核算用：N 个槽只给得出 N-1 格缓冲深度。
  static constexpr std::size_t capacity() noexcept { return PowerOfTwo - 1; }
};
\end{lstlisting}

\subsection*{满/空判别的两种约定}

两指针相等既可以解释成全空也可以解释成全满，必须选一种并付出代价。第一种是留一格：把写者追上读者的位置判为满，代价是可用容量永远是 $N-1$，好处是两个计数器足以表达全部状态，读端不需要任何附加信息。第二种是维护一个计数变量：容量能用满 $N$ 格，代价是计数与两个游标之间有三个量的相对一致性问题，在没有第三视角时无法保证同时看到三者的最新值——除非引入序列号或退回锁。本册推荐在 SPSC 上用留一格，在 MPMC 上才考虑序列号。做容量预算时，$N$ 个槽不等于 $N$ 格缓冲深度，这一点反复出错。

\subsection*{可变长负载}

控制环里多数负载是定长结构体，直接放槽里最省事。传感器数据长度可变时有两种做法。写长度前缀：槽首放一个定长宽度的长度字段，后跟数据，读端先读长度再按长度取。它把可用容量变成“槽个数”与“总字节数”两个约束，分析时两者都要算。定长槽加拷贝：负载被截断或分片到固定大小的槽，逻辑帧由若干物理槽拼起来，需要一个帧边界标记。后者更浪费内存，但索引推进是常数，最坏时间更好算。实时路径优先后者；只有当负载本身有自然上限（例如最多 64 字节）时才用前者。

\subsection*{为什么 SPSC 不需要更强的顺序保证}

关键在于谁读谁、为了什么。写者对 \Tp{tail} 的 acquire 读，是为了确保自己覆写某个槽之前读者已经把这个槽的内容取走；写者对自己 \Tp{head} 的 relaxed 读永远不需要与别人一致。写者需要保证的顺序只有一条：数据写入必须排在游标发布之前——这正是 release 存储的作用。对称地，读者 acquire 到新游标之后，对槽内容的读不会被提到发布之前。单写者、单读者，两条边，足够。

\section{MPMC：CAS 代价、序列号槽与有界数组}

多生产者或多消费者让索引更新变成竞争，于是必须引入原子读改写。

\Fn{compare\_exchange\_weak} 允许伪失败，在循环里比 \Fn{compare\_exchange\_strong} 少一条平台上的比较分支，因此官方建议就是把它放在循环里用。但在实时路径上“循环”本身就是问题：每次失败意味着这条 cache line 被另一个核拿到过，代价从几十纳秒到几百纳秒；竞争越激烈失败率越高，而失败率越高又导致更多尝试。最坏情况没有上界，因为它取决于其他核什么时候让出这条 line——那可能是另一个核上一个更高优先级的任务，也可能是内存控制器排队。

\begin{warning}[title={CAS 自旋与自旋锁上的优先级反转}]
自旋锁与 CAS 循环都要求“持有权的人很快会放手”。如果低优先级任务刚拿到 line 或持有锁就被高优先级任务抢占，所有自旋者都在为一个不放手的核消耗周期，而它永远得不到运行机会——这就是优先级反转在无锁结构上的表现，和传统锁上的一样致命，但没有现成的优先级继承机制可以救你。有界自旋加退避只缩短浪费，不解决反转。真正的解法只有两个：把临界区纳入优先级天花板协议（阻塞项如何计入响应时间界见第 10 章），或者退回 SPSC 结构，让竞争根本不发生。
\end{warning}

\subsection*{带序列号的槽}

即便用原子 CAS 抢到了某个槽的所有权，读者仍可能读到写了一半的数据：写者先更新了标志，还没写完负载。给每个槽加一个序列号，把槽变成自洽的三元组（序号、负载、序号），可以让读端判断自己读到的是一个完整写入的槽。

\begin{lstlisting}[style=cppstyle,caption={C++：带序列号的 MPMC 槽（有界队列的单元）},label={lst:ch09-seq-slot}]
#include <atomic>
#include <cstddef>

template <typename T>
struct alignas(64) SeqSlot {
  // seq 既标识本槽被使用的轮次，也用于遮蔽半写状态。
  std::atomic<std::size_t> seq{0};
  T payload{};

  // 写者：确认槽属于本轮 -> 写负载 -> 发布为"已写"标记
  bool try_write(std::size_t pos, const T& value) noexcept {
    if (seq.load(std::memory_order_acquire) != pos) {
      return false;            // 槽正被读，调用方需重试或放弃
    }
    payload = value;           // 半写状态被 seq 遮蔽，读者不会采信
    seq.store(pos + 1, std::memory_order_release);
    return true;
  }

  // 读者：确认已写 -> 取负载 -> 复核 -> 推到下一轮
  bool try_read(std::size_t pos, std::size_t step, T& out) noexcept {
    if (seq.load(std::memory_order_acquire) != pos + 1) {
      return false;
    }
    T local = payload;
    if (seq.load(std::memory_order_acquire) != pos + 1) {
      return false;            // 复核失败：读到了撕裂
    }
    out = local;
    // 下一轮的所有权标记 = 本轮 + 队列容量，由外层传入
    seq.store(pos + step, std::memory_order_release);
    return true;
  }
};
\end{lstlisting}

推进步长必须等于队列容量，由外层入队/出队逻辑计算并传入；把它写成常量会得到一个错误的队列。要点是双重复核：读负载之后再读一次序号，确认写者没有在你读的过程中再次覆写。序列号方案换来的代价是每个槽多一条原子变量与一条 cache line 的填充，内存开销通常是定长槽方案的两倍。

\subsection*{为什么有界数组优于链表}

链表在无锁语境下有两个几乎不可解决的问题。第一个是回收：一个节点被摘下来之后，你并不知道还有没有别的读者持有它的指针，于是不能释放。无锁链表要解决这一点需要 hazard pointer、引用计数或 epoch 回收，它们要么在读者侧引入额外的原子与轮询开销，要么把回收推迟到不确定时点、需要一个界定不了的“宽限期”。实时系统里最难论证的就是“不确定时点”。第二个是 ABA：一个指针值从 A 变成 B 再变回 A，CAS 认为什么都没发生；如果 A 这个地址是被释放后重新分配回来的同一个节点，判断就彻底错了。用静态池、且让槽位标识不复用（把轮次编进序号）可以缓解，但缓解不等于解决。有界数组回避了这两个问题：地址永远有效，槽位所有权由序号表达，永远不需要释放。代价是容量必须提前定，而这恰恰是我们本来就想要的那条信息。

\subsection*{四种同步形态的横向对比}

阻塞、自旋、无锁与纯原子访问这四类形态，区别不在吞吐而在最坏时间与失败模式。表 \ref{tab:ch09-sync-forms} 按可论证性把它们排开。

\begin{longtable}{>{\raggedright\arraybackslash}p{0.16\textwidth} >{\raggedright\arraybackslash}p{0.22\textwidth} >{\raggedright\arraybackslash}p{0.22\textwidth} >{\raggedright\arraybackslash}p{0.23\textwidth}}
\caption{四类同步形态的最坏时间、失败模式与适用上下文}\label{tab:ch09-sync-forms}\\
\toprule
形态 & 最坏时间性质 & 主要失败模式 & 适用上下文 \\
\midrule
\endfirsthead
\toprule
形态 & 最坏时间性质 & 主要失败模式 & 适用上下文 \\
\midrule
\endhead
\bottomrule
\endfoot
\endlastfoot
阻塞等待（信号量/条件变量） & 由持有者释放时刻决定，本身不可界；配上优先级天花板或继承后有界 & 优先级反转、等待时间不计入执行预算 & 任务上下文、非硬实时路径；中断里禁止 \\
\addlinespace
自旋锁 & 与持锁者被抢占的概率相关，未加协议时无上界 & 优先级反转、自旋消耗 CPU 反过来加剧负载 & 极短临界区且持锁段被显式界定；单核上通常不如关中断 \\
\addlinespace
无锁队列（CAS 循环） & 单次操作可界，但重试次数取决于竞争强度 & 活锁、回收问题、ABA & 多生产者或多消费者、允许显式丢弃的场合 \\
\addlinespace
SPSC/单字原子 & 一次存储或一次原子读，常数级 & 内存序用错导致读到旧值或半写值 & ISR 与任务、核与核之间最常用形态 \\
\end{longtable}

表 \ref{tab:ch09-sync-forms} 的排序逻辑是“能不能写进可调度性模型”：能写出有界阻塞项的形态才能用在硬实时路径上。多数项目的结论是：单核用关中断或 SPSC，跨核用 SPSC 风格的单向通道，真正的多对多需求尽量在设计阶段化简成多组单向通道。

\section{内存序、栅栏与 volatile 的边界}

\subsection*{acquire/release 与 seq\_cst 的取舍}

默认使用 \Tp{seq\_cst} 是安全的，但它的代价在弱序平台上不是零：顺序一致的原子读在很多架构上免费，顺序一致的原子写却常常伴随一条全屏障，在 ARM 上可能是十几条指令与若干百纳秒的存储缓冲排空。反过来，\Tp{acquire} 与 \Tp{release} 只保证“这一条数据之前的写都可见”，不要求所有观察者之间达成统一的全序。

决策规则很朴素。如果两个量之间需要所有观察者看到同一种先后顺序（典型是版本位与状态位组合、或一个开关位控制多个数据位的读取顺序），就需要 \Tp{seq\_cst} 或显式 \Fn{atomic\_fence}；如果只是“我发布的这个游标，保证带你回看我之前写的负载”，\Tp{acquire} 与 \Tp{release} 配对就够。本章的 SPSC 属于后者。需要全序的场景往往说明结构可以拆得更细，先想清楚是否真需要它。

\subsection*{relaxed 只够计数}

\Tp{memory\_order\_relaxed} 只保证原子性，不保证顺序。它可以用于统计计数、心跳、抢占次数这类“值本身有意义、与其他变量的相对顺序无意义”的量。把它用在数据发布上是错误：别人看到新计数，却可能看到你还没写进缓冲区的数据。清单 \ref{lst:ch09-spsc} 里对自己游标的那次 relaxed 读是合法的，因为只有本线程写它。

\subsection*{弱序平台上“屏障”意味着什么}

x86 的内存模型（通常称 TSO）只允许存储后移越过加载，即写缓冲让 load 先于 store 完成；它不允许 load-load 与 store-store 重排。因此在 x86 上，一条普通存储之后跟一条普通存储，别的核看到它们的顺序与发出顺序一致。很多“无锁代码怎么写都能跑通”的直觉来自这里。

ARM 与 RISC-V 是弱序：读写可以相互穿越，同一变量的两条写可以乱序可见，直到遇到显式屏障。ARM 提供 \Tp{DMB}（数据内存屏障，约束 load/store 的可见顺序）、\Tp{DSB}（更强的完成屏障，等待访问真正到达规定位置）与 \Tp{ISB}（指令屏障，用于改写寄存器或自修改代码之后）。RISC-V 用统一的 \Tp{fence} 指令，通过前后字母组合表达约束范围；\Lang{C} 与 \Lang{C++} 的原子操作在底层会展开成带回环的保留/条件序列（RISC-V 的 load-reserved 与 store-conditional，ARM 的独占加载/存储对）。结论是：“先写数据、后写标志”在弱序机器上真的会被反着看见——标志先到、数据后到，读者按新标志读到旧数据。

\subsection*{volatile 是编译器语义，不是同步}

\begin{note}[title={volatile 与 atomic 的分工}]
\Tp{volatile} 只约束编译器：每次访问都真的发出访存指令，不缓存进寄存器、不删除、不合并。它对处理器乱序、对多核可见性、对编译器把两次 volatile 访问相对于其他非 volatile 操作的移动都不提供保证，因此永远不能用来做线程同步。它有两个正当用途：其一是内存映射外设寄存器，那里“每次读都要发一次总线访问”是设备语义要求（位域读写见另一分册）；其二是与中断服务程序共享的变量——在裸机或 \DomRtos{} 语境里 ISR 与任务在同一个核上，编译器是唯一可能把值吃进寄存器的元凶，volatile 正好治它。同一个变量若需要跨核共享，就必须换成原子类型，而不是 volatile 与原子混用。
\end{note}

\begin{lstlisting}[style=cstyle,caption={C：与 ISR 共享的标志用 volatile，跨核发布用 stdatomic},label={lst:ch09-atomic-iso}]
#include <stdatomic.h>
#include <stdint.h>

/* 情形一：同一个核，任务与 ISR 共享。
   volatile 足够：需要的是"每次都重新从内存读"，
   不是"跨核看到顺序"。 */
volatile uint32_t g_tick_count;        /* 由 SysTick ISR 递增 */

void task_wait_next_tick(void) {
    uint32_t a = g_tick_count;
    while (g_tick_count == a) {        /* 必须真的重读，不能进寄存器 */
        /* 实际项目里这里应是有界等待 + 超时上报 */
    }
}

/* 情形二：跨核共享。必须用原子类型 + 内存序。
   volatile 在这里什么也不保证。 */
static _Atomic uint32_t g_shared_seq;        /* 核 0 发布 */
static uint32_t         g_shared_data[8];    /* 随序号一起发布 */

void publish_core0(void) {
    g_shared_data[0] = 0x1234u;              /* 先写数据 */
    atomic_store_explicit(&g_shared_seq, 1u,
                          memory_order_release); /* 后发布 */
}

int consume_core1(uint32_t *out) {
    if (atomic_load_explicit(&g_shared_seq,
                             memory_order_acquire) != 1u) {
        return 0;
    }
    *out = g_shared_data[0];                 /* 保证看到发布前的数据 */
    return 1;
}
\end{lstlisting}

\section{cache 行、对齐与数据放置}

\subsection*{伪共享的表现形式与检测}

伪共享的代价不体现在指令流里，因此极易被误判成“代码变慢了”。两个不相关的变量落在同一条 cache line（多数 ARMv8 与 x86 上是 64 字节），一个核写、另一个核读，或者两个核分别写各自那半条 line，就会持续触发 line 的所有权转移。从指令数上看什么都没发生，从时间上看每次访问都在跟互连和目录较劲。

检测手段：\texttt{perf stat -e cache-misses,cache-references} 看总体命中率是否异常，\texttt{perf c2c} 按 line 粒度报出出现 false sharing 的地址范围与竞争线程。裸机上没有这些工具时，把两个游标手工拆到相隔 64 字节以上的地址做 A/B 对比，是最可靠的判定：吞吐显著变化而代码语义未变，即伪共享。布局上的线索是：结构体里把不同线程写的字段混排；同一数组里高频更新项与只读项相邻；一个全局对象既有热字段又有冷字段。

对齐手段：\CPP{} 用 \Tp{alignas}（例如以 64 为参数），GCC/Clang 扩展用 \texttt{\_{\Bk}\_{\Bk}attribute\_{\Bk}\_{\Bk}((aligned(64)))}。两者都只在对象起始地址层面起作用：一条 line 可能被两个对象跨越，因此要保证独占，要么把对象大小补齐到 64 的倍数（加 padding 数组），要么让相邻对象各占整数条 line。静态数组的元素是连续排布的，逐元素对齐会把步长拉到 64 字节——这通常正是你要的，但内存占用要按这个算。

\subsection*{只读数据、常量表与噪声数据}

查找表、标定曲线、字符串常量都应落在 \texttt{.rodata}：只读页可以在多个核之间共享而不产生所有权迁移。真正要注意的是页粒度干扰——一个可写变量和一张大表同页时，写操作会让这一页在别的核上失效。把高频写的小变量与大块只读表分开排布，比逐一对齐更省事。

噪声数据（统计计数器、性能采样缓冲、调试开关）应该集中到专门区域并远离关键线。它们的共同特征是关键路径每次都要写它们一笔，于是把本不该同步的核拉进同一场 line 争夺。读写分离有两种形态：读者轮询被发布者写的 line，代价是这条 line 被多个核读走之后，写者下次写要先夺回所有权；读者订阅事件（写者完成后做一次通知），代价是通知本身（原子操作加一次跨核中断或核间邮箱）比一次读贵。当轮询频率远高于事件频率时，改成“读本地副本 + 低频同步”往往比优化原子操作有效得多。

\subsection*{打包、非对齐与共享布局}

结构体加打包属性后成员可能落在非自然边界上。对这样的对象做原子操作，多数实现要么编译期拒绝，要么生成硬件不允许的访问。更根本的问题是：原子性通常以“不跨越 line 边界的自然对齐访问”为单位承诺，跨越两条 line 的字段无法保证原子更新。因此协议结构要显式分两层：对外的字节布局可以紧凑打包，内部工作副本必须自然对齐、可原子，两者之间用显式的序列化与反序列化函数转换，绝不用指针强转。

非对齐访问的后果取决于架构。ARMv7-M 上多数非对齐访问硬件支持，但对设备类型内存的非对齐访问会产生 UsageFault；ARMv6-M 的加载存储指令不接受非对齐操作数，编译器只能拆成逐字节访问，慢且不可原子；ARMv8-M 可通过系统寄存器的非对齐控制位决定是否陷入。Cortex-A 上非对齐访问一般走硬件慢路径，但在设备内存与共享内存区域属于未定义行为。工程结论只有一句：与硬件或其他核共享的布局一律自然对齐，把对齐问题当会中断构建的错误而不是警告。

时间戳与计数器要按自身宽度自然对齐，并且如果它们是跨核读的原子量，还要满足单条 line、不跨界。8 字节时间戳在 32 位平台上不是天然原子的，两个半字的撕裂读会读出混合值；解决办法是 8 字节对齐加双字原子访问，或用“序号-值-序号”的复核读取。字节序是共享布局的另一处坑：同一份结构体在小端核与大端核（或某些外设 FIFO）之间解释不同。协议应按字节定义，用显式移位与或运算组装，并按偏移加长度声明，不要在协议边界上直接传结构体内存。

\section{静态分配与栈纪律}

\subsection*{最坏栈深度是可以算出来的}

对编译单元开启 \texttt{-fstack-usage}，每个函数产出一条自身帧大小，写在 \texttt{.su} 文件里，单位字节。把这些数据与调用图合并，就能求出每个入口（任务、ISR、外部可见函数）的最坏深度。GCC 的 \texttt{-fcallgraph-info=su} 在同一份输出里带上调用边与帧大小，配合脚本可生成调用树。工具的定位要谨慎：它们只对编译期可见的全部调用负责，函数指针、虚调用、库内部递归都在视野之外，所以结论必须搭配一条纪律：稳态路径不允许间接调用。

递归必须禁止，不是因为慢，而是因为深度由数据决定，无法静态界定。VLA（\Lang{C} 的变长数组，以及 \Lang{C++} 里作为扩展存在的那些）同理：它把栈深度和外部输入挂钩，等于在栈上装了个 \Fn{malloc}。中断栈与任务栈要分开，中断栈按“最深 ISR 加上它可能的嵌套与软件触发的级联”来定，而不是按平均。多数 \DomRtos{} 允许配置独立中断栈，或让 ISR 复用当前任务栈；后者会让中断栈预算等于任务剩余空间，一旦某个任务栈接近用满时来了中断，溢出会静默损坏相邻任务的上下文。

\begin{lstlisting}[style=cppstyle,caption={C++：静态对象池与 placement 构造（运行期零分配）},label={lst:ch09-pool-placement}]
#include <cstddef>
#include <cstdint>
#include <new>
#include <type_traits>

template <typename T, std::size_t N>
class StaticPool {
  static_assert(std::alignment_of<T>::value <= 64,
                "element alignment must fit the line stride");

  alignas(64) unsigned char storage_[N * sizeof(T)];
  unsigned char in_use_[N];
  std::size_t live_;

public:
  StaticPool() noexcept : storage_{}, in_use_{}, live_(0) {}

  // 取一个槽：只做标记与 placement 构造，没有分配器参与。
  T* acquire() noexcept {
    for (std::size_t i = 0; i < N; ++i) {
      if (in_use_[i] == 0u) {
        in_use_[i] = 1u;
        ++live_;
        // placement new：对象落在已有静态存储上。
        // 调用方必须显式 release，否则析构不会发生。
        return ::new (static_cast<void*>(&storage_[i * sizeof(T)])) T();
      }
    }
    return nullptr;            // 池耗尽：显式失败，见第 7 节
  }

  void release(T* obj) noexcept {
    if (obj == nullptr) {
      return;
    }
    obj->~T();                 // 与 placement new 成对
    const std::size_t i = static_cast<std::size_t>(
        (reinterpret_cast<unsigned char*>(obj) - storage_) / sizeof(T));
    in_use_[i] = 0u;
    --live_;
  }

  std::size_t live() const noexcept { return live_; }
  static constexpr std::size_t size() noexcept { return N; }
};
\end{lstlisting}

\subsection*{\CPP{} 特有的不可界路径}

\Fn{std::function} 的目标对象若比内联缓冲区大就会走堆，构造与赋值耗时随目标大小与是否命中内联存储而变化；它的调用还多一次间接跳转，对最坏执行时间的指令预测不利。虚函数表调用同样是间接性问题，本身不分配内存，但它让调用图不可静态闭合，栈与执行时间分析工具在此失效。异常路径需要展开信息与潜在的堆分配，处理时间高度不可界。三者共同构成“编译期看着省、运行期无法论证”的类别，因此安全关键配置通常直接开 \texttt{-fno-exceptions} 与 \texttt{-fno-rtti}，并在编码规范里禁止 \Tp{std::function} 出现在稳态路径。

全局构造的合法边界很窄：静态初始化发生在 \Fn{main} 之前，裸机或 \DomRtos{} 里它的时机取决于启动代码，可能早于外设与堆就绪，也可能在多个翻译单元之间形成不可见的顺序依赖。可以接受的形式只有常量初始化（编译期把值写进 \texttt{.rodata}，没有运行期代码）与零初始化。需要构造的对象放进显式的初始化函数，由启动序列按确定顺序调用。placement 构造的边界与之对应：只用于已存在的定长存储、成对显式析构、类型与对齐都已知，且不在 ISR 里做。

容器选型：定长数组与 \Tp{std::array} 是默认选择，上界写在类型里；环形缓冲见本章第 2 节；\Tp{std::vector}、\Tp{std::string}、\Tp{std::map} 在稳态路径上是禁忌——扩容、再平衡、节点分配都发生在最不该发生的地方。确实需要复杂索引时，应在配置阶段一次性构建，稳态只读。这一条按目标分档的判据（哪一档拿得到 \CPP{} 运行库、哪一档可以放容器）在第 1 章 \S\ref{sec:ch01-langtier}。

\begin{table}[htbp]
\centering
\caption{各类数据的内存策略（存放、分配时机与更新者）}
\label{tab:ch09-mem-policy}
\begin{tabular}{>{\raggedright\arraybackslash}p{0.19\textwidth} >{\raggedright\arraybackslash}p{0.21\textwidth} >{\raggedright\arraybackslash}p{0.21\textwidth} >{\raggedright\arraybackslash}p{0.21\textwidth}}
\toprule
数据类别 & 存放与分配 & 更新者与上下文 & 纪律要点 \\
\midrule
控制环状态量 & \texttt{.bss} 定长数组或池槽，启动期就绪 & 单一任务写，其余只读 & 单写者，跨核读走原子游标 \\
ISR 与任务交换的样本 & SPSC 环形缓冲，编译期 $2^{k}$ 容量 & ISR 写、任务读 & 满/空约定显式定义，丢弃要计数 \\
标定与查找表 & \texttt{.rodata}，链接期常量 & 配置阶段构建，稳态无人写 & 与可写变量分块，避免同页失效 \\
日志与跟踪缓冲 & 静态环，写入即忘 & 任意上下文写，后台任务读 & 环内只放定长记录，禁止格式化调用 \\
诊断与后处理 & 堆或大块静态区 & 低优先级任务 & 不参与稳态时间分析 \\
\bottomrule
\end{tabular}
\end{table}

\begin{bestpractice}[title={把静态容量变成可回归的数字}]
每个静态结构都该有一个可自动核对的边界：池的槽数、环形缓冲的 $2^{k}$、栈的字节预算。给它们配单元测试（在 \DomHost{} 上编译同一份头文件，断言 \texttt{sizeof} 与容量常量），并把 \texttt{-fstack-usage} 的输出纳入构建：任何一次让某个入口最坏深度超过上次记录的改动都应让构建失败。容量数字不写进测试，就只是注释里的期望。
\end{bestpractice}

\section{常见陷阱与失效形态}

\begin{compare}[title={内存池耗尽：静默丢帧还是复位}]
两种都可以是对的选择，但必须提前定下来并写进需求。静默丢帧要求丢失是可观测的（计数器加上层超时逻辑），并且被丢的恰好是控制量时不会让误差积累；复位或进安全态要求触发条件明确（连续耗尽，或某个关键池耗尽），并且恢复时间本身在可调度性预算内。不可接受的形态是：池耗尽时 \Fn{acquire} 返回空指针，调用方把它传给下层，或者干脆用一个从未检查的返回值继续执行。把耗尽当成一个必须被处理、被计数、被报警的输入，是这类结构的基本要求。
\end{compare}

下面五类问题在本章各节的交界处反复出现，值得单独列一遍。

\begin{enumerate}
\item 先写数据、后更新索引，但在弱序平台上没有用 release 存储：索引先到、数据后到，读者拿到旧值。这是本章最强调的一条，也是唯一在测试里几乎不能复现的一条。
\item 伪共享引起的性能回归被当成代码 bug 去重构。重构之后“变好了”，实际原因是两个变量被拆开了；如果不同时观察 cache miss 计数，你永远不知道自己做对了什么。
\item 把 \Fn{printf} 或日志留在临界区。格式化本身耗时不可界，日志后端可能有锁、可能写堆、可能写文件描述符，任何一条都足以让最坏执行时间分析作废。稳态路径上只允许往定长环里放一条二进制记录。
\item 对 volatile 目标做 \Fn{memcpy}。多数实现把源与目标当普通指针处理（甚至用向量整块移动），编译器允许把一次 volatile 访问拆成若干普通访问，从而丢掉设备要求的“每次一个访问”语义。设备访问要逐次读写。
\item 在锁或自旋等待里没有重试上限，只有忙循环。于是优先级反转从偶发变成常态，并且自旋消耗的 CPU 会推高负载，让等待更久——这两点在下一章的时间维度上还会再出现一次。
\end{enumerate}

还有一类错误是“结构对了、预算错了”：环形缓冲容量按平均帧率而不是最坏突发算，池的槽数按典型并发而不是最大在途，栈预算按最深的一次测试运行而不是静态最坏深度。本章所有结构的正确性都建立在“这些数字来自最坏情况”这一前提上；一旦数字来自平均值，代码写得再规范也无效。这与可调度性分析里的同类错误是同一个错误，第 10 章会在时间维度上正面处理它。

\subsection*{本章检查清单}

\begin{itemize}
\item 稳态路径上没有任何分配或释放调用；\texttt{.su} 输出与栈预算进入构建。
\item 每个共享结构都能回答三个问题：谁写、谁读、需要什么顺序；答案对应一个明确的 memory order。
\item 所有跨核原子量自然对齐、单 cache line、不与噪声字段同线。
\item 池耗尽、环形缓冲满、CAS 重试超限三类失败都有显式的、被测试过的处理分支。
\item 日志与格式化输出不在任何临界区里。
\end{itemize}

% ============================================================
% ---- frag_ch10.tex ----
\chapter{调度模型与可调度性分析}
\label{ch:ch10-schedulability}

上一章处理的是单个任务内部的时间与内存纪律：不分配、不自旋无界、不共享噪声数据。本章把这些量放到系统层面：许多任务同时存在，处理器只有一两个，必须回答“每个任务能不能在最坏情况下按时完成”。回答方式不是压测，而是建模与推导：给出任务模型、给出干扰来源、给出一个可以复算的判据，再用实测去校准模型里的参数。

需要强调的立场是：可调度性结论的价值不在于它证明系统能跑，而在于它把“为什么能跑”写成一组可审计的数字。任何一个数字被改掉（执行时间变长、周期变短、多了一个任务），结论立刻失效并被评审发现。这比“我们跑了三小时没问题”强得多。

\section{任务模型：周期、抖动与执行时间上界}

一个周期性任务的参数是：周期 $T$，相对截止时间 $D$（释放后多久必须完成），最坏执行时间 $C$（一次释放所需的最长处理器时间），偏移 $\phi$（第一次释放的时刻）。硬周期任务的每个参数都是硬约束；偶发任务（aperiodic）没有固定周期，但有最小间隔（带宽）与最坏执行时间，工程上通常把它折算成等效周期，或者用后台服务、带宽保留服务器把它变成可建模的形态。

任务释放的时刻并不严格落在网格上，两者之差就是到达抖动 $J$。$J$ 与 $\phi$ 是不同的东西：$\phi$ 是一次性的相位差，$J$ 是每次都存在的抖动窗口。带抖动的模型比不带抖动的模型悲观，但它符合观察——如果你的界里没有 $J$，实测的“偶发超时”就无法解释。

帧（frame）的概念在处理多速率链时很有用：一个任务实例处理的可能是上一帧的数据，于是它的截止时间比自己的周期更靠前或更靠后。端到端截止时间与单任务截止时间也不是一回事：控制链“采样到输出”跨越若干任务与队列，链上每一段都有自己的完成时刻要求。推导时必须按“哪一帧的数据在哪个时刻必须落地”来算，不能按任务自身的释放时刻算。多速率系统（内环 100 微秒、外环 1 毫秒）里，外环一次释放对应内环十次，这种帧对齐关系如果不在模型里写明，就会出现外环用的是两帧之前的电流反馈却没人报警的情况。

三个量的定义要精确：

\begin{itemize}
\item 执行时间上界是最坏执行时间，不是平均值，也不是观测到的最大值。三者关系是：平均值远小于最大观测值，最大观测值不超过真实最坏值，而真实最坏值通常只能被论证或保守估计。实测与估算工具如何配合，见本章“验证、度量与常见陷阱”一节与第 13 章。
\item 到达抖动 $J$ 是实际释放时刻与理想释放时刻的最大偏差，来自 tick 量化、上游任务本身的响应时间抖动、中断抢占。后果是等效地增加执行需求：保守处理是把 $J$ 计入响应时间界。
\item 执行抖动是同一任务相邻两次运行时间之差，主要来自 cache 与分支预测状态。它不改变 $C$ 的定义（$C$ 已是最坏），但决定实测分布的形状，因此实测数据要按上界归纳而不是按均值。
\end{itemize}

控制环里最常见的简化假设是 $D = T$，并且 $C$ 远小于 $T$。这条假设在两类情况下不成立：分帧任务（$D$ 不等于 $T$）、任务带偏移与抖动。写分析文档时应显式声明用了这条假设，并说明它为什么成立。“我们的位置环 1 毫秒周期、200 微秒截止时间”显然不满足 $D = T$，用错假设得到的结论看起来正确，实际完全无效。

还要区分三种约束：处理器时间约束（$C$ 与 $T$ 之比）、内存约束（第 9 章的静态容量）、外设约束（转换时刻、DMA 完成时刻）。前两类由调度器决定，第三类由硬件决定，端到端界必须把三者都算进去，否则会出现“任务按时结束、输出却晚了一个转换周期”。

\section{静态优先级与响应时间分析}

\subsection*{RMS 的最优性与 Liu-Layzon 判据}

在固定优先级、周期任务、$D = T$、任务间无资源竞争、彼此独立的前提下，速率单调调度（RMS，周期越短优先级越高）是最优的：任何能被某个固定优先级分配调度的任务集，也能被 RMS 分配调度。这条最优性结论经常被过度引用——它不适用于有资源共享、有偏移、$D$ 小于 $T$、有抢占开销的情形，那些情况下“按周期分配”只是常用而非最优。

Liu-Layzon 是充分判据：若 $U \le n(2^{1/n} - 1)$，则 $n$ 个独立周期任务在 RMS 下可调度。$n$ 等于 2 时界约为 0.83，随 $n$ 增大单调下降到约 0.69。两处需要注意。第一，它是充分非必要条件，不满足不代表不可调度；用它的理由是可口算，便于评审时快速判断“这个数量级的任务集还有没有余量”。第二，0.69 没有包含抢占开销、上下文切换、中断与系统任务负载，因此多数项目把“含全部开销的总负载不超过 0.5 左右”当实用门限——不是因为有定理支持，而是 0.69 到 0.5 之间那段被实测开销吃掉了。

周期集合的谐波性质会显著改变结论的可信度。当所有周期两两成整数倍（$T_i$ 整除 $T_j$ 或反之），称为谐波周期集，此时分析结果与实际行为差距小，保守界明显收紧；当周期比例落在 1 与 2 之间（例如 100 微秒与 150 微秒），低优先级任务在高优先级周期内可能被连续干扰多次，最坏情况比直觉恶劣。把非整数倍周期就近调成谐波（150 改成 200 或 100）是成本极低、收益明确的工程手段。

\subsection*{抢占与切换开销要进模型}

利用率里的每一项都应当是“处理器实际被占用的时间”。除了任务自身的 $C$，还包括：每次抢占与恢复的上下文切换（寄存器保存、栈与映射状态变化，在带 MMU 的核上可达数微秒到几十微秒）、每次定时器到期时调度器维护就绪队列的工作、每次中断入口与出口。这些开销与抢占次数成正比，因此“把周期改短、把任务拆细”会让开销项非线性上涨。工程上常见的做法是把每周期固定开销并入 $C$（即模型里的执行时间等于纯计算加一次切换加一次定时处理），并在周期缩短时重新计算总负载——只看算法复杂度而不看调度开销的成本模型，是过载的常见起点。

\subsection*{响应时间分析：迭代求解}

比利用率判据更强的是响应时间分析（RTA）。记 $hp(i)$ 为优先级高于任务 $i$ 的任务集合，最坏响应时间满足

$R_i = C_i + B_i + J_i + \sum_{j \in hp(i)} \lceil (R_i + J_j) / T_j \rceil C_j$

各项含义如下。$C_i$ 是任务自身的执行需求。$B_i$ 是优先级反转类阻塞项，来自非抢占临界区与更低优先级任务持有的资源：在有优先级继承或优先级天花板协议的实现里，它等于“所有能阻塞 $i$ 的临界区长度之和”的最大值；这个值由第 9 章讨论的同步形态决定——那里四类形态里只有能写出有界阻塞项的，才能在这一项里填一个可信的数字。$J_i$ 是到达抖动，它让第一个实例被推迟。求和项是高优先级任务在 $R_i$ 窗口内可能发出的释放次数乘以各自执行时间，向上取整因为不足一个周期也算一次释放。方程右侧含 $R_i$，因此从 $C_i + B_i + J_i$ 起步迭代，直到值不再变化（收敛，得到界）或超过截止时间（不可调度）。

\begin{warning}[title={阻塞项不能取零}]
中断服务程序的长度、不可抢占的内核临界区、禁用中断的自旋段、其他核上的长临界区都要计入 $B_i$。在许多实际系统里 $B_i$ 与 $C_i$ 同数量级：一个 20 微秒的周期任务，摊上 15 微秒的最坏中断屏蔽窗口，响应时间界接近翻倍。把 $B_i$ 取 0 只有一个前提——系统里不存在任何超过一个指令周期的不可抢占区，这在通用内核或带文件系统的 \DomRtos{} 上几乎从不成立。
\end{warning}

\subsection*{算例：电流环、位置环与通信}

三个任务的参数与迭代结果见表 \ref{tab:ch10-example}：电流环 100 微秒、位置环 1 毫秒、通信 10 毫秒。阻塞项取最坏中断屏蔽窗口 5 微秒加跨核共享临界区，到达抖动按一个 tick 的量化误差计入。

\begin{table}[htbp]
\centering
\caption{三任务参数与响应时间迭代结果（时间单位：微秒）}
\label{tab:ch10-example}
\begin{tabular}{>{\raggedright\arraybackslash}p{0.19\textwidth} >{\raggedright\arraybackslash}p{0.15\textwidth} >{\raggedright\arraybackslash}p{0.16\textwidth} >{\raggedright\arraybackslash}p{0.16\textwidth} >{\raggedright\arraybackslash}p{0.15\textwidth}}
\toprule
任务与优先级 & $T$ 与 $C$ & 阻塞与抖动 & 收敛值 $R$ & 余量 \\
\midrule
电流环（最高） & $T$ 为 100，$C$ 为 20 & $B$ 为 5，$J$ 为 0 & $R$ 为 25 & 75 \\
位置环（中） & $T$ 为 1000，$C$ 为 60 & $B$ 为 8，$J$ 为 5 & $R$ 为 93 & 907 \\
通信（最低） & $T$ 为 10000，$C$ 为 150 & $B$ 为 8，$J$ 为 10 & $R$ 为 304 & 9696 \\
\bottomrule
\end{tabular}
\end{table}

中间一步可以手算验证位置环：起值 $C + B$ 等于 68；高优先级电流环在 68 微秒窗口内释放 1 次，加 20 得 88；再计入自身抖动 5，得 93；用 93 重算窗口，电流环释放次数仍是 1 次，值不再变化，收敛。利用率是 $U$ 等于 0.20 加 0.06 加 0.015，即 0.275，远低于 $n$ 等于 3 时的 0.78，两条判据都通过。这张表真正要读的是“结论依赖哪些输入”：把最坏中断屏蔽窗口从 5 微秒改成 60 微秒，位置环立刻越界；把电流环周期从 100 改成 90，位置环窗口内的释放次数从 1 变 2，界随之跳一档。

\begin{lstlisting}[style=cstyle,caption={C：固定优先级响应时间迭代（含高优先级任务与阻塞项）},label={lst:ch10-rta-iterate}]
#include <stdbool.h>
#include <stddef.h>
#include <stdint.h>

typedef struct {
    uint32_t period_us;    /* T */
    uint32_t exec_us;      /* C：最坏执行时间 */
    uint32_t jitter_us;    /* J：到达抖动 */
    uint32_t deadline_us;  /* D：相对截止时间 */
} Task;

/* 求任务 idx 的最坏响应时间。数组前 idx 个任务是更高优先级。
   返回 true 表示收敛且不超过 deadline。 */
bool rta_worst_response(const Task *t, size_t idx, uint32_t blocking_us,
                        uint32_t *out_r) {
    const uint32_t limit = 64u;      /* 迭代上限，防死循环 */
    uint32_t r = t[idx].exec_us + blocking_us + t[idx].jitter_us;

    for (uint32_t it = 0; it < limit; ++it) {
        uint32_t demand = r;
        for (size_t j = 0; j < idx; ++j) {
            /* 窗口 r + J_j 内高优先级任务的释放次数，向上取整 */
            uint32_t w = r + t[j].jitter_us;
            uint32_t n = (w + t[j].period_us - 1u) / t[j].period_us;
            demand += n * t[j].exec_us;
        }
        if (demand == r) {           /* 不动点，已收敛 */
            *out_r = r;
            return r <= t[idx].deadline_us;
        }
        if (demand > t[idx].deadline_us) {
            *out_r = demand;         /* 越界，报告超量部分 */
            return false;
        }
        r = demand;
    }
    *out_r = r;
    return false;                    /* 未收敛，应增大上限并复查参数 */
}
\end{lstlisting}

清单 \ref{lst:ch10-rta-iterate} 的用途不是替代工具，而是让“结论怎么来的”这件事进入版本库：参数表喂进去，界算出来，任何一次改动都能被重跑。

\section{动态优先级与 EDF}

\subsection*{利用率可达 1 与它的实现代价}

最早截止时间优先（EDF）在单处理器、独立任务、允许抢占的模型下是充分的：只要 $U$ 不超过 1 就可调度。相比 RMS 的 0.69，这条界在满载场景（多轴、多环的伺服控制器）能省出实质余量。

代价在实现与失效形态两侧。实现侧：动态优先级意味着每次释放都要按截止时间排序，就绪队列的插入与取出不再是常数级的位图操作，其复杂度与任务数相关；优先级计算本身的开销要计入模型。当两个任务的绝对截止时刻相同，还需要一条确定的平局规则，否则行为随实现细节漂移。

Linux 主线的 EDF 类调度器长期停留在研究补丁状态，工程上可依赖的是内核的 \Tp{SCHED\_DEADLINE} 类：它提供带准入控制的受限 EDF（周期补充预算风格），准入时检查剩余带宽，参数由系统调用设置，运行期由内核维护预算补充。使用它的要点有三：准入失败会立即报错，这比静默过载好；每任务的 runtime、deadline、period 三元组必须与分析一致，否则内核的保证和你的保证是两回事；runtime 应填分析里的 $C$，而不是观测到的平均值——填错等于让内核按错误的预算给你带宽。

\begin{lstlisting}[style=cppstyle,caption={C++：SCHED\_DEADLINE 带宽参数设置骨架（Linux）},label={lst:ch10-sched-deadline}]
#include <sched.h>
#include <sys/syscall.h>
#include <unistd.h>

#include <cstdint>

// 用 sched_setattr 把线程设为 SCHED_DEADLINE。
// 约束：runtime 不超过 period，period 不超过 deadline，单位纳秒。
// runtime 就是你向内核承诺的最坏执行时间，应等于分析里的 C。
bool set_deadline_budget(std::int64_t pid, std::uint64_t runtime_ns,
                         std::uint64_t deadline_ns,
                         std::uint64_t period_ns) noexcept {
    sched_attr attr{};
    attr.size = sizeof(attr);
    attr.sched_policy = SCHED_DEADLINE;
    attr.sched_runtime = runtime_ns;
    attr.sched_deadline = deadline_ns;
    attr.sched_period = period_ns;
    attr.sched_flags = 0;          // 不禁止在该线程上派生子进程

    // 返回值检查要点：EINVAL 多为三元组关系不成立；
    // EBUSY 表示准入控制拒绝（该 CPU 的剩余带宽不足）；
    // EPERM 需要相应特权。
    return ::syscall(SYS_sched_setattr, pid, &attr, 0) == 0;
}
\end{lstlisting}

\subsection*{过载时的失效形态}

选 EDF 之前必须想清楚的一点是它的失效形态：过载时 EDF 不惩罚“最该被牺牲”的任务，而惩罚最接近截止时间的那个。持续过载下出现的是级联错过截止时间——一个任务超时使它的绝对截止时刻推后，推后之后它在下轮更优先，于是挤掉更多任务。这不是固定优先级下的“某个低优先级任务饿死”（那种失效范围明确、可预测、可设计成预期行为），而是全局同时劣化，症状出现得晚而扩散得快。安全关键的环因此很少完全交给 EDF。

\begin{lstlisting}[style=cppstyle,caption={C++：绝对周期任务骨架（避免相对 sleep 累积漂移）},label={lst:ch10-absolute-period}]
#include <cerrno>
#include <cstdint>
#include <ctime>

// 用绝对唤醒时刻推进周期，而不是每轮相对 sleep。
// TIMER_ABSTIME 语义：超时时刻是绝对时间，被推迟只会让本轮晚触发，
// 下一轮的目标时刻不受影响；相对 sleep 则把延迟累进后续所有轮次。
bool wait_until(std::uint64_t absolute_ns) noexcept {
    timespec ts{};
    ts.tv_sec  = static_cast<long>(absolute_ns / 1000000000ull);
    ts.tv_nsec = static_cast<long>(absolute_ns % 1000000000ull);

    const int rc = ::clock_nanosleep(CLOCK_MONOTONIC, TIMER_ABSTIME,
                                     &ts, nullptr);
    if (rc == 0) {
        return true;
    }
    if (rc == EINTR) {
        return true;   // 被信号打断：绝对时刻仍在，继续等即可
    }
    return false;      // 其余错误：目标时刻非法或时钟源配置有误
}

// 用法：网格时刻由周期累加得到，绝不使用"上一轮结束时刻 + 周期"。
//   std::uint64_t at = first_release_ns;
//   for (;;) {
//       if (!wait_until(at)) { report_wakeup_fault(); }
//       run_control_step();
//       at += period_ns;             // 相位锚定在网格上
//   }
\end{lstlisting}

\begin{compare}[title={混合模型：关键环固定优先级加旁路 EDF}]
常见折中是把硬实时部分（电流环、位置环、保护逻辑）放在固定优先级的最高层，保证过载时降级顺序明确、可静态论证；把软实时部分（通信栈、日志、状态上报）放到低优先级或带宽保留的 EDF 类里，允许其截止时间与执行需求自适应。要点是两层边界要清晰：硬层永不被软层抢占，软层可用带宽等于 1 减去硬层的 $U$。混用一个调度器又不分层，等于把两者的最差性质叠在一起。
\end{compare}

\section{抖动、偏斜与端到端延迟预算}

\subsection*{到达抖动的三个来源}

第一是 tick 量化：内核时间片粒度决定了定时器最早与最晚触发时刻之差。一个 1 毫秒周期的定时器在 4 毫秒 tick 的机器上会有接近一个 tick 的抖动；这也是使用高精度定时器与 tickless 配置时必须重新检查响应时间界的原因。第二是相对周期：用 sleep 相对时长推进，每轮的唤醒延迟都累进下一轮的绝对时刻，误差可以一直漂移；用绝对时刻（\Tp{TIMER\_ABSTIME} 形式的超时，或 \DomRtos{} 里给出下一个绝对唤醒时刻的等待接口）会把延迟吸收而不累积。清单 \ref{lst:ch10-absolute-period} 就是这个区别所在。第三是中断抢占：ISR 恰在任务被唤醒前的一刻运行，就把唤醒推后一整段 ISR 时间，这段必须同时出现在 $J$ 与 $B$ 里，不能两处都省。

\subsection*{执行抖动、偏斜与总线仲裁}

执行抖动来自微架构状态：cache 与 TLB 命中与否、分支预测、指令预取。它不改变最坏执行时间，但决定实测分布的形状——这也是为什么 P50 与最大值相差很多在实时系统里是常态，而不能用分位数做预算。

偏斜（skew）指多个任务或多个采样通道之间的相位关系。两个周期相同但相位差小于一个执行时间的任务，会持续以固定顺序抢占彼此；采样通道之间的偏斜则造成“两个量本该同一时刻取，实际差了一个转换时间”。把周期做成整数倍关系并显式安排偏移，是消除偏斜影响的常规做法。共享总线与内存控制器的仲裁会把这种相位关系进一步放大：一个通道的批量访问可能压住另一个通道的单次访问。这属于多核干扰的特例，与本章“多核”一节的问题同源。

\subsection*{端到端延迟的构成与预算}

一条典型链路是：传感或外设事件，到 ISR，到无锁队列，到任务读取，到计算，到输出写入。逐项量级（在几十到几百兆赫兹的核、配置得当的 \DomRtos{} 上的典型区间）：外设到中断响应（含硬件延迟与向量跳转）常在亚微秒到低几微秒；ISR 本身几微秒到几十微秒；任务从被唤醒到实际运行的排队时间由调度决定，可能是几微秒，也可能是一整段抢占延迟；队列与拷贝在亚微秒到几微秒；计算由算法决定；输出（写寄存器或 DMA 启动到外设生效）常是微秒级再加外设自身的转换时刻。这些数字只用来给预算划范围，项目里必须以实测与论证值替换。

\begin{bestpractice}[title={延迟预算的分配方式}]
把端到端截止时间拆成每段上界之和，并给每段一个“承诺值”而不是“观测值”。做法：按链路顺序列出各段，写入分配给它的预算，要求实现者在该预算内完成并给出证据（实测最坏值或分析界）；总预算留 20 到 30 个百分点的余量，用于未来新增项与无法建模的部分（频率调节、内存刷新、温度降频）。任何一段超预算时，评审要求的是“从别的段买回时间”或“显式吃掉总余量并记录”，而不是“平均没问题”。预算表要能自动核对：每段插桩时间戳，把实测最坏值与承诺值一起输出到测试报告。
\end{bestpractice}

\section{过载与降级}

\subsection*{负载系数与显式降级策略}

负载系数就是利用率 $U$，它既是可调度性分析的输入，也是过载定义的度量：$U$ 超过 1 意味着不存在任何调度方式能满足全部需求，此时必须回答“牺牲什么”。可选的显式策略有四类。丢帧：当前实例直接放弃，计数上报，适用于幂等的传感器输入。降采样：周期加倍、负载减半，适用于可以慢一些的外环。跳过可选段：把任务内部拆成必须段与可延时段，过载时只跑必须段，这是最平滑的降级，也是唯一能保持相位不变的策略。进安全态：当降级也不能保证正确性时，唯一正确的动作是停止输出并进入已知安全状态。

不可接受的形态是没有策略：过载时由调度器自行决定谁延后，系统表现为“某些时刻某些输出偶尔晚一点”，而这正是最难复现、最难验证的故障。把过载策略写进需求（例如“连续丢失 $n$ 帧则报警”“$U$ 超过 0.85 时自动降采样”）并做成可测试的行为，是这类系统的常识要求。

\subsection*{背压、队列深度与任务结构}

队列满时的处理（第 9 章的耗尽策略）在时间维度上就是背压：生产端要么被阻塞，要么丢弃。阻塞给生产端引入了模型里没有的等待时间，丢弃给消费端引入数据缺失，两者都要在选择时说明后果。队列深度与任务周期一起决定端到端延迟：深度越大抗突发能力越强，但最坏排队时间也越长。这个折中要显式算：$N$ 格缓冲在消费者周期 $T$ 下增加的最坏排队不超过 $N \cdot T$，或者按生产与消费速率之差算排空时间。不要让它在延迟预算里悄悄存在。

任务合并与拆分改变调度开销结构。合并减少切换次数与队列数量，但拉长单次不可抢占段，从而增大别人的 $B$；拆分让界更紧、并行度更高，但每次切换开销与每段的中断响应都要计入。判据是：如果两个段属于同一控制帧且必须同时成立，合并；如果一个段的超时不应拖垮另一个，拆分。

\subsection*{配置项对结论的影响}

内核配置与 CPU 配额会直接推翻分析结论，必须逐项核对。抢占模型（例如 \texttt{CONFIG\_PREEMPT\_RT} 与普通抢占）决定 $B$ 的量级；tick 频率（\texttt{CONFIG\_HZ}）决定 $J$ 的下界；tickless 空闲机制减少周期性 tick 干扰，但唤醒路径要重建时钟与调频状态，可能引入若干微秒到几十微秒的额外延迟，这条延迟必须进 $B$ 或唤醒项。CPU 配额同样关键：Linux 用 \texttt{sched\_rt\_runtime\_us} 与 \texttt{sched\_rt\_period\_us} 表达“每周期实时任务最多可用多少时间”，默认值意味着每个周期有约 5 个百分点的强制让出。如果分析假设实时任务可以用满整核，结论与运行行为就不一致；实时产品部署时这个配额通常被显式放开并写进运维文档，而不是保留默认值。频率调节（空闲降频、温度降频）会让同一段代码的执行时间成倍变化，所以 $C$ 要在最低允许频率下测。

\section{多核：分区、干扰与界的写法}

全局调度（所有核共享一个就绪队列，任务可迁到任何核）在多核实时上有天然困难：迁移代价（cache 与 TLB 状态丢失、目标核上的排队）难以界定；抢占次数上界在 $m$ 个核上的分析结果远差于单核，已知界会随核数迅速变松；负载不均又让任务在核间反复找位置。工程实践因此普遍选择分区（绑核）调度：任务静态分配到核，每核独立跑一个单处理器调度器，可调度性分析退化为“每核一个单核问题”。代价是利用率下降（一个核的空闲不能被另一个核的任务使用），换来的是可以论证。

同步屏障与 gang scheduling（一组任务在所有核上同一轮开始）用于必须相位对齐的场合，例如按帧同步的多轴插补。两者都把单核问题变回多核同步问题：屏障意味着最慢的核决定整组进度，gang 的轮转要求每核在同一时刻做同一件事，实现代价高、界难写，只在确实需要相位锁定的系统里用。

\begin{longtable}{>{\raggedright\arraybackslash}p{0.16\textwidth} >{\raggedright\arraybackslash}p{0.22\textwidth} >{\raggedright\arraybackslash}p{0.22\textwidth} >{\raggedright\arraybackslash}p{0.23\textwidth}}
\caption{四类调度形态的时间保证、实现要求与适用场景}\label{tab:ch10-alg}\\
\toprule
形态 & 时间与界 & 实现要求 & 适用场景 \\
\midrule
\endfirsthead
\toprule
形态 & 时间与界 & 实现要求 & 适用场景 \\
\midrule
\endhead
\bottomrule
\endfoot
\endlastfoot
RMS 固定优先级 & 利用率界约 0.69，或用响应时间迭代；过载时失效范围明确 & 优先级静态、常数级选择开销；需阻塞协议配合 & 硬实时控制环，任务数少且可论证 \\
\addlinespace
EDF 动态优先级 & 利用率界可达 1；过载时级联失效 & 每次释放按截止时刻排序，需确定平局规则 & 多速率、软实时混合，允许降级策略 \\
\addlinespace
\Tp{SCHED\_{\Bk}DEADLINE} & 受限 EDF 加带宽补充，界由内核预算保证 & 内核版本支持、准入控制、三元组填对、需特权 & 通用内核上给单个关键环留带宽 \\
\addlinespace
分区多核 & 每核一个单核问题，界可组合；跨核干扰另计 & 静态亲和性、共享资源有界、缓存或带宽划分 & 多核硬实时，实际项目的主流做法 \\
\end{longtable}

跨核干扰主要来自共享资源：共享末级缓存、内存控制器、总线仲裁、互联。典型形态是噪声邻居——一个核上的批量拷贝把内存带宽占满，另一个核上原本 3 微秒的访存变成 12 微秒，于是后者的最坏执行时间不再是它自己代码的性质，而是整个系统的性质。缓解手段包括：容量划分（缓存分组、内存控制器带宽配额）、把干扰源放到不与关键任务共享层次结构的核上、把共享资源占用当成一种资源用优先级天花板类协议管理。优先级继承与天花板在多核上需要跨核实现（等待方常以自旋而非阻塞方式等待，因为持锁者所在的核可能在跑别的任务），而自旋又消耗 CPU，把等待时间转换成别人的干扰——这正是多数项目回避跨核共享可写状态的原因。

AMP（非对称多处理，通用核与实时核并存）与 Hypervisor 分区从另一方向规避干扰：把硬实时任务放进只有它、内存区域静态划分、外设直通的处理域，可调度性分析回到“独占硬件上的单处理器问题”。代价是跨域通信必须经过显式共享内存或邮箱通道，而这些通道的延迟与吞吐要建模，且往往成为新的 $B$ 项来源。

\begin{guideline}[title={多核响应时间界怎么写才站得住}]
实际项目的界通常写成三层。其一，每个核内部用本章第 2 节的固定优先级响应时间分析，前提是任务已绑核且不迁移。其二，跨核部分只承认被明确建模的干扰项：一个有界的共享临界区、一个有界的内存访问延迟上界、或一个已划分的缓存分区，并把它们计入 $B$ 或 $C$。其三，用实测校准：在最高干扰配置下跑满负载，取实测最坏值与静态界的较大者作为最终界。任何“多核直接用单核公式”的写法都不成立；“我们绑核了所以没有干扰”同样不成立，因为共享末级缓存与内存控制器仍然共享。
\end{guideline}

\section{验证、度量与常见陷阱}

\subsection*{实测执行时间分布怎么用}

插桩测量给出的是分布：P50、P99、最大值。P50 用来发现回归趋势，P99 用来判断优化效果，最大值用来与 $C$ 比较——但最大观测值不等于最坏执行时间，它只是最坏值的一个下界。因此“可调度性结论要不要用实测最大值”的答案是：可以把它当作 $C$ 的下限依据，但不能作为唯一依据。稳妥做法是取实测最大值加上一个来自干扰模型的裕度（未测到的分支、未覆盖的最坏输入、cache 冷状态、最低频率），并说明裕度来源；若需要正式的 $C$，要依靠静态最坏执行时间分析工具或带硬件计数器的路径覆盖论证，这条线与工具定位在第 13 章展开。

\begin{warning}[title={六个反复出现的分析错误}]
(1) 用平均执行时间做预算，或把优化后的 P99 当成 $C$。(2) 忽略中断与系统任务的负载：内核线程、软中断、看门狗、遥测都消耗同一核的时间，$U$ 里没有它们，实测就出现无法解释的缺口。(3) 优先级分配与数据依赖方向相反——下游任务优先级高于上游，于是下游每轮都在等一个被它抢占的上游，形成链式等待；正确方向是让生产者（数据上游）优先级更高，除非这条反转被显式建模并计入界。(4) 把 tickless 或调频的唤醒延迟留在界外，导致实测比模型差而无法解释。(5) 在无锁结构或自旋锁里忙等，自旋消耗 CPU、推高负载，反过来让等待更久，这是第 9 章的退化路径与本章的利用率在算同一件事。(6) 算利用率只算任务执行时间，忘了每次切换与每次定时器处理的开销，也忘了把它们计入别人的阻塞项。
\end{warning}

\subsection*{度量手段的最小集合}

负载方面看 \texttt{/proc/stat}（总体与每核时间分配）与 \texttt{schedstat}（每任务累计运行时间与调度等待，需在配置里打开相应开关）。调度事件方面看 \texttt{perf sched} 的记录与汇报，它给出每次切换的等待、迁移与运行段，是识别抢占与长等待的主要来源。中断方面看向量级计数与最长服务时间。裸机与 \DomRtos{} 上的对应物是 GPIO 翻转加逻辑分析仪、内核周期计数器打点、空闲任务占用率（空闲比例等于 1 减去负载）。选哪一套不重要，重要的是所有度量都要输出最坏值与出现时刻——只输出平均值的数据对可调度性没有价值。

\begin{lstlisting}[style=shellstyle,caption={Shell：调度延迟与切换热点的采集},label={lst:ch10-measure}]
# 记录一段时间的调度事件，随后离线分析等待与切换分布
perf sched record -d 30 -- ./rt_application
perf sched latency
perf sched timehist

# 每任务累计运行时间、调度次数与等待时间（需内核暴露 schedstat）
cat /proc/PID/schedstat
grep -H . /proc/PID/sched

# 每核负载与中断分布：确认是谁在消耗那个核
mpstat -P ALL 1 10
cat /proc/interrupts
\end{lstlisting}

\subsection*{把可调度性写成可回归的门限}

分析结论应以机器可读形式进版本库：任务参数表（$T$、$C$、$D$、$J$、优先级、亲和性）、判据计算脚本、结论（每个任务的 $R$ 与余量）。任何一次改动动了这张表，构建就重算并在余量低于设定值时失败。与第 9 章的栈与容量回归门限配合，构成“时间与内存两条静态论证都在构建里跑”的最低形态。实测最坏值同理：把每个任务的上界存成基线，超出基线的改动必须显式承认。

清单 \ref{lst:ch10-rta-iterate} 与表 \ref{tab:ch10-example} 给出本章的可复算形态：给定参数、迭代、与截止时间比较。把这四步做成构建里的一条检查，可调度性就不再是文档里的段落，而是一个会失败的断言。

\subsection*{本章检查清单}

\begin{itemize}
\item 每个任务都有 $T$、$C$、$D$、$J$ 与优先级，并注明 $D = T$ 假设是否成立。
\item $C$ 来自最坏情况（实测最大值加裕度或静态界），不是平均值或分位数。
\item $B$ 有来源清单：中断屏蔽窗口、不可抢占区、跨核临界区，逐项可查。
\item 切换与定时处理开销已并入模型，周期变更时重新核算。
\item 过载策略是显式的（丢帧、降采样、跳过可选段、进安全态之一），且可测试。
\item 端到端延迟按段分配预算，实测最坏值与承诺值一并进报告。
\item 多核结论建立在绑核加有界干扰建模之上，而不是单核公式的照搬。
\end{itemize}

% ============================================================
% ---- frag_ch11.tex ----
\part{测量、确定性与工程化}

\chapter{抖动测量与运行时取证}
\label{ch:ch11-tracing}

前面几章讲的是机制：第 3 章拆过中断延迟的构成，第 4 章给了看门狗与超时的职责边界，第 8 章约束了实时抢占补丁的行为，第 10 章用可调度性分析算出过理论上界。本章谈另一件事：这些结论如何被量化地确认。工程上 ``实时性达标'' 不是形容词，而是一组可复跑的采集序列、一个写清楚的统计口径，以及一条超限即阻断的门限。缺少这三件的结论，在评审时与没测过等价。

\section{测什么：延迟、抖动、偏斜与错过率}

先把四个量区分清楚，本章所有工具的读数都挂在它们上面。

\textbf{延迟}指事件从被触发到被处理的间隔，必须写明 ``从哪到哪''。常用的四种口径互不等价：中断触发到中断服务程序首条指令入口（入口延迟）、定时器到期到线程真正获得 CPU（唤醒延迟）、请求报文发出到本机收到（链路加协议栈）、传感器采样到执行器输出（端到端）。报告里只写 ``延迟 40 微秒'' 是没有信息量的。

\textbf{抖动}是实际间隔与名义间隔的偏差。周期任务通常报三个数：峰峰值（最大值减最小值）、标准差、以及超出某个窗口的样本数。抖动的成因是 ``间隔的散布''，与延迟的均值可以完全无关：一个平均 15 微秒、偶尔 800 微秒的路径比一个稳定 60 微秒的路径危险得多。

\textbf{偏斜}描述两个时钟或两个采集点在同一时刻的读数差，是跨核、跨节点量；同一测点上的时间散布是抖动。二者混淆会导致同步类缺陷被当成调度类缺陷去修。多通道同步采集的判据只看偏斜。

\textbf{错过率}是超出门限的样本数除以总样本数，可写成比例或百万分比值。它比单个极值更适合做门限，因为它对时长不敏感、可跨轮次累加。

单位上建议同时报两种：微秒便于与硬件计数器和位时间换算，而控制环只关心 ``错过几个周期''，因此以周期数为单位的错过率更容易在跨平台比较时被读懂。换算时把任务周期记作 $t_{cycle}$，错过的连续周期数记作 $N_{miss}$。

统计口径要连同窗口一起声明：

\begin{itemize}
  \item 分位数报 P50、P99、P99.9 与 max；只报 P50 或平均值等于没测尾部。
  \item 直方图桶宽决定尾部分辨率。经验做法是让桶宽不超过关注门限的二十分之一：门限 100 微秒时用 4 或 5 微秒的桶；用 50 微秒的桶会把 120 微秒和 190 微秒的尾部并成同一格，看起来 ``没有超标''。
  \item 样本数决定分位数的稳定性。P99 至少需要几千个样本才有意义，P99.9 需要几万；单线程以 1000 微秒周期跑 60 秒只有六万个样本，看似很多，但要观测 P99.99 就远远不够。
  \item max 是极值统计量，与总时长强相关。60 秒的 max 与 4 小时的 max 不可比，也不可用短轮次的 max 去推断长轮次的上界。
\end{itemize}

\begin{warning}[title={用短轮次的 max 声明 ``最大抖动为 80 微秒''}]
极值的期望随观测时长单调增长。任何以 max 为内容的结论都必须同时给出单轮时长、轮次数与总样本数，并且把 ``本配置在 $D$ 时长内未观测到超过 $X$ 微秒'' 写成陈述式而不是上界式。真要论证上界，回到第 10 章的可调度性分析与屏蔽段枚举，测量只能证伪，不能证明。
\end{warning}

\section{Linux 侧的抖动工具族}

用户态测周期抖动的标准做法是绝对到期定时器加高优先级线程：设定下一个以 $t_{cycle}$ 为间隔的绝对到期时刻，等待，醒来后读时钟做差。差值即唤醒延迟，它包含定时器精度、中断处理、调度器决策与上下文切换。

\textbf{cyclictest} 是这一模式的实现，参数含义逐项如下。

\begin{itemize}
  \item \Tp{-p} 线程的实时优先级（\Tp{SCHED\_FIFO} 类）。基线一般用 95 或 99，与被测业务的优先级拉开一段，避免被测业务本身把测程挤掉。
  \item \Tp{-i} 名义周期，单位微秒；\Tp{-l} 循环次数；长选项 \Opt{duration} 以时长收口，二者选一个，推荐用时长，便于跨轮次对齐总样本数。
  \item \Tp{-m} 把进程全部内存锁住（\Fn{mlockall}），避免测程中发生活页与换页；不锁内存时，冷启动的几十次缺页会直接写进尾部。
  \item \Tp{-h} 输出直方图，参数是桶宽（微秒）。不带参数时的默认桶宽在不同版本上不一致，务必显式写。
  \item \Tp{-q} 抑制逐周期打印，只留汇总；\Tp{-t} 指定线程数，各线程分别统计；\Tp{-a} 把线程绑定到指定 CPU 列表，与内核的隔离核配置对齐。
  \item \Tp{-d} 起始散布：让多个线程错开最多若干微秒启动，避免所有到期时刻落在同一瞬间制造人为排队。
  \item \Tp{-S} 时钟源选择。默认用 \Tp{CLOCK\_REALTIME}，加 \Tp{-S} 改走 \Tp{CLOCK\_MONOTONIC}。测周期抖动时应选后者，因为 \Tp{CLOCK\_REALTIME} 会被时间同步的步进操作拽动，那部分偏差来自时钟维护行为而不是调度。
  \item \Tp{-r} 改用普通时间片策略运行，作为 ``实时优先级买到了什么'' 的对照；\Tp{-u} 改用相对睡眠接口，可用来验证绝对到期接口自身的开销差别。
  \item 部分版本还有把逐周期数据与直方图写入文件、以及涉及内存节点锁定一类的选项，名称在不同版本上不统一。跨版本脚本只依赖上面这些长期稳定的短选项，其余能力在现场用帮助输出核对，不要把不确定的长选项写进回归脚本。
\end{itemize}

\begin{bestpractice}[title={基线跑法：空载、加压、注入噪声三段}]
第一段在系统空载（无用户业务、控制台不刷新）下跑，得到地板值；此时若 P99 已经超门限，说明配置或硬件有问题，不必继续。第二段跑真实业务负载，这是唯一可以写进报告的数字。第三段人为注入噪声：内存与 CPU 压力、网络灌流、磁盘大量小写、以及一个高优先级抢占者。每加一类噪声就重跑一轮，尾部对哪一类敏感，成因就在哪一段。清单 \ref{lst:ch11-cyclictest} 给出的正是这套序列。
\end{bestpractice}

\begin{lstlisting}[style=shellstyle,caption={Shell：cyclictest 空载、加压与噪声注入的采集序列},label={lst:ch11-cyclictest}]
# 目标：CPU 4-7 为隔离核，业务优先级 90，测量周期 1 ms，门限 P99 = 80 us
ISO=4-7; TMO=1000; BUCKET=4; DUR=600
mkdir -p /var/log/rtbench/run07 && cd /var/log/rtbench/run07

# (1) 记录环境：内核命令行、隔离核、调速器、空闲态、中断亲和
uname -a > env.txt
cat /proc/cmdline >> env.txt
cat /sys/devices/system/cpu/isolated >> env.txt
for c in /sys/devices/system/cpu/cpu*/cpufreq/scaling_governor; do
  printf "%s %s\n" "$c" "$(cat $c)" >> env.txt
done
cat /proc/interrupts > interrupts-before.txt

# (2) 空载地板值：单线程，绑到隔离核末位，锁内存，用单调时钟
cyclictest -t1 -p95 -i"$TMO" -D"$DUR" -a7 -m -S -h"$BUCKET" -q \
  > base-quiet.txt 2>&1
mv histogram.txt hist-base-quiet.txt

# (3) 真实负载：与控制业务同优先级区间，多线程分别统计
cyclictest -t4 -p95 -i"$TMO" -D"$DUR" -d200 -a"$ISO" -m -S -h"$BUCKET" -q \
  > load-real.txt 2>&1
mv histogram.txt hist-load-real.txt

# (4) 噪声注入：每类单独加，加一类跑一轮，文件名里写清噪声来源
stress-ng --vm 2 --vm-bytes 2G --cpu 4 --timeout "$DUR" &
cyclictest -t4 -p95 -i"$TMO" -D"$DUR" -a"$ISO" -m -S -h"$BUCKET" -q \
  > load-mem.txt 2>&1; mv histogram.txt hist-load-mem.txt
iperf3 -c 10.0.0.2 -t "$DUR" -P 4 &            # 网络灌流
fio --name=dd --rw=randwrite --bs=4k --size=1G --runtime="$DUR" &

# (5) 对照：普通进程视角，看隔离与优先级到底买到了多少
cyclictest -t1 -r -i"$TMO" -D"$DUR" -a2 -m -S -h"$BUCKET" -q \
  > base-other.txt 2>&1
\end{lstlisting}

两个更专门的工具回答不同的问题。\textbf{oslat} 在用户态做忙等循环，不断读时钟，检测相邻两次读数之间被 ``拿走'' 的最长时间，因此它测的是进入延迟与屏蔽段长度，不产生周期负载，适合找最坏单次事件。它的调用形式是：\Tp{-c} 指定被测核列表，\Tp{-D} 给出时长，\Tp{-l} 给出告警阈值（单位微秒，对应的长选项名在不同版本上带 threshold 字样），超过阈值的样本连时间戳逐条打印，另有把 trace 落盘的选项。\textbf{hwlatdetect} 依赖内核的硬件延迟 tracer：被测核在关中断状态下高速轮询计数器，凡是读数跳变超过一个量化步长的事件都记下来。因为它不依赖任何中断或内核代码路径，这是唯一能把系统管理模式引起的延迟暴露出来的手段，代价是测程中该核完全不可用于业务。它的时长、阈值与报告文件都用长选项给出（\Opt{duration} 取 120、\Opt{threshold} 取 10、\Opt{report} 指向一个文本文件）。

第 8 章之后出现的 \Tp{rtla} 系列把这两类能力整合进一个前端：定时延迟子命令给出 ``中断关闭段加线程唤醒段'' 的分解直方图，并有实时统计、历史直方图、阈值捕获三种运行形态；噪声子命令给出每核噪声总量与成因归类；硬件延迟子命令对应上一段的 tracer；另有观察空闲态进出代价的子命令。相对旧工具的优势是延迟分解直接落在 trace 事件上，不必自己算差值。发行版若尚未打包，退回旧工具组合即可，结论不受影响。

实时调优守护进程提供现成配置档：实时档一次完成频率调速器、中断线程化、脏页比例与换页倾向等几十项设置，网络延迟档更保守。用 \Tp{tuned-adm profile realtime} 切换，用 \Tp{tuned-adm active} 与 \Tp{tuned-adm verify} 核对当前生效项。它的价值不在效果而在可复现：报告里写配置档名，比手写一长串 \Tp{/proc/sys} 修改更容易被复核。注意它不等于内核启动参数，\Tp{isolcpus}、\Tp{nohz\_full}、\Tp{irqaffinity} 这类只能来自引导命令行。

\begin{note}[title={工具选择取决于你怀疑哪一段}]
怀疑调度器与中断线程：cyclictest 与 rtla 的定时延迟。怀疑进入延迟或某段关中断过长：oslat。怀疑固件级的不可见延迟：hwlatdetect。怀疑具体内核代码路径：ftrace 定向抓取。四者的读数不能互相替代，也不能互相验证。
\end{note}

\begin{longtable}{>{\raggedright\arraybackslash}p{0.16\textwidth} >{\raggedright\arraybackslash}p{0.34\textwidth} >{\raggedright\arraybackslash}p{0.28\textwidth}}
\caption{抖动与取证工具覆盖的延迟段及其自身开销}\label{tab:ch11-tool-coverage}\\
\toprule
工具 & 覆盖的延迟段 & 自身对测程的扰动 \\
\midrule
\endfirsthead
\toprule
工具 & 覆盖的延迟段 & 自身对测程的扰动 \\
\midrule
\endhead
\bottomrule
\endfoot
\endlastfoot
\Tp{cyclictest} & 定时器到期到用户线程返回：唤醒加调度段 & 一到四个实时线程，构成周期负载；绑核与优先级设置本身就是结论的一部分 \\
\Tp{oslat} & 忙等读数被拿走的最大间隔：进入延迟与屏蔽段 & 忙等满占被测核，无周期负载；对被测核之外的干扰不敏感 \\
\Tp{hwlatdetect} & 一切让计数器跳变的事件，含固件级屏蔽 & 测程中该核关中断轮询，完全不可用；阈值以上的事件才可归因 \\
\Tp{rtla} & 分解视角：中断关闭段、唤醒段、进程段 & 每核一个 tracer 线程加事件缓冲，开销约百分之几 \\
ftrace & 内核路径上任意两点之间，可带栈回溯 & 事件量与缓冲大小决定开销；缓冲不足会静默丢掉尾部 \\
\Tp{trace-cmd} & 完整事件时间线，可离线看尾部 & 落盘期间开销高于纯 ftrace，事后需做符号解析 \\
\Tp{perf} & 采样视角的调度事件与硬件事件 & 采样频率线性决定开销；不适合作为逐周期证据 \\
GPIO 加逻辑分析仪 & 请求到中断入口到完成的物理端到端 & 零软件侵入，代价是需要引脚、触发条件与偏斜标定 \\
\end{longtable}

\section{ftrace 与运行时追踪：把尾部拆开}

抖动工具告诉你尾部有多长，tracer 告诉你这段时间里 CPU 在做什么。tracer 类型按用途分：\Tp{irqsoff} 记录关中断时间最长的片段并给出调用栈，是 ``屏蔽段过长'' 的直接证据；\Tp{preemptoff} 形态相同但统计禁止抢占段，在无界循环与自旋锁场景最有用；\Tp{wakeup} 与 \Tp{wakeup\_rt} 记录任务被唤醒到真正上 CPU 的间隔，后者只统计实时类，噪声远小于前者；\Tp{function} 与 \Tp{function\_graph} 分别给出调用序列与嵌套耗时视图，用来确认某条路径上的实际顺序与耗时分布。

控制面集中在追踪目录下：\Tp{current\_{\Bk}tracer} 选形态，\Tp{tracing\_on} 起停，\Tp{buffer\_{\Bk}size\_{\Bk}kb} 定每核环形缓冲，\Tp{trace} 读全局结果，\Tp{per\_{\Bk}cpu/cpuN/trace} 读单核结果，\Tp{tracing\_{\Bk}max\_{\Bk}latency} 保存当前最大值。选项目录里常用的是 \Tp{stacktrace}（为片段附栈）、\Tp{func\_{\Bk}stack\_{\Bk}trace}（为命中函数附栈）、\Tp{record-cmd}（记进程名与编号，否则栈无法归因到任务）、\Tp{context--{\Bk}info}（在缓冲头写启动环境）、\Tp{sleep\_time} 与 \Tp{display--{\Bk}graph}（调用图形态的显示控制）、\Tp{overwrite}（翻转覆盖模式，见后面的对比框）。读取 \Tp{trace} 不会清空缓冲，显式写空才清空，因此现场取证的标准动作是先停测再复制。完整路径很长，例如 \Tp{/sys/kernel/debug/tracing/per{\Bk}\_cpu/cpu7/trace}，脚本里应当用变量拼出来，而不是在正文里折断它。

\begin{lstlisting}[style=shellstyle,caption={Shell：ftrace 切换 tracer、设置过滤与抓取尾部现场},label={lst:ch11-ftrace}]
T=/sys/kernel/debug/tracing
mount -t debugfs none /sys/kernel/debug 2>/dev/null

# (1) 关掉一切、清空、设置缓冲与记录选项
echo 0 > $T/tracing_on
echo > $T/trace
echo 8192 > $T/buffer_size_kb
echo 1 > $T/options/record-cmd
echo 1 > $T/options/context-info

# (2) 找最长的关中断片段：irqsoff 自己保留最大实例
echo irqsoff > $T/current_tracer
echo 1 > $T/options/stacktrace
echo 1 > $T/tracing_on
sleep 300
echo 0 > $T/tracing_on
cp $T/tracing_max_latency /tmp/maxlat.txt
cat $T/per_cpu/cpu7/trace > /tmp/irqsoff-cpu7.txt   # per_cpu 读法不重排时间线

# (3) 只关心某条路径：调用图形态加函数白名单
echo > $T/set_ftrace_filter
echo "blk_mq_*" >> $T/set_ftrace_filter
echo function_graph > $T/current_tracer
echo 1 > $T/options/sleep_time
echo 1 > $T/tracing_on; sleep 20; echo 0 > $T/tracing_on
cat $T/trace > /tmp/graph-blk.txt

# (4) 事件视角：中断入口加调度切换，按 CPU 过滤以压低开销
echo 1 > $T/events/irq/irq_handler_entry/enable
echo 1 > $T/events/sched/sched_switch/enable
echo "cpu == 7" > $T/events/sched/sched_switch/filter
echo 1 > $T/tracing_on

# (5) 收尾：一切复位，避免 tracer 常驻改变系统行为
echo 0 > $T/tracing_on
echo nop > $T/current_tracer
echo 0 > $T/options/stacktrace
echo > $T/set_ftrace_filter
\end{lstlisting}

\Tp{trace\_pipe} 是同一份缓冲的另一个视图：读它会阻塞在管道上并把事件搬到用户态。它便于交互式观察，但每次读取都有内核到用户的复制代价，而消费速度低于生产速度时事件会堵在内核里。测程正在跑实时业务时不要用 \Tp{trace\_pipe}，改为缓冲落盘后离线分析。

\Tp{trace-cmd} 把记录与阅读从现场解耦：其 record 子命令带事件列表与输出文件名，后接一段定时（例如跑 120 秒）生成二进制数据文件，report 子命令指定输入文件事后展开为文本，底层依赖 \Lib{libtraceevent} 做事件格式解析。record 期间开销明显高于单纯 ftrace，因为它还要落盘。数据文件可直接用 KernelShark 打开，看多核甘特图与空洞；导出为浏览器可加载的时间线格式（Chrome tracing）也有现成转换手段，适合把一段 20 毫秒的现场交给同事复核。注意记录进程本身是普通优先级，落盘压力会叠加进测程。

\Tp{perf} 与调度统计提供另一视角：\Tp{perf sched record} 定时长后用 \Tp{perf sched latency} 得到每个任务的等待时间分布，适合回答 ``业务线程被谁抢占''；\Tp{perf record} 订阅 \Tp{sched:sched\_{\Bk}switch} 与 \Tp{sched:sched\_{\Bk}wakeup} 事件并用 \Tp{-C} 限定核，可与 trace 时间线交叉核对；\Tp{/proc/schedstat} 与每个进程的对应文件给出运行时间、等待时间、切换次数与时间片长度的累积计数，开销几乎为零，适合长期观察趋势。

\begin{compare}[title={tracer 即扰动：三种缓冲取舍}]
\textbf{覆盖模式}（默认）：事件写满后新事件替换旧事件，好处是尾部一定保留最近内容，坏处是启动早期与你要找的罕见现场可能被覆盖。\textbf{快照与停止模式}：命中条件即停，保住现场，代价是要预先知道触发条件（\Tp{irqsoff} 属于这一类，它内部保留最大实例）。\textbf{加大缓冲换保真}：每核缓冲翻倍，内存占用与首次访问的缺页也翻倍，预热期变长，还可能把内存压力带进结论。选择原则是先确定要保住哪一段时间，再决定模式与大小。报告里必须写出缓冲大小与是否发生丢失：缓冲不足时 ftrace 会在时间线上留下溢出提示行，看不见这行就等于结论不可信。
\end{compare}

符号解析是事后最容易出错的一步：内核需要编译进内核符号表支持，报告里出现纯地址说明现场没有保留符号表或镜像不匹配。把 \Tp{/proc/kallsyms}、未裁剪的内核镜像与调试符号一并归档，是取证流程的一部分而不是可选项。

\section{内核与硬件层的取证}

中断侧先看 \Tp{/proc/interrupts}：逐列对每个 CPU 的中断计数，测程前后各采一次做差，就能算出 ``隔离核上应当为零的外设中断'' 是否真的为零，非零的中断号是首要嫌疑对象。定向核对与调整在 \Tp{/proc/irq} 下：\Tp{/proc/irq/N/smp\_affinity}（十六进制位掩码）是期望值，\Tp{/proc/irq/N/effective\_affinity} 是内核实际接受的可投递核；两者不相等说明亲和请求被中断域或被驱动托管的中断拒绝（网卡多队列中断常常由驱动自行管理，写不进去）。若系统上有中断均衡守护进程，它会在测程中周期性重排亲和，必须先停并写进报告。

频率与低功耗状态要逐核核对现值而不是配置：\Tp{scaling\_governor} 应为性能类或固定频率调速器，\Tp{scaling\_cur\_freq} 与 \Tp{scaling\_min\_freq} 相等才说明没有中途降频；\Tp{/sys/devices/system/cpu/cpuN/cpuidle/stateM/usage} 的增量说明测程中进入了哪些空闲态，最深状态的退出延迟常常是尾部的第一个来源。频率提升开关与平台的性能倾向接口（若存在）也要一并记。

系统管理模式与机器检查级别的固件事件是取证中最难的一类：处理器进入该模式时，关中断状态与看门狗都可能被挂起，内核没有任何事件可以记录它，tracer 的时间线表现为 ``什么都没发生但时钟前进了'' 的空洞。识别手段只有两条：一是 \Tp{hwlatdetect} 或 oslat 这类不依赖中断的轮询计数，二是把固件里与功耗管理、节能、外设模拟、可靠性巡检相关的项目关掉后重跑对比。部分平台提供读取固件计数器的工具可以统计进入次数，但这些工具高度依赖厂商，报告里应写 ``已通过固件设置关闭并复测''，而不是 ``已确认不存在此类事件''。

机器检查与可纠正错误类事件在长时间尾部异常中常被忽略。装有可靠性错误记录守护进程时，用其控制命令的摘要与错误两个长选项取近期事件，或直接查它的持久化数据库；老平台上用机器检查日志工具。一次内存纠错引发的重发或降速，足以让某条路径慢上几十微秒，而它在调度层面完全无痕。

指令级取证在带跟踪宏单元与 CoreSight 的 Arm 平台可行：\Tp{perf record} 以 \Tp{cs\_etm} 事件加通配目标、用 \Tp{-C} 限定核、定时 20 秒，记录的是指令流与分支目标，事后可解出函数序列甚至基本块级路径，代价是数据量大、解码依赖完全匹配的镜像。这与调试章节讲的是同一条硬件，用途不同：一处用来单步，一处用来事后回放。

系统调用与库调用跟踪工具在实时路径上不可用。两者都依赖进程跟踪机制，被跟踪进程每次系统调用都要停下来接受一次父进程裁决，而父进程只是普通优先级；实测单调用开销在几十到几百微秒，并引入完全不可界的等待。用它们看非实时的管理进程无妨，在测程中对业务线程下手，得到的抖动全是观测者制造的。要看业务线程的系统调用，改用 ftrace 的原始系统调用事件类或 \Tp{perf trace}，它们的记录在内核里完成，不阻塞被观察者。

\section{裸机与 RTOS 侧的取证}

无操作系统与实时操作系统场景没有内核可 trace，取证靠三样东西：周期计数器、物理引脚、外部观察者。

周期计数器打点最直接。Cortex-M 的调试单元里有周期计数器寄存器，它跟随内核时钟计数、不需要额外外设，读出两个事件之间的差值除以主频就是周期数，这与第 2 章讲过的调试外设是同一个硬件，这里只把它当时间戳用。注意三点：需要先使能调试外设与计数器；读计数器本身要几条指令的代价，标定必须包含它；某些系列的计数器在低功耗模式下停止，此时读数会给出错误的短时长。

\begin{lstlisting}[style=cstyle,caption={C：DWT 周期计数器与 GPIO 翻转的组合打点},label={lst:ch11-cyccnt-gpio}]
/* 打点：进入中断服务程序立即翻转调试引脚，并用 CYCCNT 记时间戳 */
#define RT_TRACE_PIN_PORT   GPIOB
#define RT_TRACE_PIN_NUM    9u          /* 逻辑分析仪接这一根 */

volatile uint32_t dwt_trace_ts;

void trace_enable_dwt(void)
{
    CoreDebug->DEMCR |= CoreDebug_DEMCR_TRCENA_Msk;   /* 使能调试外设 */
    DWT->CYCCNT = 0u;
    DWT->CTRL |= DWT_CTRL_CYCCNTENA_Msk;              /* 使能周期计数器 */
}

void RTT_IRQHandler(void)
{
    rt_gpio_set(RT_TRACE_PIN_PORT, RT_TRACE_PIN_NUM);   /* 上升沿：请求已到 */
    uint32_t t_in = DWT->CYCCNT;
    dwt_trace_ts = t_in;
    /* ... 真正的工作 ... */
    rt_gpio_clear(RT_TRACE_PIN_PORT, RT_TRACE_PIN_NUM); /* 下降沿：处理完成 */
}
\end{lstlisting}

GPIO 翻转加逻辑分析仪是唯一对被测系统零软件扰动的端到端测量，对固件级屏蔽、对观测代码本身、对时钟漂移都免疫。标定过程要写进报告：把打点引脚与触发源同接一台仪器，测出翻转延迟与中断入口相对硬件信号的偏移，后者正是入口延迟。仪器侧注意采样率与被测周期之比（至少十倍）、通道间偏斜（同排线探头要一起做偏斜标定，否则测到的偏斜来自探头）、以及触发保持时间。

单线调试输出经仪器跟踪内存端口把字符与时间戳送到观察器，调试探针的服务程序配合其软件视图可显示任务切换与打点时间线。优势是不占串口、延迟低；劣势是带宽有限，日志一多就丢，而且探针连接会改变低功耗行为。带宽预算大致由内核时钟除以分频系数再除以每字位数决定，据此确定打点密度。

实时操作系统侧还有商业与开源的追踪扩展可选：它们在关键位置埋钩子宏（形如 \Tp{traceTASK\_SWITCHED\_IN}、\Tp{traceTASK\_SWITCHED\_OUT}、\Tp{traceBLOCKING\_ON\_QUEUE\_RECEIVE}，以及临界区进入退出的钩子），由主机端图形工具拼出每任务时间线并算出利用率。启用这类扩展要在配置里打开对应开关，代价是每次切换多几十条指令，也就是说，挂着追踪器测到的抖动不再代表未挂追踪器时的系统。负载统计可以用内核自带接口：\Fn{vTaskGetRunTimeStats} 需要打开运行时间统计开关并提供一个高频计数器的读取回调，输出每任务占用百分比；\Fn{uxTaskGetSystemState} 给出每任务状态与栈高水位，开销更小，适合长期在线。用前者时务必确认计数源不是系统节拍，否则分辨率只有一个节拍，短任务的占用全部归零。

\section{结果分析与回归门}

拿到两份直方图后先做形状对比，不做数字对比。基线与实验组的主峰宽度一致、只是右侧多出一个孤峰，说明新增了一条独立的慢路径；整体右移说明共有的前置段变长（频率、缓存、锁）；主峰内出现等间隔梳齿，通常是周期性活动（定时器、时间同步、内存回收）与测程周期拍频。桶宽不足以分辨形状时，先减小桶宽重跑，再谈结论。

把 ``哪一段被拉长'' 映射到成因，本章工具的组合大致如下：屏蔽段（关中断或禁止抢占）由 \Tp{irqsoff}、\Tp{preemptoff} 与 oslat 的读数支持；跨核迁移由调度切换事件与绑核配置核对（迁移一次通常增加一次核间中断与一段缓存冷的代价）；缺页与内存回收由锁内存是否真正生效、以及换页类事件计数判定；降频与空闲退出由频率现值与空闲态使用量增量判定；观测代码自身的影响由 ``关掉 tracer 重跑'' 这一对照实验判定。每条成因都必须既有支持它的读数，又有一条把它排除掉的对照。

门限要写成可自动判定的一段脚本加一个数字表。推荐四项：P99 上限、max 上限、超限样本计数上限（通常为零）、错过率上限。判定失败即阻断，不做 ``略微超一点可以接受'' 的软处理；确有宽限需求时把它写成显式的第二级门限，并在报告里同时列出两级结果。

\begin{lstlisting}[style=pythonstyle,caption={Python：直方图读取、分位数计算与回归门判定},label={lst:ch11-hist-gate}]
#!/usr/bin/env python3
# 用法: hist_gate.py BASELINE CURRENT --p99-max 80 --max-max 250
# 退出码: 0 通过, 1 未通过, 2 输入或口径问题
import argparse, sys
from pathlib import Path

def load_hist(path):
    """cyclictest -h 输出：每行 '桶下界 计数'，多线程以空行分隔，此处合并。"""
    b = []
    for line in Path(path).read_text(errors="ignore").splitlines():
        p = line.split()
        if len(p) == 2 and p[0].isdigit() and p[1].isdigit():
            b.append((int(p[0]), int(p[1])))
    if not b:
        sys.exit(2)                      # 一个桶都没解析到：口径或版本变了
    b.sort()
    return b

def pct(b, q):
    total = sum(c for _, c in b)
    run = 0
    for lo, c in b:
        run += c
        if run * 100 >= q * total:
            return lo
    return b[-1][0]

def summarise(name, b):
    nz = [lo for lo, c in b if c > 0]
    return {"name": name, "n": sum(c for _, c in b), "p50": pct(b, 50),
            "p99": pct(b, 99), "p999": pct(b, 999), "max": nz[-1] if nz else 0}

def main():
    ap = argparse.ArgumentParser()
    ap.add_argument("baseline"); ap.add_argument("current")
    ap.add_argument("--p99-max", type=int, required=True)
    ap.add_argument("--max-max", type=int, required=True)
    ap.add_argument("--ratio-warn", type=float, default=1.5,
                    help="P99 相对基线劣化超过此倍率即判失败")
    a = ap.parse_args()

    base = summarise("baseline", load_hist(a.baseline))
    curr = summarise("current", load_hist(a.current))
    for s in (base, curr):
        print("%-9s n=%-9d p50=%-5d p99=%-5d p99.9=%-6d max=%-5d"
              % (s["name"], s["n"], s["p50"], s["p99"], s["p999"], s["max"]))

    fail = []
    if curr["p99"] > a.p99_max:
        fail.append("p99 %d us 超过门限 %d us" % (curr["p99"], a.p99_max))
    if curr["max"] > a.max_max:
        fail.append("max %d us 超过门限 %d us" % (curr["max"], a.max_max))
    if base["p99"] and curr["p99"] > base["p99"] * a.ratio_warn:
        fail.append("p99 相对基线劣化 %d/%d 超过 %.2f 倍"
                    % (curr["p99"], base["p99"], a.ratio_warn))
    if curr["n"] < base["n"] * 0.9:
        fail.append("样本数不足九成，两次采集的口径不一致")
    for m in fail:
        print("FAIL:", m)
    sys.exit(1 if fail else 0)

if __name__ == "__main__":
    main()
\end{lstlisting}

在持续集成里跑实时性测试有三条硬约束。第一，共享执行机不可用：虚拟化层、邻居负载与云平台的电源管理都会把尾部推到与产品无关的量级，任何可复现的门限都必须跑在独占、可描述硬件上；若只有共享机可用，就只跑形状检查（直方图是否出现新峰），不跑绝对数值门限。第二，时长与预算冲突：一个能可靠覆盖 P99.9 的轮次常常要几十分钟，做法是提交级跑短轮次只做劣化检测，夜间或发布前跑长轮次与噪声注入，两级用不同门限数字并在脚本里分别命名。第三，环境漂移要能被察觉：把内核版本、启动参数、调速器、隔离核列表、中断亲和与固件版本一起落盘成工件，并在判定脚本里加一项口径检查（清单 \ref{lst:ch11-hist-gate} 里就是样本数之比那一条）。

样本量与时长之间不是线性替换关系：提高测程频率能在同样时长里得到更多样本，但更短的周期本身改变了系统行为（缓存更热、中断更密），因而得到不同的尾部。结论必须绑定周期，换周期就要重测，不能靠换算。

\section{取证报告与常见陷阱}

一次取证的证据最小集如果缺项，两个月后就无法复现，也就无法用于裁决。报告不必长，但每一项都要能被别人对着机器核出来。

\begin{center}
\begin{tabular}{>{\raggedright\arraybackslash}p{0.20\textwidth} >{\raggedright\arraybackslash}p{0.30\textwidth} >{\raggedright\arraybackslash}p{0.28\textwidth}}
\toprule
必填项 & 取值来源 & 为什么必须留 \\
\midrule
内核版本与启动命令行 & 版本查询命令与 \Tp{/proc/cmdline} & 隔离核、无节拍、中断线程化等只能来自引导参数 \\
CPU 型号、微码与频率调速器现值 & 拓扑目录与 \Tp{cpufreq} 接口 & 同一份代码在不同步进与频率策略下尾部差别可达数倍 \\
空闲态与使用量增量 & \Tp{cpuidle} 下各状态的计数字段 & 深空闲退出延迟是尾部的常见来源，需要证明它是否发生过 \\
隔离核列表与中断亲和 & \Tp{/proc/interrupts} 前后差、期望掩码与实际掩码 & 隔离核上有外设中断时，结论与调度器无关 \\
固件关键项与节能设置 & 现场截图或导出的配置 & 判断固件级屏蔽是否已被排除，以及换板后能否复现 \\
负载来源与噪声注入命令 & 脚本本体与业务版本标识 & 空载数值不能代表产品数值 \\
原始直方图与逐轮次数据 & \Tp{cyclictest} 的分桶输出与追踪数据文件 & 分位数可由原始桶重算，截图与表格不可复核 \\
工具与脚本版本 & 各工具的版本子命令、脚本校验和 & 选项语义与默认桶宽跨版本会变 \\
门限数值与判定结果 & 判定脚本的打印与退出码 & 让 ``通过'' 成为一次可重跑的机器判定而不是人工叙述 \\
\bottomrule
\end{tabular}
\end{center}

现场复现的准备同样是报告内容：说明仪器接线与触发条件、逻辑分析仪采样率与探头偏斜标定值、以及需要在什么时刻按下什么按钮。无操作系统场景尤其重要，因为固件版本号一致的板子可能熔丝配置与晶振负载都不同。

常见陷阱按代价从高到低排：

\begin{itemize}
  \item 测量程序自己跑在非隔离核上，或与业务抢占者同核，测出的是别人制造的抖动。
  \item 控制台刷新、日志守护进程、周期计划任务与图形终端留在测程里。测前降低日志级别、停掉周期任务，并把这一步写进脚本而不靠记忆。
  \item 中断均衡守护进程没停，亲和配置在测程中被改掉；要核对实际生效掩码而不是你以为写进去的值。
  \item 只看平均值与 P50，或者把一次累积统计读数当结论：累积值抹掉了时间结构，无法定位尾部。
  \item 直方图桶宽过粗，尾部并入邻近格；或反过来把桶宽压到小于时钟分辨率，读出一堆无意义的梳齿。
  \item 忘锁内存、忘改单调时钟、忘做起始散布，三个漏项分别带来缺页、时间同步步进、以及多线程到期时刻重合。
  \item 打开 tracer 后没复位，下一次运行带着常驻 ftrace 与快照缓冲，得到 ``莫名变差'' 的基线。
  \item 拿一次运行结果做发布判定，没有重复轮次，因此无法区分偶发事件与配置后果。
\end{itemize}

\begin{warning}[title={观测者即扰动源：三条自检}]
(1) 关掉所有追踪再跑一次，尾部是否仍在；(2) 把日志与统计输出的落盘移出测程，尾部是否变化；(3) 把打点代码里的计数器读取替换成常量，测程是否变快。任何一条回答不了，本次数据只能用于定性趋势，不能进门限表。
\end{warning}

本章的三张清单——采集序列（清单 \ref{lst:ch11-cyclictest}）、tracer 现场（清单 \ref{lst:ch11-ftrace}）、判定脚本（清单 \ref{lst:ch11-hist-gate}）——按顺序执行即可覆盖一次完整取证。下一章转向系统之间的时间：当抖动来自链路与时钟而不是调度器时，这里的工具只会给出 ``什么都没发生但时间前进了'' 的读数，那时需要的是总线形态与同步机制。

% ============================================================
% ---- frag_ch12.tex ----
\chapter{实时以太网、现场总线与时间同步}
\label{ch:ch12-rt-bus-tsn}

上一章的抖动来自调度器与内核路径；本章的抖动来自链路、交换与时钟。当测量给出 ``什么都没发生但时间前进了'' 的读数，而内核侧配置已经无可再动时，问题多半在总线形态或时间同步上。本章只讲总线的实时形态与同步机制：寄存器级的外设访问与逐字节的报文操作属于另一分册，位域读写这里不重复。

\section{分类框架与决定实时性的四个量}

现场总线按访问方式分四类，这个分类决定了最坏延迟能不能算出来。

\textbf{轮询主站}：一个主站按固定顺序询问每个从站，介质访问由主站序列化，因此不存在竞争。最坏延迟是轮询表中排在自己之前的所有请求加响应时间之和，可静态计算。Modbus RTU、POWERLINK 的等时相、FlexRay 动态段之外的部分都属于这一族。代价是主站成为单点，且轮询表长度直接压低周期下限。

\textbf{事件触发竞争}：节点在任意时刻可以发送，靠仲裁决定先后。CAN 是唯一广泛使用的代表：仲裁按标识符逐位进行，数值小的标识符胜出，胜出者继续占用总线，失败者自动退避重发。优点是负载低时延迟极小、无主站；缺点是最坏延迟取决于同总线上有多少更高优先级的报文恰好同时待发，这必须靠报文数据库逐帧统计。

\textbf{生产者到消费者}：节点发布数据流，交换机或主站按流转发，接收方被动订阅。DDS、OPC UA 的发布订阅模式、Sercos 的通道模型都在这一类。它的实时性由流的分隔与整形决定，与拓扑相关，需要网络基础设施配合。

\textbf{时间触发}：发送时刻由调度表预先安排，每个节点在自己的时基上执行同一张表。FlexRay 的静态段、PROFINET 等时实时通道、TSN 的门控周期都是这一类。它给出最严格的确定性上界，代价是需要全局时间同步，且调度表应由工具求解而不是手写。

无论哪一类，决定实时性的量只有四个，选型与验证都围绕它们展开：

\begin{itemize}
  \item \textbf{周期}：控制环或采样帧的重复间隔。写报告时同时给微秒与周期数，跨平台比较用后者更直观。
  \item \textbf{负载率}：一个周期内占用的传输时间除以周期长度。任何共享介质都有稳定运行上限，工业 CAN 的经验做法是把设计负载压在三分之一到二分之一之间；超过之后继续加带宽只改善平均值，不改善尾部。
  \item \textbf{最坏仲裁延迟}：竞争类总线独有的量，等于最坏情况下需要等待的帧数乘以单帧传输时间，再加帧间间隔。整形与门控类总线上，这个量由调度表里的保护带与窗口承担。
  \item \textbf{同步精度}：多节点共享时基的偏差，决定同一物理时刻在不同节点上的读数差，也就是上一章定义过的偏斜。
\end{itemize}

\begin{compare}[title={四种范式各自的上界来源}]
轮询范式的上界来自轮询表长度；竞争范式来自待发报文集合；生产者到消费者来自流的划分与交换缓冲；时间触发来自调度表与同步误差之和。写最坏延迟论证时，先确认面对的是哪一种，再决定要枚举什么。把时间触发总线当竞争总线去压负载率，或反过来给轮询总线加优先级队列，都是无效工作。
\end{compare}

\section{车载与工业串行总线的实时形态}

\textbf{CAN 2.0} 的一位由若干时间量子组成：同步段、传播段、相位段一与相位段二，采样点落在相位段一与二之间；\Tp{sjw} 即再同步补偿宽度，决定每一位允许的再同步幅度，因而决定系统能容忍多大的晶振偏差与相位裕量。配置工具里这些段以时间量子为单位给出，\Tp{sample-point} 是仲裁段的采样点比例。两者设错表现为偶发位错误或错误帧计数上升，而不是持续失败，因此很难靠一次复现定位。

标准帧携带 11 位标识符与最多 8 字节数据，算上位填充后总位数在一百四十位以内；扩展帧因标识符为 29 位再加控制位，帧长明显更长。最坏等待延迟的粗算法：设同总线上优先级高于本帧且可能同时待发的报文有 $N_{wait}$ 帧，每帧传输时间为 $t_{frame}$，帧间间隔约三个位时间，则本帧最坏要等 $N_{wait}$ 加一之后再乘以 $t_{frame}$，另加间隔。这个式子只用来看量级：精确值应按报文数据库逐帧枚举。车载里这份数据库通常是 \Tp{DBC} 文件，描述每条报文的标识符、长度、周期与归属节点，也有专门工具据此算负载与延迟分布。

\textbf{CAN FD} 保留仲裁段的原速率（车载常见 1 兆位每秒），在数据段切换为更高符号率（常见 2 到 5 兆位每秒），把数据场扩到 64 字节，并换用更长的循环冗余校验字段（17 或 21 位，取决于是否使用不兼容 ISO 的格式）。要点在于：仲裁仍然慢，因此高负载下的最坏仲裁延迟不会因为数据段提速而按比例缩小；帧变长意味着单帧占用总线更久，一条长帧能挡住更多低优先级帧。FD 适合 ``信号多、周期紧、节点数适中''，不适合 ``节点极少、报文极短、要求微秒级响应''。

\textbf{LIN} 是主从结构的低成本串行线：主站按调度表逐个发送帧头（含标识符），被指派的从站回应数据场，因此完全没有竞争，最坏延迟就是调度表上前面各槽的响应时间之和。响应槽长度与调度表由描述文件生成，主站轮询周期决定信号刷新率。局限来自物理层：波特率通常在万位每秒量级，单线、节点数少、共享地线。

\textbf{FlexRay} 把周期切成宏周期，每个宏周期包含静态段（预留时槽，时间触发）、动态段（以最小槽竞争，事件触发）与网络管理段。静态段给出可证明的上界，动态段容纳突发；通道 A 与 B 可以双通道冗余发送，再叠加时槽冗余实现容错。节点间靠帧起始边界同步，同步集群有专门参数（宏周期、静态段长度、最小槽、各类容忍度），这些参数设错的表现是随机丢帧而不是全断。

\textbf{RS-485 上的 Modbus RTU} 是轮询范式的最低成本实现：主站一次只与一个从站交换，帧界靠至少三个半字符的静默间隔识别，半双工方向切换本身要留时间。周期下限的算法是：请求字符数加响应字符数，乘以每字符位数（含起始与停止位约十一位），除以波特率，再加从站处理时间与主站轮询开销。可用波特率的上界与线长成反比：反射与偏置电路在长线上限制边沿质量，实践里距离增加十倍往往意味着可用速率降到四分之一。要测实际有效周期，抓总线上的静默间隔比看应用日志可靠。

\begin{lstlisting}[style=shellstyle,caption={Shell：CAN 与 CAN FD 的观测、压力生成与负载测量},label={lst:ch12-can}]
# (1) 仲裁段 500 kbit/s、数据段 2 Mbit/s 的 FD 配置，采样点显式给出
ip link set can0 down
ip link set can0 type can bitrate 500000 sample-point 0.875 \
    dbitrate 2000000 dsample-point 0.80 fd on
ip link set can0 up
ip -details -statistics link show can0        # 核对生效参数与错误计数

# (2) 观测：-ta 给相对时间戳，只收自己关心的标识符区间
candump -ta can0,100:200 > /tmp/can-prio.txt &
candump can0 > /tmp/can-all.txt               # 全量抓，用于负载统计

# (3) 压力生成：固定标识符与数据长度，间隔 200 us，共 20000 帧
cangen can0 -I 320 -D 8 -g 200 -L 20000
cangen can0 -I 1FFFFFFF -g 100 -L 5000        # 29 位标识符即扩展帧，看仲裁影响

# (4) 收发一致性与负载估算
cansequence -p can0                           # 逐个标识符递增发送
cansequence -r can0                           # 接收并检查序号连续性
busload can0                                  # 按帧间隔反推负载率
cansniffer can0                               # 只显示发生变化的字节

# (5) 事后核算：一行一帧，帧数除以观测秒数得帧率，再乘每帧位数除带宽
wc -l < /tmp/can-all.txt
# 错误帧与总线关闭状态看驱动统计，抓包文本里没有这些信息
ip -details -statistics link show can0 | grep -E 'state|berr|error'
\end{lstlisting}

\section{工业以太网的实时形态}

这一节只讲实时形态与同步，报文格式与设备编程不在本册范围。

\textbf{EtherCAT} 的报文由主站发出，逐跳穿过全部从站：每个从站在帧经过时读出属于自己的数据、写入自己的数据，只引入一个转发延迟（硬件实现下可达微秒以下量级），不做存储转发。因此循环时间下界由三段相加：帧在线上走完整个拓扑的时间、各从站转发延迟之和、主站侧的调度与操作系统延迟。前两段在有硬件转发时基本可算，第三段属于上一章的取证范围。分布式时钟让每个从站用自己的硬件时间戳记录帧经过时刻，与各段传播延迟联立求解，把各节点的采样与输出锁到同一时基上，同步精度可达亚微秒。设计要点：拓扑顺序影响周期，线缆长度进入周期计算，线性与树形拓扑的延迟不同，热插拔会改变上界。

\textbf{PROFINET} 分两个实时等级。实时等级走标准交换与优先级，不保证延迟上界，适用于多数离散量与中等速率过程数据；等时实时等级在同步周期里为设备预留收发时刻，需要支持该等级的网络组件与统一参考时钟，并组成同步域（域内共享同一份循环计划，由控制器给出参考）。等时等级的循环时间下界由每周期可发送数据量与转发设备数量共同约束，规划要靠组态工具算出的通信映像。较新的路线是把它承载在 TSN 之上，作为一份通信配置档存在。

\textbf{POWERLINK} 用明显的两相结构：等时相里管理站点先发出发布帧（承载控制数据），随后按时槽依次订阅各节点数据，每个时槽独占发送窗口；异步相才允许标准以太网帧以竞争方式发送。等时相的时槽表决定确定性上界，异步相的拥塞不会污染等时相。全网用一个周期广播的时间信息对齐本地时基。

\textbf{Sercos III} 用通道模型：高优先级与低优先级通道再加一个用于监控数据的安全通道，实时通道由专用网络接口芯片处理，其余走标准以太网；同步基于精确时间协议族。环网冗余用双变压器结构实现，链路失效时的切换时间要计入周期预算。

\textbf{CC-Link IE} 以千兆环形网为骨干，把周期性刷新（所有节点的共享数据区自动交换）与消息通信分开，刷新周期是可计算量；后续的 TSN 版本叠加了时间感知整形与同步，形态与其他 TSN 方案对齐。

\begin{note}[title={为什么通用交换机做不到工业实时}]
上述方案共同的硬件前提是：转发与时间戳在交换或接口芯片里完成，而不是在主机的协议栈里。通用交换机只提供存储转发加尽力而为队列，缓冲深度不可见、排队延迟取决于当时全网流量，因此没有任何一条路径能给出上界；它也缺少分布式时钟所需的 ``在物理层记录帧边界'' 能力。把现场设备接到家用交换机上 ``跑通了''，与 ``最坏延迟可界'' 是两件不同的事。
\end{note}

\begin{longtable}{>{\raggedright\arraybackslash}p{0.15\textwidth} >{\raggedright\arraybackslash}p{0.23\textwidth} >{\raggedright\arraybackslash}p{0.24\textwidth} >{\raggedright\arraybackslash}p{0.17\textwidth}}
\caption{总线与方案的实时特性对比：周期、负载、同步精度与硬件要求}\label{tab:ch12-bus-traits}\\
\toprule
总线或方案 & 周期与负载形态 & 抖动与同步精度 & 硬件要求 \\
\midrule
\endfirsthead
\toprule
总线或方案 & 周期与负载形态 & 抖动与同步精度 & 硬件要求 \\
\midrule
\endhead
\bottomrule
\endfoot
\endlastfoot
CAN 2.0 & 毫秒级周期，设计负载压在三分之一到二分之一 & 低负载时亚毫秒，高负载时尾部不可预测；无原生时基 & 控制器加收发器即可，成本最低 \\
CAN FD & 同上，但单帧更长、数据段更快 & 与 CAN 同量级，仲裁段仍然慢 & 需 FD 控制器与收发器，两端参数必须一致 \\
LIN & 主站轮询表决定，典型毫秒到十几毫秒 & 槽位固定，抖动来自主站调度；无跨节点同步 & 单线收发器，微控制器自带外设 \\
FlexRay & 宏周期内静态槽加动态段，可到几十微秒 & 帧起始同步优于微秒量级；支持双通道冗余 & 需专用控制器与调度表求解工具 \\
EtherCAT & 循环时间可低至几十微秒到 1 毫秒 & 分布式时钟下亚微秒级同步与抖动 & 需带硬件转发的从站芯片与专用主站栈 \\
PROFINET 等时级 & 同步周期可低至数百微秒 & 同步域内微秒级；普通实时等级无保证 & 需支持的交换组件与组态工具生成的映像 \\
POWER{\Bk}LINK & 等时相时槽加异步相，毫秒或亚毫秒 & 时槽决定抖动；周期时间信息对齐到微秒量级 & 需支持的开关芯片与主站实现 \\
TSN 以太网 & 由门控表决定，窗口可到几十微秒 & 域内亚微秒，随级联跳数衰减 & 需支持相应任务组的网卡与交换机 \\
普通 TCP 或 UDP & 应用层自定，链路层无保证 & 抖动来自排队与重传，可达毫秒 & 通用网卡与交换机，生态最好 \\
\end{longtable}

\section{TSN 任务组：同步、整形、抢占与冗余}

时间敏感网络不是一种协议，而是一组任务组提供的能力，必须按需组合。

\textbf{802.1AS}（广义精确时间协议）解决时基。它建立在精确时间协议机制上，用二层组播收发同步报文，每个端口测链路延迟（两侧请求应答时间戳差的一半，且两侧时钟必须来自同一个硬件时间戳单元），级联时每一跳引入自己的误差。级联精度衰减是硬约束：每增加一跳，误差预算要加上该设备的时间戳粒度、收发路径不对称与频率偏差，因此同步域的规模与级联深度必须写进设计，不能假设 ``全网都微秒级''。设备的物理层时钟是这些时间戳的来源，主机侧通过字符设备接口读取，内核用 \Tp{SO\_TIMESTAMPING} 让应用同时拿到软件与硬件时间戳。较新的修订版还支持非锚定从时钟与更好的路径延迟处理，规划时按实际芯片支持的版本核对。

\textbf{802.1Qbv}（时间感知整形）在每台设备的每个输出端口放一张门控列表：表项指定哪些队列打开以及保持多久，整表循环，起点由基准时刻给出并以同步时间为参照。控制流被安排在自己的窗口里发送，窗口长度决定该队列能承载的字节数，其余流量共用剩下的窗口。要点是保护带：若窗口起点处还在传输上一队列遗留的长帧，本窗口的第一批报文会被推迟，所以保护带要覆盖最长帧的传输时间。一千五百字节在一百兆位链路上约需一百二十微秒。

\textbf{802.1Qav}（基于信用的流量整形）给高优先级流一份 ``单位时间可积累与消耗多少发送额度'' 的信用，从而把突发限制在预算内，而不必精确到时刻。它不保证发送时刻，只保证突发不会长期挤压别人，适合视频与大容量非控制流量。

\textbf{802.1Qbu} 与 \textbf{802.3br} 提供帧抢占：普通帧允许被截断成碎片传输，立即类队列的短帧可以插进碎片之间，接收端按碎片重组。关键量是最长碎片对应的延迟，也就是抢占生效前必须发完的那部分的传输时间；它把保护带从整个最长帧降到一个碎片，是在百兆链路上做亚百微秒窗口的前提。前提是链路全双工、两端都支持，中间任何不支持该能力的桥都会让抢占静默失效。

\textbf{802.1CB} 做帧复制与消除：同一报文沿两条不相交路径各发一份，收端按序列号丢弃后到的重复帧。上界取两路中较好的一路，单点失效不中断通信，代价是带宽翻倍与乱序处理。

\textbf{802.1Qci} 做逐流过滤与监管：以流句柄（源、目的、标识符加类别）识别流，为每条流单独配置门、阻塞时长与额度，把不属于它的报文在入口丢弃。它是防止一个节点的错误压垮全网的机制，也是流级隔离的基础。

\textbf{802.1Qcc} 解决配置分发：集中式模式下由网络配置实体收集能力、求解调度并下发；分布式模式下各节点自行协商。它取代了早期的流预留思路，是全网能对齐门控表的现实路径。\textbf{802.1BA} 则把同步、整形与流预留的组合写成音频视频应用档，说明这些任务组本来就不是互相独立的选项。

\begin{warning}[title={时间源切换时的跳变会污染控制环}]
同步域里的时间主从关系变化（更高优先级的主时钟上线、级联路径改变、时钟域根切换）会让本地时基用若干周期去收敛，期间可能出现步进或大幅频率调整。若门控窗口用绝对基准时刻表达，一次跳变就让所有设备的窗口错开，表现为突发丢包而不是缓慢劣化。防护办法：同步侧限制步进幅度并优先用频率微调；调度侧在同步状态未就绪时不放实时流；验证时必须人为拔线触发一次切换，观察门控是否重排。
\end{warning}

一个共存设计例：控制环 500 微秒周期、单帧 200 字节、要求零排队；两路视频各 2 兆比特每秒；尽力而为的数据采集。在一百兆位链路上，200 字节的传输时间约十六微秒。门控表可以这样安排：队列三（控制）每 500 微秒打开二十微秒；剩余窗口里队列零（视频，走信用整形，额度按链路速率的三分之一）与队列一、队列二（采集与背景）轮转；保护带取最长帧传输时间，约一百二十微秒；若启用帧抢占，则只取碎片对应的时间。由于保护带比控制窗口本身长得多，控制窗口的相对位置必须避开上一帧残留，通常做法是把保护带放在窗口之前而不是之后。设计必须用脚本或求解工具验证三件事：窗口不重叠、额度足够、跨设备使用同一个基准时刻。

\begin{guideline}[title={协议选型是同步、整形、冗余三件事一起选}]
选定一种总线形态时，同时回答三个问题：时间基准从哪来、谁在什么时刻能发、单点失效后怎么办。三个答案必须来自同一套可支持的机制。只买 ``支持 TSN'' 的网卡，而不同时确认它实现了哪几个任务组、交换机是否支持同样的组合、驱动与固件是否真的暴露这些能力，等于只答了半个问题。三件拆开买，通常会在集成阶段暴露不兼容。
\end{guideline}

\begin{center}
\begin{tabular}{>{\raggedright\arraybackslash}p{0.16\textwidth} >{\raggedright\arraybackslash}p{0.33\textwidth} >{\raggedright\arraybackslash}p{0.28\textwidth}}
\toprule
任务组 & 提供的能力 & 落地时的约束 \\
\midrule
802.1AS & 域内统一时基与链路延迟测量，时间戳来自硬件 & 精度随跳数衰减；需要网卡与交换机都支持 \\
802.1Qbv & 按时间表开关队列，给出发送窗口 & 需要保护带；全网基准时刻必须一致 \\
802.1Qav & 用信用限制突发，保住别人的带宽 & 只给速率上界，不给时刻上界 \\
802.1Qbu 与 802.3br & 允许高优帧插入普通帧的碎片之间 & 需全双工且全链路支持，否则静默退化 \\
802.1CB & 双路径复制与去重，单点失效不断流 & 带宽翻倍，去重逻辑要按序列号实现 \\
802.1Qci & 逐流识别、监管与入口丢弃 & 需要每条流的句柄定义，流数受设备表项限制 \\
802.1Qcc & 集中式或分布式下发调度与流参数 & 依赖工具链，手工维护极易错 \\
802.1BA & 把上述能力组合成可引用的应用档 & 档内组合须与芯片实际能力逐条核对 \\
\bottomrule
\end{tabular}
\end{center}

\begin{lstlisting}[style=yamlstyle,caption={设备树：以示意形式声明端口的 TSN 队列调度策略},label={lst:ch12-dts-queues}]
/* 示意：属性名与取值在不同厂商的绑定里差别很大，落地前必须核对
   内核的绑定期望文件与驱动实际解析的字段，不要照抄。 */
&mdio0 {
    switch0: ethernet-switch@0 {
        compatible = "vendor,tsn-dsa";
        ports {
            #address-cells = <1>;
            #size-cells = <0>;
            port@3 {
                reg = <3>;
                ethernet = <&ethernet0>;
                tsn {
                    queue-count = <8>;
                    queue3 {                        /* 控制环：时间感知整形 */
                        schedule = "qbv";
                        base-time-ns = <0>;
                        gate-states = <0x08 20000>, <0x07 480000>;
                        guard-band-ns = <120000>;
                    };
                    queue0 {                        /* 视频：信用整形 */
                        schedule = "qav";
                        idle-slope-bps = <25000000>;
                        send-slope-bps = <100000000>;
                    };
                    queue1 { schedule = "express"; };        /* 帧抢占立即类 */
                    queue7 { schedule = "best-effort"; };
                };
            };
        };
    };
};
\end{lstlisting}

\section{主机侧实现与硬件依赖}

主机能不能做到硬件级实时发送，取决于网卡暴露了什么。第一步是查询：\Tp{ethtool -T eth0} 的输出会列出是否支持硬件发送与接收时间戳、是否支持按时间发送、可用的时钟（例如物理层时钟与 \Tp{CLOCK\_TAI}）、以及时间戳格式。看不到硬件发送时间戳能力，后面所有配置都退化成软件排队。

发送侧的核心是时间触发发送队列（\Tp{etf}）：它把报文按到达顺序排队，然后在指定时刻交给驱动，可以由软件实现，也可以卸载到网卡的硬件门控表。它需要与多队列根规则配合，把不同优先级导到不同队列，再在目标队列上挂这个规则。

\begin{lstlisting}[style=shellstyle,caption={Shell：ptp4l 与 phc2sys 建立时基，并用 tc 挂载按时刻发送的队列},label={lst:ch12-ptp-etf}]
# (1) 确认网卡能力：硬件时间戳与按时间发送；卸载开关改前留档
ethtool -T eth0
ethtool -k eth0 > /tmp/offload-before.txt
ip -details link show eth0

# (2) 建立域内时基：二层 gPTP 配置档，日志留在控制台便于看状态机
#      配置路径按发行版不同，linuxptp 上游自带一份 gPTP 示例配置
ptp4l -i eth0 -2 -f /etc/ptp4l.d/gptp.cfg -m -l LOG_INFO &
phc2sys -s eth0 -c CLOCK_TAI -O 0 -m -w &        # 网卡 PHC 作来源，喂系统侧
# 反方向（用系统时钟喂网卡 PHC）只应在时间根节点上启用，两者勿同时开

# (3) 队列与整形：多队列根规则，把目标流量导到队列 1，队列 1 挂 etf
tc qdisc replace dev eth0 parent root handle 100 multiq
tc qdisc add dev eth0 parent 100:1 etf clockid TC_CLOCK_TAI \
    offload deadline 2000000 interval 500000
tc filter add dev eth0 protocol ip parent 100: prio 1 u32 \
    match ip protocol 17 0xff match ip dport 40000 0xffff flowid 100:1
# 不同 iproute2 版本里 etf 的子选项名略有差异，用 tc qdisc add 的用法帮助核对

# (4) 关掉会破坏时间戳与排队语义的卸载与节能项
ethtool -K eth0 tso off gso off gro off tx-checksum-ip-generic off
ethtool -A eth0 autoneg off rx off tx off
ethtool --set-eee eth0 eee off                   # 节能以太的唤醒延迟是抖动源
cat /sys/class/net/eth0/statistics/tx_timeout    # 发送复位计数应恒为零
ethtool -S eth0 > /tmp/nic-stats-after.txt
\end{lstlisting}

按时间发送的套接字选项（\Tp{SO\_TXTIME}）的逐包用法是：先在套接字上打开该选项，再在发送控制信息里携带一个六十四位无符号的绝对发送时刻（单位纳秒，取自队列所绑定的那个时钟）。内核对这条路径的检查较严，通常要求数据报套接字、明确绑定设备、并且优先级能映射到挂了发送队列的那个队列上；不满足时错误从 \Fn{sendmsg} 直接返回，而不是静默退化成普通发送，这是最常见的调试陷阱。

\begin{lstlisting}[style=cppstyle,caption={C++：用 SO\_TXTIME 把一个报文安排在绝对时刻发出},label={lst:ch12-so-txtime}]
// 依赖：内核支持按时间发送，目标队列上已挂 etf 并绑定同一时钟
#include <arpa/inet.h>            // htons
#include <cstdint>
#include <cstring>
#include <ctime>
#include <linux/if_ether.h>       // ETH_P_IP
#include <net/if.h>               // struct ifreq, IFNAMSIZ
#include <netpacket/packet.h>     // struct sockaddr_ll
#include <string>
#include <sys/socket.h>

struct TxHandle {
    int fd{-1};
    std::string ifname;
};

bool tx_open(TxHandle &h, const char *ifname, int prio)
{
    h.ifname = ifname;
    h.fd = ::socket(AF_PACKET, SOCK_DGRAM | SOCK_CLOEXEC, htons(ETH_P_IP));
    if (h.fd < 0) return false;

    int on = 1;
    if (::setsockopt(h.fd, SOL_SOCKET, SO_TXTIME, &on, sizeof(on)) != 0)
        return false;           // 不支持就退出，不要退回普通发送
    if (::setsockopt(h.fd, SOL_SOCKET, SO_PRIORITY, &prio, sizeof(prio)) != 0)
        return false;           // 优先级决定报文落到哪个队列

    ifreq b{};
    std::strncpy(b.ifr_name, ifname, IFNAMSIZ - 1);
    return ::setsockopt(h.fd, SOL_SOCKET, SO_BINDTODEVICE, &b, sizeof(b)) == 0;
}

bool tx_send_at(TxHandle &h, const sockaddr_ll &dst, uint64_t t_ns,
                const void *data, size_t len)
{
    char ctl[CMSG_SPACE(sizeof(uint64_t))];
    iovec io{const_cast<void *>(data), len};
    msghdr m{};
    m.msg_name = const_cast<sockaddr_ll *>(&dst);
    m.msg_namelen = sizeof(dst);
    m.msg_iov = &io;
    m.msg_iovlen = 1;
    m.msg_control = ctl;
    m.msg_controllen = sizeof(ctl);

    cmsghdr *cm = CMSG_FIRSTHDR(&m);
    cm->cmsg_level = SOL_SOCKET;
    cm->cmsg_type = SCM_TXTIME;
    cm->cmsg_len = CMSG_LEN(sizeof(uint64_t));
    *reinterpret_cast<uint64_t *>(CMSG_DATA(cm)) = t_ns;   // 绝对发送时刻

    const ssize_t r = ::sendmsg(h.fd, &m, 0);
    return r == static_cast<ssize_t>(len);
}

uint64_t tx_now_ns()
{
    timespec ts{};
    ::clock_gettime(CLOCK_TAI, &ts);      // 时钟必须与队列的 clockid 一致
    return static_cast<uint64_t>(ts.tv_sec) * 1000000000ull + ts.tv_nsec;
}
\end{lstlisting}

其余主机侧选项各有代价。XDP 的用户态套接字把收发搬到映射的填充环与完成环上，绕过协议栈但仍由中断或轮询驱动网卡，收益是延迟散布明显收窄，代价是要自己处理报文生命周期，并且需要网卡支持这一模式。用户态协议栈框架走得更远：彻底消除中断与协议栈，端到端延迟可稳定在几微秒到几十微秒，代价是独占整个端口与一个满占的核，功耗与可用性代价都实打实；而且它与内核的时间戳、同步机制脱钩，需要另外打通。内核与用户态共享收发环的折中模式保留内核可见性又减少拷贝，但要确认网卡与驱动版本都支持，否则会无声退化。

内核默认的多优先级队列（例如通用分段、优先队列类规则）只做 ``按类别排队''，不看时间：它们能让短队列排在长队列之前，却无法保证某个报文在某个时刻发出。这与时间触发发送队列的区别是本质性的，也是 ``加了优先级却仍然抖'' 的常见解释。

验证必须双端做。发送端比较实际硬件发送时间戳与计划时刻之差；接收端用 \Tp{SO\_TIMESTAMPING} 的硬件接收时间戳记录入帧边界，并与本机物理层时钟读数比对；\Tp{ptp4l} 与 \Tp{phc2sys} 的日志给出同步状态收敛过程与偏移量。最可靠的是双网卡对测：一台机器按门控表发送并记下硬件发送时间戳，另一台用硬件接收时间戳记录，两者处在同一时钟域，差值直接给出端到端抖动，不依赖任何软件层假设。

\section{上层中间件与选型矩阵}

\textbf{OPC UA 承载在 TSN 上}把原来走 TCP 或普通 UDP 的信息模型改成发布订阅：发布方用统一架构的二进制编码把数据集写进组播或以太类型报文，订阅方按流句柄接收，于是每条流可以一一对应一个逐流监管点与一个门控窗口。要点有两个：数据集划分要与门控窗口大小一致，否则一个窗口的剩余时间被无关字段填满；加密与安全层放在哪一层，决定这段延迟还能不能计算。

\textbf{DDS} 以实时发布订阅协议为线格式，靠组播做发现与传输。服务质量策略里与时间相关的是截止期、活性、传输优先级、延迟预算，以及目的端顺序（按源时间戳还是按接收时刻）。它提供了让应用自己声明实时性需求的表达方式，但底层排队与时间戳仍在内核与网卡上：策略写对了不会让通用交换机变成确定性的。

由此得到本章最重要的判断：选型时同时决定同步、整形与冗余，中间件只是把这三个决定表达出来。

\begin{longtable}{>{\raggedright\arraybackslash}p{0.15\textwidth} >{\raggedright\arraybackslash}p{0.21\textwidth} >{\raggedright\arraybackslash}p{0.24\textwidth} >{\raggedright\arraybackslash}p{0.16\textwidth}}
\caption{按周期、抖动、同步精度、规模、生态、认证与成本的选型对照}\label{tab:ch12-selection}\\
\toprule
需求画像 & 周期与抖动 & 同步与冗余 & 规模、生态、认证与成本 \\
\midrule
\endfirsthead
\toprule
需求画像 & 周期与抖动 & 同步与冗余 & 规模、生态、认证与成本 \\
\midrule
\endhead
\bottomrule
\endfoot
\endlastfoot
车身门窗与座椅，毫秒级 & LIN：毫秒到十几毫秒，槽位固定 & 无跨节点同步，靠主站重发兜底 & 节点个位数，生态极成熟，成本最低，认证要求低 \\
车载动力域，百微秒到毫秒 & CAN 或 CAN FD：低负载亚毫秒，高负载尾部不可界 & 无原生时基；可做报文级双发 & 数十节点，工具链完备，现场案例最多 \\
车载域控，几十微秒 & 以太网加门控：窗口可到几十微秒 & 域内亚微秒；帧复制可做双路径 & 需支持相应任务组的交换芯片，认证路径仍在积累 \\
工厂运动控制，125 到 500 微秒 & FlexRay：静态槽可证上界 & 帧起始同步优于微秒；双通道 & 依赖调度表工具，节点规模有限，成本高 \\
伺服轴与高速 IO，几十微秒 & EtherCAT：循环时间可低至几十微秒 & 分布式时钟亚微秒；冗余需计入周期 & 器件与主站栈绑定较深，生态大，成本中等 \\
过程控制与安全回路 & PROFINET 等时级或 POWERLINK：数百微秒到毫秒 & 同步域或周期时间信息，微秒量级 & 需专用交换组件，行业认证资料齐全 \\
混合承载（控制加视觉加采集） & TSN：按门控表逐项分配窗口与额度 & 同步、整形、复制三件一起配 & 硬件与配置复杂度最高，需要求解工具与专职人员 \\
毫秒级、单跳、可重传 & 普通 UDP 加应用层超时 & 不需要共享时基 & 通用网卡与交换机，成本与人力最低 \\
\end{longtable}

最后一条值得展开。满足以下全部条件时，普通 UDP 就是合理选择：周期在毫秒级以上且抖动预算还有几十微秒以上余量；不存在需要同一时刻采样的多节点；一跳或两跳的简单拓扑，链路上没有背景大流量；应用层有明确的超时、幂等或序号处理；网络负载率长期低于四分之一。反之，只要出现 ``必须知道另一台设备在何时采样''，或者 ``尾部延迟会造成机械损伤''，就必须动用同步与整形机制，而不是继续调优内核。

\section{常见陷阱与现场验证清单}

\textbf{交换机不透明}。普通交换的存储转发延迟随帧长变化，缓冲深度不可读，队列积压只在事后统计里以丢弃计数出现。后果是延迟完全不可界，而且劣化过程安静。判定办法是把链路缩到一跳直连再测：若直连达标而经过交换后劣化，问题就在设备侧。

\textbf{物理层与线缆}。自协商在链路建立期反复重训，节能以太的空闲低功耗模式唤醒要几十微秒，长线加劣质屏蔽造成边沿退化与重传。测程前应把速率强制固定、关掉节能模式与链路休眠，并把这几项写进上一章报告模板里的环境项。

\textbf{卸载只关一半}。报文分段与接收聚合会改变帧的实际边界与时序，硬件时间戳在部分卸载路径上不可用或指向错误位置。常见错误是关了发送分段却留着接收聚合，于是接收时间戳看起来正常而整体延迟被拉长；或者关掉校验和卸载后网卡进入软件兼容路径，性能与时间戳能力同时改变。要逐项核对卸载开关，并在改前改后各测一轮。

\textbf{把 \Tp{TCP\_NODELAY} 当实时手段}。它只取消小包合并，不消除重传、流控与排队。在共享链路上开它，尾部延迟反而可能变差，因为小包数量上升。真正需要即时发送时，用 UDP 加发送时刻控制。

\textbf{广播与未知单播风暴}。转发表项老化或目的地址未学习时，帧被泛洪到全部端口，实时流的窗口被这些无关帧填满。表现为某个瞬间全网抖动，事后极难复现。对策是静态表项、入口过滤，以及把组播地址提前注册进转发表；同时检查诊断类工具是否留在生产网段里。

\textbf{错误计数被忽略}。CAN 的仲裁丢失与位错误计数、以太网的物理层与循环冗余校验计数、交换机的丢弃与缓冲溢出计数，是唯一能证明 ``抖动来自链路而不是软件'' 的证据。它们大多是增量计数器，必须在测程前后各取一次做差；只取一次读数没有意义，这与上一章对中断计数做差是同一个方法。

\begin{bestpractice}[title={现场验证最小清单}]
(1) 网卡能力核对：硬件时间戳与按时间发送能力、卸载开关的改前改后对照留档。(2) 时基核对：同步日志中偏移与均值收敛，同时记录物理层时钟与系统时钟的差值。(3) 门控核对：两端各发一轮，比较硬件发送时间戳与计划时刻。(4) 端到端：双网卡对测，统计连续一小时的抖动直方图并按上一章的口径判定。(5) 故障注入：拔掉一条冗余链路、上线第二个时间主源、给尽力而为流量灌满带宽，逐项确认实时流窗口未被侵占。(6) 留证：所有原始时间戳文件、配置导出与判定脚本一并归档。
\end{bestpractice}

到这里，分册第三部分的方法闭环已经完整：先测出抖动，再判断它发生在调度、内核、硬件还是链路，最后把它写成可自动阻断的门限。

% ============================================================
% ---- frag_ch13.tex ----
\chapter{微架构、编译器与确定性陷阱}
\label{ch:ch13-determinism-pitfalls}

前面各章已经把时间打点、中断延迟、看门狗、DMA、调度与总线讲成了可测量的对象；本章处理另一半问题：同一段代码、同一块芯片，为什么这一轮跑了 40 µs 而下一轮跑了 180 µs，以及为什么换一个小版本的编译器之后，最坏值整体抬高了三成。答案几乎总是落在三个层面上：存储层次的状态、控制流的可预测性、以及编译器对源码语义的改写。这三层共同决定时间是否可界，而可界性与实测很快是两件完全不同的事。

本章口径与前几章保持一致：先用打点与追踪把抖动来源归因到具体的微架构事件，再用静态手段把不会走到那条慢路径变成可检查项。凡涉及中断到达、调度响应、总线仲裁的机制，前几章已经建立，这里只点名不复述。

\section{存储层次的时间代价}
\label{sec:ch13-memhier}

\subsection{命中与缺失的量级差别}

在现代嵌入式处理器上，取指与取数的代价相差可达一到两个数量级。以典型的 Cortex-M7 与多数应用处理器核为例：L1 数据命中在个位数周期量级，L1 缺失走到紧耦合的 L2 或低延迟内部总线时是十几到几十周期，走到外部 SDRAM 或 QSPI Flash 且需要跨突发拆分时则是数百周期；如果这次缺失又赶上 AXI 上的多主仲裁等待或刷新周期，还要叠加前几章测过的总线抖动。指令侧同理：一段从未执行过的代码首次取指会把整条 I-cache 填充路径走满，这就是冷启动的第一轮循环常常是全片最慢一轮的原因。

对硬实时任务而言，关键量不是平均值而是这一轮究竟缺失了多少条 cache line。一次缺失的代价可以粗略看成：被触及的字节数除以 line 大小，再乘以每条 line 的填充周期，最后加上总线等待。因此估算最坏执行时间时，正确的建模单位是被触及的 line 条数，而不是迭代次数——前者与微架构状态相关但可以界住，后者会掩盖 line 数随输入变化这一事实。

\subsection{组相联、映射冲突与一个真实的实例}

组相联缓存把地址映射到固定的组：组号由地址中间若干位决定。Cortex-M7 的 L1 数据缓存是 16 KB、4 路、32 字节 line，按容量除以 line 大小再除以路数得 128 组，决定组号的位宽是 7 位，于是地址每相差 4096 字节（32 乘 128）就落回同一个组。这不是理论游戏：如果两个数组恰好相差 4096 字节，又被同一个循环按相同下标访问，它们的元素会争抢同一组的 4 个位置；此时循环里再放一个同样 4096 间隔的查找表，冲突就超出路数，缓存反复淘汰、反复填充，最坏路径长度成倍增加。这类同一缓冲区的两个区落在同一组的现象，在直接映射的老式缓存与部分 DSP 核上更严重，因为那里路数只有 1。

\begin{lstlisting}[style=cstyle,caption={C：4 KB 间隔引起的组冲突，以及错相填充后的写法},label={lst:ch13-alias}]
/* 反例：两个缓冲区恰好相差 4096 字节（32 B line x 128 set），
 * 同一个循环同时访问，二者争抢同一组，路数不足 -> 反复缺失 */
#define BUF_BYTES   1024u
uint8_t bad_in[BUF_BYTES];
uint8_t bad_out[BUF_BYTES];   /* 链接后极可能与 bad_in 相差 4096 的整数倍 */

/* 正例：错相半个 cache line，使两个数组落入不同的组集合。
 * 32 B line，取 16 B 偏移即可打散；填充区不参与任何功能 */
#define ALIAS_PAD   16u
uint8_t fix_in[BUF_BYTES];
uint8_t pad[ALIAS_PAD];       /* 只为改变后续数组的组分布 */
uint8_t fix_out[BUF_BYTES];

/* 常量表若很大，优先放紧耦合内存或只读存储器，并显式按 line 对齐，
 * 避免表头跨两条 line 造成首轮多取一次 */
const uint16_t sin_table[256] __attribute__((aligned(32)));
\end{lstlisting}

错相与填充是最廉价的手段，代价只是几十字节 RAM。更系统的做法是把关键缓冲区集中到一段连续、已知对齐的区域内，让组分布可以被静态枚举；这条纪律应当与第 9 章的内存纪律合并成同一份检查项，避免两套布局原则互相打架。

\subsection{预取器：平均变好、最坏变坏}

硬件预取器观察访问流并把后续 line 提前拉进来，顺序访问的吞吐量因此常常翻倍。但对硬实时代码它是一把双刃剑：预取占用总线带宽与缓存容量，把无关数据挤进有限的路数里，从而抬高别人或同一任务下一阶段的最缺失计数；某些核的预取器遇到跨页访问时会产生一次被丢弃的预取，白白占用队列。更麻烦的是，预取行为依赖历史状态，界分析工具很难给出保守而可用的上界。可行的处理是三选一：为关键核关闭数据预取（部分核有实现定义的控制寄存器或调试访问口，需在板级验证其确实生效）；把访问模式改成工具可推断的固定步长；或者接受收益，同时把最坏情形的额外缺失项显式写进预算，并按第 11 章的口径复核。

\subsection{TLB、大页与地址翻译}

带 MMU 的核上，虚拟地址到物理地址还要过一层 TLB，一次 TLB 缺失要走页表，代价常常高于一次数据缓存缺失。因此实时任务的地址工作集应当小且连续：Linux 加 PREEMPT\_RT 上常见做法是把共享内存与关键环形缓冲放在大页上，减少表项数量并抬高单次覆盖范围；访问跨页边界也要小心，某些核的单次翻译无法覆盖跨越页缝的两半。多核 SMP 上还应注意，页表变更可能伴随一次 TLB 失效的核对过程，这段时间对实时性是隐藏的，且不会出现在任何应用层打点里。

\subsection{写缓冲、内存类型与 WC 上的读回退化}

写操作先进入写缓冲，随后按突发出栈，这使连续写比随机写便宜得多；代价是别的观察者（DMA、另一个核）看到写入的时刻并不确定，这正是第 5 章一致性讨论的根因之一。内存类型决定这段行为的边界：write-back 允许合并与延迟、读可命中缓存；write-through 每次写出但读仍可命中；write-combining 允许把整条 line 合并成一次不可拆的写出，但读会强制清空合并缓冲，退化到逐字节的设备访问速度。WC 的坑正在于读回：驱动里如果混用合并写与调试性读回，一次看似无害的读可以造成数百周期惩罚，在慢速外设上还要再加几个量级。Linux 侧确实需要显式同步时使用 \Fn{smp\_wb\_invalidate\_all}，且只应出现在非热路径；更稳的做法是根本不在 WC 区域上读。

\begin{longtable}{>{\raggedright\arraybackslash}p{0.20\textwidth}>{\raggedright\arraybackslash}p{0.22\textwidth}>{\raggedright\arraybackslash}p{0.22\textwidth}>{\raggedright\arraybackslash}p{0.19\textwidth}}
\caption{微架构特性对平均时间与最坏时间的影响及对策}\label{tab:ch13-microarch}\\
\toprule
特性 & 对平均时间 & 对最坏时间 & 工程对策 \\
\midrule
\endfirsthead
\toprule
特性 & 对平均时间 & 对最坏时间 & 工程对策 \\
\midrule
\endhead
\bottomrule
\endfoot
\endlastfoot
L1 命中与缺失 & 平均被大量命中拉低 & 每条缺失按填充周期累加，是最坏值主体 & 按触及的 line 条数建模；热路径局部性优先 \\
组冲突与路数不足 & 影响小范围热点 & 超出路数即呈阶跃式恶化 & 缓冲区错相、集中布局、静态枚举组分布 \\
硬件预取 & 顺序流吞吐显著上升 & 抢占带宽与路数，抬高他人的缺失 & 关闭或收敛访问模式；额外缺失写进预算 \\
TLB 与页边界 & 命中后近乎免费 & 一次页表遍历可达数百周期 & 大页、地址工作集小且连续、避免跨页 \\
写缓冲与合并 & 连续写吞吐高 & 观察者可见时刻不确定 & 只在固定边界做显式同步 \\
WC 区域读回 & 纯写时极快 & 一次读清空合并缓冲，惩罚巨大 & 关键区域设为 write-back 或强序 \\
\end{longtable}

\section{紧耦合内存、锁定与可界性}
\label{sec:ch13-tcm}

紧耦合内存（TCM，某些厂商称 CCM、DTCM、ITCM）挂在核的专用端口上，访问周期固定且通常只有一到几个周期，不参与缓存协议，也不与总线上的其他主设备争抢。对微秒级的控制环，这是最有效的确定性手段：把关键环的代码放 ITCM、把它的状态与滤波系数放 DTCM，最坏路径就从取决于缓存状态变成取决于指令条数。实现上靠段属性完成，段与区域由链接脚本提供。

\begin{lstlisting}[style=cstyle,caption={C：把关键环代码与数据钉到 ITCM 与 DTCM，并显式对齐},label={lst:ch13-tcmsec}]
/* 由链接脚本提供 .itcm / .dtcm（或 .ccm）区域；这里只负责把符号放进去 */
#define ITCM_FN  __attribute__((section(".itcm"), aligned(32), noinline))
#define DTCM_VAR __attribute__((section(".dtcm"), aligned(32)))

/* 关键环：整段落在 ITCM，运行期不会被替换；32 字节对齐即一条缓存行 */
ITCM_FN static float loop_step_impl(float err, float kp, float ki)
{
    /* 定点或单精度浮点均可；禁止在此调用任何间接层 */
    return err * kp + ki;
}

DTCM_VAR static float s_integral;        /* 状态放 DTCM，读写周期恒定 */
DTCM_VAR static float s_coef[4];         /* 若为编译期常量，改放只读段更省 RAM */
\end{lstlisting}

上面的定义里刻意保留了 \Tp{noinline}：钉到 ITCM 的函数若被内联进一个位于 Flash 的调用者，代码就按调用者所在段放置，钉定随之失效。这一点在 \S\ref{sec:ch13-compiler} 还会遇到。

缓存锁定（按路锁定或按线锁定）是退一步的方案：若平台提供，可把关键代码或表钉在若干路里保证永不被替换。锁定通常要求先清理缓存、再配置、后使能，且必须与 MPU 或内存类型配置配合，否则 DMA 写入依然会造成观察不一致。可用的路数很少，一般只够锁住一两个关键环与一两张大表，因此它只能作为整体有界方案的一环，不能指望它救下整个固件。

MPU（Cortex-M 一类）与 MMU 属性（Cortex-A、x86）决定一段地址的缓存性、共享性与访问权限。把关键环的只读表标记为只读且可缓存、把设备区标记为强序或设备类型，是最基本的配置。这里有条容易被忽视的原则：属性本身就是界分析模型的输入。如果一段区域的缓存性可以被运行期改写（例如调试代码临时关闭缓存），最坏时间上界就不再成立，除非能证明该改写不会出现在实时窗口内。

{\bf 为什么实测快不等于可界。} 冷启动、缓存被后台任务冲刷、多主设备抢占总线、预取器的历史状态——这些条件在压测里通常不会同时出现，于是实测得到的 P99 与真正的上界之间存在一条无法靠延长观测时间弥补的缝隙。可界性关心的是能否列举所有慢路径并证明其上限，它要求代码结构满足一簇可机械检查的性质：无递归、循环有静态可证的上界、无动态分配、热路径不跨越不可界的外部调用、没有未定义行为。第 10 章的可调度性论证需要这些上界作为输入；拿统计量喂给调度论证，等于把责任推给运气。

\begin{guideline}[title={关键路径的存储布局三条纪律}]
第一，代码与它的数据同域：热路径函数与它访问的缓冲区要么一起进紧耦合内存，要么一起留在可缓存区，避免跨越两套延迟模型。第二，布局要能被静态枚举：段归属、对齐、缓冲区间隔写进链接脚本并在 map 文件里复核，评审时能一眼看完。第三，任何为实测更快而做的布局调整（开启预取、加大缓存、放宽属性）都必须在最坏路径口径下重算一遍，算不动就回退。
\end{guideline}

\section{控制流时间抖动的来源}
\label{sec:ch13-ctrlflow}

分支预测把顺序执行的代价压低，代价只在预测失败时暴露：一次误预测要清空流水线并丢弃投机结果，周期数在十几到几十的量级，深度流水的核上更高。因此硬实时代码里的抖动常常不来自做了什么，而来自走了哪条路。数据相关的分支最危险：分支走向由输入决定时，最坏路径就取决于输入分布，而这正是压测最难穷举的部分。典型表现是同一函数的执行时间在两个相邻数值附近跳变，例如比较阈值刚好落在边界，或 CRC 循环里的条件位随数据翻转。

对策是把数据相关的控制流换成与输入无关的形式：让比较结果参与算术而不参与跳转，或者用索引恒定的查表代替串行比较。

\begin{lstlisting}[style=cppstyle,caption={C++：用选择与查表替代数据相关分支，并消除热路径的间接层},label={lst:ch13-hotpath}]
// 反例：分支走向由采样值决定，最坏路径取决于输入序列，难以界定
float gain_bad(float v) {
    if (v > 4.0f) { return v * 1.5f; }   // 误预测十几到几十周期
    if (v > 2.0f) { return v * 1.2f; }
    return v;
}

// 正例 (1)：饱和处理。写成两两比较加选择，编译器通常降为选择指令
// 或位运算掩码，路径与输入无关（若目标无选择指令，可用 mask = -(int32_t)cond）
float sat_branchless(float x, float lim) {
    const float a = (x > lim) ? lim : x;
    return (a < -lim) ? -lim : a;
}

// 正例 (2)：查表。索引由比较结果的整数和给出，代价恒定
static const float kTab[3] = {1.0f, 1.2f, 1.5f};
float gain_table(float v) {
    const int idx = (int)(v > 2.0f) + (int)(v > 4.0f);
    return v * kTab[idx];
}

// 反例：热路径上的间接层。虚调用与类型擦除的可调用对象都要多跳一次指针，
// 间接跳转还会撞上预测器的弱项，并落入侧信道缓解措施的加收范围
struct Mode { virtual float step(float v) = 0; };
float loop_bad(Mode& m, const float* p, size_t n) {
    float acc = 0.0f;
    for (size_t i = 0; i < n; ++i) { acc += m.step(p[i]); }  // 每轮一次间接调用
    return acc;
}

// 正例：模板参数在编译期把类型定死，循环内不留间接跳转
template <class Op>
float loop_good(const Op& op, const float* p, size_t n) {
    float acc = 0.0f;
    for (size_t i = 0; i < n; ++i) { acc += op(p[i]); }
    return acc;
}
\end{lstlisting}

\subsection{跳转表、间接分支与侧信道缓解}

\Tp{switch} 的生成策略影响很大：分支少且取值稀疏时编译成比较链，分支多且取值密集时编译成跳转表。前者是若干顺序分支，路径长度与比较次数成正比但分支模式稳定；后者是一次间接跳转，代价恒定却更难预测，而且落在间接分支缓解措施的加收范围内。现代核为了堵住推测执行侧信道（Spectre 一类），对间接分支与返回引入了额外序列化：x86 上的 retpoline 与 IBRS/IBPB 一类控制，ARM 上针对预测失败的硬件或软件措施。这些机制把间接调用与返回的固定开销抬高了可观的一段，并让栈回溯的代价更难静态推断。工程结论很直接：热路径里不要出现间接分支，包括函数指针表、虚函数、符号延迟绑定，以及通过绑定包装出来的可调用对象；确实需要多态时用编译期多态，或者把跳转点放到环的外层，保证每轮最多一次。若工具链提供禁止跳转表的开关（GCC 一族的 \Tp{-fno-jump-tables}），可以用它把稀疏 \Tp{switch} 钉成比较链，代价是链长变成需要显式设计的对象。

\subsection{循环展开、流水与到达的扰动}

展开减少循环开销并增加指令级并行的机会，但也改变分支模式与访存的时间局部性，并可能把一次缺失摊到更多指令上；更关键的是，展开后的代码与源码结构差距变大，界分析模型必须按展开后的实际控制流图建立。中断与异常到达带来的抖动属于另一类来源：响应时刻的不确定性已在第 3 章度量过，这里的纪律是热路径内部的关中断窗口长度固定且可数，窗口里不做外部调用；一次到达会让流水入口重来，其代价在界分析里要按每个窗口一次计入。测量时如果为了省事关掉中断或拉长打点自身，观察到的分布就会与真实不符——第 2 章讲过的观测污染在微架构层面尤其敏感。

\begin{warning}[title={异常表与栈展开是隐藏的控制流代价}]
在开启 \Lang{C++} 异常、或带 \Tp{EH\_FRAME} 与 \Tp{ARM\_EXIDX} 一类展开信息的固件里，异常路径的代价不体现在任何源码行上：展开器要遍历表项、做二分查找、可能访问未被缓存的只读段。即便你从不抛出，表的存在也会增大镜像、改变段布局，并让静态分析的函数边界变得模糊。硬实时代码通常的做法是关闭异常与 RTTI，把恢复语义交给显式返回值与错误码；若必须保留（例如依赖的中间件要求），就要把展开路径本身纳入最坏时间预算。
\end{warning}

\section{编译器的双面性}
\label{sec:ch13-compiler}

同一份源码，编译选项决定它变成什么样的时间序列。下面按收益与风险成对出现的方式过一遍常用几组。

优化等级与可调试性。\Tp{-Og} 保留变量与调用结构，便于源码级单步；\Tp{-O2} 会消掉变量、合并函数、重排基本块，此时断点与日志的位置都不可信。调试期的结论不能外推到发布镜像，这条常识在实时场景里还要再加一句：抖动形态也不同步，用 \Tp{-Og} 观察到的最大间隔不能用来论证 \Tp{-O2} 的行为。

自动向量化。它把标量循环变成宽寄存器操作，吞吐量上升的同时，把每次迭代一条流水的直觉换成一次宽访问加若干数据搬运。后果有两个：访存的粒度与对齐要求变严，未对齐或跨界访问会退化成标量序列；执行时间分布改变，短循环与长循环的收益完全不成比例。GCC 一族用 \Tp{-ftree-vectorize} 控制（\Tp{-O2} 起通常已启用），用 \Tp{-fvect-cost-model} 调整激进度；源码侧可写 \Tp{\#pragma GCC ivdep} 声明本次迭代间无依赖，让缺少别名信息时的向量化仍可进行。若某个环要求逐周期可解释，把它标成不向量化比解释向量化结果省事。

链接期优化与函数边界消失。\Tp{-flto} 把跨文件的调用变成同一编译单元里的内联，收益是性能与更小的镜像，代价是函数这个分析单位不再存在：界分析工具、静态检查、以及任何按函数边界建模的手段都失去输入，map 文件与符号表也不再对应源码结构。确实要用 LTO，就得为关键环留出可检查的边界：把它们放进独立的静态库、以禁止内联与独立段的方式构建，或在链接前后各做一次反汇编比对，确认关键环的实际指令与静态分析所用模型一致。任何只在发布构建里出现的代码形态，都是不可界性的来源。

内联与 I-cache 压力。内联消除调用开销，同时撑大指令工作集；多任务共享指令缓存时，一次过度内联可能让热路径跨出锁定的路数，最坏值反而变差。折中手段是显式控制，但两个方向的强度\textbf{并不对称}：\Tp{noinline} 是硬保证，标了它，函数就一定还是界分析与静态检查能拿到的那个独立实体；\Tp{always\_inline} 只是尽力而为加一次申报失败——函数体在调用点不可得（跨翻译单元且未开 LTO）、调用点与定义的目标属性不匹配、经由函数指针或虚派发抵达，展开都会落空，落空时 GCC 报编译错误、Clang 降级成警告、MSVC 的 \Tp{\_\_forceinline} 文档直接写明不保证。所以它换来的是可诊断性（编译器必须尝试，做不到就告诉你做不到），不是``关键环里没有调用''的确定性；后者只能在发布镜像上反汇编确认调用点是否真的存在，属性名本身进不了证据链。

浮点重结合。\Tp{-ffast-math} 及其子开关（\Tp{-funsafe--{\Bk}math--{\Bk}optimizations}、\Tp{-ffinite--{\Bk}math--{\Bk}only}、\Tp{-fassociative--{\Bk}math}、\Tp{-fno--{\Bk}signed--{\Bk}zeros}）允许编译器重排加法顺序、假设不存在非有限值、忽略有符号零。浮点加法不满足结合律，同一控制律在不同重排下会得到不同的尾数，结果随编译器版本与目标选项变化；它还常常破坏针对非有限值的防护逻辑，而这种防护在安全相关代码里本身就是功能。

\begin{compare}[title={保留异常与 RTTI，还是关闭它们}]
保留时：镜像里有展开表与类型信息，抛出是一条真实存在的时间路径，代价必须进入模型；动态类型查询可用，中间件兼容性最好；静态分析的边界被隐藏路径污染。关闭时（\Tp{-fno-exceptions} 与 \Tp{-fno-rtti}）：镜像显著缩小，类型查询与运行时类型识别不可用，虚调用仍在但不再携带类型查询数据；时间模型只剩显式调用图，可界性论证可以直接落在函数粒度上。判据是错误处理方式而不是库偏好：只要恢复语义靠返回值与错误码表达，关闭就是纯收益。
\end{compare}

间接层的代价。关闭异常与 RTTI 并不会消除类型擦除带来的间接调用，也不会消除虚表的指针追逐；带堆分配的可调用对象包装还会违反无动态分配的纪律。需要可调用对象时优先用模板参数、放在只读段并保证跳转次数固定的函数指针表，或编译期多态。若编译器提供抑制符号插交的开关（GCC 一族 \Tp{-fno-semantic-interposition}），在静态链接的固件里收益有限，不必依赖。

\begin{lstlisting}[style=cstyle,caption={C：局部钉定：用属性覆盖内联与优化等级，以及屏障的正确形态},label={lst:ch13-pin}]
/* 1) 关键环内的极短函数：请求内联，展开结果要在产物上核对 */
#define HOT __attribute__((always_inline))
static HOT float clamp01(float x) { return x < 0.0f ? 0.0f : (x > 1.0f ? 1.0f : x); }

/* 2) 环外的辅助函数：显式禁止内联，保住函数边界供静态分析使用 */
__attribute__((noinline)) static void commit_diag(uint32_t code);

/* 3) 局部覆盖优化等级：整片 -O2，但这一支要求逐指令可解释
 *    （s_kp / s_ki / s_int 为已定义的系数与状态）*/
__attribute__((optimize("O1"), noinline))
static float pid_step_impl(float err) { return s_kp * err + s_ki * s_int; }

/* 4) 编译器屏障的正确形态：禁止跨此点的访存重排，本身不产生指令 */
#define COMPILER_BARRIER() __asm__ __volatile__("" ::: "memory")

/* 5) MMIO 读写：volatile 加编译器屏障；硬件顺序要靠专门的顺序指令 */
static inline uint32_t reg_rd(volatile const uint32_t* a) {
    uint32_t v = *a; COMPILER_BARRIER(); return v;
}
static inline void reg_wr(volatile uint32_t* a, uint32_t v) {
    COMPILER_BARRIER(); *a = v;
}
/* 与 DMA 描述符交换时，编译器屏障不足以建立硬件顺序：
 * ARM 用 dsb / isb 一类专用指令，x86 用 mfence / sfence 或内核提供的内建；
 * 空的内联汇编只约束编译器，不要用它假装约束了总线 */

/* 6) 把不可能路径告知优化器：静态前提，不产生运行时代码 */
static void fill(uint16_t* d, uint32_t n) {
    if (n > 64u) { __builtin_unreachable(); }   /* 前提：n 已由调用方钳位 */
    for (uint32_t i = 0; i < n; ++i) { d[i] = (uint16_t)i; }
}
\end{lstlisting}

\section{可复现构建与工具链锁定}
\label{sec:ch13-reprobuild}

时间行为是编译器版本的函数。同一份源码在不同 GCC 小版本下可能得到不同的寄存器分配、不同的内联决策、不同的展开深度；这些都改变分支模式与缓存工作集，实测里最常见的表现就是 P99 不动而最坏值涨三成。因此工具链锁定不是洁癖，而是可界性论证的前提：锁定编译器版本与补丁串（把编译器自身的版本选项完整输出、以及工具链容器的镜像摘要记进构建记录）、锁定 binutils 与链接脚本、锁定所有影响代码生成的环境变量、锁定构建环境的区域设置与时区。构建系统上把工具链封进带摘要的容器镜像，在流水线里用摘要而不是浮动标签引用它；再把编译器与关键库的版本指纹、源码提交号、关键段校验和写进固件只读区，运行时可由版本查询读出，现场设备不必拆开就能核对。

\begin{lstlisting}[style=cmakestyle,caption={CMake：硬实时目标的编译开关成组与可复现构建要点},label={lst:ch13-cmake}]
# 关键环：确定性相关的开关成组出现，禁止在别处单独追加
add_library(ctrl_core STATIC ctrl/pid.c ctrl/filter.c ctrl/loop.c)

target_compile_options(ctrl_core PRIVATE
  -O2                      # 发布级；调试用另一份构建树，不用 -Og 的结论外推
  -g3                      # 保留宏与内联信息，便于事后反汇编比对
  -fno-exceptions -fno-rtti
  -fno-unwind-tables -fno-asynchronous-unwind-tables  # 不留隐藏展开表
  -fno-stack-protector     # 避免栈金丝雀引入额外访存
  -fno-strict-aliasing     # 与寄存器别名混用时取保守一侧
  -ffp-contract=off        # 禁止融合成 FMA，控制律结果不随选项漂移
  -fno-fast-math           # 显式否定，防止上游 preset 带入
  -fno-jump-tables         # 稀疏 switch 钉成比较链（按需）
  -fconserve-stack
  -Wall -Wextra -Wshadow -Wconversion -Wcast-qual -Werror
)

target_link_options(ctrl_core PRIVATE
  -Wl,--gc-sections        # 需配合 -ffunction-sections -fdata-sections
  -Wl,-Map,ctrl_core.map
  -Wl,--build-id=sha1
)

# 可复现：固定随机种子、抹掉路径前缀、关闭着色
target_compile_options(ctrl_core PRIVATE
  "-frandom-seed=$<TARGET_FILE_NAME:ctrl_core>"
  "-fdiagnostics-color=never"
  "$<$<CONFIG:Release>:-fmacro-prefix-map=${CMAKE_SOURCE_DIR}=>")

# 关键环不被过程间优化改写
set_target_properties(ctrl_core PROPERTIES INTERPROCEDURAL_OPTIMIZATION OFF)
\end{lstlisting}

\subsection{段回收、链接顺序与镜像形态}

未用段回收依赖把每个函数与每份数据切进独立 section，而回收结果对链接顺序敏感：库的先后次序改变，被保留的符号集合就可能改变，进而改变段的地址与缓存组分布。因此链接顺序要写进构建脚本而不是靠命令行习惯，map 文件必须作为构建产物归档。要特别警惕因为没被引用所以被回收的隐式行为：某张安全相关的表如果只在异常路径被引用，一次回收就可能把它删掉而无人察觉；这类符号要用保留属性或显式引用钉住，并在门禁脚本里核对它仍然存在。

\subsection{尺寸与 map 回归}

把 \Tp{size} 的分段输出、关键段的起止地址、关键环的反汇编摘要作为构建门禁项：任一超出容差即判失败并要求人工确认。这个做法的价值在于它把编译器换了变成一次显式评审，而不是上线后从抖动里反推。反汇编层面的比对可以用固定规则提取关键环的指令条数、分支数与访存次数，形成三元组指纹；三者同时不变，微架构层面的时间模型才允许复用。

\begin{bestpractice}[title={把工具链变更当成一次设计变更}]
换编译器小版本、换链接器、改优化预设、升级 SDK 里的启动文件，都要走与代码变更同级或更严的流程：独立的构建指纹对比、关键环反汇编比对、按第 11 章口径重测延迟与抖动分布，并把结论写进变更单。允许顺手升级工具链的团队，最终无法解释任何一次最坏值的抬升。
\end{bestpractice}

\section{静态检查与时间界分析}
\label{sec:ch13-static-wcet}

\subsection{规则集与工具定位}

MISRA C:2012 是嵌入式 C 最常用的强制规则集，它明确禁止标准库的动态内存分配函数，其余与时间相关的条目集中在指针算术、整数转换与未定义行为上；它没有一条通用规则禁止递归，实践中递归由项目补充规则或调用图深度限制来管。CERT C 与 CERT C++ 以安全编码声明组织，覆盖内存、整数、并发与异常，适合作为规则映射的另一套坐标。工具侧的分工大致是：\Lib{clang-tidy} 在编译数据库上跑可扩展检查（含 bugprone、cert、misc 等检查组），适合日常增量；\Lib{cppcheck} 独立解析、不需要完整编译环境，适合在没有目标工具链的机器上先过一遍，并可加载 MISRA 类附加组件；GCC 自带的 \Tp{-fanalyzer} 做跨过程路径分析，擅长发现未初始化、越界与双重释放，代价是编译时间与误报。这些工具的共同定位是把不可界的结构找出来，而不是证明时间。

\subsection{界分析工具的输入要求}

学术与商用的 WCET 界分析工具（aiT、Bound-T 一类产品与若干开源流程）给出的是上界而不是统计量，代价是对输入有硬性要求。典型清单如下：可解析的二进制或中间表示，加不含歧义的段与符号映射；处理器与存储层次模型（缓存容量、相联度、line 大小与替换策略，部分工具还要求紧耦合内存与外设访问的延迟表）；每个循环的上界，无法自动推断时必须由标注、不变式或断言形式提供；不可行路径的证据（谓词、输入范围或人工断言），否则结果被无用路径支配而失去意义；中断与抢占的处理方式，要么建模每次到达的代价，要么证明窗口内不可到达；以及排除未定义行为——程序一旦可能进入未定义行为，任何上界都不再成立。这也解释了为什么可界要写成一簇可机械检查的项：无递归、无动态分配、循环上限可证、热路径无外部间接调用、无未定义行为。

\begin{lstlisting}[style=cstyle,caption={C：把可界性写成可检查项：显式上限、静态断言与不可界结构的排除},label={lst:ch13-bounded}]
/* (1) 循环上限用编译期常量，便于界分析推断并与静态断言核对 */
#define MAX_TAPS   16u
#define MAX_FRAME  64u
_Static_assert(MAX_TAPS * sizeof(float) <= 256u, "滤波状态必须一屏可枚举");

float fir_fixed(float x, const float* c) {
    float acc = 0.0f;
    for (uint32_t i = 0; i < MAX_TAPS; ++i) { acc += c[i] * x; }  /* 上界：16 次 */
    return acc;
}

/* (2) 长度来自外部时先钳位，再把钳位值作为前提交给优化器 */
static uint32_t clamp_len(uint32_t n) { return n > MAX_FRAME ? MAX_FRAME : n; }
int32_t sum_frame(const int16_t* p, uint32_t n) {
    const uint32_t k = clamp_len(n);
    if (k > MAX_FRAME) { __builtin_unreachable(); }  /* 新版本可用 assume 属性表达同一前提 */
    int32_t s = 0;
    for (uint32_t i = 0; i < k; ++i) { s += p[i]; }
    return s;
}

/* (3) 评审中直接否掉的形态（对照检查清单逐项打勾）：
 *   递归调用自身            -> 改为定长迭代或显式栈数组
 *   动态内存分配函数         -> 一律来自静态对象池
 *   热路径上的格式化输出     -> 写入无锁事件环，离线消费（第 11 章）
 */

/* (4) 断言的取舍：静态断言留在编译期；运行期检查不返回而是进失效闩锁 */
#define RT_CHECK(cond) do { if (!(cond)) { fault_latch(__LINE__); } } while (0)
\end{lstlisting}

清单里 \Fn{fault\_latch} 的语义与第 4 章的失效处理一致：不返回、不参与正常路径的时间模型。这给出一个实用判据——热路径里的错误处理要么代价恒定，要么干脆不存在。

\subsection{控制流图与人工复核的配合}

即便不使用正式的界分析工具，控制流图的思路仍然有效：对关键环取出反汇编，画出基本块与边，标注每个循环的静态上界，逐条检查是否存在由输入决定的分支与跳出循环的隐藏出口（例如宏里的提前返回、或被到达处理改写的路径）。把这张图连同上界结论作为评审材料归档，再配合第 10 章的可调度性论证，就构成这一环在给定预算内一定走得完的最小证据链。

\section{硬件辅助观测与确定性设计检查清单}
\label{sec:ch13-hw-aux-checklist}

性能计数器（PMU）给出缓存缺失、分支误预测、流水线停滞等事件计数，配合第 2 章的打点可以把这一轮慢了归因到多了十二次 L1 数据缺失。Linux 侧的 \Lib{perf} 能直接读通用事件（\Tp{cache-misses}、\Tp{branch-misses}），具体核另有分层与访问类型的专用事件；裸机侧要自己写协处理器或系统寄存器访问，并注意计数器组的复用限制。

对 PMU 的信任边界必须说清楚。第一，计数器可能被复用或近似计数（部分核在超出可用计数器时按比例采样），据此得到的缺失次数带误差。第二，特权执行环境会让计数与真实延迟脱节：x86 的 SMM（系统管理模式）在操作系统与虚拟化层之下执行，任何用户态与内核态的测量都看不见它，却实实在在推迟实时任务，这是最坏值抖不出解释的经典来源，只能在 BIOS 或固件层面收敛（关闭不必要的 SMI 源，或把相关外设移出实时核）；ARM 侧的安全世界与虚拟化维护同理。第三，PMU 中断与性能采样本身会引入抖动，应作为被观察对象而不是背景。第四，SMT 让两个逻辑核共享取指、缓存与执行端口，邻居的访问模式直接决定本道的最坏延迟；硬实时核通常要么关闭 SMT，要么靠绑核与隔离保证同一物理核上没有第二负载，而这要靠第 6 至第 8 章的部署纪律兑现。

指令级追踪（ARM CoreSight ETM、RISC-V 追踪扩展与若干厂商实现）能还原执行过的指令序列与分支走向，是回归比对最有力的工具：同一激励下路径应当逐分支一致，一旦展开深度、向量化或内联决策改变，追踪流会先于时间戳暴露差异。代价是追踪带宽与缓冲区占用，通常只在标定与非产线构建里开启，并且要先确认追踪本身不改变实时性。

\begin{longtable}{>{\raggedright\arraybackslash}p{0.40\textwidth}>{\raggedright\arraybackslash}p{0.24\textwidth}>{\raggedright\arraybackslash}p{0.18\textwidth}}
\caption{确定性设计检查清单（发布前逐项签核）}\label{tab:ch13-checklist}\\
\toprule
检查项 & 判据出处 & 证据形式 \\
\midrule
\endfirsthead
\toprule
检查项 & 判据出处 & 证据形式 \\
\midrule
\endhead
\bottomrule
\endfoot
\endlastfoot
关键环代码与数据同在紧耦合内存或已锁定的路数 & \S\ref{sec:ch13-tcm} & map 文件的段归属记录 \\
热路径无数据相关分支，或已改为选择与查表 & \S\ref{sec:ch13-ctrlflow} & 反汇编基本块枚举 \\
热路径无间接调用：无虚函数、无类型擦除包装、无函数指针 & \S\ref{sec:ch13-compiler} & 调用图与符号表 \\
所有循环有编译期可证上界，无递归 & \S\ref{sec:ch13-static-wcet} & 静态检查报告加人工标注 \\
无动态分配，缓冲区尺寸与对齐固定 & \S\ref{sec:ch13-static-wcet} & 规则集零违例加 map \\
write-combining 与设备区域上不存在读回 & \S\ref{sec:ch13-memhier} & 内存类型表加代码评审 \\
编译开关成组且记录在构建树，关键环关闭 LTO & \S\ref{sec:ch13-compiler} & 构建配置原文 \\
工具链版本、容器摘要与构建指纹已写入固件 & \S\ref{sec:ch13-reprobuild} & 版本查询读出值 \\
镜像尺寸、map 差异与关键环指令三元组过门 & \S\ref{sec:ch13-reprobuild} & 流水线产物比对记录 \\
PMU 缺失计数与预算一致，SMM 与安全世界已收敛 & \S\ref{sec:ch13-hw-aux-checklist} & 事件计数与预算表 \\
热路径无日志、无格式化输出、无断点与调试钩子 & \S\ref{sec:ch13-static-wcet} & 静态检查加追踪回放 \\
生产构建的断言策略与失效闩锁一致 & 第 4 章口径 & 失效注入记录 \\
\end{longtable}

\subsection{常见陷阱}

把 \Fn{printf}、串口日志或断点留在热路径里，是本章所有问题中最贵也最常见的一种：一次格式化几十到数百周期，还会拖入锁、堆与缓冲区，串口更慢几个量级；正确形态是把观测面交给无锁事件环，离线消费。生产固件里开启 \Fn{assert} 等于把错误路径变成时间路径，而且断言表达式本身可能有副作用。全局构造（\Lang{C++} 静态对象初始化）在进入 \Fn{main} 之前按不易控制的顺序执行，若其中含外设时序敏感的准备，就与启动窗口纠缠；做法是把时间敏感的准备显式移进带序号的初始化阶段。大块 \Fn{memcpy} 常被编译器换成宽拷贝，把整段数据推进缓存并冲刷热路径的工作集——对确定敏感的系统，要么按固定小分块拷贝，要么走 DMA 加显式同步（第 5 章），要么使用平台提供的非临时存储。最后再强调一次：\Tp{-ffast-math} 让控制律的数值结果随编译器与目标选项变化，任何需要复现现场行为的产品都不该带着它发布。

\begin{warning}[title={把统计改善当成确定性提升}]
压测里看到平均执行时间下降，通常来自缓存与预取的收益，而最坏值往往原地不动甚至更差。判定标准只能来自两件事：结构上是否消除了不可界的路径（静态），以及最坏情形是否被复测并落在预算内（动态，按第 11 章口径）。只有后一件而没有前一件，就是赌运气。
\end{warning}

本章把时间不稳定拆成可分别处置的三层：存储层次、控制流、编译器改写。下一章转向流程侧——当这些结论需要被认证、被测试、被交接时，怎样组织证据，以及怎样为一个具体产品选到合适的平台档位。

% ============================================================
% ---- frag_ch14.tex ----
\chapter{功能安全、在环测试与选型收束}
\label{ch:ch14-safety-hil}

前面十三章解决的是怎样把时间做出来并证明它；本章处理剩下的三件事：怎样让这套结论被认证接受、怎样在真硬件与台架上系统地暴露问题、以及怎样为一个具体产品挑到合适的平台档位。本章的主线只有一条：标准约束的首先是流程与证据，代码层面的要求（语言子集、无动态分配、覆盖度指标）都是这条主线的派生物。理解了这一点，就不会出现代码写得很干净但认证不过的困惑。

\section{标准与等级：约束流程还是约束代码}
\label{sec:ch14-standards}

IEC 61508 是功能安全的横向基础标准，把安全完整性等级划成 SIL1 到 SIL4，等级越高对危险失效的每小时平均概率要求越严格，SIL4 的量级要到 10 的负九到负十次方。它同时给出硬件度量（失效概率、安全失效分数、诊断覆盖率）与软件度量（结构复杂度、需求与测试的完整性），并用低、中、高三档安全失效分数决定软件技术措施的严格程度。ISO 26262 把这套方法汽车化，等级称 ASIL A 到 ASIL D，由危害分析与风险评估按严重度、暴露概率与可控性推出。ASIL 允许分解：一个 ASIL D 的要求可以拆成两条相互独立、分别满足 ASIL B(D) 与 ASIL C(D) 的措施，前提是能证明共因失效被控制住——这条规则直接改变架构选择，因为 `` ``两条独立通道'' `` 往往比 `` ``一条通道做到 D'' `` 更容易落地，但也把独立性证据（电源、时钟、内存、软件分区）变成了硬任务。

IEC 62304 面向医疗器械，按软件造成的危害可能性把软件项分成 A、B、C 三类，类别决定文档与验证深度，并要求与风险管理流程联动。EN 50128 覆盖铁路控制与保护系统的软件，沿用 SIL0 到 SIL4（SIL0 表示不带安全功能），对开发、验证与维护的独立性有明确的角色要求。DO-178C 是机载软件的适航指南，按失效后果从 A（灾难性）到 E（无安全影响）分五级，等级不同需要满足的目标条数差别很大；它以证据为中心，Level A 要求做到对象代码层的 MC/DC（修正条件/判定覆盖），较低等级接受语句与分支覆盖，并要求需求追溯到代码与测试、测试结果追溯到需求。配套的 DO-330 定义工具鉴定等级（TQL 分档，从严到宽），DO-332 与 DO-333 则给出操作系统与已鉴定服务的复用路径——这就是 `` ``RTOS 的认证包在机载项目里值钱'' `` 的原因。

编码语言层面的实际约束通过方法与工具的可信性间接落地，业界收敛出的共识做法是：使用受限的语言子集（\Lang{C} 与 \Lang{C++} 各有常用子集规则，如 MISRA 系列），MISRA C:2012 明确禁止标准库的动态内存分配函数；不使用递归，也不让循环上界由外部输入直接驱动；避免未定义行为；覆盖度按等级要求做到语句、分支直至 MC/DC；并把静态检查与测试纳入配置管理。

\begin{longtable}{>{\raggedright\arraybackslash}p{0.17\textwidth}>{\raggedright\arraybackslash}p{0.16\textwidth}>{\raggedright\arraybackslash}p{0.24\textwidth}>{\raggedright\arraybackslash}p{0.24\textwidth}}
\caption{主要标准的作用域、等级与对代码的直接要求}\label{tab:ch14-standards}\\
\toprule
标准 & 等级体系 & 约束重心 & 对代码与测试的直接要求 \\
\midrule
\endfirsthead
\toprule
标准 & 等级体系 & 约束重心 & 对代码与测试的直接要求 \\
\midrule
\endhead
\bottomrule
\endfoot
\endlastfoot
IEC 61508 & SIL1 到 SIL4 & 流程、硬件度量与安全生命周期 & 软件措施随安全失效分数分档；结构化验证与工具可信性 \\
ISO 26262 & ASIL A 到 D，可分解 & 项级与硬件级安全分析联动 & ASIL D 要求 MC/DC；语言子集、无动态分配、独立性证据 \\
IEC 62304 & 软件项 A、B、C 类 & 风险管理驱动文档深度 & 类别越高验证越严；缺陷与变更须可追溯 \\
EN 50128 & SIL0 到 SIL4 & 独立验证角色与维护 & 需求到测试四向追溯；结构复杂度度量 \\
DO-178C & 目标 A 到 E & 目标证据与审定可审计性 & Level A 需 MC/DC 与对象代码覆盖；DO-330 工具鉴定 \\
\end{longtable}

\begin{warning}[title={追溯断裂是认证现场最常见的直接失败原因}]
评审现场不会因为某行代码复杂而否掉项目，而是因为 `` ``这条需求找不到实现、这个函数找不到需求、这次变更没有对应的回归记录'' ``。追溯一旦断裂，即使代码质量很高也要重做证据链。把追溯矩阵当成与源码同等重要的产物，并在流水线里检查它的完整性（每个安全需求至少一条测试用例编号、每个偏离至少一条审批记录），是唯一可扩展的做法。
\end{warning}

\section{流程、可追溯与工具鉴定}
\label{sec:ch14-process}

四向追溯指需求到设计、设计到代码、代码到测试的正向链，以及从任一节点回查其上游的反向链。落地形式不重要，重要的是可自动化检查：需求条目有唯一编号，代码侧用注释或专用属性引用，测试用例引用同一编号，脚本扫描三个集合的差集并把结果贴在合并请求上。变更控制与配置管理的核心是任何进入证据链的东西都有版本：源码、工具链容器摘要、脚本、台架自身的固件版本、被试件的构建指纹（第 13 章的做法），缺一不可。

工具鉴定回答的问题是你凭什么相信工具的输出。编译器把源码变成不是人所写的东西；静态分析工具既可能漏报也可能误报；测试框架与被测件共享运行时。判断是否需要鉴定看两点：工具能否引入或掩盖错误，以及人工检查能否替代该工具。据此，编译器与其输出验证方法需要证据（版本锁定、已知缺陷记录、反汇编抽查），静态检查工具的规则配置需要验证记录，测试框架与桩需要自测。可复现构建在这里发挥第二重价值：它让 `` ``这一版二进制确实由这一版源码与这一版工具链产出'' `` 成为可核查的事实，而不是一句承诺。

编码规范的落地靠三件东西：规则集进流水线（新增违例阻断，历史违例清单化管理）、偏离必须书面审批并带失效日期、评审记录可查。MISRA 的偏离机制正是为此设计的：偏离不是隐藏问题的手段，而是把判断显式化的地方。

缺陷分类与安全分析决定代码结构。FMEA 逐项列举失效模式与影响，FTA 从顶层危害向下用逻辑门串起原因，共因失效分析（\Lang{ISO} 26262 里的依赖性失效分析）关注看似独立实则共享的东西：同一段内存、同一个时钟、同一份工具链、同一个人写的两个模块。对应的代码结构措施是可见的：监控者与被监控者不共享堆与全局状态，两个通道不共用同一个第三方库版本，诊断代码与被诊断代码分开构建。硬实时系统还多一项时间上的共因——所有通道都受同一次总线仲裁延迟或同一次系统管理模式影响，这类失效只能靠第 12 章与第 13 章的隔离手段消解。

\section{测试金字塔在嵌入式上的形态}
\label{sec:ch14-tiers}

主机侧单元测试承担最大基数。做法是把硬件层抽象成薄接口，用桩或模拟对象替换：\Lib{Unity} 提供断言与运行器，\Lib{CMock} 依头文件生成模拟对象与期望，\Lib{Ceedling} 把三者与构建、覆盖率采集串起来；代码混用 \Lang{C++} 时，\Lib{googletest} 的 \Tp{TEST\_F} 与 \Tp{MOCK\_METHOD} 更顺手。这一层负责算法分支、边界值、状态机的非法转移，并确立同一份控制代码在桩上跑通这一事实。

\begin{lstlisting}[style=cstyle,caption={C：Unity 风格的单元测试与硬件层桩化（模拟对象由 CMock 依头文件生成）},label={lst:ch14-unity}]
#include "unity.h"
#include "mock_adc.h"         /* CMock 依 adc.h 生成：Expect / Ignore / StubWithCallback */
#include "mock_gpio.h"
#include "ctrl_loop.h"

#define TRACE_LEN 64u

static uint16_t s_adc_seq[TRACE_LEN];
static uint32_t s_adc_calls;

/* 回调签名与被测函数一致，尾部多一个调用序号，可当激励指针用 */
static uint16_t adc_stub(uint8_t ch, int cmock_num_calls) {
    (void)ch;
    s_adc_calls += 1u;
    return s_adc_seq[(uint32_t)cmock_num_calls % TRACE_LEN];
}

void setUp(void) {
    s_adc_calls = 0u;
    MockAdc_Init();
    MockGpio_Init();
}

void tearDown(void) {
    MockAdc_Verify();          /* 未满足的期望在这里失败，而不是被静默忽略 */
    MockAdc_Destroy();
    MockGpio_Verify();
    MockGpio_Destroy();
}

void test_loop_step_reads_exactly_one_sample(void) {
    for (uint32_t i = 0; i < TRACE_LEN; ++i) { s_adc_seq[i] = (uint16_t)(100u + i); }

    /* 整个用例只允许读一次配置寄存器：多读说明热路径混进了 MMIO */
    Gpio_ReadConfig_ExpectAndReturn(0x12u);
    Adc_ReadChannel_StubWithCallback(adc_stub);

    float out = 0.0f;
    for (uint32_t k = 0; k < 200u; ++k) { out = ctrl_loop_step(); }

    TEST_ASSERT_FLOAT_WITHIN(0.02f, 4.0f, out);
    TEST_ASSERT_EQUAL_UINT32_MESSAGE(200u, s_adc_calls, "每次控制步必须恰好一次采样");
}

int main(void) {
    UNITY_BEGIN();
    RUN_TEST(test_loop_step_reads_exactly_one_sample);
    return UNITY_END();
}
\end{lstlisting}

SIL（软件在环）把同一份控制代码放进模型环境：被控对象用数学模型或离线回放数据代替，输入输出与时间基准由桩提供，步长固定。它的价值在于能跑几千小时的仿真与参数扫描，也能把真实系统里难以复现的工况（传感器缓慢漂移、执行器老化）廉价造出来。SIL 的局限同样明确：时序、缓存、总线、中断全都不存在，任何关于时间的结论都不能从 SIL 得出。HIL（硬件在环）把被试件换成真实控制器与真实固件，台架负责被控对象仿真、故障注入与时间测量：商用实时仿真机承担高动态模型（电机、电力电子、整车动力学），低成本台架则用一台带 PREEMPT\_RT 的主机加 FPGA 或微控制器协处理器产生 PWM 与编码器信号并执行注入。HIL 是唯一能同时给出功能对不对与时间够不够的层级。再往上是全机集成、路试与场试，处理接口组合、环境应力与真实工况。

\begin{longtable}{>{\raggedright\arraybackslash}p{0.14\textwidth}>{\raggedright\arraybackslash}p{0.24\textwidth}>{\raggedright\arraybackslash}p{0.22\textwidth}>{\raggedright\arraybackslash}p{0.21\textwidth}}
\caption{单元、SIL、HIL 与全机各层级的覆盖对象、成本与不可代替性}\label{tab:ch14-tiers}\\
\toprule
层级 & 覆盖对象 & 典型设备与成本 & 不可代替之处 \\
\midrule
\endfirsthead
\toprule
层级 & 覆盖对象 & 典型设备与成本 & 不可代替之处 \\
\midrule
\endhead
\bottomrule
\endfoot
\endlastfoot
主机单元 & 算法分支、边界值、状态机非法转移 & 普通开发机与容器，边际成本近乎零 & 大规模用例与变异测试只在此层负担得起 \\
SIL & 控制律与模型的联合行为、参数扫描 & 仿真工作站；主要成本是模型开发工时 & 长时运行与危险工况的零风险重复 \\
HIL & 真实固件的功能、时序与故障响应 & 实时仿真机或自研台架加接口板与被试件 & 真实中断、总线、缓存与执行器闭环 \\
全机与路试 & 系统集成、环境应力、真实工况 & 样机、场地与安全装备，成本最高 & 台架上不存在的接口组合与人为因素 \\
\end{longtable}

\begin{compare}[title={SIL 与 HIL 的判据边界}]
SIL 回答逻辑是否正确：同一份控制代码在模型闭环里能否让被控量收敛，在参数边界与异常输入下是否保持预期的状态迁移。它的时间是虚拟步长，与真实执行时间没有对应关系。HIL 回答逻辑与时间是否同时成立：真实的采样时刻、真实的中断到达、真实的总线负载下，闭环仍然稳定且期限不丢。工程推论是：任何时序类缺陷（抖动、超时、错过窗口）只有在 HIL 与全机上成立才算结论；任何覆盖率与参数空间探索只有在单元与 SIL 层才负担得起。两层必须用同一份被测源码与同一组需求编号连接，否则会出现 SIL 全绿而 HIL 从未跑过同一条用例的假安全。
\end{compare}

覆盖度指标要按用途选工具。语句与分支覆盖用编译器自带的插桩：GCC 一族加 \Tp{-coverage}，产出 \Tp{gcov} 可读的数据，再用 \Lib{lcov} 采集、\Tp{genhtml} 出报告；\Lang{LLVM} 侧用 \Tp{-fprofile-instr-generate} 加 \Tp{-fcoverage-mapping}，配合 \Tp{llvm-profdata} 合并原始计数与 \Tp{llvm-cov} 出报告。MC/DC 要求判定内每个条件都能被证明独立影响结果，所需的最小用例集通常由商用的嵌入式测试与覆盖分析工具在插桩与用例生成时提供，通用开源覆盖率工具不直接给出该度量。测试有效性用变异测试评估：向源码注入小改动（改比较符、改常量、删语句），若现有用例杀不死变异体，说明覆盖是形式上的。面向 \Lang{C} 与 \Lang{C++} 的变异工具已有可用实现（在中间表示层与在源码层各有取向），这一手段在安全等级高的模块上尤其值得投入，因为它把 `` ``分支都走到了'' `` 变成 `` ``分支走到的方式有判别力'' ``。

\section{故障注入与鲁棒性验证}
\label{sec:ch14-fault-injection}

传感器侧要注入的是失效与漂移：开路、短路到电源或地、读数冻结、缓变偏置、噪声放大。做法是在模拟前端串可编程源与继电器矩阵，或在数字链路上改桩——SIL 里改模型参数，HIL 里改仿真机输出通道。判据不是有没有报错，而是进入安全状态的时间与路径：第 4 章定义的窗口在这里被实测。

总线侧注入丢帧、延迟、乱序、重复与错误帧。主机侧的以太网与 TSN 链路用流量整型：\Tp{tc} 命令的 \Tp{netem} 排队规则可叠加延迟及其抖动分布、丢包率、乱序与重复比例，用来验证第 12 章的配置在偏离理想时系统是否仍然守得住期限。CAN 侧用 SocketCAN 的虚拟接口做纯软件注入，真实总线则用支持错误帧的接口卡或专用注入工具产生错误帧与总线关闭，配合继电器做拔线与短接。电源侧用可编程电源做跌落、缓升缓降、瞬断与抖动，验证复位与上电时序、以及参数存储的完整性。内存位翻转：带 ECC 或奇偶校验的平台通常提供注错寄存器或保留区，可注入单比特与多比特错误，观察纠正与检出报告是否按预期产生、诊断覆盖率是否达标；Linux 侧可经内核的错误检测与纠正子系统读取计数，核对它是否被日志与看门狗正确反映。

时间相关失效是本章特有的科目：故意让某个周期超时（拖慢任务、增加负载、注入延迟），验证策略是丢帧、降级还是进入安全状态，以及降级是否可观测。错过期限之后的行为必须有需求条目，否则测试无从判定通过。看门狗与安全状态到达时间的实测沿用第 4 章的机制与第 11 章的判据口径：多次测量、报告分布、把上界与预算比对，不以一次最坏值定结论。

\begin{lstlisting}[style=pythonstyle,caption={Python：HIL 台架的故障注入调度脚本（传感器失效、总线延迟、电源跌落三类场景编排）},label={lst:ch14-hil}]
"""场景 = 注入时刻 + 激励动作 + 复位动作 + 期望判据；时刻一律用台架统一时间基准。"""
import time

SCENARIOS = [
    dict(name="sensor_open", at=1.20, dur=0.50,
         act=lambda b: b.ai_force(0, "open"),                 # 模拟前端强制开路
         reset=lambda b: b.ai_release(0),
         expect=lambda r: r.safe_state_ms <= 20.0 and r.fault_code == 0x21),
    dict(name="eth_delay_5ms", at=3.00, dur=2.00,
         # 在实时链路的以太网口上叠加 5 ms 均值、2 ms 抖动的延迟与 1% 丢包
         act=lambda b: b.tc("qdisc add dev eth0 root netem delay 5ms 2ms loss 1%"),
         reset=lambda b: b.tc("qdisc del dev eth0 root"),
         expect=lambda r: r.deadline_miss == 0 and r.period_p99_ms <= 1.05),
    dict(name="psu_dip_9v", at=6.00, dur=0.05,
         act=lambda b: b.psu_sweep(12.0, 9.0, 0.005),          # 缓降到 9 V
         reset=lambda b: b.psu_set(12.0),
         expect=lambda r: r.reboots == 0 and r.nvram_crc_ok),
]

def bench_loop(bench, recorder, scenarios, rounds=3):
    failed = []
    for n in range(rounds):
        for sc in scenarios:
            bench.restart_firmware()          # 每轮从头初始化，避免残留状态
            t0 = time.monotonic()
            bench.start_model(setpoint=sc.get("setpoint", 4.0))
            while time.monotonic() - t0 < sc["at"]:
                time.sleep(0.001)
            sc["act"](bench)
            t_hit = time.monotonic()
            time.sleep(sc["dur"])
            sc["reset"](bench)                                 # 无论判据如何都必须复位
            result = recorder.collect(since=t_hit, fields=["safe_state_ms", "deadline_miss",
                                                           "period_p99_ms", "fault_code",
                                                           "reboots", "nvram_crc_ok"])
            recorder.dump_csv("%s_r%d.csv" % (sc["name"], n), result)
            if not sc["expect"](result):
                failed.append((sc["name"], n, result))
    return failed

if __name__ == "__main__":
    from hil_bench import Bench, Recorder     # 台架驱动与采集：封装仿真机、电源与 tc 调用
    with Bench(profile="rig-a") as b, Recorder(b) as rec:
        bad = bench_loop(b, rec, SCENARIOS)
        for row in bad:
            print(row)
        raise SystemExit(1 if bad else 0)
\end{lstlisting}

\section{持续集成、发布与现场取证}
\label{sec:ch14-ci-release}

交叉编译矩阵的规模很快超过直觉：板型乘以工具链版本乘以配置档位。流水线任务的划分应当以产物可辨识为原则——每个组合产出带指纹的镜像、map 文件、尺寸报告、静态检查结果，并保留可供事后比对的原始反汇编。实时性回归绝不能跑在共享执行器上：虚拟化的时钟源、邻居负载与迁移会让抖动数据失去意义。做法是把少数独占硬件注册成自建执行器并打标签，任务用并发锁保证同一块板子同时只被一条流水线占用，门限沿用第 11 章的口径与第 10 章的预算。

\begin{lstlisting}[style=yamlstyle,caption={YAML：交叉编译矩阵与实时性回归的门限设置},label={lst:ch14-ci}]
name: fw-matrix
on:
  pull_request: {}
  push: { branches: [main] }

jobs:
  build:
    strategy:
      fail-fast: false
      matrix:
        board: [ctrl-m7, gateway-a53, vision-a72]
        toolchain: ["gcc-12.3", "gcc-13.2"]   # 版本随容器摘要一起锁定
    container:
      image: ghcr.io/example/xbuild@sha256:0000000000000000000000000000000000000000000000000000000000000000
    steps:
      - uses: actions/checkout@v4
      - run: cmake -B build -DBOARD=${{ matrix.board }} -DCMAKE_BUILD_TYPE=Release
      - run: cmake --build build
      - run: python3 tools/map_gate.py build/ctrl_core.map build.elf  # 尺寸、段归属、指令三元组
      - run: cppcheck --error-exitcode=1 --enable=warning,style,performance sources
      - uses: actions/upload-artifact@v4
        with: { name: fw-${{ matrix.board }}-${{ matrix.toolchain }}, path: build/dist/ }

  hil-latency:
    needs: build
    runs-on: [self-hosted, rig-a, exclusive]  # 独占台架，绝不使用共享执行器
    timeout-minutes: 90
    steps:
      - uses: actions/download-artifact@v4
        with: { name: fw-ctrl-m7-gcc-13.2, path: dist/ }
      - run: python3 tools/flash.py --port usb0 dist/firmware.bin
      - run: python3 tests/hil/latency_run.py --hours 2 --max-miss 0 --p99-budget-us 120
      - run: python3 tools/latency_gate.py report.json  # 门限来自第 11 章口径与第 10 章预算

  conformance:
    needs: build
    runs-on: [self-hosted, rig-b, exclusive]
    steps:
      - run: python3 tests/hil/fault_inject.py --scenario-file scenarios/nightly.json
\end{lstlisting}

发布侧的更新与回滚按平台规模分两种形态。带 MMU 的应用处理器上通常是 A/B 双分区：引导程序维护一小块带校验的元数据，记录活动槽、尝试次数与镜像摘要；启动后由健康判据（是否进入正常运行、是否完成一次有效心跳）确认新镜像，否则计数递增并在超限后回落到旧槽。防降级依赖单调版本计数，计数落在一次性熔丝或受保护的持久区，签名要同时覆盖镜像内容与计数值。更新过程断电必须是原子的：镜像写入非活动槽期间活动槽始终可启动，切换只发生在元数据的一次带校验写入上；双覆盖一类方案额外用一个小的交换表描述槽位重叠。微控制器侧常用双 bank 交换或补丁式更新，交换动作本身要在掉电保护区里留下可恢复的记录。

\begin{lstlisting}[style=cstyle,caption={C：A/B 槽选择与回滚判据（元数据带校验，尝试超限即回落安全态）},label={lst:ch14-aboot}]
#define SLOT_A            0u
#define SLOT_B            1u
#define BOOT_MAGIC        0x424f4f54u
#define MAX_BOOT_ATTEMPTS 3u

typedef struct {
    uint32_t magic, active_slot, attempt_count, image_crc, hdr_crc;
} BootMeta;   /* 落在一块独立的小容量持久区；写入前先擦除，写入后回读校验 */

/* 以下原语由板级支持包与安全存储提供 */
static bool meta_load(BootMeta* m);
static bool meta_store_atomic(const BootMeta* m);  /* 一次带校验的写入，掉电可回到旧值 */
static bool image_ok(uint32_t slot);               /* 校验长度与摘要签名 */
static uint32_t image_ver(uint32_t slot);          /* 从镜像头读版本 */
static uint32_t nv_monotonic_ver(void);            /* 一次性计数，防降级下限 */
static void start_image(uint32_t slot);
static void watchdog_short_period(void);

void boot_select(void) {
    BootMeta m;
    if (!meta_load(&m) || m.magic != BOOT_MAGIC) { m.active_slot = SLOT_A; m.attempt_count = 0u; }

    const uint32_t min_ver = nv_monotonic_ver();
    if (!image_ok(m.active_slot) || image_ver(m.active_slot) < min_ver) {
        const uint32_t other = SLOT_A + SLOT_B - m.active_slot;   /* 试另一槽 */
        if (image_ok(other) && image_ver(other) >= min_ver) { m.active_slot = other; }
        else { goto safe_mode; }
    }
    if (m.attempt_count >= MAX_BOOT_ATTEMPTS) { goto safe_mode; } /* 反复失败即不再重试 */

    m.attempt_count += 1u;                    /* 先记账再启动：崩溃也不会丢这次尝试 */
    if (!meta_store_atomic(&m)) { goto safe_mode; }
    start_image(m.active_slot);
    return;

safe_mode:
    watchdog_short_period();   /* 停在可被外部服务恢复的安全态，等待重新刷写 */
    for (;;) { }
}

/* 应用侧：进入正常运行并完成一次有效心跳后，把尝试次数清零 */
void boot_confirm_ok(void) {
    BootMeta m;
    if (meta_load(&m)) { m.attempt_count = 0u; (void)meta_store_atomic(&m); }
}
\end{lstlisting}

现场取证要求产品在失效之后仍然能被读懂：日志分区采用环形覆盖并保留时间基准与复位原因，RAM 里的事件缓冲在复位后由启动代码搬进持久区，崩溃转储（应用处理器上的核心转储与寄存器组，微控制器上的故障状态寄存器加栈快照）要有可导出的通道，并且必须与归档的镜像版本对应，否则事后无法符号化。取证机制本身要遵守第 11 章关于观测污染的结论：不得改变被观察对象的实时性。

\section{选型收束：四档平台与三个贯穿实例}
\label{sec:ch14-platform-choice}

选型的落点是把六个判据按优先级排出来：需要的响应延迟量级、外设与软件栈需求、内存与成本、认证要求、团队能力、长期维护与供货。其中延迟量级几乎总是第一个筛子，因为它决定其余选项是否可用。

\begin{longtable}{>{\raggedright\arraybackslash}p{0.16\textwidth}>{\raggedright\arraybackslash}p{0.22\textwidth}>{\raggedright\arraybackslash}p{0.22\textwidth}>{\raggedright\arraybackslash}p{0.21\textwidth}}
\caption{四档平台的判据矩阵}\label{tab:ch14-platform}\\
\toprule
档位 & 可达延迟量级与依据 & 主要代价与风险 & 适合的认证与团队画像 \\
\midrule
\endfirsthead
\toprule
档位 & 可达延迟量级与依据 & 主要代价与风险 & 适合的认证与团队画像 \\
\midrule
\endhead
\bottomrule
\endfoot
\endlastfoot
裸机前后台加中断 & 微秒级；路径由人穷举，最坏时间可直接论证 & 中间件要自带；规模变大后结构易失控 & 认证强度高、团队小、功能单一 \\
RTOS（\Lib{FreeRTOS}、\Lib{Zephyr}、\Lib{ThreadX}、\Lib{RTEMS}） & 十微秒到百微秒级抢占响应；取决于优先级与第 6 至第 8 章的纪律 & 配置复杂、驱动质量参差；需认证时依赖安全包 & 中等到高强度认证；需要栈与生态 \\
\Lang{Linux} 加 PREEMPT\_RT & 亚毫秒到毫秒级最坏延迟；依赖绑核与隔离 & 全量鉴定不现实；后台负载影响需长期管控 & 复杂栈、需网络与文件；团队有内核调优能力 \\
混合临界（Hypervisor 分区加 RTOS 域） & 实时域可达微秒级；隔离带来的额外代价小且可界定 & 调试与工具链门槛高；板级资源需求上升 & 安全与非安全共存；需复用已鉴定组件 \\
\end{longtable}

单核还是多核取决于隔离需求而非算力。多核上的实时性并不自动更好：共享 L2、内存控制器与总线的耦合会把一个核的抖动传给另一个核，除非采用 AMP 式分区（每核独立职责、明确的共享内存协议、按第 9 章与第 13 章处理一致性），或者 SMP 加严格绑核与中断亲和。总线选择同理：TSN 或工业以太网只在需要同步多设备、跨机柜确定性与诊断时才值得付出配置成本；单个控制环内部用板级直连更简单可靠。

\subsection{实例一：100 µs 周期的电机控制环}

需求是 100 µs 周期的电流环，含两个观测器、SVPWM 输出与故障保护，安全要求包含快速降磁与堵转保护。推导链是：100 µs 周期意味着端到端可用预算只有几十微秒，留四分之一裕度后约为 60 µs（其余留给传感与输出，见 \S\ref{sec:ch14-acceptance} 的分段口径），这个量级排除了任何带文件系统与通用网络的操作系统，也排除了依赖软件浮点或长延迟外设访问的设计。选定带硬件浮点与紧耦合内存的单核实时微控制器，控制环跑在由硬件定时器触发的最高优先级上下文里，ADC 与 PWM 同步触发并经 DMA 送入固定缓冲区。平台档位选裸机加一个轻量 RTOS：环本身用前后台式的硬实时上下文实现以保证没有调度开销，诊断、通信与状态管理交给 RTOS 的低优先级任务，于是第 8 章的负载隔离问题被限制在非关键任务上。代码按第 13 章把环函数与系数表钉进 ITCM 与 DTCM、关闭异常与 RTTI、消除环内数据相关分支。PWM 与故障输入用硬件保护单元（死区发生器的硬件故障输入直接关闭输出），使最坏保护时间不依赖软件。认证上这一档最容易给证据：调用图小、无动态分配、结构复杂度低。

\subsection{实例二：车载 CAN 网关心}

需求是八个 CAN 通道加两条 CAN FD、若干诊断功能与一条非关键的上云链路，路由延迟目标在亚毫秒，承担安全相关路由的通道按 ASIL B 论证，非关键的通信功能按 ASIL A 分区。这个规模下纯裸机会在通道组合与诊断状态机上失控，而应用处理器的启动时间、内存占用与鉴定负担都过头。选定多核或带锁步的实时微控制器加 RTOS：路由任务按通道优先级划分，每通道独立缓冲与静态内存池，中断只做取帧与时间戳（第 2 章与第 3 章的口径），表查询与转发在任务级完成；上云与诊断跑在同一芯片的低优先级核或域里，用硬件存储保护单元把两者内存隔开。总线侧不需要 TSN，但必须有高精度时间戳与帧序号，供路试复现。测试策略按本章层级：主机侧用桩替换 CAN 控制器寄存器，SIL 用录制回放覆盖路由表全组合，HIL 用双总线注入丢帧、错误帧与总线关闭，并实测最坏转发延迟。

\subsection{实例三：1 ms 级的工业视觉与控制}

需求是一路相机做缺陷检测，同时以 1 ms 周期驱动执行器，还要跑协议栈、配方存储与上位机接口。控制环需要微秒到亚毫秒的确定性，视觉需要内存、网络与算力，两类判据互相冲突，单档平台都不合适。选定混合临界：应用处理器上用 Hypervisor 分出实时域与非实时域，实时域跑轻量 RTOS 承载 1 ms 环与 IO 任务并独占若干核与外设（网卡队列、定时器），非实时域跑 \Lang{Linux} 做视觉、存储与协议；两域之间用共享内存加无锁环（第 9 章的纪律）与显式门限传递结果。网络侧值得上 TSN：相机与控制共网时需要时间同步与流量整形，第 12 章的配置与验证在这里成为前置条件。认证上采用组合论证——实时域按安全包构建，非实时域作为不贡献安全的分区并被监控，监控由独立看门狗与心跳通道实现。

\begin{bestpractice}[title={先定延迟量级，再谈平台偏好}]
按最坏响应延迟要求把候选砍到一档，再看外设与软件栈是否可用，最后用认证要求与团队能力校准可行性。反过来从熟悉的平台出发找理由，几乎一定会在集成后期付出重做架构的代价。
\end{bestpractice}

\section{上线前实时性验收与常见陷阱}
\label{sec:ch14-acceptance}

验收的核心工具是延迟预算分解表：把端到端路径拆成有限段，逐段给预算、指定实测方法与上限，任何一段无法实测就退回设计。表的读法是每段上限之和留出裕度后必须覆盖系统级要求；裕度的分配本身要有理由，而不是统一取一半。

\begin{longtable}{>{\raggedright\arraybackslash}p{0.22\textwidth}>{\raggedright\arraybackslash}p{0.18\textwidth}>{\raggedright\arraybackslash}p{0.22\textwidth}>{\raggedright\arraybackslash}p{0.19\textwidth}}
\caption{1 ms 控制路径的端到端延迟预算分解示例（数值为方法示意，须以实测替换）}\label{tab:ch14-budget}\\
\toprule
路径段 & 预算上限 & 实测方法 & 主要影响因素 \\
\midrule
\endfirsthead
\toprule
路径段 & 预算上限 & 实测方法 & 主要影响因素 \\
\midrule
\endhead
\bottomrule
\endfoot
\endlastfoot
传感器与调理到转换就绪 & 80 µs & 触发脚与转换完成脚同时打点 & 采样保持、前端带宽 \\
转换就绪到中断服务开始 & 60 µs & 第 3 章的最坏响应测量法 & 优先级、屏蔽窗口、总线等待 \\
中断到任务恢复运行 & 50 µs & 时间戳配对加第 11 章的分布口径 & 抢占、绑核、缓存冷启 \\
控制律执行 & 300 µs & 环首尾打点并叠加静态上界 & 第 13 章的存储与控制流项 \\
输出到执行器动作 & 120 µs & PWM 边沿与电流波形对照 & 驱动级延迟、硬件保护阈值 \\
总线传输到下一节点 & 150 µs & 报文时间戳与总线分析仪 & 仲裁、负载、整形配置 \\
余量与抖动吸收 & 140 µs & 由其余段实测上界反推 & 后台任务、日志、电源 \\
\end{longtable}

验收执行的最小口径：在独占硬件上连续运行足够时长并覆盖全部工况（以小时计的长测加数千轮场景重放，时长与轮次写进验收文档并成为门禁项），统计口径与图表沿用第 11 章；失败判定预先定义（任一期限丢失、任一安全状态到达超窗、任一尺寸或 map 门禁失败均判不过）；失败后的回退路径明确（回到上一发布镜像并保留现场取证产物）；遗留风险逐条记录，包含未验证的工况、被接受的偏离与其缓解措施，并由具备权限的角色签核。

\begin{guideline}[title={上线前逐项签核的验收清单}]
(1) 追溯矩阵完整且脚本可复核，每个安全需求至少一条自动化用例。(2) 静态检查零新增违例，偏离记录均在有效期内。(3) 覆盖度达到等级要求，关键模块另用变异测试验证用例有效性。(4) 工具链与容器摘要锁定，构建指纹已写入固件并通过尺寸、map 与指令三元组门禁。(5) 实时性回归在独占台架完成，端到端预算逐段有实测上界与裕度。(6) 故障注入场景全部通过，安全状态到达时间在窗口内。(7) 更新与回滚经过断电、降级攻击与尝试超限三类测试。(8) 取证通道可导出，且取证机制本身不改变实时性。(9) 非实时任务、后台日志与通信负载在满配下参与回归。(10) 遗留风险清单有签核与处置计划。
\end{guideline}

\subsection{常见陷阱}

第一类是证据与代码脱节：文档改了而代码没改、代码重构后追溯没重建，审定与评审时直接判失败，而且往往牵连整批交付物。第二类是只测快乐路径：用例的输入分布与真实工况重合度高，边界、并发与失效路径集中在少数几条上，覆盖数字漂亮但判据从未被独立触发。第三类是把主机测试通过率当成硬件结论：单元与 SIL 没有中断、没有缓存、没有总线，它们的绿色不代表时序成立，跨层级的结论必须重新测量。第四类是忽略非实时任务与后台日志对验收的影响：验收跑精简配置，现场跑满配，多出来的日志、监控与通信足以改变 P99 与最坏值，因此回归必须带全部后台负载。第五类是流水线用共享执行器得出假的抖动结论：虚拟化时钟、邻居负载与迁移产生的抖动被误读成代码退化，进而做出错误的架构决策。第六类是把看门狗当成万灵药而不实测到达时间，这与第 4 章的结论直接冲突。最后一类常见于选型：为将来可能的需求提前上更高一档的平台，结果团队能力与维护成本吃掉了裕度，实时性反而不如低一档的设计——第 6 章至第 8 章的所有复杂度都源自多活的东西，而复杂度本身就是最大的不确定性来源。

\begin{note}[title={把不可界当成需求缺陷处理}]
测试中出现无法解释的尾部延迟时，不要用延长测量来解决。第 13 章检查清单里任何一项不满足，都意味着时间模型不完整，正确的动作是回到设计层面补齐：把慢路径找出来、证明其上限、或者把它从关键路径上移走。延长观测只是让无法解释的东西藏得更久。
\end{note}

本章把前面十三章的技术结论放进流程与决策的框架：可界性来自结构与纪律，证据来自可复现的构建与独占的测量，选型来自延迟量级与团队能力的交集。全书的落点不是某一块芯片或某一个开关，而是这条能反复执行的判断链——先让最坏情况可列举，再让列举结果可测量，最后让测量证据可审计。

% ============================================================
% ---- frag_apA.tex ----
\appendix

\chapter{术语与缩写}

\label{ap:apA-glossary}

本附录收录本册正文实际用到的术语与缩写，按主题分成七组。
释义不追求词典式完备，只说明该条目在硬实时语境下为何重要、它以何种方式影响延迟与可界定性。
全称一律采用标准文本或上游文档中的原始写法，大小写与硬件手册保持一致；
同一概念在不同平台叫法不同时，正文取使用最广的一个，本表补充别名。
数值均为标准或器件手册给出的上限值，若与所用具体器件不符，以器件手册为准。

\section{硬件与微架构}

表~\ref{tab:apA-hw} 给出核内寄存器、中断控制器与存储结构相关条目。

\begin{longtable}{@{} >{\raggedright\arraybackslash}p{0.20\textwidth} >{\raggedright\arraybackslash}p{0.28\textwidth} >{\raggedright\arraybackslash}p{0.35\textwidth}@{}}
\caption{硬件与微架构术语}\label{tab:apA-hw} \\
\toprule
术语 & 全称 & 释义 \\
\midrule
\endfirsthead
\toprule
术语 & 全称 & 释义 \\
\midrule
\endhead
\bottomrule
\endfoot
\endlastfoot
\Tp{MMIO} & Memory-Mapped I/O & 外设寄存器映射进统一地址空间，实时路径上的访问次数与宽度必须固定，否则时序随编译器结果漂移。 \\
\Tp{CSR} & Control and Status Register & 控制状态寄存器，读写常伴随特权级切换，切换开销需计入最坏执行时间。 \\
\Tp{ISR} & Interrupt Service Routine & 中断服务程序，其长度与嵌套深度直接构成中断延迟，实时代码要求可静态界定其执行时间。 \\
\Tp{NMI} & Non-Maskable Interrupt & 不可屏蔽中断，不能被普通屏蔽位压下，多用于看门狗超时与硬件致命故障上报。 \\
\Tp{NVIC} & Nested Vectored Interrupt Controller & Cortex-M 的中断控制器，提供优先级分组、尾链与延迟抢占以缩短中断响应。 \\
\Tp{GIC} & Generic Interrupt Controller & Arm 通用中断控制器，负责中断路由、优先级仲裁与目标核分配。 \\
\Tp{SGI} & Software Generated Interrupt & 核间软件中断，是多核同步与远程唤醒的硬件通道，其投递延迟计入调度响应。 \\
\Tp{VTOR} & Vector Table Offset Register & 中断向量表基址寄存器，把向量表搬进紧耦合内存可减少异常入口取指缺失。 \\
\Tp{BASEPRI} & Base Priority Mask Register & 按优先级阈值屏蔽中断，用于以恒定开销进入极短临界区，比全局关中断更易界定。 \\
\Tp{PENDSV} & Pendable Service Request & 可挂起系统服务异常，内核常把上下文切换放在其中执行，以避开不可靠的栈状态。 \\
\Tp{DWT} & Data Watchpoint and Trace & 数据断点与跟踪单元，其周期计数器常用作实测最坏执行时间的最小侵入手段。 \\
\Tp{ETM} & Embedded Trace Macrocell & 指令跟踪宏单元，输出指令流供最坏路径与缓存行为分析使用。 \\
\Tp{CoreSight} & Arm debug and trace architecture & Arm 调试与跟踪体系，跟踪缓冲写回内存会干扰被测任务时序，测量时需扣除该开销。 \\
\Tp{TCM} & Tightly Coupled Memory & 紧耦合内存，访问延迟恒定且不经过缓存，适合放置向量表与实时任务栈。 \\
\Tp{ITCM} & Instruction Tightly Coupled Memory & 指令紧耦合内存，保证取指不受数据访问争用影响。 \\
\Tp{DTCM} & Data Tightly Coupled Memory & 数据紧耦合内存，常用于共享数据与直接内存访问的双缓冲。 \\
\Tp{CCM} & Core Coupled Memory & 核耦合内存，仅本核可达、不参与总线仲裁的私有存储区。 \\
\Tp{MPU} & Memory Protection Unit & 内存保护单元，在无内存管理单元的平台上以区域属性实现空间隔离，越界即刻触发异常。 \\
\Tp{MMU} & Memory Management Unit & 内存管理单元，提供地址转换与页属性，但转换失败即缺页，是主要延迟源之一。 \\
\Tp{TLB} & Translation Lookaside Buffer & 页表转换缓存，缺失引发的页表遍历可让单次访问耗时放大数十倍。 \\
\Tp{WC} & Write Combining & 写合并缓冲，适合顺序写入，但读回语义弱，实时驱动需显式刷缓冲才可见。 \\
\Tp{WB} & Write Back & 回写缓存，写命中只改缓存行，脏行回写时刻不定，会造成偶发长访问。 \\
\Tp{WT} & Write Through & 直写，写操作即时到达外设或内存，时序可预测但带宽代价高。 \\
\Tp{SMT} & Simultaneous Multithreading & 同时多线程，同核线程共享流水级使执行时间互相影响，硬实时配置一般关闭。 \\
\Tp{PMU} & Performance Monitoring Unit & 性能计数单元，用于量化缓存缺失与总线争用，但读取计数本身有成本。 \\
\Tp{SEU} & Single Event Upset & 单粒子翻转，位翻转可落在计数器或状态机，安全平台需纠错与周期校验。 \\
\Tp{ECC} & Error Correcting Code & 纠错码，单比特纠错、双比特检错，纠错与上报耗时需纳入实时预算。 \\
\Tp{BOR} & Brown-Out Reset & 欠压复位，电源跌落导致重启，恢复时间应计入故障容忍时间间隔。 \\
\Tp{PVD} & Programmable Voltage Detector & 可编程电压检测，在欠压复位之前产生早期告警中断，以便有序进入安全状态。 \\
\Tp{WFI} & Wait for Interrupt & 等待中断的停核指令，唤醒延迟与中断延迟叠加，硬实时路径常限制其深度。 \\
\Tp{WFE} & Wait for Event & 等待事件指令，需配合发送事件类指令使用，构成低开销的核间等待通道。 \\
\Tp{STOP} & Stop mode & 深度休眠模式，唤醒需重建时钟树，恢复时间远超一般硬实时截止期。 \\
\Tp{XIP} & eXecute In Place & 片上闪存直接执行，读速受限且擦写时常需暂停取指，直接放大抖动。 \\
\Tp{DMA} & Direct Memory Access & 直接内存访问，搬运不占核，但总线争用会拉长实时任务的访问时间。 \\
\Tp{FPU} & Floating-Point Unit & 浮点单元，核内或外置之分决定上下文大小，进而影响上下文切换耗时。 \\
\Tp{GPIO} & General Purpose Input/Output & 通用输入输出，控制回路的输入读取与输出置位需要固定访问次数。 \\
\Tp{EXTI} & External Interrupt or Event Controller & 外部中断控制器，边沿捕获到中断送达之间的传播延迟计入传感器响应。 \\
\Tp{AArch64} & Arm 64-bit Execution State & Arm 64 位执行状态，其异常模型与定时器指令决定实时移植的工作量。 \\
\end{longtable}

\section{时间、计数器与同步}

表~\ref{tab:apA-time} 覆盖时钟源、时间语义与延迟度量三类条目。
读一个时钟的代价、它能否被回拨、是否计入休眠时间，就是选择它的全部理由。

\begin{longtable}{@{} >{\raggedright\arraybackslash}p{0.20\textwidth} >{\raggedright\arraybackslash}p{0.28\textwidth} >{\raggedright\arraybackslash}p{0.35\textwidth}@{}}
\caption{时间、计数器与同步术语}\label{tab:apA-time} \\
\toprule
术语 & 全称 & 释义 \\
\midrule
\endfirsthead
\toprule
术语 & 全称 & 释义 \\
\midrule
\endhead
\bottomrule
\endfoot
\endlastfoot
\Tp{SysTick} & System Tick Timer & Cortex-M 的 24 位递减系统定时器，常作节拍源，位宽小需处理重装窗口。 \\
\Tp{Generic Timer} & Arm architected system timer & Arm 架构级定时器，各核从同一计数源取时，免除总线定时器带来的锁竞争。 \\
\Tp{CNTPCT} & Counter-timer Physical Count Register & 物理计数器寄存器，读取即得当前周期，是裸机上最经济的单调时钟源。 \\
\Tp{CNTFRQ} & Counter-timer Frequency Register & 计数器频率寄存器，运行时读取它可避免把频率硬编码进单位换算。 \\
\Tp{CNTVCT} & Counter-timer Virtual Count Register & 虚拟计数器，受偏移寄存器影响，虚拟化环境取时须确认偏移为零。 \\
\Tp{TSC} & Time Stamp Counter & x86 每核周期计数寄存器，非恒定频率核上的取值不可直接用于跨核比较。 \\
\Tp{HPET} & High Precision Event Timer & 高精度事件定时器，常作后备时钟源，读取需寄存器访问，开销远高于周期计数。 \\
\Tp{clocksource} & kernel clock source & 内核取时后端，实时系统需选定无副作用且频率稳定的源并禁止运行期切换。 \\
\Tp{PHC} & PTP Hardware Clock & 网卡硬件时钟，在收发瞬间打戳，是软件栈抖动之外的稳定时间基准。 \\
\Tp{PTP} & Precision Time Protocol & IEEE 1588 精密时间协议，局域网内主从同步可达亚微秒量级。 \\
\Tp{gPTP} & generalized Precision Time Protocol & 泛化精密时间协议，即 IEEE 802.1AS 为桥式网络规定的同步方案。 \\
\Tp{IEEE 802.1AS} & Timing and Synchronization for Time-Sensitive Applications & 桥式网络的同步标准，规定对等延迟测量与级联主从关系。 \\
\Tp{Announce} & PTP Announce Message & 主钟宣告报文，收敛期主从选举引起的不确定会短时破坏时间基准。 \\
\Tp{pdelay} & peer delay measurement & 邻居延迟测量，链路上下行不对称会使估计偏大，需在截止期预算中留余量。 \\
\Tp{NTP} & Network Time Protocol & 网络时间协议，毫秒量级，适合管理面而不是控制回路的时间基准。 \\
\Tp{PPS} & Pulse Per Second & 秒脉冲信号，作为硬件绝对参考，用于标定软件时间戳的偏差。 \\
\Tp{TAI} & International Atomic Time & 国际原子时间，不含闰秒的连续标度，控制回路应使用它而非墙钟。 \\
\Tp{UTC} & Coordinated Universal Time & 协调世界时，存在闰秒插入，跨越该时刻的差值计算需专门处理。 \\
\Tp{CLOCK\_MONOTONIC} & POSIX monotonic clock & 不受墙钟回拨影响的单调时钟，适合测量间隔，但不计入休眠时间。 \\
\Tp{CLOCK\_{\Bk}MONOTONIC\_{\Bk}RAW} & raw monotonic clock & 不受频率校正影响的原始计数，长期比对应改用经校正的时钟。 \\
\Tp{CLOCK\_REALTIME} & POSIX real-time clock & 墙钟，可被步进或回拨，用它表达绝对截止期存在语义风险。 \\
\Tp{CLOCK\_BOOTTIME} & suspend-aware monotonic clock & 计入挂起时间的单调时钟，适合与休眠共存的长期定时。 \\
\Tp{RTC} & Real-Time Clock & 日历时钟，掉电由电池维持，仅用于上电后设定初始墙钟。 \\
\Tp{jitter} & timing jitter & 事件实际时刻相对理想时刻的离散度，硬实时看最坏值而不是均值。 \\
\Tp{skew} & clock skew & 两个时钟在同一时刻的读数差，同步的目标是压缩其最坏值。 \\
\Tp{drift} & clock drift & 频率偏差随时间累积，决定两次同步之间允许的最大间隔。 \\
\Tp{holdover} & clock holdover & 参考丢失后靠本地振荡器守时，其斜率决定可用的时间窗口。 \\
\Tp{latency} & 延迟 & 事件从发出到被观测的耗时，实时语境必须区分最好、典型与最坏。 \\
\Tp{response time} & 响应时间 & 从任务就绪到完成的总耗时，是可调度性分析的直接对象。 \\
\Tp{deadline} & 截止期 & 任务必须完成的时间点，错失即在安全功能层面构成失效。 \\
\Tp{tick} & OS tick & 周期节拍，节拍越长定时粒度越粗，最小可测时间误差随之增大。 \\
\Tp{HZ} & kernel tick frequency & 内核节拍频率，提高它可缩小定时粒度，但增加节拍自身开销。 \\
\Tp{timestamp} & 时间戳 & 与事件绑定的时刻记录，硬件打戳优于软件打戳，后者含采集抖动。 \\
\end{longtable}

\section{中断、调度与 RTOS}

表~\ref{tab:apA-rtos} 列出调度算法、中断处理机制与内核统计条目。

\begin{longtable}{@{} >{\raggedright\arraybackslash}p{0.20\textwidth} >{\raggedright\arraybackslash}p{0.28\textwidth} >{\raggedright\arraybackslash}p{0.35\textwidth}@{}}
\caption{中断、调度与 RTOS 术语}\label{tab:apA-rtos} \\
\toprule
术语 & 全称 & 释义 \\
\midrule
\endfirsthead
\toprule
术语 & 全称 & 释义 \\
\midrule
\endhead
\bottomrule
\endfoot
\endlastfoot
\Tp{FIFO} & First In First Out & 先到先服务，就绪队列按入队顺序执行，长任务可使后到的任务等待。 \\
\Tp{RR} & Round Robin & 同优先级轮转，时间片长度决定同优先级任务的最坏等待时间。 \\
\Tp{EDF} & Earliest Deadline First & 最早截止期优先，动态优先级，利用率可达 1，但每次释放都要重新比较截止期。 \\
\Tp{RMS} & Rate Monotonic Scheduling & 周期越短优先级越高的静态优先级算法，是硬实时可调度性分析的基础模型。 \\
\Tp{DM} & Deadline Monotonic & 相对截止期越短优先级越高，当截止期不等于周期时其判据优于按周期排序。 \\
\Tp{CBS} & Constant Bandwidth Server & 以固定带宽为后备任务分配运行时，防止非关键任务耗尽关键预算。 \\
\Tp{Polling Server} & 轮询服务 & 周期性为异步实时任务注入执行时间，其注入周期本身即为额外延迟。 \\
\Tp{utilization bound} & 利用率判据 & 静态优先级的充分判据，任务数增加时固有利用率上界趋近约 0.69。 \\
\Tp{schedulability} & 可调度性 & 存在某个调度使全部截止期都被满足这一性质，须在实现之前给出证明。 \\
\Tp{RTA} & Response Time Analysis & 迭代求解最坏响应时间并与截止期比较，比利用率判据更紧且允许利用率超界。 \\
\Tp{WCET} & Worst-Case Execution Time & 最坏执行时间，一切实时计算的输入，取值必须覆盖缓存与分支的最差情形。 \\
\Tp{preemption} & 抢占 & 高优先级就绪即剥夺运行任务，抢占边界是共享资源竞争窗口的主要来源。 \\
\Tp{preemption point} & 抢占点 & 允许被抢占的程序位置，调度器只在安全点切换以保证数据结构一致。 \\
\Tp{priority inversion} & 优先级反转 & 低优先级任务持锁阻塞高优先级任务，反转时长无法仅由代码本身界定。 \\
\Tp{priority inheritance} & 优先级继承 & 持锁者临时继承等待者的优先级，缓解反转但锁链传递仍会延长阻塞。 \\
\Tp{priority ceiling} & 优先级天花板 & 取锁时把优先级升到该锁的最高引用优先级，可一次性消除反转与嵌套。 \\
\Tp{tail chaining} & 尾链 & 前一中断服务未完即直接进入下一入口，省去出栈入栈，降低中断风暴期开销。 \\
\Tp{late EOI} & late end of interrupt & 延迟发送中断结束，允许同优先级中断嵌套，代价是栈深度不确定性增大。 \\
\Tp{APIC} & Advanced Programmable Interrupt Controller & x86 中断控制器，局部件负责每核排队，投递延迟随拓扑层级增加。 \\
\Tp{MSI} & Message Signaled Interrupt & 消息信号中断，以内存写代替引脚请求，减少中断共享与线路延迟。 \\
\Tp{threaded IRQ} & 线程化中断 & 把中断后半部放进可调度的内核线程，使其在高优先级任务面前可被抢占。 \\
\Tp{softirq} & 软中断 & 中断上下文中的延迟处理，不保证串行化，实时路径须注意其抢占行为。 \\
\Tp{tasklet} & 任务项 & 基于软中断的串行化工作单元，同一实例不并发，适合短小收尾。 \\
\Tp{workqueue} & 工作队列 & 进程上下文的异步执行框架，默认优先级低，实时场景需专用高优先级队列。 \\
\Tp{cooperative priority} & 协作优先级 & 不主动抢占他人的线程类别，用于极短且绝不可被打断的操作序列。 \\
\Tp{time slice} & 时间片 & 同优先级轮转的配额长度，片越小响应越快但切换占比越高。 \\
\Tp{context switch} & 上下文切换 & 保存与恢复寄存器及栈指针，其耗时须计入响应时间与切换次数预算。 \\
\Tp{run-time stats} & 运行时间统计 & 逐任务执行时间记录，依赖高频计数源，测量自身会引入额外开销。 \\
\Tp{heap\_1} & minimal heap implementation & 只分配不释放的堆实现，零碎片、行为可预测，最适合硬实时产品。 \\
\Tp{heap\_5} & growable heap implementation & 可合并空闲块的堆实现，遍历空闲链的时长不确定，实时路径不应调用。 \\
\Tp{tickless idle} & 无节拍空闲 & 空闲期停节拍以省电，唤醒需补偿节拍，实时应用中禁用以免引入非周期延迟。 \\
\Tp{CPU load} & CPU 负载 & 单位时间非空闲占比，负载升高时最坏响应时间随抢占次数放大。 \\
\Tp{deadline miss} & 截止期错失 & 未能在截止期前完成，系统须能检测并逐次上报每一次错失。 \\
\Tp{overrun} & 超时执行 & 周期任务执行跨越下一释放点，需要显式的丢失计数与恢复策略。 \\
\Tp{interrupt latency} & 中断延迟 & 从请求到入口指令的耗时，包含屏蔽窗口与仲裁排队时间。 \\
\Tp{interrupt nesting} & 嵌套深度 & 同时活动的中断层数，决定中断栈的最坏需求量。 \\
\Tp{stack overflow check} & 栈溢出检测 & 在切换点校验栈水印，开销固定，需与最坏栈深评估配套使用。 \\
\end{longtable}

\section{并发、内存与语言}

表~\ref{tab:apA-conc} 覆盖无锁结构、内存序、屏障与语言子集条目。

\begin{longtable}{@{} >{\raggedright\arraybackslash}p{0.20\textwidth} >{\raggedright\arraybackslash}p{0.28\textwidth} >{\raggedright\arraybackslash}p{0.35\textwidth}@{}}
\caption{并发、内存与语言术语}\label{tab:apA-conc} \\
\toprule
术语 & 全称 & 释义 \\
\midrule
\endfirsthead
\toprule
术语 & 全称 & 释义 \\
\midrule
\endhead
\bottomrule
\endfoot
\endlastfoot
\Tp{SPSC} & Single Producer Single Consumer & 单生产者单消费者环形缓冲，仅凭两个索引即可无锁，耗时可静态界定。 \\
\Tp{MPSC} & Multi Producer Single Consumer & 多生产者单消费者队列，入队需原子比较交换，重试次数影响最坏路径。 \\
\Tp{MPMC} & Multi Producer Multi Consumer & 多生产者多消费者队列，实现复杂且最坏步数难以界定。 \\
\Tp{ABA} & ABA problem & 比较交换无法区分值回到旧值的过程变化，实时系统常用索引加版本号规避。 \\
\Tp{lock-free} & 无锁 & 保证系统整体推进，但不保证单个线程在有限步内完成，最坏时间仍可能不可界。 \\
\Tp{wait-free} & 等待自由 & 每个操作都在有界步数内完成，硬实时并发结构应优先选择这一类实现。 \\
\Tp{obstruction-free} & 阻塞自由 & 线程独占运行即可在有限步完成，仅提供前进性，不提供延迟上界。 \\
\Tp{fence} & 内存栅栏 & 约束访存可见顺序的指令，插入位置决定跨核数据交接是否正确。 \\
\Tp{memory\_{\Bk}order\_{\Bk}relaxed} & 宽松序 & 只保证原子性不保证顺序，适合统计计数，不能用于发布订阅式交接。 \\
\Tp{memory\_{\Bk}order\_{\Bk}acquire} & 获取序 & 其后的访问不会早于本次读，消费端借此看到生产端已写好的数据。 \\
\Tp{memory\_{\Bk}order\_{\Bk}release} & 释放序 & 其前的访问不会晚于本次写，生产端借此完成数据发布。 \\
\Tp{seq\_cst} & sequentially consistent & 最强顺序，代价是额外屏障，实时路径只在必要时使用。 \\
\Tp{DMB} & Data Memory Barrier & Arm 数据内存屏障，只排序不保证完成，写外设寄存器常需配合同步屏障。 \\
\Tp{DSB} & Data Synchronization Barrier & 数据同步屏障，等待此前访存真正完成，常用于启动直接内存访问之前。 \\
\Tp{ISB} & Instruction Synchronization Barrier & 指令同步屏障，冲刷流水线，用于改写向量表或控制寄存器之后。 \\
\Tp{sfence} & store fence & x86 存储栅栏，保证内存映射写按序到达设备。 \\
\Tp{false sharing} & 伪共享 & 不同核写同一缓存行的不同字节导致持续互相失效，按行对齐即可消除。 \\
\Tp{cache line} & 缓存行 & 一致性管理的最小单位，典型宽度 64 字节，实时结构体常显式按其填充。 \\
\Tp{object pool} & 对象池 & 预分配固定大小对象并循环复用，分配与释放耗时恒定，取代动态申请。 \\
\Tp{static allocation} & 静态分配 & 编译期确定存储，无碎片、无失败路径，是硬实时的默认选择。 \\
\Tp{memory fragmentation} & 内存碎片 & 动态分配使空闲块离散化，最坏申请时长随之不确定，是常见隐患。 \\
\Tp{placement new} & 定位 new & 在指定缓冲上构造对象，使生命周期与存储位置分离，析构需自行安排。 \\
\Tp{volatile} & volatile 限定 & 禁止优化掉读写，用于内存映射输入输出与中断共享标志，但不提供任何原子性。 \\
\Tp{atomic} & 原子操作 & 不可分割的读改写，实时中要求访问宽度受硬件单次访问保证。 \\
\Tp{data race} & 数据竞争 & 同一位置并发访问且至少一个为写且无同步，属未定义行为，与竞态条件不同。 \\
\Tp{deadlock} & 死锁 & 循环等待锁造成永久阻塞，硬实时用固定锁序与有界等待根除。 \\
\Tp{lock ordering} & 锁序 & 全局固定的取锁顺序，可在检测模式下断言验证，是消除死锁的主要手段。 \\
\Tp{RCU} & Read-Copy-Update & 读端几乎无开销的同步机制，写端回收延迟且依赖宽限期，实时常需卸载其收尾工作。 \\
\Tp{MISRA C} & Motor Industry Software Reliability Association C & 面向安全关键编码的规则集，核心目的是抑制未定义行为与不可界定的控制流。 \\
\Tp{language subset} & 语言子集 & 安全标准允许的构造集合，实用子集多排除异常、动态类型与运行期分配。 \\
\Tp{exceptions} & 异常 & 栈展开的耗时无界且与执行路径耦合，硬实时子集普遍禁用。 \\
\Tp{RTTI} & Run-Time Type Information & 运行期类型信息，体积与耗时无界，安全子集常禁用。 \\
\Tp{Security By Contract} & 契约式安全 & 以接口契约与静态约束替代运行期检查，避免为兜底而引入无界开销。 \\
\Tp{strict aliasing} & 严格别名 & 不同类型指针不得在同址读写，违反会使编译器生成错误顺序，驱动层需显式处置。 \\
\Tp{inline assembly} & 内联汇编 & 直达指令的控制手段，需正确书写约束，用于消除不确定的优化行为。 \\
\Tp{undefined behavior} & 未定义行为 & 标准不约束结果的构造，程序不含它是一切最坏时间分析的前提。 \\
\end{longtable}

\section{Linux 实时栈}

表~\ref{tab:apA-linux} 覆盖内核配置、线程模型、调优参数与测量工具条目。

\begin{longtable}{@{} >{\raggedright\arraybackslash}p{0.20\textwidth} >{\raggedright\arraybackslash}p{0.28\textwidth} >{\raggedright\arraybackslash}p{0.35\textwidth}@{}}
\caption{Linux 实时栈术语}\label{tab:apA-linux} \\
\toprule
术语 & 全称 & 释义 \\
\midrule
\endfirsthead
\toprule
术语 & 全称 & 释义 \\
\midrule
\endhead
\bottomrule
\endfoot
\endlastfoot
\Tp{PREEMPT\_RT} & full preemptible real-time model & 自 5.12 起进入主线的完全抢占模型，硬中断与自旋锁转为可调度的内核线程或锁。 \\
\Tp{NO\_HZ\_FULL} & tickless userspace & 单任务运行核上停周期性节拍，消除节拍带来的固定抖动源。 \\
\Tp{nohz\_full} & boot parameter & 指定哪些核进入无节拍用户态，需与核隔离同时设置，核数超界会被忽略。 \\
\Tp{rcu\_nocbs} & boot parameter & 把读复制更新的回调卸载到专用线程，避免隔离核上出现回收收尾工作。 \\
\Tp{isolcpus} & boot parameter & 把核从调度域摘出，普通任务不再落入，代价是须显式布置实时任务。 \\
\Tp{irqaffinity} & boot parameter & 设定中断默认落核集合，把设备中断集中到管家核是最小延迟做法。 \\
\Tp{cpuset} & control group cpu affinity controller & 按组约束可用核与内存节点，实时组的边界应稳定而不随负载变化。 \\
\Tp{cgroup} & control group & 资源受控分组机制，带宽限制设置过强会直接拉长实时任务的到达间隔。 \\
\Tp{RT throttling} & realtime runtime quota & 防止实时任务饿死内核的配额机制，超配额即被暂停，形成较长时间无法运行。 \\
\Tp{mlockall} & lock all memory & 一次性锁定当前与未来映射，杜绝运行期缺页，失败必须显式处理。 \\
\Tp{hugepage} & huge page & 大页减少转换条目与缺页次数，预留固定页池后行为可预测。 \\
\Tp{THP} & transparent huge pages & 自动大页的后台合并与回收是不可控延迟源，实时系统常关闭或设为延迟模式。 \\
\Tp{swap} & 交换 & 页被换出后访问代价达磁盘量级，实时核上必须彻底禁用。 \\
\Tp{SCHED\_FIFO} & realtime FIFO policy & 无时间片的实时策略，同优先级须显式让出，误用可使整核不可用。 \\
\Tp{SCHED\_DEADLINE} & deadline scheduling policy & 内核内的最早截止期优先与带宽预留实现，按运行时间、周期与截止期声明资源。 \\
\Tp{governor} & CPU frequency governor & 调频策略，实时常用固定频率的策略以消除升频延迟。 \\
\Tp{C-state} & idle state & 空闲状态进入与退出各有延迟，硬实时核常限制为最浅的睡眠状态。 \\
\Tp{P-state} & performance state & 标称性能状态，最低状态对应最保守的最坏执行时间假设。 \\
\Tp{HWP} & hardware managed performance & 硬件自主调频，介入时刻不完全受软件控制，评估时需额外留余量。 \\
\Tp{Turbo Boost} & 睿频 & 短时超频后回落，回落使同一段代码先后耗时不一致，测量须固定频率。 \\
\Tp{SMI} & System Management Interrupt & 系统管理中断，对内核不可见地抢占运行，是最大单点延迟源之一。 \\
\Tp{SMM} & System Management Mode & 系统管理模式，在独立地址空间执行，测量工具难以归因其中的耗时。 \\
\Tp{tuned} & adaptive system tuning daemon & 配置管理守护进程，以实时方案统一设置内核参数并降低后台服务干扰。 \\
\Tp{cyclictest} & cyclic latency test tool & 周期唤醒延迟测量工具，报告每核最坏延迟与延迟直方图。 \\
\Tp{oslat} & user space latency tool & 仅测用户态线程延迟，排除内核调度噪声，用于验证隔离配置效果。 \\
\Tp{hwlatdetect} & hardware latency detector & 以忙轮询发现系统管理中断类硬件停顿的内核模块与配套工具。 \\
\Tp{rtla} & realtime analysis tool & 内核自带的实时分析套件，包含定时延迟与噪声观测两类应用。 \\
\Tp{timerlat} & timer latency application & 测量定时线程在硬中断、内核线程与用户态各级的唤醒延迟分布。 \\
\Tp{ftrace} & function tracer & 内核跟踪框架，事件与函数记录自带时间戳，用于事后延迟分解。 \\
\Tp{trace-cmd} & trace interface utility & 采集与打包跟踪数据的命令行工具，便于长时间录制与离线分析。 \\
\Tp{KernelShark} & trace visualiser & 图形化事件时序浏览，用于定位长延迟窗口前后的干扰来源。 \\
\Tp{perf} & performance counters tool & 采样计数工具，采样中断本身干扰时序，实时测量须压低采样频率。 \\
\Tp{schedstat} & scheduler statistics & 调度统计接口，含等待与切换耗时，可作负载与干扰的复核数据。 \\
\Tp{inotify} & inode change notification & 文件系统事件通知接口，不保证送达时限，实时路径不应依赖它。 \\
\Tp{O\_DIRECT} & direct I/O flag & 绕过页缓存直接传输，减少缓存管理，但要求地址对齐且缓冲区常驻。 \\
\Tp{io\_uring} & io uring interface & 基于共享环形队列的低开销异步接口，提交与完成可绑定硬件时间戳。 \\
\Tp{AF\_XDP} & eXpress Data Path socket family & 网卡直达用户空间的零拷贝收发通道，绕开协议栈以缩短网络延迟。 \\
\Tp{DPDK} & Data Plane Development Kit & 用户态轮询式网络栈，以独占核换取微秒以下的处理延迟。 \\
\Tp{ETF} & earliest departure time first qdisc & 按绝对发送时刻出队的排队规则，配合发送时刻选项实现周期性发送。 \\
\Tp{SO\_TXTIME} & socket option for transmit time & 让报文在指定的硬件时刻发出，是确定性网络出口的关键接口。 \\
\Tp{SO\_TIMESTAMPING} & socket timestamping option & 在收发路径打软件或硬件时间戳，用于回路延迟测量与时间同步。 \\
\end{longtable}

\section{总线、网络与工业协议}

表~\ref{tab:apA-bus} 给出车载与工业现场总线、确定性以太网及相关标准条目。

\begin{longtable}{@{} >{\raggedright\arraybackslash}p{0.20\textwidth} >{\raggedright\arraybackslash}p{0.28\textwidth} >{\raggedright\arraybackslash}p{0.35\textwidth}@{}}
\caption{总线、网络与协议术语}\label{tab:apA-bus} \\
\toprule
术语 & 全称 & 释义 \\
\midrule
\endfirsthead
\toprule
术语 & 全称 & 释义 \\
\midrule
\endhead
\bottomrule
\endfoot
\endlastfoot
\Tp{CAN} & Controller Area Network & 多主仲裁总线，标称位速率上限 1 Mb/s、载荷 8 字节，非破坏性仲裁使高优先级标识报文先行。 \\
\Tp{CAN FD} & CAN with Flexible Data-Rate & 数据段可提速并把载荷扩展到 64 字节，改善单帧信息量与回路周期。 \\
\Tp{SAE J1939} & CAN higher layer for heavy vehicles & 商用车高层协议，多帧传输使端到端延迟随报文数线性增长。 \\
\Tp{LIN} & Local Interconnect Network & 低成本主从总线，传输时刻由调度表决定，不宜承担硬实时闭环。 \\
\Tp{FlexRay} & flexible automotive real-time bus & 周期与事件混合的时触总线，槽位预先静态规划因而延迟有界。 \\
\Tp{ARINC 429} & avionics digital data bus & 航电点到点单向总线，字间隔固定，带宽上限直接决定刷新周期。 \\
\Tp{SPI} & Serial Peripheral Interface & 主从同步串行接口，速率由主机时钟决定，适合片上传感器高速采样。 \\
\Tp{I2C} & Inter-Integrated Circuit & 双线共享总线，时钟拉伸与仲裁使传输时刻不精确，不应放在关键路径。 \\
\Tp{UART} & Universal Asynchronous Receiver Transmitter & 异步串行接口，无共享时钟，字符级延迟受波特率限制。 \\
\Tp{SDIO} & Secure Digital Input Output & 存储与外设一体接口，内部擦除使访问时延出现长尾。 \\
\Tp{RS-485} & balanced multipoint standard & 差分多点物理层，常承载现场总线协议，方向切换开销计入回路周期。 \\
\Tp{Modbus TCP} & Modbus over Industrial Ethernet & 主从式工业以太网协议，无时触机制，请求响应延迟随网络负载波动。 \\
\Tp{EtherCAT} & Ethernet for Control Automation Technology & 帧经过每个节点就地转发并处理，周期可达微秒级并自动补偿传播延迟。 \\
\Tp{DC} & EtherCAT distributed clocks & 分布式时钟机制，把各节点时钟对齐到主站，实现运动控制级同步。 \\
\Tp{PROFINET RT} & Real-Time class & 标准以太网上的实时类别，靠优先级排队，不保证硬截止期。 \\
\Tp{PROFINET IRT} & Isochronous Real-Time class & 预留时槽的等时通信类别，面向运动控制的抖动要求。 \\
\Tp{POWERLINK} & Ethernet POWERLINK & 主站广播加等时同步的实时以太网，循环周期可短至数百微秒。 \\
\Tp{Sercos III} & open real-time Ethernet interface & 实时以太网运动与自动化接口，提供带时间同步的服务通道。 \\
\Tp{CC-Link IE} & Industrial Ethernet open network & 工业以太网开放网络，循环传输面向控制器之间的数据交换。 \\
\Tp{TSN} & Time-Sensitive Networking & IEEE 802.1 下的确定性以太网标准集合，融合同步、调度、整形与冗余。 \\
\Tp{IEEE 802.1Qbv} & Enhancements for Scheduled Traffic & 时间感知整形，按周期门控队列，保证关键流窗口内不被抢占。 \\
\Tp{IEEE 802.1Qav} & Forwarding and Queueing Enhancements & 基于信用的整形，限制突发流量占用带宽以保护关键流。 \\
\Tp{IEEE 802.1Qbu} & Frame Preemption & 帧抢占，允许高优先级短帧打断长帧发送以压低最大延迟。 \\
\Tp{IEEE 802.1Qci} & Per-Stream Filtering and Policing & 逐流过滤与管制，在入口阻止超预算的流破坏整体确定性。 \\
\Tp{IEEE 802.1CB} & Frame Replication and Elimination for Reliability & 帧复制与消除，用冗余路径在不增加延迟的前提下提高可用性。 \\
\Tp{IEEE 802.1Qcc} & Stream Reservation Protocol Enhancements & 流预留的集中式与分布式配置控制模型，决定上线前能否完成预算验证。 \\
\Tp{PRP} & Parallel Redundancy Protocol & 双网并行发送并按序列号去重，故障切换延迟为零但带宽翻倍。 \\
\Tp{HSR} & High Availability Seamless Redundancy & 环上双向发送的无缝冗余协议，适用于工业环形拓扑。 \\
\Tp{OPC UA} & Open Platform Communications Unified Architecture & 跨层数据与服务模型，实时性能取决于所选传输绑定与订阅参数。 \\
\Tp{DDS} & Data Distribution Service & 以数据为中心的发布订阅规范，提供截止期、可靠性与流控策略。 \\
\Tp{AUTOSAR} & Automotive Open System Architecture & 汽车软件架构标准，其 OS 提供时触与周期任务模型和静态调度表。 \\
\Tp{PWM} & Pulse Width Modulation & 脉宽调制输出，周期与占空比精度取决于定时器分辨率与中断抖动。 \\
\Tp{HIL} & Hardware In the Loop & 被控对象由实时仿真机充当，仿真步长与总线周期必须严格对齐。 \\
\Tp{SIL test} & Software In the Loop & 控制软件在通用计算机上与模型闭环，注意与安全完整性等级使用同一缩写。 \\
\end{longtable}

\section{安全、质量与流程}

表~\ref{tab:apA-safety} 列出功能安全标准、覆盖度量、硬件指标与更新机制条目。

\begin{longtable}{@{} >{\raggedright\arraybackslash}p{0.20\textwidth} >{\raggedright\arraybackslash}p{0.28\textwidth} >{\raggedright\arraybackslash}p{0.35\textwidth}@{}}
\caption{安全、质量与流程术语}\label{tab:apA-safety} \\
\toprule
术语 & 全称 & 释义 \\
\midrule
\endfirsthead
\toprule
术语 & 全称 & 释义 \\
\midrule
\endhead
\bottomrule
\endfoot
\endlastfoot
\Tp{IEC 61508} & Functional Safety of Electrical/{\Bk}Electronic/{\Bk}Programmable Electronic Safety-Related Systems & 基础功能安全标准，按安全完整性等级规定软硬件与开发过程要求。 \\
\Tp{SIL} & Safety Integrity Level & 安全完整性等级，分 1 至 4 共四级，等级越高要求的验证强度与硬件指标越严。 \\
\Tp{ISO 26262} & Road Vehicles Functional Safety & 道路车辆功能安全标准，按 ASIL A 至 ASIL D 组织开发要求与措施。 \\
\Tp{ASIL} & Automotive Safety Integrity Level & 汽车安全完整性等级，由危害分析与风险评估结合严重度、暴露概率与可控性导出。 \\
\Tp{IEC 62304} & Medical Device Software Life Cycle Processes & 医疗器械软件生命周期标准，按软件安全分级规定开发与维护活动。 \\
\Tp{EN 50128} & Railway Applications Communication Signalling and Processing Systems Software & 铁路信号软件标准，定义软件完整性等级与相应的验证深度。 \\
\Tp{DO-178C} & Software Considerations in Airborne Systems Certification & 机载系统软件认证标准，以目标、软件问题级别与审查方式组织证据。 \\
\Tp{DO-254} & Design Assurance Guidance for Airborne Electronic Hardware & 机载电子硬件设计保证指南，与机载软件标准配套使用。 \\
\Tp{DO-330} & Software Tool Qualification Considerations & 工具鉴定指南，给出工具鉴定等级的判定方法。 \\
\Tp{TQL} & Tool Qualification Level & 工具鉴定等级，取决于工具对安全的影响程度以及其错误被后续检查发现的难易。 \\
\Tp{SWCL} & software criticality level & 机载软件问题级别，决定目标数量与结构覆盖要求的高低。 \\
\Tp{MC/DC} & Modified Condition/Decision Coverage & 修正条件判定覆盖，要求每个条件都能独立影响判定结果，最高等级常用它。 \\
\Tp{statement coverage} & 语句覆盖 & 最低的覆盖度量，不足以证明实时分支在运行期确实可达。 \\
\Tp{branch coverage} & 分支覆盖 & 每个判定的真假分支均被执行，与修正条件判定覆盖的差距在于条件独立性。 \\
\Tp{FMEA} & Failure Mode and Effects Analysis & 失效模式与影响分析，自下而上枚举失效并给出检测与缓解措施。 \\
\Tp{FTA} & Fault Tree Analysis & 故障树分析，自上而下用逻辑门组合基本事件以求顶事件概率。 \\
\Tp{CCF} & Common Cause Failure & 共因失效，冗余通道共享设计、时钟或电源时，失效并非相互独立。 \\
\Tp{HARA} & Hazard Analysis and Risk Assessment & 危害分析与风险评估，是确定安全等级与实时预算的起点。 \\
\Tp{FFI} & Freedom From Interference & 免干扰，混合关键性系统用分区、内存保护与时间预算证明低等级失效不波及高等级。 \\
\Tp{FTTI} & Fault Tolerant Time Interval & 故障容忍时间间隔，从失效发生到进入安全状态所允许的最大时间，实时预算据此制定。 \\
\Tp{safe state} & 安全状态 & 失效后主动进入的状态，定义它并证明可达是实时设计的一部分。 \\
\Tp{fail-safe} & 失效安全 & 检测到失效即停止功能，代价是可用性下降，但结果对人是安全的。 \\
\Tp{fail-operational} & 失效可运行 & 失效后维持降级运行，通常需要冗余通道与免干扰论证。 \\
\Tp{fail-silent} & 失效静默 & 失效后停止对外输出错误数据，由其它节点判断其缺席。 \\
\Tp{safe library} & 安全库 & 经过验证可直接使用的库，使用前需确认目标频率与输入域与本项目匹配。 \\
\Tp{static analysis} & 静态分析 & 不执行代码查找缺陷与规则违背，工具输出需经假阳性处理后才可用于鉴定。 \\
\Tp{sanitizer} & runtime instrumentation checker & 编译期插桩的运行时检查器，代价极高，只用于开发期与验收测试。 \\
\Tp{mutation testing} & 变异测试 & 向代码注入缺陷以检验测试有效性，是覆盖数字之外的质量证据。 \\
\Tp{traceability} & 可追溯性 & 需求、代码与用例之间的双向映射，工具可辅助但无法自动保证语义正确。 \\
\Tp{PMHF} & Probability of Mission Threatening Hardware Failure & 威胁任务的硬件失效概率，最高等级要求平均每小时不超过 10 FIT。 \\
\Tp{SPFM} & Single Point Fault Metric & 单点故障度量，衡量单点失效被安全机制覆盖的比例。 \\
\Tp{LFM} & Latent Fault Metric & 潜在故障度量，评估双重失效中第二个失效被诊断出的程度。 \\
\Tp{FIT} & Failures In Time & 失效单位，每 10 亿小时一次，把硬件目标数字落到可核对的量纲。 \\
\Tp{IWDG} & Independent Watchdog & 独立看门狗，由独立低频时钟驱动，主时钟失效后仍能复位。 \\
\Tp{WWDG} & Window Watchdog & 窗口看门狗，喂狗只允许落在指定时间窗内，因而能同时发现过早与过晚两种异常。 \\
\Tp{task watchdog} & 任务看门狗 & 由软件为每个任务登记活性凭证，全部登记被按期刷新才允许喂硬件狗。 \\
\Tp{OTA} & Over The Air & 空中远程升级，更新窗口与截止期冲突需专门设计回退时机。 \\
\Tp{A/B partition} & dual bank update & 双分区更新，新镜像写入备用区，校验通过才切换，失败可回退到原分区。 \\
\Tp{code signing} & 代码签名 & 对固件签名并验签后才允许写入备用区，须与防降级配合才构成纵深防御。 \\
\Tp{anti-rollback} & 防降级 & 拒绝安装版本号低于当前值的镜像，防止攻击者退回已知有漏洞的旧版本。 \\
\end{longtable}

% ============================================================
% ---- frag_apB.tex ----
\chapter{API 与参数速查}

\label{ap:apB-api-quickref}

本附录把正文里反复出现的工具开关、API 与配置项集中成可抄的表。按栈分组：交叉工具链与构建、链接脚本、FreeRTOS、Zephyr、Linux 实时、测量取证工具、TSN 与总线参数、编码纪律检查项。表中``实时注意点''一栏优先写会造成抖动、溢出或未定义行为的那一条，而不是通用用法说明。

\section{交叉工具链与构建}

工具前缀由目标三元组决定，脚本里应把前缀做成变量而不是硬编码。编译开关的选取直接决定 WCET 分析的可信度：优化等级每上一档，最坏路径就可能换一条。

\begin{longtable}{>{\raggedright\arraybackslash\footnotesize}p{0.24\textwidth} >{\raggedright\arraybackslash\footnotesize}p{0.30\textwidth} >{\raggedright\arraybackslash\footnotesize}p{0.29\textwidth}}
\caption{交叉工具链与关键编译/链接开关速查}\label{tab:apB-toolchain} \\
\toprule 名称 & 作用 & 实时注意点 \\ \midrule
\endfirsthead
\toprule 名称 & 作用 & 实时注意点 \\ \midrule
\endhead
\bottomrule \endfoot
\endlastfoot
\Tp{arm-none-eabi-gcc} & Cortex-M/R 与 ARMv7-A 裸机 GCC 主驱动 & 换主版本前先比 \Tp{arm-none-eabi-gcc -dumpmachine} 与随附 newlib 版本，二者一起决定 ABI \\
\Tp{arm-none-eabi-g++} & 同上，C++ 前端 & 异常与 RTTI 关掉后容器和 \Tp{dynamic\_cast} 在链接期才报错，别只测编译 \\
\Tp{aarch64-none-elf-*} & 64 位 bare-metal 前缀族（gcc/objcopy/size 同构） & AArch64 裸机缺省 startfile 不适用，向量表与 \Tp{TCR/MAIR} 要自写 \\
\Tp{riscv64-unknown-elf-gcc} & RV32/RV64 bare-metal ELF 工具链 & \Tp{-march} 与 \Tp{-mabi} 必须成对（如 rv32imafdc 配 ilp32d），错配会静默退回软浮点 \\
\Tp{clang} \Opt{target=arm-none-eabi} & LLVM 前端做同一目标 & 需配 \Tp{-ffreestanding} 与显式 sysroot；libstdc++ 头配 newlib 运行时是常见 ABI 事故源 \\
\Tp{arm-none-eabi-objcopy} & ELF 转 bin/hex，裁剪与填充 & \Tp{--gap-fill 0xFF} 与 \Tp{--pad-to} 决定空洞字节；\Tp{.bss} 不落在镜像里，别指望上电清零 \\
\Tp{arm-none-eabi-size} & 输出 text/data/bss 尺寸（\Tp{-A} 为 SysV 格式） & 每次改优化等级或加开关后对比三段，data+bss 才是 RAM 预算 \\
\Tp{arm-none-eabi-nm} & 列符号与大小（\Tp{-S --size-sort}） & 看不到栈上临时对象，栈预算要另用 \Tp{-fstack-usage}；回溯地址也靠它给出的区间 \\
\Tp{arm-none-eabi-readelf} & \Tp{-a}、\Tp{-l}、\Tp{-S} 看段表与属性 & 核对 \Tp{Tag\_CPU\_arch} 与浮点 ABI；属性不一致时混链不报错只跑飞 \\
\Tp{-mcpu=cortex-m7} & 选定内核，隐含调度模型与可用指令 & 只改 \Tp{-mcpu} 不改 \Tp{-mfpu} 时浮点配置不跟着变；FPU 未使能则浮点走慢路径 \\
\Tp{-march=armv7e-m} & 指令集基线（clang 常要求显式给） & 与 \Tp{-mcpu} 冲突时以显式 \Tp{-march} 为准，易产出不可执行编码 \\
\Tp{-mfpu=fpv5-d16} & 声明 FPU 形态 & Cortex-M4/M7 常见组合是 \Tp{fpv4-sp-d16} 或 \Tp{fpv5-sp-d16}，写错只在用到双精度时炸 \\
\Tp{-mfloat-abi=hard} & 浮点参数走 FPU 寄存器 & 与 \Tp{softfp} 混链会``链接通过、结果错''；中断里用硬浮点须确保上下文保存 FPU 寄存器 \\
\Tp{-Os} & 以尺寸为主的优化 & 尺寸赢、时间不稳；热函数可用函数属性单独升到 \Tp{-O2}，但要重新量最坏执行时间 \\
\Tp{-Og} & 保留可调试性的优化 & 只给调试版；\Tp{-Og} 版与发布版路径不同，抖动数据不可互相引用 \\
\Tp{-O2} & 通用性能档 & 会展开循环、合并条件判断，最坏路径可能变长，需复查循环体有界性 \\
\Tp{-ffunction-sections} & 每个函数独立成段 & 必须与 \Tp{-Wl,--gc-sections} 配对，否则只增段不减肥 \\
\Tp{-fdata-sections} & 每个全局独立成段 & 让链接脚本能把高频数据单独放进 TCM/DTCM，粒度越细越要写清归属 \\
\Tp{-flto} & 跨编译单元内联与死代码消除 & 未被 \Tp{KEEP()} 保护的向量表与注册段可能被吞；\Tp{-fstack-usage} 报告不再逐 TU 可比 \\
\Tp{-fno-exceptions} & 去掉异常表与 landing pad & 栈展开表消失后最坏栈深度可静态估；残留 \Tp{throw} 在链接期报错 \\
\Tp{-fno-rtti} & 去掉 typeinfo 与 vtable 运行时信息 & 关掉后 \Tp{dynamic\_cast} 与 \Tp{typeid} 不可用，HAL 抽象要改编译期分派 \\
\Tp{-fno-tree-vectorize} & 禁止自动向量化（NEON/SIMD/MVE） & 让执行时间随数据量近似线性，便于最坏时间推导；真要并行就手写并单独验证 \\
\Tp{-fstack-usage} & 每函数输出 .su 文件给出栈用量 & 不含内联与间接调用开销，LTO 后失真；取包络要再留余量 \\
\Tp{-{\Bk}fstack-{\Bk}protector-{\Bk}strong} & 栈金丝雀检查 & 函数进出各多几条指令，开销可测且稳定；与 \Tp{-{\Bk}Os} 同用会吃掉尺寸收益 \\
\Tp{-ffast-math} & 允许重排浮点、放松 IEEE 语义 & 结果依赖编译顺序；控制回路的浮点算法不要用 \\
\Tp{-Wall} \Tp{-Wextra} & 常规诊断集 & 实时代码还要盯未初始化与隐式转换，缺省级别看不全 \\
\Tp{-Wconversion} & 隐式窄化与符号转换告警 & 寄存器与位域读写处告警量大，建议只对驱动目录开启 \\
\Tp{-Werror} & 告警即错误 & 只在 CI 固定工具链版本时开，避免升级编译器后整树阻塞 \\
\Tp{-std=c17} & C 语言基线 & 需要 \Tp{atomic\_int}、\Tp{static\_assert} 时别退回旧标准 \\
\Tp{-std=c++20} & C++ 基线 & \Tp{consteval} 与 \Tp{concepts} 可把参数校验搬到编译期，运行期零成本 \\
\Tp{-T script.ld} & 指定链接脚本 & 换芯片先改脚本；脚本缺符号会让 \Tp{.data} 拷贝悄悄失效 \\
\Tp{-Wl,--gc-sections} & 丢弃未被引用的段 & 被向量表或注册宏引用的段必须 \Tp{KEEP()}，否则当垃圾回收 \\
\Tp{-Wl,-Map=firmware.map} & 输出段与符号布局 & 每次发布存档 map，回溯故障地址要靠它定位函数区间 \\
\Tp{-{\Bk}Wl,-{\Bk}-{\Bk}print-{\Bk}memory-{\Bk}usage} & 链接末尾打印各区占用 & 只看百分比不够，向量表所在 Flash 页的余量要单独盯 \\
\Tp{-nostartfiles} & 不用标准启动文件 & 自备向量表与 \Fn{\_start}，\Tp{.data} 与 \Tp{.bss} 的初始化全部自己写 \\
\Tp{-specs=nosys.specs} & 为无系统调用桩补齐 libc 依赖 & \Fn{printf} 走半主机或空实现，实时路径绝不能依赖 \\
\Tp{-specs=nano.specs} & newlib-nano 缩小 libc & 浮点格式需额外 \Tp{-u \_printf\_float}，且内部缓冲更大 \\
\end{longtable}

\section{链接脚本与段}

链接脚本决定符号是否落在正确的存储区，也决定启动代码能否把数据搬到位。第一列是脚本里的关键字、段名或符号名。

\begin{longtable}{>{\raggedright\arraybackslash\footnotesize}p{0.24\textwidth} >{\raggedright\arraybackslash\footnotesize}p{0.30\textwidth} >{\raggedright\arraybackslash\footnotesize}p{0.29\textwidth}}
\caption{链接脚本关键字段与启动符号速查}\label{tab:apB-linker} \\
\toprule 名称 & 作用 & 实时注意点 \\ \midrule
\endfirsthead
\toprule 名称 & 作用 & 实时注意点 \\ \midrule
\endhead
\bottomrule \endfoot
\endlastfoot
\Tp{MEMORY} & 声明物理区域（FLASH/RAM/TCM/OTP）与长度 & 长度按数据手册写并扣掉保留页，写大就会静默重叠 \\
\Tp{SECTIONS} & 描述输出段布局与归属 & 段的书写顺序即布局顺序，有对齐要求的段要显式 \Fn{ALIGN} \\
\Tp{AT$>$} & 让加载地址与运行地址（VMA）分离 & 只给 VMA 而漏 \Tp{AT$>$flash} 时 \Tp{.data} 初值不在镜像里，上电读到的是垃圾 \\
\Tp{$>$flash} & 当前段归入名为 flash 的 MEMORY 区 & 若段内含后续要写的数据或函数指针表，就不能放 flash \\
\Tp{$>$ram} & 当前段归入 ram 区 & 与 \Tp{AT$>$flash} 联用时启动代码必须负责搬运 \\
\Tp{KEEP()} & 强制保留段，即使未被引用 & 向量表、注册段、异常表都要包起来；开了 \Tp{--gc-sections} 尤其关键 \\
\Tp{.isr\_vector} & 中断/异常向量表，须位于 0 或 VTOR 指向处 & 表搬到中段却没设 VTOR 是最典型的一进中断就 HardFault \\
\Tp{.text} & 代码与指令常量 & 把热函数用段属性单独收进 TCM 前，先确认调用者也在同区 \\
\Tp{.rodata} & 只读常量、字符串、查找表 & 大查找表跨 cache 行会带来命中率方差，可按访问频率分区 \\
\Tp{.data} & 已初始化可写数据 & 需要启动代码从加载地址拷到运行地址；漏拷时读到的初值全不可信 \\
\Tp{.bss} & 零初始化可写数据 & 只占 RAM 不占镜像，启动必须清零，否则旧内容被当有效数据用 \\
\Tp{.\_user\_heap\_stack} & ST 风格脚本里对 heap 与栈合区的检查段 & 与 \Tp{\_estack} 一起改，两处不一致时只在溢出才暴露 \\
\Tp{\_estack} & 栈顶（高地址，向下增长） & 太小立刻崩，太大侵占 heap；应从栈用量包络反推 \\
\Tp{\_sidata} & \Tp{.data} 在 flash 中的起始（源）地址 & 由链接器生成，名字要与启动代码引用完全一致 \\
\Tp{\_sdata} 与 \Tp{\_edata} & \Tp{.data} 运行区间界 & 拷贝按半开区间写，闭区间会多拷一个字 \\
\Tp{\_sbss} 与 \Tp{\_ebss} & \Tp{.bss} 区间界 & 按字清零时若 \Tp{.bss} 未 4 对齐会漏尾字节 \\
\Tp{PROVIDE}() & 提供可被覆盖的符号默认值 & 注册表边界符号用 PROVIDE，避免重复定义 \\
\Tp{ASSERT}() & 链接期检查条件，失败即中断构建 & 用来断言向量表所在区、heap 与栈不重叠、余量下限，比事后调查看得早 \\
\end{longtable}

最小骨架的要点是``两区三段两界一断言''：

\begin{lstlisting}[style=cstyle,caption={链接脚本最小骨架要点（区域、段、启动符号、断言）},label={lst:apB-linker-skel}]
MEMORY {
  FLASH (rx)  : ORIGIN = 0x08000000, LENGTH = 1024K
  RAM   (rwx) : ORIGIN = 0x20000000, LENGTH = 256K
}
_estack = ORIGIN(RAM) + LENGTH(RAM);
SECTIONS {
  .isr_vector : { KEEP(*(.isr_vector)) } > FLASH
  .text : { *(.text*) *(.rodata*) } > FLASH
  .data : AT(ADDR(.text) + SIZEOF(.text)) {
      _sdata = .; *(.data*) _edata = .; } > RAM
  .bss : { _sbss = .; *(.bss*) *(COMMON) _ebss = .; } > RAM
  ._user_heap_stack : { . = ALIGN(8); heap_end = .;
      . = . + _Min_Heap_Size; . = . + _Min_Stack_Size;
      _end = .; } > RAM
  ASSERT(heap_end < _estack, "heap overlaps stack")
}
\end{lstlisting}

\begin{note}[title={常见误用清单}]
一是把可写段放进 FLASH 区，链接通过、运行写失败；二是只给 VMA 不给 \Tp{AT$>$}，\Tp{.data} 初值缺失；三是新增注册段却没有 \Tp{KEEP()}，配合 \Tp{--gc-sections} 静默消失；四是 \Tp{\_estack} 与 \Tp{\_Min\_Stack\_Size} 各改一半，栈预算被算两遍或零遍；五是 heap 与栈共用一段却没有 \Fn{ASSERT}，只在压测时才炸。
\end{note}

\section{FreeRTOS API 与配置宏}

FreeRTOS 里带 \Tp{FromISR} 后缀的版本只可在中断中调用，普通版本在中断里调用属于未定义行为。配置宏要成组改：动了 tick 频率，就必须重算所有以 tick 表达的超时。

\begin{longtable}{>{\raggedright\arraybackslash\footnotesize}p{0.24\textwidth} >{\raggedright\arraybackslash\footnotesize}p{0.30\textwidth} >{\raggedright\arraybackslash\footnotesize}p{0.29\textwidth}}
\caption{FreeRTOS API 与配置宏速查}\label{tab:apB-freertos} \\
\toprule 名称 & 作用 & 实时注意点 \\ \midrule
\endfirsthead
\toprule 名称 & 作用 & 实时注意点 \\ \midrule
\endhead
\bottomrule \endfoot
\endlastfoot
\Fn{xTaskCreate} & 动态创建任务，返回 pdTRUE/pdFALSE & 栈以``字''为单位；对象来自 heap，创建失败要在启动阶段处理，实时期不再建任务 \\
\Fn{xTaskCreate{\Bk}Static} & 静态创建任务，调用者提供栈与 TCB & 无分配器依赖，适合安全关键路径；需开 \Tp{configSUPPORT\_{\Bk}STATIC\_{\Bk}ALLOCATION} \\
\Fn{vTaskDelete} & 删除任务并回收其资源 & 被删任务持有的 mutex 与队列不会自动释放，资源归属要在设计里写明 \\
\Fn{vTaskSuspend} & 挂起任务 & 与阻塞的区别是恢复无就绪延迟；挂起时若持有 mutex 会让高优先级任务反转 \\
\Fn{vTaskResume} & 恢复挂起任务 & 若任务其实处于阻塞态则恢复无效，表现为``按了没反应'' \\
\Fn{vTaskDelay} & 相对超时让出 CPU & 粒度是 tick，实际落在 $n-1$ 到 $n+1$ 个 tick 区间，不要用来做微秒级定时 \\
\Fn{xTaskDelayUntil} & 按绝对时刻做周期延后 & 周期任务首选；入参时刻被更早的抖动拖后时需自行处理错过周期，否则相位漂移 \\
\Fn{uxTask{\Bk}Get{\Bk}Stack{\Bk}HighWaterMark} & 返回历史最小剩余栈（字） & 必须在压出峰值后读；配合栈溢出检测更可靠 \\
\Fn{vTaskGetRunTimeStats} & 打印每任务 CPU 占比与直方图 & 依赖运行统计开关与高频计数源，自身有开销，量产版常关掉 \\
\Tp{taskYIELD} & 请求重新调度 & 只在临界区外生效，临界区内置 pending，退出才切换 \\
\Fn{xSemaphoreGiveFromISR} & 中断里释放信号量 & 需配 \Tp{portYIELD\_FROM\_ISR} 决定是否立刻切换；二值语义不累计次数，会丢事件 \\
\Fn{xQueueSendFromISR} & 中断里入队 & 队列满返回 pdFALSE，中断里必须处理返回值，否则静默丢数据 \\
\Fn{vTask{\Bk}Notify{\Bk}Give{\Bk}FromISR} & 中断里给单任务发计数通知 & 比信号量快且省对象；只适合一个接收者，任务侧等待时要显式指定清除哪些位 \\
\Tp{portYIELD\_FROM\_ISR} & 退出中断前触发上下文切换 & 漏掉它会把响应延迟放大到下一个 tick，实测就是毫秒级毛刺 \\
\Fn{xSemaphore{\Bk}Create{\Bk}Binary} & 建立二值信号量 & 初值为空，需先释放一次才可取；不累计事件次数 \\
\Fn{xSemaphoreCreateMutex} & 建立互斥量（带优先级继承） & 不可在中断里使用；继承只在确实发生阻塞时生效 \\
\Fn{xSemaphore{\Bk}Create{\Bk}Counting} & 建立计数信号量 & 最大值即容量，溢出直接丢事件；用作 ISR 到任务的事件计数时容量按峰值算 \\
\Fn{xSemaphoreTake} & 获取，超时 0 表示轮询 & \Tp{portMAX\_DELAY} 是无超时，实时路径上会掩盖逻辑死锁 \\
\Fn{xSemaphoreGive} & 释放 & 已达上限时返回失败，循环 give 不会``攒出''更多许可 \\
\Fn{xEventGroupSetBits} & 置位事件组若干位 & 有 \Tp{FromISR} 版本；多位组合的原子性只保证在一次调用之内 \\
\Fn{xEventGroupWaitBits} & 等待若干位（与/或、可选自动清零） & 判位与清零是一步，之后再读会看到别的任务刚置的位 \\
\Fn{xTimerCreate} & 建软件定时器，回调在 timer 任务里执行 & 回调时长决定所有定时器的精度，绝不在回调里干重活 \\
\Fn{xTimerStart} & 启动或重装定时器 & 命令是排队给 timer 任务的，命令队列满则不生效，必须看返回值 \\
\Fn{xQueueCreate} & 建立队列（拷贝语义） & 开销是长度乘元素字节数；大结构体改为传指针入队 \\
\Fn{xQueueSend} & 队尾入队（另有 \Fn{xQueueSendToBack}、\Fn{xQueueSendToFront}） & 队满会阻塞或超时，超时后不要假定旧数据仍与帧对齐 \\
\Fn{xQueueReceive} & 出队 & 与 \Fn{xQueuePeek} 的差别是否消费；peek 后不消费会造成流水线滞后 \\
\Tp{configCPU\_CLOCK\_HZ} & 内核计数频率 & 与实际系统时钟不同步时，tick 与延时同时失准 \\
\Tp{configTICK\_RATE\_HZ} & 系统 tick 频率（常取 1000） & 提高 tick 让延迟分辨更细，但空闲中断开销线性上升 \\
\Tp{configMAX\_PRIORITIES} & 优先级数量 & 数值越大优先级越高；每多一级多一套就绪链表 \\
\Tp{configMINIMAL\_{\Bk}STACK\_{\Bk}SIZE} & idle 等内部任务的最小栈（字） & 加了钩子或运行统计后 idle 栈需求上升，是常见隐性溢出点 \\
\Tp{configTOTAL\_HEAP\_SIZE} & 供 heap\_N 使用的总堆字节数 & 只在启动期建对象；碎片情况靠 \Fn{xPortGetFreeHeapSize} 一类接口看不全 \\
\Tp{configUSE\_PREEMPTION} & 抢占式/协作式开关 & 置 0 即协作式，长任务会独占 CPU，需保证每处主动让出 \\
\Tp{configUSE\_TIME\_SLICING} & 同优先级时间片轮转 & 关掉后同优先级任务可能饿死；硬实时一般不让关键任务同优先级 \\
\Tp{configUSE\_MUTEXES} & 使能互斥量与优先级继承 & 关掉后 mutex 创建失败，改用二值信号量等于放弃继承 \\
\Tp{configUSE\_{\Bk}COUNTING\_{\Bk}SEMAPHORES} & 使能计数信号量 & 不开则计数信号量创建断言失败 \\
\Tp{configUSE\_IDLE\_HOOK} 与 \Tp{configUSE\_TICK\_HOOK} & 空闲钩子与 tick 钩子 & tick 钩子在中断上下文运行，禁止阻塞与日志 \\
\Tp{configUSE\_TICKLESS\_IDLE} & 低功耗停 tick & 唤醒时间计入延迟；软件定时器与绝对延时的分辨力随之变化 \\
\Tp{configCHECK\_{\Bk}FOR\_{\Bk}STACK\_{\Bk}OVERFLOW} & 栈溢出检测，1/2 两档 & 靠边界字节标记，一次压栈跨过整个缓冲区的``静默溢出''查不到 \\
\Tp{configUSE\_{\Bk}MALLOC\_{\Bk}FAILED\_{\Bk}HOOK} & 分配失败钩子 & 失败后系统多半继续错误运行，钩子里应进安全态而不是随意重启 \\
\Tp{configGENERATE\_{\Bk}RUN\_{\Bk}TIME\_{\Bk}STATS} & 运行时间统计 & 需要额外高频时基与两个配套宏定义，缺一个则统计恒 0 \\
\Tp{configASSERT} & 配置期与运行期断言入口 & 缺省空实现等于把 API 误用变成静默失败，至少接到日志与安全循环 \\
\Tp{heap\_1.c} & 只分配不释放，最小最确定 & 永不碎片、分配时间可静态给出；适合一次性建好全部对象的项目 \\
\Tp{heap\_2.c} & 可释放，首次适配且不合并 & 碎片永久存在，长周期运行的延迟分布不可预测 \\
\Tp{heap\_3.c} & 仅最坏情况优先（包装 libc \Fn{malloc}） & 与其他模块同时使用 libc \Fn{malloc} 时互相蚕食 \\
\Tp{heap\_4.c} & 首次适配并合并空闲块 & 分配耗时随碎片变化，实时路径只允许启动期使用 \\
\Tp{heap\_5.c} & 多存储区合并管理，策略同 heap\_4 & 可把 TCM 与主 RAM 分块注册，注意各区对齐与栈区分离 \\
\end{longtable}

\section{Zephyr}

Zephyr 的对象多为``先定义、编译期登记''，运行期创建走另一套接口。线程优先级分协作类与抢占类，数值语义与 FreeRTOS 相反，移植时最容易记错。

\begin{longtable}{>{\raggedright\arraybackslash\footnotesize}p{0.24\textwidth} >{\raggedright\arraybackslash\footnotesize}p{0.30\textwidth} >{\raggedright\arraybackslash\footnotesize}p{0.29\textwidth}}
\caption{Zephyr API、构建与配置项速查}\label{tab:apB-zephyr} \\
\toprule 名称 & 作用 & 实时注意点 \\ \midrule
\endfirsthead
\toprule 名称 & 作用 & 实时注意点 \\ \midrule
\endhead
\bottomrule \endfoot
\endlastfoot
\Tp{K\_THREAD\_DEFINE} & 编译期定义静态线程（栈、入口、优先级、周期） & 参数按位置生效，错位一个就串到别的字段；\Tp{K\_NO\_SCHEDULE} 表示不做周期让出 \\
\Tp{K\_THREAD\_STACK\_DEFINE} & 为线程声明静态栈区 & 以字节计，还要加 ISR 栈与接口调用余量；溢出通常是随机寄存器值 \\
\Fn{k\_thread\_create} & 运行期创建线程 & 需开动态线程支持；实时期建线程会引入不确定分配 \\
\Fn{k\_thread\_priority\_set} & 修改线程优先级 & 立即改变就绪顺序，临时提权后务必改回 \\
\Fn{k\_sleep} & 按 \Tp{K\_MSEC} 一类相对时长休眠 & 粒度取决于 tick 频率；亚 tick 的睡眠会被上取整 \\
\Fn{k\_busy\_wait} & 忙等微秒（参数为 k\_useconds\_t） & 不做上下文切换，占住 CPU；超过几个微秒就该改成等待对象 \\
\Fn{k\_usleep} & 微秒级睡眠 & 在 tickless 平台上未必比 \Fn{k\_sleep} 细，取决于时基；别用来做高精度节拍 \\
\Tp{k\_work} & workqueue 上的一次性工作项 & \Fn{k\_work\_submit} 只排队，实际在系统工作线程优先级上执行，与提交者不同 \\
\Tp{k\_work\_delayable} & 可延时、可重排的工作项 & \Fn{k\_work\_schedule} 与 \Fn{k\_work\_reschedule} 的区别在于是否重置已有延时 \\
\Fn{k\_work\_queue\_start} & 启动专用 workqueue（自带栈与优先级） & 未启动就往该队列 submit 是典型的静默不执行 \\
\Tp{k\_sem} & 计数信号量 & \Fn{k\_sem\_take} 的超时用 \Tp{K\_MSEC} 或 \Tp{K\_FOREVER}；中断里可 give 不可 take \\
\Tp{k\_mutex} & 互斥量（带优先级继承） & \Fn{k\_mutex\_lock} 不可重入，同线程二次加锁等于自锁 \\
\Tp{k\_msgq} & 定长消息队列（拷贝语义） & 满则按超时丢弃；容量单位是消息数而不是字节 \\
\Tp{k\_fifo} & 单链表 FIFO（传指针） & 元素需 4 字节对齐，链表头由内核占用，复用已释放节点会破坏链表 \\
\Tp{k\_pipe} & 字节流管道，生产与消费可异步 & 需要 get 与 put 两个环形缓冲，尺寸与写者数不匹配会出现假``满'' \\
\Tp{k\_mem\_slab} & 固定尺寸块池 & 实时路径取块用超时 0；块数按最坏并发算 \\
\Tp{k\_timer} & 软件定时器，回调在系统时钟中断上下文 & 回调必须极短且不得阻塞，长活转成 work item \\
设备驱动注册宏 & 老接口 \Tp{DEVICE\_DEFINE} 已由 \Tp{Z\_DEVICE\_*} 宏族取代，按状态与初始化优先级登记实例 & 具体宏名随版本变化，写之前先确认树内仍在；登记顺序决定依赖链，优先级填错会得到``已登记但外设未就绪'' \\
\Tp{west build -b board} & 为目标板配置并构建 & board 名与设备树必须同源；沿用旧 build 目录会带上上一次的 \Tp{.config} \\
\Tp{west build -p always} & 清理后全量重建 & \Tp{-p auto} 只按需清理；改了 Kconfig 默认值却没 prune 时，症状是``改了没生效'' \\
\Tp{CONFIG\_MAIN\_STACK\_SIZE} & main 线程栈字节数 & main 里有静态大数组或深层调用时这项最先撑不住 \\
\Tp{CONFIG\_{\Bk}NUM\_{\Bk}PREEMPT\_{\Bk}PRIORITIES} & 可抢占优先级数量 & 与协作类数量共同决定范围，越界值会被静默夹取 \\
\Tp{CONFIG\_{\Bk}COOP\_{\Bk}THREAD\_{\Bk}PRIORITIES} & 协作线程优先级数（旧名，部分版本仍用） & 版本相关，迁移时先 grep 树内定义，不要照网络写法直接抄 \\
\Tp{CONFIG\_TIME\_SLICING} & 同优先级轮转开关 & 关掉后同优先级任务只能靠主动让出，饿死风险上升 \\
\Tp{CONFIG\_{\Bk}ARCH\_{\Bk}HAS\_{\Bk}USERSPACE} & 架构能力位，由 Kconfig 推导 & 手工置位而硬件不支持会立刻崩；用户真正打开的是下一项 \\
\Tp{CONFIG\_USERSPACE} & 打开用户/内核隔离与对象权限检查 & 系统调用边界有固定开销；硬实时线程通常留内核态，只把接口层放用户态 \\
\Tp{CONFIG\_THREAD\_ANALYZER} & 导出线程与栈用量 & 需要输出通道配合，采样本身要计入预算 \\
\Tp{CONFIG\_TICKLESS\_IDLE} & tickless 低功耗 & 唤醒源延迟与最小睡眠误差要实测，长睡还会改变高精度定时器语义 \\
\Tp{CONFIG\_EVENT\_LOGGER} & 结构化事件日志后端 & 后端在低优先级线程写入，实时线程只允许非阻塞提交 \\
\Tp{status = \textquotedbl okay\textquotedbl} & 设备树里启用节点 & 多数节点缺省 \Tp{disabled}，驱动登记了也不实例化，症状是按名取设备返回空 \\
\Tp{interrupts} & 描述中断号、类型与触发方式 & 边沿/电平写错表现为丢中断或中断风暴；多 cell 时顺序依绑定文档 \\
\Tp{clocks} & 声明外设所需时钟源与分频 & 父时钟频率与驱动假设不一致时，波特率与周期全错却不报错 \\
\Tp{dmas} & 声明 DMA 通道与请求号 & 通道被两个节点复用会在完成回调上互相踩，吞吐出现偶发重复 \\
\end{longtable}

\section{Linux 实时 API 与开关}

Linux 侧的实时性由三件事共同决定：调度策略、内存驻留、中断落点。任一处没配好，另外两处的收益会被吃掉。

\begin{longtable}{>{\raggedright\arraybackslash\footnotesize}p{0.24\textwidth} >{\raggedright\arraybackslash\footnotesize}p{0.30\textwidth} >{\raggedright\arraybackslash\footnotesize}p{0.29\textwidth}}
\caption{Linux 实时 API、启动参数与检查点速查}\label{tab:apB-linuxrt} \\
\toprule 名称 & 作用 & 实时注意点 \\ \midrule
\endfirsthead
\toprule 名称 & 作用 & 实时注意点 \\ \midrule
\endhead
\bottomrule \endfoot
\endlastfoot
\Fn{sched\_setscheduler} & 设置策略与静态优先级（FIFO/RR/OTHER） & \Tp{SCHED\_FIFO} 下 100\% 忙等的错误任务会锁死该核，必须有超时或看门狗；需 \Tp{CAP\_SYS\_NICE} \\
\Fn{sched\_setattr} & 通用属性设置，含 deadline 的 runtime、deadline、period & 准入检查不过直接 EBUSY；总带宽受内核上限约束，改策略前先读回现值 \\
\Fn{clock\_gettime} & 读时钟（\Tp{CLOCK\_MONOTONIC}、\Tp{CLOCK\_BOOTTIME}、\Tp{CLOCK\_REALTIME}、\Tp{CLOCK\_TAI}） & 算时长别用会被步进的 \Tp{CLOCK\_REALTIME}；与网卡时间戳对齐用 \Tp{CLOCK\_TAI} \\
\Fn{clock\_nanosleep} & 绝对或相对睡眠（\Tp{TIMER\_ABSTIME}） & 周期任务用绝对时刻，相对睡眠会累积漂移；被信号打断后要重算剩余 \\
\Fn{timerfd\_create} & 定时器变成可读 fd，接入 poll/epoll & \Tp{TFD\_TIMER\_CANCEL\_ON\_SET} 在墙钟被改时给出取消事件；fd 设非阻塞并保证内存驻留 \\
\Fn{mlockall} & \Tp{MCL\_CURRENT} 加 \Tp{MCL\_FUTURE} 锁全部地址空间 & 缺 \Tp{MCL\_FUTURE} 时新映射仍会页错误；受 \Tp{RLIMIT\_MEMLOCK} 限制，失败必须硬退出 \\
\Fn{posix\_memalign} & 对齐分配（2 的幂） & 分配后仍需 \Fn{mlock} 并逐页预触碰，否则首次访问仍是页错误 \\
\Fn{mmap} & 匿名或文件映射，可控增长与驻留 & \Tp{MAP\_POPULATE} 提前建页，避开首次访问页错误；大页路径还要挂载对应文件系统 \\
\Tp{madvise(MADV\_HUGEPAGE)} & 建议改用透明大页 & TLB 命中率上升但缺页路径变长；后台整理会引入毫秒级尖峰 \\
\Fn{pthread\_{\Bk}setaffinity\_{\Bk}np} & 线程绑核 & 绑到被隔离的核才有效；同核上的 worker 与软中断仍会插队 \\
\Fn{eventfd} & 轻量计数事件通知，跨线程或内核到用户 & 二值与阻塞语义要选对，否则一次读只清一个计数 \\
\Fn{signalfd} & 把信号变成 fd 读取，取代异步处理函数 & 必须先把对应信号屏蔽，否则仍走默认处理 \\
\Fn{io\_uring\_setup} & 建立提交与完成环做异步 IO & 内核轮询线程模式会占 CPU 且优先级低，硬实时侧慎用；缓冲区需注册并常驻 \\
\Tp{isolcpus} & 启动期把指定核排除在通用调度外 & 现代内核建议与 cpuset 配合；隔离不等于中断也不落在该核 \\
\Tp{nohz\_full} & 指定核上减少周期性 tick & 只在核数足够时安全，需与 RCU offload、cpuset 组合；效果是消掉 tick 唤醒 \\
\Tp{rcu\_nocbs} & 把 RCU 回调 offload 到专门线程 & 未 offload 时被隔离核仍会被 RCU 唤醒，这是``绑核仍有毛刺''最常见的原因 \\
\Tp{irqaffinity} & 启动期把 IRQ 默认落点排除实时核 & 与运行时写 \Tp{/proc/irq/N/smp\_affinity} 可互相覆盖，要在开机脚本里统一 \\
\Tp{idle=poll} & 空闲轮询不进低功耗态 & 消掉唤醒延迟但发热与功耗显著上升，只用于可接受降级的专用机 \\
\Tp{processor.max\_cstate} & 限制最深 C 态 & 深 C 态唤醒是微秒级；编号与平台相关，先看 \Tp{cpuidle} 目录再取值 \\
\Tp{intel\_idle.max\_cstate=0} & 让 intel\_idle 只用浅 C 态 & 与上一项同时设置时以驱动参数为准，别期望叠加生效 \\
\Tp{nowatchdog} 或 \Tp{tsc=nowatchdog} & 关闭 NMI watchdog & 省掉每核 PMU 计数与 NMI 注入；代价是硬死锁无法被 NMI 检出 \\
\Tp{.../{\Bk}cpufreq/{\Bk}scaling\_{\Bk}governor} & 调频策略，实时侧用 \Tp{performance} & 只改 governor 不足以消除频率切换毛刺，还要处理 boost 与 cpuidle \\
\Tp{/{\Bk}sys/{\Bk}fs/{\Bk}cgroup/{\Bk}cpuset/{\Bk}...} & cpuset 限定可用核与内存节点 & 与 \Tp{isolcpus} 互补；cgroup v2 下遗漏内存节点会引入跨 NUMA 访存 \\
\Tp{/proc/interrupts} & 每核每 IRQ 计数 & 看两帧差值判断 IRQ 是否落在实时核；计数不涨说明亲和性没生效 \\
\Tp{/{\Bk}proc/{\Bk}irq/{\Bk}N/{\Bk}smp\_{\Bk}affinity} & 单 IRQ 的 CPU 掩码（十六进制写法） & 多字掩码顺序易写反；被标记为 managed 的 IRQ 不可改回 \\
\Tp{/{\Bk}proc/{\Bk}sys/{\Bk}kernel/{\Bk}sched\_{\Bk}rt\_{\Bk}runtime\_{\Bk}us} & RT 任务总带宽上限（缺省 950000/1000000） & 置 -1 关闭限流，误写的忙等就会拖死系统；关闭前先想清楚 \\
\Tp{.../debug/tracing/} & ftrace 实例目录（\Tp{current\_tracer}、\Tp{trace}、\Tp{buffer\_size\_kb}、\Tp{tracing\_on}） & 缓冲太小会丢最坏样本；取 trace 自身有开销，测量配置要单独存档 \\
\end{longtable}

\begin{lstlisting}[style=shellstyle,caption={隔离实时核并采样一轮的典型命令},label={lst:apB-linux-deploy}]
# 1) 内核命令行（bootloader 下发）
isolcpus=2-5 nohz_full=2-5 rcu_nocbs=2-5 irqaffinity=0
# 2) 把外设中断赶出实时核
echo f > /proc/irq/48/smp_affinity
# 3) 频率封顶
for c in /sys/devices/system/cpu/cpu[2-5]; do
  echo performance > $c/cpufreq/scaling_governor
done
# 4) 采样
cyclictest -t4 -p99 -i1000 -l300000 -m -h100 -q -a2-5
\end{lstlisting}

\section{测量与取证工具参数}

先定``测什么、在哪一层测''，再选工具。用户态唤醒延迟、硬件事件时间戳、总线周期是三条独立证据链，只有一条时结论不成立。

\begin{longtable}{>{\raggedright\arraybackslash\footnotesize}p{0.24\textwidth} >{\raggedright\arraybackslash\footnotesize}p{0.30\textwidth} >{\raggedright\arraybackslash\footnotesize}p{0.29\textwidth}}
\caption{延迟测量与取证工具参数速查}\label{tab:apB-tools} \\
\toprule 名称 & 作用 & 实时注意点 \\ \midrule
\endfirsthead
\toprule 名称 & 作用 & 实时注意点 \\ \midrule
\endhead
\bottomrule \endfoot
\endlastfoot
\Tp{cyclictest} & 用户态周期唤醒延迟基准，输出 min/avg/max 与直方图 & 结论只在当前机器与当前负载下成立；测前关掉图形与其他测试 \\
\Tp{-t N} & 线程数，每线程测一个 CPU & 线程数应等于待测核数，多出的线程本身是污染源 \\
\Tp{-p} 与 \Opt{priority} & 设定测量线程的 RT 优先级 & 必须高于被测负载优先级，否则测到的是排队时间而不是唤醒时间 \\
\Tp{-i} & 周期（ms） & 周期过小会让测量自身变成负载，常用 100 到 1000 \\
\Tp{-l} & 循环次数 & 看长尾要``小周期乘长时间''，样本不足时 max 只是下限 \\
\Tp{-m} & 测试前锁定内存 & 不加时页错误混入结果，冷启动样本被当成常态 \\
\Tp{-h} & 直方图上限（us） & 超上限计入溢出桶，长尾被压扁 \\
\Tp{-q} & 只打最终汇总 & 长跑可配合按时长汇总的选项，单位随版本略有差异 \\
\Tp{-a} & CPU 亲和掩码 & 要与 \Tp{isolcpus} 一致，否则线程被搬到 housekeeping 核 \\
\Tp{-c} & 时钟源选择 & 单调与墙钟两类结果不同，跨次对比必须固定 \\
\Tp{-S} & 信号方式唤醒的对称测量模式 & 便于把中断路径与唤醒路径分开分析 \\
\Opt{relative} 一类选项 & 选择相对或绝对睡眠语义 & 绝对语义更接近真实周期任务；选项名随版本略有差异，以 \Opt{help} 输出为准 \\
\Tp{oslat} & 专用核上的最坏单次延迟（polling 判据） & 阈值分桶用于量化罕见大值；至少留一个非隔离核给系统 \\
\Tp{hwlatdetect} & 检测硬件源（SMI、管理引擎）引起的延迟 & 检测期自身屏蔽中断，只在验证机跑；用于判定``非软件可解释''的毛刺 \\
\Tp{rtla timerlat} & 内核内测量唤醒路径，含 top、hist、detect 子命令 & 与 cyclictest 的差别是给出阶段拆分；detect 可在超阈值时自动抓 trace \\
\Tp{rtla oslat} & rtla 版 polling 最坏延迟 & 按桶统计，阈值取自基线而不是拍脑袋 \\
\Tp{trace-cmd record} & 抓取 ftrace 事件到 trace.dat & \Tp{-e} 逐类开事件（irq、sched、syscalls 一族）；开太多会溢出缓冲，丢掉最坏样本 \\
\Tp{trace-cmd record -p} & 用插件格式化输出 & 插件在写盘路径上，长 trace 建议先 \Tp{-o} 落文件再离线解析 \\
\Tp{perf stat} & 计数（cycles、instructions、cache-misses、page-faults） & PMU 计数与 NMI watchdog 争资源，关 watchdog 后数值更稳 \\
\Tp{perf record} & 采样调用栈与事件 & 采样频率过高会自身引入抖动；实时代码优先 tracepoint 而不是采样 \\
\Tp{perf sched latency} & 从调度事件复原等待与执行时间 & 需先用配套的 record 采集；符号缺失时指定同一次构建的 vmlinux \\
\Tp{ethtool -T} & 查询网卡时间戳能力（硬件/软件、PTP 过滤器） & 显示支持硬件发送打点，不代表所有队列都支持，需实测 \\
\Tp{ethtool -L} & 设置队列（channel）数 & 多队列会把同一条流分到不同中断；可预测性优先时宁可单队列 \\
\Tp{ptp4l} & 802.1AS/1588 协议栈守护进程 & 走硬件时间戳的开关必须打开；配置里的传输方式与时钟选择决定精度量级 \\
\Tp{phc2sys} & 在 PHC 与系统时钟之间同步 & 同步方向要与 \Tp{ptp4l} 一致，双向同时开会互相拉扯 \\
tsntool 一类工具 & 部分 SDK 提供的 TSN 配置命令行（门控、帧抢占） & 并非各发行版自带；不可用时改用 \Tp{tc} 与驱动接口，不要照搬不存在的选项名 \\
\Tp{tc qdisc ... netem} & 注入延迟、抖动、丢包、乱序做鲁棒性测试 & 分布参数只用于测试台，注入量与随机种子要写进报告 \\
\Tp{tc qdisc ... etf} & 启用按时刻发送（\Tp{clockid}、\Tp{delta}、\Tp{offload}） & 需要驱动支持发送时刻 offload；纯软件模式时 delta 要吸收软件滞后 \\
\Tp{tc qdisc ... taprio} & 下发 802.1Qbv 门控（\Tp{base-time}、\Tp{queues}、\Tp{schedule}） & 条目时间以纳秒表达并依赖同一时钟基准；base-time 设在未来才不与下发时延竞争 \\
\Tp{candump} & 显示 CAN 接口报文 & \Tp{-t} 时间戳格式变化会改变解析；终端打印本身就是负载 \\
\Tp{candump -l} & 落日志文件便于回放 & 写慢速文件系统会引入总线背压，测试台应落 tmpfs \\
\Tp{cansequence} & 发送递增序列统计丢序与丢帧 & 高负载时会灌满总线，掩盖错误帧统计 \\
\end{longtable}

\begin{note}[title={测量三查}]
一查时钟：所有工具是否共用同一时间基准（TAI、单调、开机时间）。二查负载：被测机上是否还跑着 offload 线程、图形、日志写盘。三查样本量与上限：直方图溢出桶非零时，max 只是下限而不是最坏值。三查任一不过，结论只报趋势不报数值。
\end{note}

\section{TSN 与总线关键参数}

\begin{longtable}{>{\raggedright\arraybackslash\footnotesize}p{0.24\textwidth} >{\raggedright\arraybackslash\footnotesize}p{0.30\textwidth} >{\raggedright\arraybackslash\footnotesize}p{0.29\textwidth}}
\caption{CAN、CAN-FD、EtherCAT 与 TSN 关键参数速查}\label{tab:apB-bus} \\
\toprule 名称 & 作用 & 实时注意点 \\ \midrule
\endfirsthead
\toprule 名称 & 作用 & 实时注意点 \\ \midrule
\endhead
\bottomrule \endfoot
\endlastfoot
\Tp{brt} & 位时间预分频，决定一个 tq 的长度 & 与晶振误差、时钟源档位耦合，改时钟树必须重算整套位时间 \\
\Tp{tq} & 位时间的最小刻度单位 & 一位由若干 tq 组成，总数与采样点位置共同决定时钟容差 \\
\Tp{sample-point} & 采样时刻在位内的百分比 & 常取 75 到 87.5 个百分点；逼近上限时相位余量小，长总线易出错误帧 \\
\Tp{prop-seg} & 传播段，吸收往返延迟 & 直接决定可挂总线长度，压缩它就等于缩短线长上限 \\
\Tp{phase-seg1} 与 \Tp{phase-seg2} & 相位缓冲段，用于硬同步与再同步 & phase-seg2 参与再同步幅度上限，配小则抖动吸收能力差 \\
CAN 仲裁段波特率 & 125k、250k、500k、1M 一类取值 & 速率与线长权衡；1 Mbps 只适合短支线，负载率超过 70\% 时长帧延迟急剧上升 \\
CAN-FD 数据段速率 & 仲裁段低速、数据段高速（如 500k 配 2M） & 两段位时间参数各自独立，只改一段会让采样点错配；帧变长还要同步放大接收缓冲 \\
EtherCAT 循环时间 & 主站一轮遍历全部从站的周期 & 与从站数量和转发延迟近似线性；低于硬件最小周期表现为随机丢同步 \\
EtherCAT 分布式时钟 & 从站用硬件时间戳对齐锁存输入 & 偏移与抖动未收敛时各站采样时刻不一致，控制环等效引入采样噪声 \\
\Tp{cycleTime} 与 \Tp{taskInterval} & 802.1Qbv 中描述门控周期与配置下发节奏的量 & 所有时间量以纳秒表达且必须同源时钟，单位错一个数量级就把周期打散 \\
802.1Qbv 门控列表 & 各优先级队列在周期内的开启时长与状态 & 队列数与优先级映射要一致；一帧跨两个条目时会被推迟到下一个开启窗口 \\
802.1Qbv base-time & 门控周期的绝对起点 & 设在未来某一刻并配合周期滚动；写``现在''会与下发时延竞争 \\
802.1AS 同步精度 & 全网时间基准对齐（gPTP） & 量级为亚微秒到几十纳秒，取决于硬件打点与透明时钟；无硬件打点只能到几十微秒 \\
\Tp{SO\_TXTIME} & 让报文按指定发送时刻出网卡，错误经扩展数据回报 & 需驱动与整形硬件配合；时钟与网卡 PHC 不同源时包被丢弃并报错误 \\
\Tp{SO\_TIMESTAMPING} & 请求接收或发送的硬件/软件时间戳 & 时间戳以控制消息异步到达，需独立接收路径；漏读会让时间戳队列堆积 \\
\end{longtable}

\section{实时循环与内存纪律的编码检查项}

最后一张表是评审清单：判据写成可静态检查或可开机自检的形式，症状写``第一次失败会长什么样''，便于快速对号。

\begin{longtable}{>{\raggedright\arraybackslash\footnotesize}p{0.24\textwidth} >{\raggedright\arraybackslash\footnotesize}p{0.30\textwidth} >{\raggedright\arraybackslash\footnotesize}p{0.29\textwidth}}
\caption{实时循环与内存纪律编码检查项}\label{tab:apB-discipline} \\
\toprule 检查项 & 判据 & 违反后的症状 \\ \midrule
\endfirsthead
\toprule 检查项 & 判据 & 违反后的症状 \\ \midrule
\endhead
\bottomrule \endfoot
\endlastfoot
周期路径内禁止动态分配 & 硬实时调用链上无 \Fn{malloc}、\Tp{new}、\Fn{k\_heap\_malloc}，靠调用图或静态检查核对 & 前几小时正常，之后偶发毫秒级停顿，或分配失败分支被走到 \\
语言子集与目标档位一致 & 规范里每条语言禁令都写明它属于哪一档（第 1 章 \S\ref{sec:ch01-langtier}），并按产物核对：段表无 \Tp{.ARM.exidx}、未定义符号无 \Tp{\_\_cxa\_*} 与 \Tp{\_Znwj} & 上电第一次输出或第一次分配落在最坏路径；小容量器件上功能段被展开表与格式化实现挤掉 \\
周期路径内禁止无界循环 & 每个循环有常量上界并在入口断言 & 观测不到的最长周期；外部输入触发时表现为卡住一个周期甚至任务饿死 \\
周期路径内禁止日志与格式化 & 热路径只写定长环形缓冲，格式化在低优先级侧 & 一次串口输出就是几百微秒到毫秒，直方图出现整齐的长尾 \\
ISR 时长上界 & 每个 ISR 有实测上界并计入周期预算，中断内无阻塞与格式化打印 & 关中断窗口叠加后 tick 丢失，表现为周期任务相位突变 \\
关中断窗口有界 & 临界区上界写入设计文档并被 trace 实测 & 中断延迟累积，外部事件在最坏情况下整窗口丢失 \\
ISR 与任务共享数据 & 单写者环形缓冲，或 volatile 与原子访问并明确内存序 & 读到半个更新的结构体；数值离谱却复现不了 \\
阻塞调用清单 & 热路径无 sleep 与无超时等待，所有超时都是确定值 & 任务永久阻塞或看门狗复位；CPU 满负荷却不产出 \\
栈预算 & 由栈用量文件或高水位反推，含 ISR 栈与最坏嵌套深度 & 小概率随机 HardFault 或跑飞，换编译器版本后现象改变 \\
内存驻留与首触碰 & 启动期完成页面触碰、页表建立与缓存预热 & 第一个周期显著慢于稳态，被误判为``预热问题'' \\
CPU 亲和与中断落点 & 实时核上无通用任务、无内核线程、无 RCU 回调 & 绑核之后仍有毛刺，且只在别的核跑重活时出现 \\
时钟源与步进 & 只用单调或 TAI 类时钟算时长，显式处理步进事件 & 时间同步调整系统时钟后出现负时长或超大间隔 \\
错误处理不引发重启风暴 & 失败路径进安全态并保持可观测，退避而非立即重试 & 一次瞬时故障导致反复复位，日志只剩同一串启动记录 \\
外部依赖有超时上界 & 所有总线与网络读取带超时，失败可降级 & 对端停机时本地停在等数据上，整条流水线停摆 \\
观测通道不影响被测 & trace 与统计有容量和速率预算，可整体关闭 & 开 trace 后延迟反而变好（丢样本）或明显变差（探针开销） \\
\end{longtable}

最后一句纪律：表中的数值（tick 频率、采样点百分比、优先级、超时）都是常见取值而非推荐值。任何被引用的数值都必须由当前这台机器、当前这份负载测出来，并把工具版本、内核命令行、\Tp{.config} 或 map 文件与数据一起存档，否则下次无法确认测的还是同一件事。

% LocalWords:  arm objcopy readelf riscv clang freertos zephyr west cyclictest
% LocalWords:  oslat hwlatdetect rtla phc ethtool candump taprio skbprio

% ============================================================
% ---- frag_apC.tex ----
\chapter{选型矩阵与排障手册}

\label{ap:apC-selection}

本附录是可翻查的速查件：先给平台与组件的选型矩阵与验收门限，再给按症状索引的排障条目，最后是上线前的勾选项。它不复述正文第 14 章的论证过程，只给出结论、判据与可执行动作。使用顺序建议是：立项时用第一节定档与设门限，出问题时按第二节的症状分组查，交付前按第三节逐条勾。

\section{平台与组件选型矩阵}

\subsection{平台档位对比}

\begin{guideline}[title={先定延迟档位，再谈平台偏好}]
把控制周期与最大允许偏差写成数字，档位就定了；在此之前的平台争论多半是口味之争。先用表~\ref{tab:apC-platform} 排除不可行档，再用内存下限、认证先例与团队能力三项收窄到一到两个候选。
\end{guideline}

\begin{longtable}{>{\raggedright\arraybackslash\footnotesize}p{0.11\textwidth}>{\raggedright\arraybackslash\footnotesize}p{0.16\textwidth}>{\raggedright\arraybackslash\footnotesize}p{0.16\textwidth}>{\raggedright\arraybackslash\footnotesize}p{0.16\textwidth}>{\raggedright\arraybackslash\footnotesize}p{0.10\textwidth}>{\raggedright\arraybackslash\footnotesize}p{0.11\textwidth}}
\caption{平台档位与九个选型维度对比}\label{tab:apC-platform} \\
\toprule
平台 & 可达延迟量级与最坏抖动 & 可界性论证难度与静态分配 & 外设与协议栈生态、认证先例 & 内存下限 & 团队与人才、许可与成本 \\
\midrule
\endfirsthead
\toprule
平台 & 可达延迟量级与最坏抖动 & 可界性论证难度与静态分配 & 外设与协议栈生态、认证先例 & 内存下限 & 团队与人才、许可与成本 \\
\midrule
\endhead
\bottomrule
\endfoot
\endlastfoot
裸机前后台 & 环路周期决定；中断响应可到亚微秒，抖动源为关中断段与嵌套深度，无队列排队项 & 最易论证：控制流可枚举，最坏路径由代码审查加计时实测；天然全静态 & 外设靠厂商 HAL 手写，协议栈自搭；无 RTOS 认证议题，但整机仍需自证 & RAM 数 KB 可用 & 上手最低，规模化后维护成本转移到人；无许可费 \\
FreeRTOS & tick 常为 1\,ms，硬边沿靠外设中断直驱可达几十微秒；抖动为中断响应加调度器临界区 & 低到中：调度器路径短、可用静态内存池；最坏执行时间需自建计时证据 & 生态最广（各家板级支持包与驱动样例多）；官方内核存在可借用的安全鉴定版本 & 十几 KB 只读存储、数 KB RAM & 人才最多；MIT 类许可，商用免版税 \\
Zephyr & 与 FreeRTOS 同档，无滴答模式加高分辨率定时器可做几十微秒；抖动取决于驱动是否线程化 & 中：设备树与抽象层引入间接性，边界要按配置逐项核；支持静态堆与零分配配置 & 协议栈（蓝牙、网络、CAN）现成度最高；认证以第三方安全鉴定分支为路径 & 数十 KB 起，含协议栈则上百 KB & 需接受上游节奏；Apache 类许可，无版税 \\
ThreadX 类商用 RTOS & 微秒级到几十微秒，事件组与字节池路径短且被长期实测 & 低：源码可读、行为文档化，厂商可提供时序与安全论证件 & 中间件（USB、图形、网络）成熟；航空航天与工业功能安全先例多 & 数 KB RAM & 需商业授权与版税；合同可含技术支持与认证件交付 \\
RTEMS & 几十微秒量级，周期任务与优先级带可精确配置 & 中：面向任务与周期服务器的模型清晰，文档可引用；移植层要自证 & 板级支持包数量有限但覆盖冷门架构；有航空与核电使用先例可参考 & 数百 KB 级 & 社区维护，宽松许可；招到熟悉的人最难 \\
NuttX & 毫秒级为主，POSIX 线程加驱动框架，微秒级需大量裁剪验证 & 中到高：POSIX 兼容层路径长，边界要逐子系统收缩 & 类 Unix 驱动模型、Shell 与应用生态；认证先例薄 & 数百 KB RAM & 许可友好；适合既要 POSIX 抽象又必须小内存的折中 \\
Linux 加 \Tp{PREEMPT\_RT} & 分位值可到几十微秒，最大值尖刺受系统管理中断、缺页与回收影响，常到数百微秒 & 高：内核路径不可完全枚举，需靠隔离、限流与实测统计建立置信 & 生态最强（协议栈、文件系统、工具链）；认证存在但范围严格受限 & 数十 MB 起，含服务更多 & 会写 Linux 的人多，会做实时调优的人少；许可合规要评估 \\
Hypervisor 分区混合临界 & 硬实时分区可达微秒到几十微秒，通用分区为毫秒级；跨分区通信增加尾延迟 & 高：需论证分区调度与内存隔离，形式化程度最高 & 可同时跑小型 RTOS 与 Linux；商用微内核有航空与功能安全先例 & 每分区独立预留，总量为各分区之和 & 需要虚拟化与认证双背景团队；授权与定制费用高 \\
\end{longtable}

\textbf{按延迟量级选档。}(1) \emph{亚微秒级硬实时}（电机电流环、硬件保护联锁、单粒子事件响应）：只有裸机直驱外设中断，或把该环路完全放进硬件（定时器、PWM、霍尔捕获）才有意义，RTOS 在此档只承担管理角色。(2) \emph{几十微秒}（多轴位置环、快速安全联锁、高采样传感）：FreeRTOS、裁剪后线程化驱动的 Zephyr、ThreadX 类商用核、RTEMS 同档；Linux 在此档只能靠独占核加关闭全部可变源，且必须接受偶发数百微秒尾。

(3) \emph{亚毫秒}（姿态控制、关节总线、视觉触发流水线）：小型 RTOS 稳态有余；Linux 加 \Tp{PREEMPT\_RT} 在做完隔离、绑核、关节能后可达标，验收看分位值而非平均值。(4) \emph{毫秒级软实时}（人机界面、日志、参数下发、非安全相关协调）：任何档都行，此时决策依据换成生态、人才与许可；把这类工作从实时核挪到管理核是更常见的正解。

\subsection{时间源与计数}

\begin{longtable}{>{\raggedright\arraybackslash\footnotesize}p{0.15\textwidth}>{\raggedright\arraybackslash\footnotesize}p{0.32\textwidth}>{\raggedright\arraybackslash\footnotesize}p{0.31\textwidth}}
\caption{时间源与计数手段的判据与禁忌}\label{tab:apC-clocksrc} \\
\toprule
选项 & 选型判据 & 禁忌 \\
\midrule
\endfirsthead
\toprule
选项 & 选型判据 & 禁忌 \\
\midrule
\endhead
\bottomrule
\endfoot
\endlastfoot
DWT 周期计数器（Cortex-M） & 需要亚微秒分辨率的单板内测量；与内核同频、无额外硬件；打戳开销为几条指令 & 深度睡眠后计数丢失；多核不共享；被调试器占用时读数不可靠，不能作为产品时间戳唯一来源 \\
SysTick 加滴答计数 & 调度与超时只需要毫秒粒度；全系统单一时基 & 用它评估微秒级抖动；滴答溢出回绕未处理；把滴答直接当日志时间戳上报 \\
通用定时器或性能计数寄存器 & 需要较长区间且独立于调度器的计数；性能分析场景 & 各核计数未对齐就跨核相减；低功耗状态下时钟门控使计数停摆 \\
x86 时间戳计数器 & 用户态可直读、单条指令；恒定频率平台适用 & 异构核频点不同或可变频；跨核未经同步即比较；把不变频保证之外的读数当时间基准 \\
网卡硬件打戳（\Tp{SO\_{\Bk}TIMESTAMPING}） & 跨设备时间必须与网络时基一致；需要亚微秒级收发起点 & 用软件打戳的时钟去对齐硬件打戳的时钟；读 \Tp{/{\Bk}dev/{\Bk}ptp0} 而不检查它是否正在步进 \\
IEEE 1588 或 gPTP 从时钟 & 多节点需要共同时基；允许周期性步进或频率牵引 & 控制环路读取被步进的时间戳；未定义失锁保持时间与牵引速率上下界 \\
\end{longtable}

\subsection{同步与中断方案}

\begin{longtable}{>{\raggedright\arraybackslash\footnotesize}p{0.15\textwidth}>{\raggedright\arraybackslash\footnotesize}p{0.32\textwidth}>{\raggedright\arraybackslash\footnotesize}p{0.31\textwidth}}
\caption{同步与触发方案的判据与禁忌}\label{tab:apC-sync} \\
\toprule
选项 & 选型判据 & 禁忌 \\
\midrule
\endfirsthead
\toprule
选项 & 选型判据 & 禁忌 \\
\midrule
\endhead
\bottomrule
\endfoot
\endlastfoot
定时器到期中断驱动任务 & 周期固定且需要硬件级抖动下限；把触发点交给硬件而非调度器 & 在中断体做超过几微秒的处理；到期后仍靠软件补偿累积漂移 \\
外部事件捕获（中断线、输入捕获） & 事件时刻本身是数据（编码器、过零、联锁）；需要锁存时刻与电平分离 & 未定义边沿选择与抖动抑制；把捕获值当计数用；不清挂起位就退出 \\
每源独立中断线 & 需要严格区分优先级且各源到达率不同 & 在引脚紧张的封装上照搬，导致线不得不复用后改成轮询 \\
共享中断线加状态回读 & 引脚与向量数量受限时；挂在同一线上的源都能在常数时间内枚举 & 状态读取顺序与清除顺序颠倒；未对无归属中断计数报警；一个慢设备让整线排队 \\
线程化中断（\Tp{request\_{\Bk}threaded\_{\Bk}irq}、二级处理） & 处理需要阻塞、需要被抢占、或耗时不可压缩 & 把延迟敏感的第一级动作也放进线程；未给中断线程设定实时优先级与亲和 \\
硬件时间触发表或同步链（SYNC、级联定时器） & 多通道需要同相位更新（三相 PWM、多轴同时发令） & 忽略主从启动顺序；未考虑跨芯片传播延迟；把软件写寄存器的时刻当同步时刻 \\
自旋加忙等轮询总线 & 等待时长有界且小于一次调度往返（数十微秒内）；已运行在专用核上 & 无界等待；等待对象是可能挂起的外部节点；在关中断状态下自旋 \\
\end{longtable}

\subsection{缓冲与进程间形态}

\begin{longtable}{>{\raggedright\arraybackslash\footnotesize}p{0.15\textwidth}>{\raggedright\arraybackslash\footnotesize}p{0.32\textwidth}>{\raggedright\arraybackslash\footnotesize}p{0.31\textwidth}}
\caption{缓冲与进程间形态的判据与禁忌}\label{tab:apC-ipc} \\
\toprule
选项 & 选型判据 & 禁忌 \\
\midrule
\endfirsthead
\toprule
选项 & 选型判据 & 禁忌 \\
\midrule
\endhead
\bottomrule
\endfoot
\endlastfoot
阻塞队列（互斥量加条件变量） & 允许等待、端到端延迟非关键、需要容量背压 & 放在硬实时路径上（唤醒延迟不可界）；跨优先级共用一个队列造成排队反转 \\
自旋锁加短临界区 & 临界区可在数十条指令内完成；持有者不会睡眠 & 锁内做输入输出、拷贝或分配；多核高争用；在无抢占核上长时间持有 \\
单生产者单消费者无锁环 & 一个采集线程到一个处理线程的经典形态；索引与数据同处一处可核对 & 容量非 2 的幂却做掩码取模；生产端无释放序、消费端无获取序；两端互相踢同一缓存行 \\
多生产者多消费者无锁队列 & 确实多对多且不能接受中心锁；实现经过形式化或长时间验证 & 为省事而选它；重复使用问题未处理；实时路径上带内存回收的通用实现 \\
共享内存加时间戳（序号锁存、双缓冲、三缓冲） & 数据量大、读侧可容忍重试或旧值、需要零拷贝 & 读侧不校验序号前后一致；标志与数据落在同一缓存行；未定义数据陈旧上限却用于控制 \\
消息中间件或 DDS 类发现机制 & 跨进程、跨主机、拓扑可变，需要主题级过滤 & 把发现与序列化延迟算进控制周期；在小内存实时节点上跑完整栈 \\
\end{longtable}

\subsection{总线与网络选型}

\begin{longtable}{>{\raggedright\arraybackslash\footnotesize}p{0.13\textwidth}>{\raggedright\arraybackslash\footnotesize}p{0.20\textwidth}>{\raggedright\arraybackslash\footnotesize}p{0.16\textwidth}>{\raggedright\arraybackslash\footnotesize}p{0.12\textwidth}>{\raggedright\arraybackslash\footnotesize}p{0.18\textwidth}}
\caption{现场总线与工业以太网对比}\label{tab:apC-bus} \\
\toprule
方案 & 可用控制周期与节点数 & 同步精度 & 是否需专用硬件 & 生态开放性 \\
\midrule
\endfirsthead
\toprule
方案 & 可用控制周期与节点数 & 同步精度 & 是否需专用硬件 & 生态开放性 \\
\midrule
\endhead
\bottomrule
\endfoot
\endlastfoot
CAN & 毫秒级为主，典型 1 到 4\,ms；数十节点，负载上限由仲裁决定 & 各节点自行计时，无全局时基，精度与帧往返及负载相关 & 控制器加收发器，普遍内置 & 标准公开，扩展依赖各家应用层 \\
CAN-FD & 亚毫秒可做到；仍受仲裁负载约束，节点数与 CAN 同 & 同上；数据段变短使抖动占比下降 & 需支持 FD 的控制器与收发器 & 标准公开，工具链与 CAN 兼容 \\
LIN & 数毫秒到十毫秒级；一条线一对主从，节点少 & 由主节点调度表决定，无亚毫秒保证 & 最简收发器 & 规范公开，成本最低 \\
FlexRay & 微秒级周期与静态时槽；数十节点 & 同步周期明确，可达微秒级 & 需专用控制器与芯片组 & 规范开放但器件与工具来源集中，成本高 \\
EtherCAT & 百微秒级（典型 125 微秒到 1\,ms）；可挂数百从站 & 分布式时钟可到亚微秒对齐 & 从站需专用交换芯片 & 规范公开，从站芯片供应商有限 \\
PROFINET IRT & 亚毫秒同步周期；节点数受规划限制 & 硬同步窗口，精度到微秒级 & 需支持该特性的交换机与控制器 & 受控生态，工程依赖厂商规划工具 \\
POWER{\Bk}LINK & 约 1\,ms 周期并可选等时阶段；上百节点 & 广播轮询给出确定窗口 & 需实时主站协议栈（可用软件实现） & 存在开源主站栈，标准化程度好 \\
Sercos III & 亚毫秒；轴类设备常见，节点数中等 & 微秒级生产者消费者模型 & 需专用从站芯片与授权 & 规范由协会管理，需会员与认证 \\
普通以太网加 TSN & 取决于配置，亚毫秒可规划 & 时基同步到亚微秒，转发有界取决于调度配置 & 需支持时间感知整流的交换机与网卡 & 标准公开，实际能力与芯片实现差距大 \\
普通 UDP 加 gPTP 时间戳 & 毫秒到亚毫秒尽力而为；节点数无硬上限 & 时间戳可比到微秒，投递时刻无保证 & 不需要 & 完全开放，调试最容易 \\
\end{longtable}

\textbf{普通 UDP 加时间戳就够用的边界。}同时满足四点时不必引入专用实时以太网：(1) 该流量是\emph{观测与记录}而非闭环控制，丢一帧只降低信息密度而不改变被控对象状态；(2) 接收端动作是幂等的（最新值覆盖、快照刷新），不存在必须按序按帧执行的命令流；(3) 时间戳来自与数据面一致的时基，下游只做相关与排序，不做精确触发；(4) 网络域可控：链路长度、交换机与广播流量在文档里写死，且已实测满载下的丢帧率与尾延迟。反过来，只要出现下列任一情形就换总线或上 TSN：闭环周期小于 1\,ms 且跨节点、需要同相位多轴发令、故障时必须可界地送达、或网络与其他业务共享且无法隔离。

\subsection{度量与验收门限}

\begin{longtable}{>{\raggedright\arraybackslash\footnotesize}p{0.18\textwidth}>{\raggedright\arraybackslash\footnotesize}p{0.24\textwidth}>{\raggedright\arraybackslash\footnotesize}p{0.22\textwidth}>{\raggedright\arraybackslash\footnotesize}p{0.13\textwidth}}
\caption{建议门限、测量手段与采样时长}\label{tab:apC-gate} \\
\toprule
指标 & 建议上限或门限 & 测量工具 & 采样时长 \\
\midrule
\endfirsthead
\toprule
指标 & 建议上限或门限 & 测量工具 & 采样时长 \\
\midrule
\endhead
\bottomrule
\endfoot
\endlastfoot
分位值 99 周期偏差 & 不超过控制周期的 5\%，且绝对值不超过设计标注的最小余量 & \Tp{cyclictest} 直方图、板内引脚打戳加逻辑分析仪 & 至少 10 万个周期，覆盖一个完整工况循环 \\
最大抖动（尾值） & 单独给数值并说明触发条件；不把平均值当上限 & 同上；关注直方图尾部桶而非仅报最大值 & 不少于 30 分钟，交付前建议 8 小时 \\
错过截止时间计数 & 安全相关任务为 0；非安全相关给出连续窗口内允许个数与恢复动作 & 任务内截止时间检查、事件缓冲计数 & 每个工况至少 1 小时连续 \\
喂狗裕量 & 实际喂狗周期落在超时窗内，裕量不少于 30\%；最长关中断段远小于超时 & 复位原因寄存器、备份区记录的心跳位图 & 全生命周期，含启动与掉电恢复 \\
每核 CPU 负载 & 实时核稳态不高于 50\% 至 60\%，留出突发与测量余量 & 进程占用率统计、空闲计数法、\Tp{mpstat} & 最坏工况连续 1 小时 \\
栈水位 & 历史最高使用不高于分配量的 70\% 至 75\%；逐任务统计 & \Fn{uxTask{\Bk}Get{\Bk}Stack{\Bk}HighWaterMark} 类接口、填充图案加金丝雀、静态估算 & 覆盖全部代码路径的压力测试之后读取 \\
丢帧率与误码 & 观测类给百分比上限；控制类给连续丢失不超过 N 帧 & 序号与校验统计、协议栈错误计数寄存器、抓包 & 与部署寿命相关，至少 24 小时 \\
中断到处理起点延迟 & 给出分位值与最大值两个数字，与中断预算一致 & 逐中断号计数配合打戳、\Tp{ftrace} 事件、引脚翻转 & 每种中断独立测，不少于 10 万次 \\
\end{longtable}

\section{排障手册}

下面按症状分组。每组先跑\emph{同一台硬件、同一二进制}的对照实验，一次只动一个变量；所有处置都要能用同一条测量手段复核，否则等于没修。取证最小动作集见清单~\ref{lst:apC-forensics}。

\begin{lstlisting}[style=shellstyle,caption={症状发作时先抓的最小证据集},label={lst:apC-forensics}]
date; uname -a; cat /proc/cmdline; cat /proc/interrupts
cat /sys/devices/system/cpu/isolated
grep -E "Model name|siblings|cpu cores" /proc/cpuinfo
cat /sys/devices/system/cpu/cpu0/cpufreq/scaling_governor
chrt -p 1234; taskset -pc 1234; cat /proc/1234/smaps_rollup
dmesg | tail; journalctl --since "-10min" --no-pager
\end{lstlisting}

\subsection{周期任务抖动变大与最大值尖刺}

这一组的共同前提是均值没变、只是尾部变肥。尾部问题不要去查平均耗时，直接按最大值与分位数定位，再按证据逐个排除可变源。

\begin{longtable}{>{\raggedright\arraybackslash\footnotesize}p{0.12\textwidth}>{\raggedright\arraybackslash\footnotesize}p{0.18\textwidth}>{\raggedright\arraybackslash\footnotesize}p{0.18\textwidth}>{\raggedright\arraybackslash\footnotesize}p{0.19\textwidth}>{\raggedright\arraybackslash\footnotesize}p{0.11\textwidth}}
\caption{抖动与尖刺的症状处置}\label{tab:apC-jitter} \\
\toprule
症状 & 先看的三个证据 & 常见成因 & 处置 & 复核方法 \\
\midrule
\endfirsthead
\toprule
症状 & 先看的三个证据 & 常见成因 & 处置 & 复核方法 \\
\midrule
\endhead
\bottomrule
\endfoot
\endlastfoot
整体抬升且尾部肥厚 & 连续两次 \Tp{/{\Bk}proc/{\Bk}interrupts} 的每行增量落在哪个核；任务的调度类与优先级（\Tp{chrt -{\Bk}p}）；直方图桶宽与上界 & 中断均衡服务把外设中断挪进隔离核；线程实际退回普通分时；桶过粗造成视觉放大 & 停用并屏蔽 \Tp{irqbalance} 服务；逐中断号把亲和写为管理核（\Tp{smp\_{\Bk}affinity\_{\Bk}list}）；固定桶宽后重测 & 重启后中断仍不在隔离核上；同一桶宽下 8 小时最大值不回升 \\
最大值尖刺与负载无关 & 当前频率的采样曲线；各空闲状态进入计数的增量；\Tp{turbostat} 的驻留统计 & 进入深度空闲态后的唤醒开销；调速器降频后首个周期爬坡 & 频率调速器设为性能模式；逐个把空闲状态的禁用位写 1；内核参数固定最大空闲状态与禁用等待指令；裸机则关深度睡眠 & 空闲状态计数不再增长；尖刺幅度降到目标内 \\
每隔固定时间一次尖刺 & 尖刺时刻与周期定时器列表对照；日志服务落盘时间戳；脏页回写阈值 & 周期服务（时间同步、日志轮转、缓存回写）与实时任务同核；RCU 回调集中执行 & 把周期服务错相并绑到管理核；回写阈值调低并限流；用 \Tp{nohz\_full} 与 \Tp{rcu\_nocpus} 把回调移出实时核 & 停用该服务一小时做对照，尖刺随之消失 \\
尖刺与打印时刻吻合 & 内核消息时间戳；串口或终端是否作为控制台；单进程写字节数 & 控制台同步输出；日志格式化在实时路径；小记录频繁强制落盘 & 测试期把控制台级别压到最低或去掉控制台；日志改环形缓冲加异步批量落盘；裸机串口发送改直访内存或做抽样丢弃 & 用引脚翻转量测打印路径占用时间，改前后对照 \\
只有刚启动那几秒抖 & 主缺页计数（进程统计文件字段）；常驻集增长曲线；锁页调用返回值 & 缺页与首次触碰：堆内存未预分配，代码与数据未进内存 & 启动时锁住当前与将来页面（\Tp{MCL\_FUTURE}）；所有缓冲在初始化阶段显式写一遍；大块改静态分配或带预填的映射；小型 RTOS 用静态池并预先建栈 & 缺页增量为 0；启动后 1 分钟与 1 小时曲线一致 \\
偶发单核尖刺，关掉超线程就好 & 每核线程数与兄弟逻辑核列表；尖刺核与其兄弟核上的负载；单实例双实例对照 & 超线程共享一级缓存与执行单元，邻居是日志或工作线程 & 以物理核为粒度隔离（兄弟核一并处理）或在固件里关超线程；把关键任务与其中断线程放同一物理核 & 开关超线程各跑 8 小时的对照数据 \\
毫秒级空洞，内核日志干净 & \Tp{turbostat} 的系统管理中断计数增量；固件事件日志与固件版本；同型号不同板对照 & 固件抢占（风扇、可靠性巡检、传统串口仿真）；深度睡眠组合状态 & 在固件设置里关闭非必要系统管理功能与传统仿真；把系统管理中断增量写进验收条款；必要时换无此机制的处理器 & 计数增量为 0 的时段内尖刺不再出现 \\
CPU 空闲但中断响应变慢 & 伪中断计数；中断线程是否常态处于运行态；抢占模型配置是否真生效 & 长时间关中断的驱动临界区；抢占模型退化为自愿抢占；外设直访内存抢总线 & 确认使用抢占实时内核；把输入输出移出锁；临时打开线程化中断以定位慢驱动；把直访内存带宽计入预算 & 硬件延迟探测工具与设优先级、不设优先级两组测量差值收窄 \\
抖动随任务数线性增长 & 就绪队列长度分布；每秒上下文切换次数；最长等待时间的归属 & 全局锁或单一队列把并发串行化；优先级反转（用了无继承的信号量）；线程池过小 & 每设备独立线程加优先级继承互斥量；队列按优先级分片；确认使用互斥量而非计数信号量 & 切换次数不再随负载线性；最长等待时间有界且被记录 \\
\end{longtable}

\subsection{中断变慢或不进中断}

\begin{longtable}{>{\raggedright\arraybackslash\footnotesize}p{0.12\textwidth}>{\raggedright\arraybackslash\footnotesize}p{0.18\textwidth}>{\raggedright\arraybackslash\footnotesize}p{0.18\textwidth}>{\raggedright\arraybackslash\footnotesize}p{0.19\textwidth}>{\raggedright\arraybackslash\footnotesize}p{0.11\textwidth}}
\caption{中断路径的症状处置}\label{tab:apC-isr} \\
\toprule
症状 & 先看的三个证据 & 常见成因 & 处置 & 复核方法 \\
\midrule
\endfirsthead
\toprule
症状 & 先看的三个证据 & 常见成因 & 处置 & 复核方法 \\
\midrule
\endhead
\bottomrule
\endfoot
\endlastfoot
中断计数完全不涨 & 时钟门控与使能位；中断控制器的使能与挂起寄存器；引脚实际电平（仪器） & 外设时钟未开；触发源未配置（边沿方向反、电平与脉冲选择错）；控制器使能位未置 & 按时钟树、外设配置、控制器使能、总使能的顺序逐项核对；逻辑分析仪同时看事件线与计数 & 计数随注入事件线性增长；断线即不增长 \\
计数涨但服务不过来、事件丢失 & 同一线上各源计数分布；状态寄存器被读取的次数；标志清除与读取的先后 & 共享线逐源轮询导致排队；先清标志后读状态吞掉同帧事件；无归属中断未报警 & 先读状态与数据再清挂起位；把慢设备移到独立线；对进入但无人认领的次数计数并设阈值报警 & 连续百万帧注入无丢失计数 \\
进了中断但整体延迟高 & 中断入口到退出的周期计数；浮点压栈配置位；栈深实测 & 硬件自动浮点压栈（惰性状态保存）；中断体里做日志与拷贝 & 中断体只清标志并入队；关闭惰性状态保存或保证中断体不碰浮点；把工作下移到线程 & 引脚打戳量测响应到退出时间，改前后对照 \\
只有某些中断慢 & 优先级分组寄存器值；同组内的子优先级排布；数值语义说明 & 把子优先级当抢占用；数值方向理解反（数值越小越优先）；关键源与噪声源同组 & 显式设定分组，抢占位只留必要位数；把关键源置于更高抢占级；在代码里写死语义注释 & 人为制造并发中断，观察嵌套与抢占顺序符合预期 \\
换了启动方式后中断全不进 & 向量表偏移寄存器实际值；链接输出的段起始地址；向量表内容与期望是否一致 & 向量表未重定位到应用起始；偏移不满足对齐要求；引导程序跳转不干净 & 按对齐要求设置向量表偏移；核对链接脚本中向量段与复位入口；启动早期即完成重定位 & 遍历触发每一个中断，逐一确认计数；核对复位原因 \\
冷启动或唤醒后头几次极慢 & 首次处理耗时与第 100 次的差；缓存使能顺序；锁相位标志读取 & 指令与数据缓存未热；闪存预取未生效；倍频未锁定就跑；首次缺页 & 唤醒路径显式等待锁定；把向量表与热点函数放入快速存储；上电即开缓存；启动阶段做一次预热调用 & 首次与稳态差值落在阈值内；重复唤醒一百次做统计 \\
\end{longtable}

\subsection{偶发复位与看门狗咬}

\begin{warning}[title={在中断里喂狗等于关掉监督}]
只要喂狗动作挂在中断或单一任务入口，看门狗就只能证明``时钟在跑''，不能证明系统在推进。正确形态是：各关键任务置心跳位，监督体确认位图齐全后才喂，并把位图快照写入复位后仍保留的备份存储。
\end{warning}

\begin{longtable}{>{\raggedright\arraybackslash\footnotesize}p{0.12\textwidth}>{\raggedright\arraybackslash\footnotesize}p{0.18\textwidth}>{\raggedright\arraybackslash\footnotesize}p{0.18\textwidth}>{\raggedright\arraybackslash\footnotesize}p{0.19\textwidth}>{\raggedright\arraybackslash\footnotesize}p{0.11\textwidth}}
\caption{复位与看门狗的症状处置}\label{tab:apC-reset} \\
\toprule
症状 & 先看的三个证据 & 常见成因 & 处置 & 复核方法 \\
\midrule
\endfirsthead
\toprule
症状 & 先看的三个证据 & 常见成因 & 处置 & 复核方法 \\
\midrule
\endhead
\bottomrule
\endfoot
\endlastfoot
狗咬但主循环看似正常 & 喂狗调用点位置；狗最终挂在哪个事件上；复位原因寄存器 & 中断里喂狗掩盖任务卡死；心跳位图未被真正检查 & 改为位图齐全才喂；多域级联狗；复位后第一时间打印并留存上次位图 & 注入死循环能可靠被咬；正常时不误咬 \\
窗口狗在轻负载下偶发咬 & 喂狗时刻相对窗口下界的分布；窗口边界配置；计数与时钟分频换算 & 喂得过早（落在允许窗口之外）；窗口宽度计算用错分频 & 只在窗口内固定位置喂；启用提前喂狗中断做统计；把窗口换算写成可核对常量 & 百万次喂狗时刻全部落在窗口内 \\
复位指向软件故障或未定义指令 & 复位标志与错误状态寄存器组合；出错地址；栈使用高水位 & 栈溢出触发保护；递归或异常路径未测；函数指针被踩 & 开启栈保护与警戒区域；逐任务栈水位实测；加递归深度上限 & 每任务栈水位保留 25\% 以上余量；灌满测试可定位 \\
复位与电机或继电器动作同时 & 供电跌落检测标志；复位时刻与外部通断电波形双通道对齐；地线走向 & 电源压降、地弹、共模干扰超过复位门限 & 加隔离、缓启动与独立电源域；把复位计数与外部事件时间戳一并归档 & 台架重复开关十万次以上复位计数为 0 \\
高温或振动后出现位翻转 & 纠错错误计数与地址；是否为单比特可纠正；温度记录 & 软错误；描述符与配置块无校验；时序余量不足 & 关键结构加纠错或校验和；周期性巡检并回写；必要时降频或提高电压余量 & 巡检日志与温度联合分析，错误率降到预算内 \\
长临界区导致误咬 & 关中断段的进入与退出打戳；全局搜索关中断调用点；狗超时设置 & 关中断里做大块拷贝、浮点或闪存写入；闪存写期间总线停顿 & 给关中断段设硬上限（数十微秒）；长操作分片；超时与最长关中断段保持安全比例 & 24 小时最长关中断段统计低于预算 \\
时钟切换窗口内跑飞 & 切换序列与手册一致性；切换期间的计数频率；锁定标志读取 & 切到未锁定的源；切换中分频与电源配置未同步更新 & 严格按序列切换并查询锁定；切换在关中断且位于快速存储中的函数里完成；切换后回读频率 & 反复切换百万次无异常；频率计核对每个阶段 \\
\end{longtable}

\subsection{数据错但时间正常}

\begin{longtable}{>{\raggedright\arraybackslash\footnotesize}p{0.12\textwidth}>{\raggedright\arraybackslash\footnotesize}p{0.18\textwidth}>{\raggedright\arraybackslash\footnotesize}p{0.18\textwidth}>{\raggedright\arraybackslash\footnotesize}p{0.19\textwidth}>{\raggedright\arraybackslash\footnotesize}p{0.11\textwidth}}
\caption{数据完整性问题的症状处置}\label{tab:apC-data} \\
\toprule
症状 & 先看的三个证据 & 常见成因 & 处置 & 复核方法 \\
\midrule
\endfirsthead
\toprule
症状 & 先看的三个证据 & 常见成因 & 处置 & 复核方法 \\
\midrule
\endhead
\bottomrule
\endfoot
\endlastfoot
直访内存缓冲读到旧值 & 直访内存目标地址与读侧地址；数据缓存是否开启；区域起始地址是否按行对齐 & 未做缓存维护，处理器写在缓存行而外设读的是内存旧行；起始非对齐使整行操作越界 & 接收前失效整行对齐区间、发送前清写；或把缓冲放入不可缓存区域（用保护区域配置） & 加序号与校验和，百万帧零不匹配 \\
标志更新了但数据半新半旧 & 标志与缓冲的地址是否落在同一缓存行；写入顺序；反汇编是否重排 & 同一行内标志与数据互相牵扯；无发布顺序保证，编译器与硬件都可重排 & 用对齐使标志与数据分行；发布用释放序、消费用获取序；跨核路径显式核对 & 反汇编确认屏障存在；满载长跑无半包 \\
双缓冲或序号偶发错乱 & 序号是否单调且不会回绕到有效值；读侧是否重试；写侧个数 & 单缓冲或双缓冲在慢读者下被覆盖；读侧未做前后可信检查 & 改用序号锁存（读前后各取一次，要求为偶数且相等）或三缓冲；序号位宽足够 & 一写多读压力下逐帧校验无撕裂 \\
单字节对、多字节错 & 结构是否紧凑布局；原子类型的对齐要求与是否需要锁；生成代码是否拆成多次访问 & 紧凑结构成员上做原子操作（可能根本不原子）；对齐不满足；误用内存拷贝语义 & 原子只用于对齐良好的标量；线上格式用紧凑布局加显式序列化；对每个原子类型断言免锁性质 & 单元与压力测试；跨编译器版本复测 \\
多核之间读到乱序 & 内存序使用点清单；屏障指令是否出现；运行平台是否为弱序 & 用最宽松弛序做指针发布；缺屏障；把必要的顺序保证当开销删掉 & 发布用释放序、消费用获取序，必要时用顺序一致；寄存器访问前后加数据同步屏障 & 在弱序架构（ARM）上验证而非只在 x86；反汇编核对 \\
加了易失限定还是不收敛 & 有多少共享变量仅靠易失限定；是否存在两两组合的``看起来原子''写法 & 易失限定不保证原子性也不保证顺序，却被当同步手段 & 易失限定只保留于外设寄存器访问；其余改原子类型或加锁；加静态检查规则禁止共享易失变量 & 静态扫描结果为零；行为在优化级别切换下稳定 \\
跨设备字段错位 & 抓取的原始字节与协议定义逐字段对照；填充与偏移；是否有显式打包函数 & 直接把内存结构整体拷贝上线；大小端与对齐差异 & 定义线上格式并显式做字节序转换与打包；加版本号与长度字段；双向交叉测试 & 两种端序交叉收发全部通过；版本不匹配能被拒绝 \\
\end{longtable}

\subsection{吞吐上不去与负载虚高}

\begin{longtable}{>{\raggedright\arraybackslash\footnotesize}p{0.12\textwidth}>{\raggedright\arraybackslash\footnotesize}p{0.18\textwidth}>{\raggedright\arraybackslash\footnotesize}p{0.18\textwidth}>{\raggedright\arraybackslash\footnotesize}p{0.19\textwidth}>{\raggedright\arraybackslash\footnotesize}p{0.11\textwidth}}
\caption{吞吐与负载问题的症状处置}\label{tab:apC-throughput} \\
\toprule
症状 & 先看的三个证据 & 常见成因 & 处置 & 复核方法 \\
\midrule
\endfirsthead
\toprule
症状 & 先看的三个证据 & 常见成因 & 处置 & 复核方法 \\
\midrule
\endhead
\bottomrule
\endfoot
\endlastfoot
满载但吞吐不涨 & 热点函数占比；自旋与等待类函数占比；等待时长分布 & 忙等自旋无退避；锁争用中心；单线程瓶颈被误算成全核负载 & 有界自旋后转阻塞或让出；改用优先级继承锁或单生产者单消费者环；把中心锁分片 & 空闲等待时间与占用率对照；吞吐随核数可扩展曲线 \\
每帧耗时逐段增大 & 各阶段拷贝次数与字节数；缓冲归属是否转移；内存池水位 & 解析、过滤、日志各拷一次；接口按值传大结构 & 传句柄或索引配合对象池；原位解析；给每帧拷贝字节数设上限并入指标 & 每帧拷贝字节数指标下降；周期耗时同步下降 \\
日志一开就掉性能 & 日志与格式化函数在调用栈中占比；串口波特率与等待时间；单帧日志字节数 & 同步串口输出；格式化在实时路径；全量打印 & 去掉生产环境控制台；改二进制记录加离线解码；抽样与限流；发送用直访内存 & 关闭日志与开启日志的对照差值落在预算内 \\
闪存磨损与回写尖刺并存 & 每帧写入字节乘以频率的写量估算；强制落盘次数；分区写放大 & 小记录频繁落盘；文件系统日志与擦除块放大 & 环形缓冲批量写；追加式掉电安全格式；把写预算写进设计并核算寿命 & 48 小时写入量不超过预算；落盘次数与抖动相关性消失 \\
任务很多、切换很密 & 每秒上下文切换数；每任务实际工作量；排队时长 & 任务切得过细；滴答频率取得过高 & 同源工作合并进同一周期链；降低滴答或用无滴答模式；精确触发交给硬件定时器 & 切换数下降且分位延迟不劣化 \\
\end{longtable}

\subsection{换编译器小版本就变差}

\begin{note}[title={把产物可追溯性当成前置条件}]
时序回归必须建立在可比较的产物上：保存每个版本的链接地址映射文件、镜像摘要与工具链摘要，并把``地址映射差异超阈值即阻断''做成自动门。否则问题会被归因给编译器，实际是变更不可追溯。
\end{note}

\begin{longtable}{>{\raggedright\arraybackslash\footnotesize}p{0.12\textwidth}>{\raggedright\arraybackslash\footnotesize}p{0.18\textwidth}>{\raggedright\arraybackslash\footnotesize}p{0.18\textwidth}>{\raggedright\arraybackslash\footnotesize}p{0.19\textwidth}>{\raggedright\arraybackslash\footnotesize}p{0.11\textwidth}}
\caption{工具链变更的症状处置}\label{tab:apC-toolchain} \\
\toprule
症状 & 先看的三个证据 & 常见成因 & 处置 & 复核方法 \\
\midrule
\endfirsthead
\toprule
症状 & 先看的三个证据 & 常见成因 & 处置 & 复核方法 \\
\midrule
\endhead
\bottomrule
\endfoot
\endlastfoot
小版本升级后尾延迟变差 & 两版本地址映射文件对比；热点函数大小与段偏移；缓存行与组分布是否变化 & 内联与对齐变动使热点跨行跨组；数据布局改变使并发访问落在同一行 & 把映射差异纳入回归门限；关键函数显式对齐；保留基线产物做交替测试 & 同板交替运行两版本各 8 小时，差异可复现且被解释 \\
只有发布版变差 & 优化档位对照（关闭全部优化、调试档、常规档、激进档）；逐项关闭单特性；未定义行为报警 & 自动向量化与激进内联改变路径；严格别名下的未定义行为被暴露；断言被去掉 & 用逐函数调优替代整档降级；先修未定义行为；关键路径保留断言与运行时检查 & 特性开关矩阵结果记录归档；修复后高档优化稳定 \\
开了跨阶段优化就不同 & 链接期跨模块内联清单；被合并或被删的符号；函数指针可达性 & 跨模块内联改变时序；段清理把必要符号去掉；链接顺序变动引起布局变动 & 跨阶段优化版本与测量结论绑定；对入口与中断表加保留属性；基线与验证版本一致 & 反汇编确认中断入口路径；符号表核对 \\
换了 C 库或板级支持包 & 版本锁定清单与镜像摘要；分配器行为差异；驱动补丁记录 & 分配策略与缓冲默认值变化；驱动补丁改了等待路径；库版本升级带来隐性拷贝 & 工具链、内核、库三者版本随产物归档；禁止自动拉取最新版本；升级走单独评审 & 用归档工具链重建并核对摘要；升级前后同条件对照 \\
\end{longtable}

\subsection{现场复现难}

\begin{lstlisting}[style=cppstyle,caption={发布版可保留的最小取证设施：事件环},label={lst:apC-evidence}]
// 事件环放在不被清零的备份段，复位后仍可读回
struct Event { uint32_t tick; uint16_t id; uint16_t aux; };
constexpr uint32_t kRingMask = 63u;
alignas(64) Event g_ring[kRingMask + 1];
alignas(64) std::atomic<uint32_t> g_head;   // 与数据分行

void log_event(uint32_t tick, uint16_t id, uint16_t aux) {
    g_ring[g_head.load(std::memory_order_relaxed)] = {tick, id, aux};
}
\end{lstlisting}

\begin{longtable}{>{\raggedright\arraybackslash\footnotesize}p{0.12\textwidth}>{\raggedright\arraybackslash\footnotesize}p{0.18\textwidth}>{\raggedright\arraybackslash\footnotesize}p{0.18\textwidth}>{\raggedright\arraybackslash\footnotesize}p{0.19\textwidth}>{\raggedright\arraybackslash\footnotesize}p{0.11\textwidth}}
\caption{复现与取证的症状处置}\label{tab:apC-repro} \\
\toprule
症状 & 先看的三个证据 & 常见成因 & 处置 & 复核方法 \\
\midrule
\endfirsthead
\toprule
症状 & 先看的三个证据 & 常见成因 & 处置 & 复核方法 \\
\midrule
\endhead
\bottomrule
\endfoot
\endlastfoot
现场与台架都不复现 & 复位后是否还留有事件环与原因寄存器内容；现场与台架的工况差异清单；统计口径是否一致 & 没有留取证设施，证据在复位瞬间被清空；现场工况（温度、供电、负载）无法在台架复现 & 发布版保留环形事件缓冲与复位前快照；建立现场工况复现表并按表扫描 & 人为制造一次复位，确认整条证据链可读回 \\
持续集成结果与本地差很多 & 运行环境（是否共享或虚拟化主机、邻居负载）；频率调节与节能状态；结果附带的运行环境指纹 & 虚拟化时间被偷走；共享主机邻居干扰；调速与空闲态未锁 & 使用独占硬件；锁定调速器并禁空闲态；隔离核；每次结果附环境指纹与内核命令行 & 同二进制在独占机与共享机各跑一轮，差值可量化 \\
直方图看不出问题 & 桶宽与上界设置；样本总量；报告用的是最大值还是分位数 & 桶过粗使尾部被合并；样本量远小于事件稀有度；只报平均值 & 桶宽取 1 微秒并让上界覆盖关注范围；样本量按周期数定（至少十万级）；同时报分位数与最大值 & 用两种桶宽重跑，尾部形状一致 \\
只在某温度或某批次出现 & 温升曲线与失效时刻的关联；多板同序列结果分布；供电边界记录 & 时序余量在边界；器件批次差异；用单点结论掩盖分布 & 温箱梯度扫描加供电边界扫描；多板并行同序列；结论写成裕量分布而非通过与否 & 边界条件下全部板通过长时间连续运行；裕量数据入库 \\
问题与``是否接调试器''相关 & 有无调试器两组的差异；断点与单步对缓存和等待状态的影响；追踪接口是否占用引脚 & 调试器改变时序与缓存状态；引脚被追踪功能占用；功耗与热状态不同 & 用无侵入手段（引脚打戳、事件环、追踪缓冲）替代断点；把无调试器作为验收配置 & 无调试器配置下复现，并与有调试器的结果对照 \\
\end{longtable}

\section{上线前检查清单}

\subsection{构建产物}

\begin{itemize}
\item 锁定编译器版本、板级支持包版本、内核与库版本，并把三者摘要随产物归档。
\item 确认构建可重复：同一份源码与工具链能重建出可核对的镜像摘要。
\item 归档链接地址映射文件与符号表，作为后续回归比较的基线。
\item 确认发布版本的优化档位与实际测量时使用的一致，不存在``测的是低档、发的是高档''。
\item 检查关键段的对齐与放置（中断向量表、热点代码、直访内存缓冲）与链接脚本一致。
\item 确认看门狗、时钟与引脚配置以固化配置进入产物，而不是靠上电默认值。
\item 去掉开发期开关：控制台日志、测试注入钩子、临时提权。
\end{itemize}

\subsection{启动与配置}

\begin{itemize}
\item 确认引导到应用的跳转路径完整，向量表偏移与对齐满足要求。
\item 核对时钟树：源、倍频锁定检查、分频与外设实际频率一致，切换有明确序列。
\item 检查外设与中断控制器的使能顺序，早期初始化不得访问未开时钟的外设。
\item 锁定电源与复位相关配置：掉电检测阈值、复位原因读取并留档。
\item 确认启动期完成页面预触碰与内存锁定，稳态不再发生缺页。
\item 记录启动时间与首帧输出时间，纳入验收数据。
\end{itemize}

\subsection{时间与调度}

\begin{itemize}
\item 确认实时进程的调度类与优先级，且在重启后仍生效（启动脚本或单元文件里显式设置）。
\item 锁定中断亲和：外设中断不落在隔离核，关键任务的中断线程与任务同物理核。
\item 确认隔离以物理核为粒度，超线程兄弟核一并处理。
\item 锁定频率调速器为性能模式并禁用深度空闲状态，核对该设置在掉电后仍生效。
\item 确认时间源选择：打戳分辨率、跨核可用性与溢出回绕处理。
\item 检查周期任务的相位关系，避免同相叠加造成排队峰。
\item 记录分位延迟与最大值，作为基线数据入库。
\end{itemize}

\subsection{内存与并发}

\begin{itemize}
\item 确认实时路径无动态分配：所有缓冲在初始化阶段建立。
\item 检查每任务与每线程栈的水位测量结果，保留余量不低于既定配额的四分之一。
\item 确认直访内存缓冲的缓存维护（接收前失效、发送前清写）或位于不可缓存区域。
\item 检查共享标志与共享数据不在同一缓存行，且各有明确归属者。
\item 确认跨核发布用释放序、消费用获取序，寄存器访问前后有屏障。
\item 清点易失限定用途：仅用于外设寄存器，未被当同步手段。
\item 检查队列满、内存池耗尽、分配失败路径的处理动作与报警方式已定义。
\end{itemize}

\subsection{外设与总线}

\begin{itemize}
\item 确认中断标志读取与清除顺序：先读状态与数据，再清挂起位。
\item 检查总线负载预算：最坏负载下的尾延迟与丢帧实测数据入库。
\item 确认节点地址、拓扑与线缆长度在文档中写死，并与现场一致。
\item 检查共享中断线上每个源都能被常数时间枚举，并有未认领计数报警。
\item 确认外部事件捕获的边沿选择与抖动抑制配置。
\item 检查总线错误恢复策略：错误计数阈值、恢复动作、恢复时间上下界。
\end{itemize}

\subsection{监督与失效}

\begin{itemize}
\item 确认喂狗条件为``关键任务心跳位图齐全''，不在中断里无条件喂。
\item 核对看门狗超时与最长关中断段的比例余量，并实测最长关中断段。
\item 检查窗口看门狗的喂狗时刻分布全部落在允许窗口内。
\item 确认复位原因与复位前快照在重启后可读，并写入持久记录。
\item 定义掉电、失同步、传感器断链、总线静默四类失效的检测时间与反应动作。
\item 执行一次故障注入（任务卡死、总线静默、供电跌落），确认监督路径如设计动作。
\end{itemize}

\subsection{测量与证据}

\begin{itemize}
\item 确认验收指标为数字，且含分位数与最大值两个口径。
\item 锁定测量手段、采样时长与桶宽，并把这些参数随结果一起归档。
\item 测量环境与目标部署环境的差异清单化（节能、超线程、后台服务、调试器）。
\item 确认发布版保留事件环形缓冲与打戳能力，取证不依赖断点。
\item 跑满一轮长时间连续运行（建议 8 小时以上，交付级建议 24 小时），数据入库。
\item 与上一版本的地址映射和分位延迟做回归对照，差异有解释。
\end{itemize}

\begin{bestpractice}[title={一次只动一个变量}]
排障时同时改两处配置，即使症状消失也无法归因，修复不可迁移到别的板子。把每次变更、对应的测量口径与结果一起记录，形成可复用的因果条目。
\end{bestpractice}

本手册没有一条``通用最优解''，只有在本机、本工况、本版本下测出来的数字。实时性的成立不来自某个漂亮的算法，也不来自某个被迷信的操作系统，而来自档位选择有据、可变源被逐个关闭、指标被长期测量、失效路径被注入验证。任何一次改善，若不能用同一测量手段复现，就尚未完成；任何一次异常，若复位后无证据留存，就只能等它再次发生。把测量与流程固定下来，交付的才是可预期的时间行为。

\end{document}
